\documentclass[11pt, french, openany]{report}

% preamble.tex — BTS FED GCF — Cours complet 2 ans
\usepackage[T1]{fontenc}
\usepackage[utf8]{inputenc}
\usepackage[french]{babel}
\usepackage{geometry}
\geometry{a4paper, margin=2.5cm}

% Mathématiques et unités
\usepackage{amsmath, amssymb, mathtools}
\usepackage{siunitx}
\sisetup{locale=FR, output-decimal-marker={,}}

% Tableaux et figures
\usepackage{booktabs, tabularx, multirow}
\usepackage{graphicx, float}
\usepackage{tikz}
\usetikzlibrary{shapes.geometric, arrows.meta, positioning}

% Mise en forme
\usepackage{xcolor}
\definecolor{btsprimary}{RGB}{0, 82, 147}
\definecolor{btssecondary}{RGB}{220, 237, 255}
\definecolor{btsaccent}{RGB}{255, 140, 0}

\usepackage{tcolorbox}
\tcbuselibrary{breakable, skins}

% Boîtes pédagogiques
\newtcolorbox{definition}[1][]{
  colback=btssecondary, colframe=btsprimary,
  fonttitle=\bfseries, title=Définition,#1
}
\newtcolorbox{methode}[1][]{
  colback=orange!10, colframe=btsaccent,
  fonttitle=\bfseries, title=Méthode,#1
}
\newtcolorbox{attention}[1][]{
  colback=red!8, colframe=red!60!black,
  fonttitle=\bfseries, title=Attention,#1
}
\newtcolorbox{exemple}[1][]{
  colback=green!8, colframe=green!50!black,
  fonttitle=\bfseries, title=Exemple résolu,#1
}

% En-têtes et pieds de page
\usepackage{fancyhdr}
\pagestyle{fancy}
\fancyhf{}
\fancyhead[L]{\small\textcolor{btsprimary}{BTS FED — Option GCF}}
\fancyhead[R]{\small\textcolor{btsprimary}{\leftmark}}
\fancyfoot[C]{\thepage}
\renewcommand{\headrulewidth}{0.4pt}

% Liens
\usepackage{hyperref}
\hypersetup{
  colorlinks=true, linkcolor=btsprimary,
  urlcolor=btsprimary, citecolor=btsprimary,
  pdftitle={Cours BTS FED — Option GCF},
  pdfauthor={Généré par agent Claude Code}
}

% Numérotation des exercices
\newcounter{exercice}[chapter]
\newcommand{\exercice}[1]{%
  \stepcounter{exercice}%
  \subsection*{Exercice \theexercice\ — #1}%
}


\begin{document}

% ============================================
% PAGE DE TITRE
% ============================================

\begin{titlepage}
    \begin{tikzpicture}[remember picture, overlay]
        % Bande latérale gauche en FedBlue
        \node[anchor=north west, inner sep=0pt] at (current page.north west) {
            \begin{tcolorbox}[
                enhanced, colback=FedBlue, colframe=FedBlue, arc=0pt, outer arc=0pt,
                width=1.5cm, height=\paperheight, boxrule=0pt, left=0pt, right=0pt, top=0pt, bottom=0pt
            ]
            \end{tcolorbox}
        };
        
        % Texte rotaté sur la bande latérale
        \node[rotate=90, anchor=west, white] at ([xshift=0.85cm, yshift=2cm]current page.south west) {
            \sffamily\Huge\bfseries RÉFÉRENTIEL
        };
        
        % Contenu principal
        \node[anchor=north west, text width=14cm] at ([xshift=3cm, yshift=-4cm]current page.north west) {
            % Titre département/diploma
            {\huge\sffamily\color{FedBlue} \textbf{BTS FLUIDES ÉNERGIES DOMOTIQUE}}\\[0.5cm]
            {\large\sffamily OPTION A : GÉNIE CLIMATIQUE ET FLUIDIQUE}\\[2.5cm]

            % Titre principal en gros
            {\fontsize{45}{50}\sffamily\bfseries\selectfont RÉFÉRENTIEL}\\[0.3cm]
            {\fontsize{45}{50}\sffamily\bfseries\selectfont PROFESSIONNEL}\\[0.8cm]
            {\color{FedOrange}\rule{10cm}{3pt}}\\[1cm]

            % Sous-titre
            {\large\sffamily\color{FedBlue}\bfseries Compétences et activités professionnelles}
        };
        
        % Informations au bas
        \node[anchor=south west] at ([xshift=3cm, yshift=2cm]current page.south west) {
            \sffamily
            \begin{tabular}{ll}
                \textbf{Établissement :} & Lycée Gaudier Brzeska \\
                \textbf{Enseignant :} & M. Dylan PIMONT \\
                \textbf{Module :} & Référentiel professionnel BTS FED
            \end{tabular}
        };
        
        % Décoration inférieure
        \fill[FedOrange] (current page.south west)
            rectangle ([yshift=8pt] current page.south east);
    \end{tikzpicture}
\end{titlepage}

\newpage

% ============================================
% TABLE DES MATIÈRES
% ============================================

\tableofcontents
\newpage

% ============================================
% CHAPITRES
% ============================================

% À partir d'ici, ajouter les chapitres avec \input{chapitres/XX-nom.tex}

% ============================================================
%  Chapitre 00 — Analyse dimensionnelle et remise à niveau
%  BTS FED option GCF — Cours complet
%  Créé le 2026-04-07
% ============================================================

\chapter{Analyse dimensionnelle et remise à niveau mathématique}

% ============================================================
\section{Objectifs pédagogiques}
% ============================================================

\subsection*{Compétences GCF mobilisées}

\begin{itemize}
  \item \textbf{C1-1} — Exploiter des sources d'information : récupérer et interpréter des
        données techniques (fiches produit, tableaux, courbes).
  \item \textbf{C3-1} — Analyser un système : identifier les paramètres et les grandeurs
        physiques pertinentes.
  \item \textbf{C3-3} — Dimensionner : maîtriser les transformations d'équations et les
        conversions d'unités.
  \item \textbf{C7-4} — Conduire des mesures : interpréter les résultats en cohérence avec
        les unités employées.
\end{itemize}

\subsection*{Savoirs associés GCF}

\begin{itemize}
  \item \textbf{S1} — Mathématiques appliquées : calcul littéral, résolution d'équations,
        exploitation de graphiques. (Niveau GCF\,: 3)
  \item \textbf{S4} — Sciences physiques : grandeurs, unités SI et dérivées, analyse
        dimensionnelle. (Niveau GCF\,: 2)
\end{itemize}

% ============================================================
\section{Introduction}
% ============================================================

L'analyse dimensionnelle est une méthode fondamentale en sciences de l'ingénieur permettant
de vérifier la cohérence physicienne d'une équation ou d'un calcul. En génie climatique, on
rencontre fréquemment des formules complexes mettant en jeu des puissances (W), des débits
(m³/h, kg/s), des pressions (Pa, bar), des températures (K, °C)…

Ce chapitre fournit les outils pour :
\begin{itemize}
  \item \textbf{Vérifier l'homogénéité} d'une formule avant de l'appliquer ;
  \item \textbf{Convertir des unités} avec assurance et rigueur ;
  \item \textbf{Transformer des équations} pour exprimer une inconnue en fonction d'autres
        paramètres ;
  \item \textbf{Exploiter des graphiques et abaques} : lecture de courbes, interpolation,
        extrapolation dans un réseau d'isobares ou d'isothermes.
\end{itemize}

La maîtrise de ces concepts garantit la fiabilité des dimensionnements et la détection rapide
d'erreurs.

% ============================================================
\section{Système international d'unités (SI)}
% ============================================================

\subsection{Les sept unités de base}

Le Système International (SI) repose sur sept unités de base, dont sont dérivées toutes les
autres :

\begin{table}[h]
  \centering
  \sffamily
  \caption{Unités de base du SI}
  \label{tab:unite-base}
  \begin{tabular}{|l|c|l|}
    \hline
    \textbf{Grandeur} & \textbf{Symbole} & \textbf{Unité SI} \\
    \hline
    Longueur & $\ell$, $L$, $d$ & \si{\meter} (m) \\
    Masse & $m$ & \si{\kilogram} (kg) \\
    Temps & $t$ & \si{\second} (s) \\
    Intensité de courant & $I$ & \si{\ampere} (A) \\
    Température thermodynamique & $T$ & \si{\kelvin} (K) \\
    Quantité de matière & $n$ & \si{\mole} (mol) \\
    Intensité lumineuse & $I_v$ & \si{\candela} (cd) \\
    \hline
  \end{tabular}
\end{table}

En génie climatique, les trois premières unités (longueur, masse, temps) sont prédominantes,
associées à la température thermodynamique.

\subsection{Unités dérivées courantes en CVC}

De nombreuses grandeurs dérivées des sept unités de base sont d'usage courant. Le tableau
ci-dessous en liste les plus importantes :

\begin{table}[h]
  \centering
  \sffamily
  \caption{Unités dérivées du SI usuelles en génie climatique}
  \label{tab:unite-derivee}
  \begin{tabular}{|l|c|c|l|}
    \hline
    \textbf{Grandeur} & \textbf{Symbole} & \textbf{Unité SI} & \textbf{Expression en base} \\
    \hline
    Vitesse & $v$ & \si{\meter\per\second} (m/s) & $\mathrm{m \cdot s^{-1}}$ \\
    Accélération & $a$ & \si{\meter\per\second\squared} (m/s²) & $\mathrm{m \cdot s^{-2}}$ \\
    Force & $F$ & \si{\newton} (N) & $\mathrm{kg \cdot m \cdot s^{-2}}$ \\
    Pression & $P$ & \si{\pascal} (Pa) & $\mathrm{kg \cdot m^{-1} \cdot s^{-2}}$ \\
    Énergie, travail & $E$, $W$ & \si{\joule} (J) & $\mathrm{kg \cdot m^2 \cdot s^{-2}}$ \\
    Puissance & $\dot{Q}$, $P$ & \si{\watt} (W) & $\mathrm{kg \cdot m^2 \cdot s^{-3}}$ \\
    Débit massique & $\dot{m}$ & \si{\kilogram\per\second} (kg/s) & $\mathrm{kg \cdot s^{-1}}$ \\
    Débit volumique & $\dot{V}$ & \si{\meter\cubed\per\second} (m³/s) & $\mathrm{m^3 \cdot s^{-1}}$ \\
    \hline
  \end{tabular}
\end{table}

% ============================================================
\section{Analyse dimensionnelle — Principe d'homogénéité}
% ============================================================

\subsection{Définition}

\begin{definition}
  Une équation physique est \textbf{dimensionnellement homogène} si, et seulement si, tous
  les termes de l'équation ont la même dimension. Chaque terme doit s'exprimer dans les mêmes
  unités (ou équivalentes via les unités dérivées).
\end{definition}

\textbf{Principe cardinal} : \emph{On ne peut ajouter ou soustraire que des grandeurs de même
dimension. On ne peut comparer que des grandeurs de même dimension.}

\subsection{Méthode — Vérifier l'homogénéité}

Soit l'équation de bilan thermique d'un échangeur :
\begin{equation}
  \dot{Q} = \dot{m}_f \cdot c_p \cdot \Delta T
  \label{eq:bilan-thermal}
\end{equation}

où :
\begin{itemize}
  \item $\dot{Q}$ : puissance thermique échangée (\si{\watt}, W)
  \item $\dot{m}_f$ : débit massique du fluide (\si{\kilogram\per\second}, kg/s)
  \item $c_p$ : capacité thermique massique (\si{\joule\per\kilogram\per\kelvin}, J/(kg·K))
  \item $\Delta T$ : différence de température (\si{\kelvin}, K)
\end{itemize}

\vspace{0.5cm}

\textbf{Analyse de dimension} :

\begin{align*}
  [\dot{Q}] &= \mathrm{W} = \mathrm{kg \cdot m^2 \cdot s^{-3}} \\
  [\dot{m}_f \cdot c_p \cdot \Delta T] &= [\mathrm{kg/s}] \times [\mathrm{J/(kg \cdot K)}] \times [
  \mathrm{K}] \\
  &= \mathrm{(kg \cdot s^{-1})} \times \mathrm{(kg^{-1} \cdot m^2 \cdot s^{-2} \cdot K^{-1})} \times
  \mathrm{K} \\
  &= \mathrm{kg \cdot s^{-1} \cdot kg^{-1} \cdot m^2 \cdot s^{-2} \cdot K \cdot K^{-1}} \\
  &= \mathrm{m^2 \cdot s^{-3}} = \mathrm{W}
\end{align*}

$\Rightarrow$ L'équation \eqref{eq:bilan-thermal} est dimensionnellement homogène. ✓

\subsection*{Exemple 1 — Détection d'erreur}

Supposons une formule proposée (erronée) :
\begin{equation*}
  P_{\text{pompe}} = \dot{m} + \Delta P
\end{equation*}

où $P_{\text{pompe}}$ est en W, $\dot{m}$ en kg/s et $\Delta P$ en Pa.

\[
  [\mathrm{W}] \stackrel{?}{=} [\mathrm{kg/s}] + [\mathrm{Pa}]
\]

Les dimensions ne correspondent pas :
\begin{itemize}
  \item Membre de gauche : $\mathrm{kg \cdot m^2 \cdot s^{-3}}$
  \item Membre de droite (1) : $\mathrm{kg \cdot s^{-1}}$
  \item Membre de droite (2) : $\mathrm{kg \cdot m^{-1} \cdot s^{-2}}$
\end{itemize}

$\Rightarrow$ La formule est \textbf{fausse}. La bonne formule est : $P_{\text{pompe}} = \dot{V}
\times \Delta P$

% ============================================================
\section{Conversions d'unités}
% ============================================================

\subsection{Principes généraux}

Une conversion d'unité consiste à exprimer une grandeur dans une unité différente, sans en
modifier la valeur physique. Trois approches sont possibles :

\begin{enumerate}
  \item \textbf{Utiliser un facteur de conversion} (rapport dimensionnel)
  \item \textbf{Passer par l'unité SI} de référence
  \item \textbf{Utiliser des tableaux prédéfinis}
\end{enumerate}

\subsection{Conversions courantes en génie climatique}

\subsubsection{Puissance et énergie}

\begin{table}[h]
  \centering
  \sffamily\small
  \caption{Conversions courantes — Énergie et puissance}
  \label{tab:conv-energie}
  \begin{tabular}{|l|c|l|}
    \hline
    \textbf{De} & \textbf{Vers} & \textbf{Facteur de conversion} \\
    \hline
    \si{\watt} (W) & \si{\kilowatt} (kW) & $\times 10^{-3}$ \\
    \si{\joule} (J) & \si{\kilo\watt\hour} (kWh) & $\times 2{,}778 \times 10^{-7}$ \\
    kWh & \si{\joule} (J) & $\times 3{,}6 \times 10^6$ \\
    Wh & kWh & $\times 10^{-3}$ \\
    cal & \si{\joule} (J) & $\times 4{,}1868$ \\
    kcal & \si{\joule} (J) & $\times 4186{,}8$ \\
    \hline
  \end{tabular}
\end{table}

\textbf{Exemple — Conversion W $\leftrightarrow$ kW} :

Une chaudière gaz produit une puissance de $35~\mathrm{kW}$.

\[
  P = 35 \text{ kW} = 35 \times 10^3 \text{ W} = 35\,000 \text{ W}
\]

\subsubsection{Pression}

\begin{table}[h]
  \centering
  \sffamily\small
  \caption{Conversions courantes — Pression}
  \label{tab:conv-pression}
  \begin{tabular}{|l|c|l|}
    \hline
    \textbf{De} & \textbf{Vers} & \textbf{Facteur de conversion} \\
    \hline
    \si{\pascal} (Pa) & bar & $\times 10^{-5}$ \\
    bar & \si{\pascal} (Pa) & $\times 10^5$ \\
    bar & kgf/m² & $\times 10\,197{,}16$ \\
    \si{\pascal} (Pa) & mH\textsubscript{2}O & $\times 1{,}02 \times 10^{-4}$ \\
    mH\textsubscript{2}O & \si{\pascal} (Pa) & $\times 9\,810$ \\
    bar & mH\textsubscript{2}O & $\times 10{,}2$ \\
    \hline
  \end{tabular}
\end{table}

\textbf{Exemple — Conversion Pa $\leftrightarrow$ bar} :

Un détendeur hydrodynamique supporte une perte de charge $\Delta P = 25\,000~\mathrm{Pa}$.

\[
  \Delta P = 25\,000 \text{ Pa} = 25\,000 \times 10^{-5} \text{ bar} = 0{,}25 \text{ bar}
\]

\subsubsection{Débit volumique et débit massique}

\begin{table}[h]
  \centering
  \sffamily\small
  \caption{Conversions courantes — Débits}
  \label{tab:conv-debit}
  \begin{tabular}{|l|c|l|}
    \hline
    \textbf{De} & \textbf{Vers} & \textbf{Facteur de conversion} \\
    \hline
    \si{\meter\cubed\per\hour} (m³/h) & \si{\meter\cubed\per\second} (m³/s) & $\times 1/3\,600 \approx 2{,}778 \times 10^{-4}$ \\
    \si{\meter\cubed\per\second} (m³/s) & m³/h & $\times 3\,600$ \\
    L/s & m³/h & $\times 3{,}6$ \\
    m³/h & L/s & $\times 1/3{,}6 \approx 0{,}2778$ \\
    m³/h & L/min & $\times 1000/60 \approx 16{,}67$ \\
    \hline
  \end{tabular}
\end{table}

\textbf{Relation débit massique / débit volumique} :

\[
  \dot{m} = \rho \cdot \dot{V}
\]

Avec $\rho$ densité (kg/m³).

\textbf{Exemple — Conversion m³/h $\leftrightarrow$ m³/s} :

Un extracteur d'air a un débit volumique $\dot{V} = 800~\mathrm{m^3/h}$.

\begin{align*}
  \dot{V} &= 800 \text{ m³/h} = 800 \times \frac{1}{3\,600} \text{ m³/s} \\
  &= 800 \times 2{,}778 \times 10^{-4} \text{ m³/s} \\
  &= 0{,}222 \text{ m³/s}
\end{align*}

Avec une densité air $\rho = 1{,}2~\mathrm{kg/m^3}$, le débit massique est :

\[
  \dot{m} = 1{,}2 \times 0{,}222 = 0{,}267 \text{ kg/s}
\]

\subsubsection{Température}

\begin{table}[h]
  \centering
  \sffamily\small
  \caption{Conversions courantes — Température}
  \label{tab:conv-temp}
  \begin{tabular}{|l|c|l|}
    \hline
    \textbf{De} & \textbf{Vers} & \textbf{Formule} \\
    \hline
    Celsius (°C) & Kelvin (K) & $T\,[\mathrm{K}] = \theta\,[\mathrm{°C}] + 273{,}15$ \\
    Kelvin (K) & Celsius (°C) & $\theta\,[\mathrm{°C}] = T\,[\mathrm{K}] - 273{,}15$ \\
    Celsius (°C) & Fahrenheit (°F) & $T\,[\mathrm{°F}] = \theta\,[\mathrm{°C}] \times \frac{9}{5} + 32$ \\
    \hline
  \end{tabular}
\end{table}

\textbf{Remarque importante} : Il convient de distinguer \textbf{température absolue} (K, utilisée
en thermodynamique) et \textbf{écart de température} (K ou °C, utilisé pour les différences).

Une température $\theta = 20~\mathrm{°C}$ correspond à $T = 293{,}15~\mathrm{K}$.

Un écart $\Delta T = 5~\mathrm{K} \leftrightarrow \Delta \theta = 5~\mathrm{°C}$.

% ============================================================
\section{Transformation d'équations et résolution}
% ============================================================

\subsection{Principes de manipulation algébrique}

La résolution d'équations repose sur l'application rigoureuse d'opérations mathématiques
identiques aux deux membres de l'égalité.

\textbf{Règles fondamentales} :
\begin{itemize}
  \item Additionner/soustraire le même terme aux deux membres
  \item Multiplier/diviser par le même terme non nul aux deux membres
  \item Élever à la même puissance (avec attention aux racines)
  \item Appliquer une fonction (log, exp, sin, …) aux deux membres
\end{itemize}

\subsection{Exemple 1 — Équation linéaire simple}

\textbf{Problème} : Exprimer la température de sortie d'un radiateur en fonction des autres
paramètres.

Bilan énergétique du radiateur :

\[
  \dot{Q} = \dot{m} \cdot c_p \cdot (T_{\mathrm{in}} - T_{\mathrm{out}})
\]

Chercher $T_{\mathrm{out}}$ :

\begin{align*}
  \dot{Q} &= \dot{m} \cdot c_p \cdot (T_{\mathrm{in}} - T_{\mathrm{out}}) \\
  \frac{\dot{Q}}{\dot{m} \cdot c_p} &= T_{\mathrm{in}} - T_{\mathrm{out}} \quad \text{(diviser par
  $\dot{m} \cdot c_p$)} \\
  T_{\mathrm{out}} &= T_{\mathrm{in}} - \frac{\dot{Q}}{\dot{m} \cdot c_p}
\end{align*}

\subsection{Exemple 2 — Équation non linéaire}

\textbf{Problème} : Exprimer le diamètre d'une conduite en fonction du débit et de la vitesse.

Formule de continuité :

\[
  \dot{V} = v \cdot A = v \cdot \pi \cdot \frac{d^2}{4}
\]

Chercher $d$ :

\begin{align*}
  \dot{V} &= v \cdot \pi \cdot \frac{d^2}{4} \\
  \frac{4 \cdot \dot{V}}{v \cdot \pi} &= d^2 \quad \text{(isoler $d^2$)} \\
  d &= \sqrt{\frac{4 \cdot \dot{V}}{v \cdot \pi}} \\
  d &= 2 \cdot \sqrt{\frac{\dot{V}}{v \cdot \pi}}
\end{align*}

\subsection{Exemple 3 — Équation avec logarithme}

\textbf{Problème} : Exprimer le nombre d'unités de transfert en fonction de l'efficacité d'un
échangeur.

Relation DTLM (Différence de Température Logarithmique Moyenne) :

\[
  \dot{Q} = U \cdot A \cdot \Delta T_{\mathrm{lm}}
\]

avec $\Delta T_{\mathrm{lm}} = \dfrac{\Delta T_1 - \Delta T_2}{\ln(\Delta T_1 / \Delta T_2)}$

Pour certaines configurations, on utilise : $\varepsilon = 1 - \exp(-\mathrm{NUT})$

Chercher $\mathrm{NUT}$ en fonction de $\varepsilon$ :

\begin{align*}
  \varepsilon &= 1 - \exp(-\mathrm{NUT}) \\
  \varepsilon - 1 &= - \exp(-\mathrm{NUT}) \\
  1 - \varepsilon &= \exp(-\mathrm{NUT}) \quad \text{(retirer les signes)} \\
  \ln(1 - \varepsilon) &= -\mathrm{NUT} \quad \text{(appliquer ln)} \\
  \mathrm{NUT} &= -\ln(1 - \varepsilon)
\end{align*}

% ============================================================
\section{Exploitation de graphiques et abaques}
% ============================================================

\subsection{Lecture directe}

Certaines données sont fournies graphiquement (courbes, nomogrammes, tables). Savoir extraire
une valeur nécessite :

\begin{itemize}
  \item \textbf{Identifier les axes} : abscisse (x), ordonnée (y), unités
  \item \textbf{Localiser le point} correspondant à la donnée cherchée
  \item \textbf{Lire la valeur} avec la précision appropriée à l'échelle
  \item \textbf{Respecter les unités} du graphique
\end{itemize}

\textbf{Exemple — Propriétés de l'eau} :

Un abaque standard donne $\rho_{\mathrm{eau}}$ et $\nu_{\mathrm{eau}}$ en fonction de la
température.

À $T = 60°C$ : on lit $\rho = 983{,}2~\mathrm{kg/m^3}$ et $\nu = 4{,}74 \times
10^{-7}~\mathrm{m^2/s}$.

\subsection{Interpolation linéaire}

Lorsque le point cherché se situe \textit{entre} deux points connus, on fait une interpolation
linéaire :

\[
  y = y_1 + (y_2 - y_1) \times \frac{x - x_1}{x_2 - x_1}
\]

\textbf{Exemple — Déduction d'une valeur intermédiaire} :

Tableau de viscosité cinématique de l'air (à 1 atm) :

\begin{center}
  \sffamily
  \begin{tabular}{|c|c|}
    \hline
    $T\,[\mathrm{°C}]$ & $\nu\,[\times 10^{-5}~\mathrm{m^2/s}]$ \\
    \hline
    0 & 13{,}28 \\
    20 & 15{,}89 \\
    40 & 18{,}97 \\
    \hline
  \end{tabular}
\end{center}

Chercher $\nu$ à $T = 25°C$ :

\begin{align*}
  \nu(25) &= 15{,}89 + (18{,}97 - 15{,}89) \times \frac{25 - 20}{40 - 20} \\
  &= 15{,}89 + 3{,}08 \times \frac{5}{20} \\
  &= 15{,}89 + 0{,}77 \\
  &= 16{,}66 \times 10^{-5}~\mathrm{m^2/s}
\end{align*}

\subsection{Lecture sur un réseau de courbes (diagramme de Mollier, …)}

Les diagrammes complexes (psychrométrique, Mollier, Ashrae) permettent la détermination
simultanée de plusieurs propriétés.

\textbf{Méthode générale} :
\begin{enumerate}
  \item Localiser deux propriétés connues (ex : T et HR)
  \item Suivre les isoçorbes/isobares/isothermes appropriées
  \item Lire les propriétés secondaires au point d'intersection (ex : humidité absolue,
        enthalpie)
  \item Exprimer le résultat dans l'unité du diagramme
\end{enumerate}

% ============================================================
\section{Exercices résolus}
% ============================================================

\subsection*{Exercice 1 — Vérification d'homogénéité}

\textbf{Énoncé} : Vérifier l'homogénéité dimensionnelle de la formule de perte de charge en
conduite :

\[
  \Delta P = \lambda \cdot \frac{\ell}{d} \cdot \frac{\rho \cdot v^2}{2}
\]

où $\lambda$ est un coefficient sans dimension (coefficient de frottement), $\ell$ la longueur
(m), $d$ le diamètre (m), $\rho$ la masse volumique (kg/m³), $v$ la vitesse (m/s).

\textbf{Solution} :

\begin{align*}
  [\Delta P] &= \mathrm{Pa} = \mathrm{kg \cdot m^{-1} \cdot s^{-2}} \\
  [\lambda \cdot \frac{\ell}{d} \cdot \frac{\rho \cdot v^2}{2}] &= [1] \times [\mathrm{m/m}]
  \times \frac{[\mathrm{kg/m^3}] \times [\mathrm{m/s}]^2}{[1]} \\
  &= [\mathrm{kg \cdot m^{-3} \cdot m^2 \cdot s^{-2}}] \\
  &= [\mathrm{kg \cdot m^{-1} \cdot s^{-2}}] \\
  &= [\mathrm{Pa}]
\end{align*}

$\Rightarrow$ L'équation est homogène. ✓

\subsection*{Exercice 2 — Conversion de débits}

\textbf{Énoncé} : Une pompe de circulation a un débit nominal $Q = 2\,400~\mathrm{L/min}$.
Exprimer ce débit en m³/h et en m³/s.

\textbf{Solution} :

\begin{align*}
  Q &= 2\,400 \text{ L/min} \times \frac{1~\mathrm{m^3}}{1\,000~\mathrm{L}} \times
  \frac{60~\mathrm{min}}{1~\mathrm{h}} \\
  &= 2\,400 \times \frac{60}{1\,000}~\mathrm{m^3/h} \\
  &= 144~\mathrm{m^3/h}
\end{align*}

Puis :

\begin{align*}
  Q &= 144~\mathrm{m^3/h} \times \frac{1~\mathrm{h}}{3\,600~\mathrm{s}} \\
  &= 0{,}04~\mathrm{m^3/s}
\end{align*}

\subsection*{Exercice 3 — Résolution d'équation}

\textbf{Énoncé} : On mesure une consommation énergétique annuelle $E = 45~\mathrm{kWh}$. La
puissance moyenne est estimée à $\bar{P} = 3{,}5~\mathrm{kW}$. Combien de temps (en heures)
l'équipement a-t-il fonctionné ?

\textbf{Solution} :

Relation : $E = \bar{P} \cdot t$

Isoler $t$ :

\[
  t = \frac{E}{\bar{P}} = \frac{45}{3{,}5} = 12{,}86~\mathrm{h}
\]

Soit environ \textbf{12 h 52 min}.

\subsection*{Exercice 4 — Interpolation sur un abaque}

\textbf{Énoncé} : Soit l'abaque des propriétés de l'eau (pression atm.) :

\begin{center}
  \sffamily
  \begin{tabular}{|c|c|c|}
    \hline
    $T\,[\mathrm{°C}]$ & $\rho\,[\mathrm{kg/m^3}]$ & $c_p\,[\mathrm{kJ/(kg \cdot K)}]$ \\
    \hline
    20 & 998{,}2 & 4{,}181 \\
    40 & 992{,}2 & 4{,}179 \\
    60 & 983{,}2 & 4{,}185 \\
    \hline
  \end{tabular}
\end{center}

Déterminer $\rho$ et $c_p$ à $T = 50°C$ par interpolation linéaire.

\textbf{Solution} :

Pour $\rho$ (entre 40°C et 60°C) :

\begin{align*}
  \rho(50) &= 992{,}2 + (983{,}2 - 992{,}2) \times \frac{50 - 40}{60 - 40} \\
  &= 992{,}2 + (-9) \times \frac{10}{20} \\
  &= 992{,}2 - 4{,}5 \\
  &= 987{,}7~\mathrm{kg/m^3}
\end{align*}

Pour $c_p$ (entre 40°C et 60°C) :

\begin{align*}
  c_p(50) &= 4{,}179 + (4{,}185 - 4{,}179) \times \frac{50 - 40}{60 - 40} \\
  &= 4{,}179 + 0{,}006 \times 0{,}5 \\
  &= 4{,}182~\mathrm{kJ/(kg \cdot K)}
\end{align*}

% ============================================================
\section{Diagrammes et schémas de synthèse}
% ============================================================

\subsection{Les sept dimensions fondamentales du SI}

La figure~\ref{fig:dimensions-si} présente les sept grandeurs de base du Système International
et leurs unités associées.

\begin{figure}[h]
  \centering
  \begin{tikzpicture}[
      node distance = 1.4cm and 1.8cm,
      dim/.style  = {draw, rounded corners=4pt, minimum width=3.2cm, minimum height=0.9cm,
                     font=\sffamily\small, align=center, thick},
      base/.style = {dim, fill=FedBlue!15, draw=FedBlue},
      center/.style = {draw, circle, fill=FedBlue, text=white, font=\sffamily\bfseries\small,
                       minimum size=1.8cm, align=center},
      arr/.style  = {-Stealth, thick, FedBlue!70}
    ]
    % Centre
    \node[center] (SI) {SI\\Dimensions};

    % 7 nœuds autour
    \node[base, above=1.6cm of SI]               (M) {\textbf{Masse} $[M]$\\\si{\kilogram} (kg)};
    \node[base, above right=1.1cm and 2.2cm of SI](L) {\textbf{Longueur} $[L]$\\\si{\meter} (m)};
    \node[base, right=2.8cm of SI]               (T) {\textbf{Temps} $[T]$\\\si{\second} (s)};
    \node[base, below right=1.1cm and 2.2cm of SI](I) {\textbf{Courant} $[I]$\\\si{\ampere} (A)};
    \node[base, below=1.6cm of SI]               (th){\textbf{Température} $[\Theta]$\\\si{\kelvin} (K)};
    \node[base, below left=1.1cm and 2.2cm of SI] (N) {\textbf{Quantité de matière} $[N]$\\\si{\mole} (mol)};
    \node[base, left=2.8cm of SI]                (J) {\textbf{Intensité lumineuse} $[J]$\\\si{\candela} (cd)};

    % Flèches
    \foreach \n in {M,L,T,I,th,N,J}
      \draw[arr] (SI) -- (\n);
  \end{tikzpicture}
  \caption{Les sept dimensions fondamentales du Système International d'unités}
  \label{fig:dimensions-si}
\end{figure}

\subsection{Méthode d'homogénéité — Schéma de raisonnement}

La figure~\ref{fig:methode-homogeneite} illustre la démarche pas-à-pas pour vérifier
l'homogénéité d'une équation physique (exemple : bilan thermique $\dot{Q} = \dot{m}\,c_p\,\Delta T$).

\begin{figure}[h]
  \centering
  \begin{tikzpicture}[
      node distance=0.55cm,
      step/.style  = {draw, rounded corners=4pt, fill=FedBlue!12, draw=FedBlue,
                       text width=8cm, minimum height=0.85cm, align=left,
                       font=\sffamily\small, inner sep=6pt},
      result/.style= {draw, rounded corners=4pt, fill=FedGreen!20, draw=FedGreen,
                       text width=8cm, minimum height=0.85cm, align=left,
                       font=\sffamily\small, inner sep=6pt},
      err/.style   = {draw, rounded corners=4pt, fill=FedOrange!20, draw=FedOrange,
                       text width=8cm, minimum height=0.85cm, align=left,
                       font=\sffamily\small, inner sep=6pt},
      arr/.style   = {-Stealth, thick, FedBlue}
    ]
    \node[step]   (s1) {\textbf{Étape 1} — Écrire l'équation : $\dot{Q} = \dot{m}\cdot c_p\cdot\Delta T$};
    \node[step, below=of s1] (s2) {\textbf{Étape 2} — Remplacer chaque grandeur par sa dimension :
      $[\dot{Q}]=\mathrm{W}=\mathrm{kg\cdot m^2\cdot s^{-3}}$};
    \node[step, below=of s2] (s3) {\textbf{Étape 3} — Développer le membre droit :
      $[\dot{m}][c_p][\Delta T] = (\mathrm{kg\cdot s^{-1}})(\mathrm{m^2\cdot s^{-2}\cdot K^{-1}})(\mathrm{K})$};
    \node[step, below=of s3] (s4) {\textbf{Étape 4} — Simplifier et regrouper les exposants};
    \node[result, below=of s4] (ok) {\textcolor{FedGreen!60!black}{\textbf{Résultat :}}
      $\mathrm{kg\cdot m^2\cdot s^{-3} = kg\cdot m^2\cdot s^{-3}}$ \quad \textbf{Équation homogène} \checkmark};

    \draw[arr] (s1) -- (s2);
    \draw[arr] (s2) -- (s3);
    \draw[arr] (s3) -- (s4);
    \draw[arr] (s4) -- (ok);

    % Annotation à droite de s4
    \node[err, right=1.2cm of s4, text width=4.2cm]
      (warn) {\textcolor{FedOrange!80!black}{\textbf{Si différent :}} équation fausse — vérifier coefficients ou formule};
    \draw[-Stealth, dashed, FedOrange, thick] (s4.east) -- (warn.west);
  \end{tikzpicture}
  \caption{Démarche de vérification d'homogénéité dimensionnelle (exemple thermique)}
  \label{fig:methode-homogeneite}
\end{figure}

\subsection{Grandeurs physiques courantes en CVC — Rappel visuel}

\begin{figure}[h]
  \centering
  \begin{tikzpicture}[
      col/.style={draw=FedBlue, fill=FedBlue!8, rounded corners=3pt,
                  minimum width=2.8cm, minimum height=0.75cm,
                  font=\sffamily\footnotesize, align=center},
      hdr/.style={draw=FedBlue, fill=FedBlue, text=white, rounded corners=3pt,
                  minimum width=2.8cm, minimum height=0.8cm,
                  font=\sffamily\small\bfseries, align=center},
      colO/.style={col, draw=FedOrange, fill=FedOrange!10},
      hdrO/.style={hdr, draw=FedOrange, fill=FedOrange},
      colG/.style={col, draw=FedGreen, fill=FedGreen!10},
      hdrG/.style={hdr, draw=FedGreen, fill=FedGreen},
      node distance=0pt
    ]
    % Colonne 1 — Thermique
    \node[hdr]  (h1) at (0,0)        {Thermique};
    \node[col, below=of h1] (r11)    {Puissance : \si{\watt} (W)};
    \node[col, below=of r11] (r12)   {Énergie : \si{\kilo\watt\hour} (kWh)};
    \node[col, below=of r12] (r13)   {Temp. : \si{\kelvin} / \si{\degreeCelsius}};
    \node[col, below=of r13] (r14)   {$c_p$ : J/(kg·K)};

    % Colonne 2 — Hydraulique
    \node[hdrO] (h2) at (3.2,0)      {Hydraulique};
    \node[colO, below=of h2] (r21)   {Pression : \si{\pascal} (Pa)};
    \node[colO, below=of r21] (r22)  {Débit vol. : m³/h};
    \node[colO, below=of r22] (r23)  {Débit mass. : kg/s};
    \node[colO, below=of r23] (r24)  {Vitesse : m/s};

    % Colonne 3 — Aéraulique / géométrie
    \node[hdrG] (h3) at (6.4,0)      {Géométrie/Air};
    \node[colG, below=of h3] (r31)   {Longueur : \si{\meter} (m)};
    \node[colG, below=of r31] (r32)  {Surface : \si{\meter\squared} (m²)};
    \node[colG, below=of r32] (r33)  {Volume : \si{\meter\cubed} (m³)};
    \node[colG, below=of r33] (r34)  {Masse vol. : kg/m³};
  \end{tikzpicture}
  \caption{Grandeurs physiques et unités courantes en génie climatique (CVC)}
  \label{fig:grandeurs-cvc}
\end{figure}

% ============================================================
\section{Données terrain — Extraits de catalogues fabricants}
% ============================================================

Ce chapitre prépare l'étudiant à exploiter des données techniques réelles. Les extraits
ci-dessous sont représentatifs des catalogues industriels Grundfos (circulateurs) et
Daikin (pompes à chaleur) — compétence \textbf{C1-1 : exploiter des sources d'information}.

\subsection{Extrait catalogue Grundfos — Circulateur MAGNA3 25-100}

La gamme \textbf{Grundfos MAGNA3} est représentative des circulateurs haute efficacité
(classe IE5) utilisés en chauffage et climatisation. Les données suivantes correspondent
aux conditions d'essai EN\,1151 (eau à \SI{60}{\celsius}, $\rho = \SI{983}{\kilogram\per\metre\cubed}$).

\begin{table}[H]
  \centering\sffamily\small
  \caption{Extrait performances Grundfos MAGNA3 25-100 — Points caractéristiques (eau 60\,°C, EN\,1151)}
  \label{tab:grundfos-magna3}
  \begin{tabular}{|c|c|c|c|c|}
    \hline
    \textbf{Débit} $\dot{V}$ & \textbf{HMT} & \textbf{Puiss. absorbée} $P_1$ & \textbf{Rendement} $\eta$ & \textbf{Vitesse} $n$ \\
    (\si{\metre\cubed\per\hour}) & (\si{\meter} CE) & (\si{\watt}) & (\%) & (tr/min) \\
    \hline
    0{,}0 & 10{,}0 &  85 &  0{,}0 & 3\,200 \\
    0{,}5 &  9{,}5 &  92 & 14{,}3 & 3\,100 \\
    1{,}0 &  8{,}6 & 105 & 22{,}6 & 2\,950 \\
    1{,}5 &  7{,}4 & 118 & 25{,}8 & 2\,750 \\
    2{,}0 &  6{,}0 & 128 & 25{,}8 & 2\,550 \\
    2{,}5 &  4{,}5 & 132 & 23{,}5 & 2\,350 \\
    3{,}0 &  2{,}6 & 128 & 16{,}9 & 2\,100 \\
    3{,}2 &  1{,}8 & 120 & 13{,}2 & 2\,000 \\
    \hline
  \end{tabular}
\end{table}

\textbf{Rappel — Conversion HMT $\leftrightarrow$ pression différentielle :}
\begin{equation}
  \Delta P\,[\si{\pascal}] = \rho \cdot g \cdot H_{MT}\,[\si{\meter}]
  \quad\Rightarrow\quad
  \Delta P \approx 983 \times 9{,}81 \times H_{MT} \approx 9\,644 \times H_{MT}
  \label{eq:hmt-dp}
\end{equation}

\begin{exemple}[title={Application — Lecture catalogue Grundfos MAGNA3}]
\textbf{Contexte :} Un circuit de chauffage présente une perte de charge totale
$\Delta P_{circuit} = \SI{45\,000}{\pascal}$ pour un débit nominal $\dot{V} = \SI{1{,}5}{\metre\cubed\per\hour}$.

\textbf{1 — Conversion en mètres de colonne d'eau (vérification dimensionnelle) :}
\[
  H_{circuit} = \frac{\Delta P}{\rho \cdot g}
  = \frac{45\,000\,[\si{\pascal}]}{983\,[\si{\kilogram\per\metre\cubed}] \times 9{,}81\,[\si{\metre\per\second\squared}]}
  = \frac{45\,000}{9\,644} \approx \SI{4{,}67}{\meter}~\text{CE}
\]

\textbf{Homogénéité :} $\dfrac{\mathrm{Pa}}{\mathrm{kg/m^3 \cdot m/s^2}}
= \dfrac{\mathrm{kg \cdot m^{-1} \cdot s^{-2}}}{\mathrm{kg \cdot m^{-2} \cdot s^{-2}}}
= \mathrm{m}$ \checkmark

\textbf{2 — Vérification du point de fonctionnement :}

À $\dot{V} = 1{,}5\,\mathrm{m^3/h}$, la MAGNA3 25-100 développe $H_{MT} = 7{,}4\,\mathrm{m}$ CE.
La pompe peut alimenter ce circuit ($7{,}4 > 4{,}67\,\mathrm{m}$).
Le point de fonctionnement réel (intersection courbe pompe / résistance réseau) se situe
entre $\dot{V} = 1{,}5$ et $\SI{2{,}0}{\metre\cubed\per\hour}$.

\textbf{3 — Rendement au point de fonctionnement :}
$\eta \approx 25{,}8\,\%$ — valeur typique d'un petit circulateur de chauffage.
\end{exemple}

\subsection{Extrait catalogue Daikin — Pompe à chaleur Altherma 3 R (air/eau)}

La gamme \textbf{Daikin Altherma 3 R} est représentative des PAC air/eau résidentielles.
Les COP sont mesurés selon EN\,14511 ; ils \textbf{diminuent} quand la température extérieure
baisse ou que la température de départ augmente.

\begin{table}[H]
  \centering\sffamily\small
  \caption{Extrait performances Daikin Altherma 3 R 8\,kW — COP selon conditions d'essai EN\,14511}
  \label{tab:daikin-cop}
  \begin{tabular}{|l|c|c|c|c|}
    \hline
    \textbf{Conditions} & $T_{ext}$ & $T_{départ}$ & \textbf{Puiss. calorifique} & \textbf{COP} \\
    & (\si{\degreeCelsius}) & (\si{\degreeCelsius}) & (\si{\kilo\watt}) & (—) \\
    \hline
    A7/W35 (nominale)      & $+7$  & 35 & 8{,}0 & 4{,}65 \\
    A7/W45                 & $+7$  & 45 & 8{,}1 & 3{,}40 \\
    A2/W35 (intermédiaire) & $+2$  & 35 & 7{,}3 & 3{,}80 \\
    A$-$7/W35 (hiver froid)  & $-7$  & 35 & 6{,}0 & 2{,}90 \\
    A$-$10/W35 (hiver froid) & $-10$ & 35 & 5{,}4 & 2{,}60 \\
    A$-$15/W35 (grand froid) & $-15$ & 35 & 4{,}5 & 2{,}20 \\
    A7/W55 (haute temp.)   & $+7$  & 55 & 8{,}5 & 2{,}55 \\
    \hline
  \end{tabular}
\end{table}

\textbf{Clé de lecture :} la notation \textbf{A7/W35} signifie
$T_{air\_ext} = +7\,\si{\degreeCelsius}$ / $T_{eau\_départ} = 35\,\si{\degreeCelsius}$.

\begin{exemple}[title={Application — Vérification dimensionnelle avec données Daikin}]
\textbf{Contexte :} Maison de $120\,\mathrm{m^2}$ à Strasbourg (zone H1c,
$T_{ext,base} = \SI{-11}{\celsius}$, $DJU = 3\,100$).
Besoin de pointe : $\Phi_{ch} = \SI{6{,}0}{\kilo\watt}$.

\textbf{1 — Puissance électrique absorbée en pointe ($T_{ext} = -15\,\si{\celsius}$, COP = 2,20) :}
\[
  P_{el,pointe} = \frac{\Phi_{ch}}{COP}
  = \frac{6{,}0\,[\si{\kilo\watt}]}{2{,}20\,[-]}
  = \SI{2{,}73}{\kilo\watt}
  \quad\text{Homogénéité :}\ \frac{\mathrm{kW}}{1} = \mathrm{kW}\ \checkmark
\]

À $-15\,\si{\celsius}$, l'Altherma 3R 8\,kW ne délivre que $4{,}5\,\mathrm{kW}$
$< 6{,}0\,\mathrm{kW}$ requis. \textbf{Appoint électrique nécessaire} ou modèle 12\,kW.

\textbf{2 — Consommation annuelle ($SCOP_{annuel} = 3{,}6$) :}
\[
  W_{el,annuel} = \frac{\Phi_{ch} \times DJU \times 24}{SCOP \times \Delta T_{base}}
  = \frac{6{,}0 \times 3\,100 \times 24}{3{,}6 \times 30}
  = \frac{446\,400}{108}
  \approx \SI{4\,133}{\kilo\watt\hour\per\year}
\]

\textbf{Homogénéité :} $\dfrac{[\si{\kilo\watt}] \times [\text{jours}] \times [\si{\hour\per\day}]}{[1] \times [\si{\kelvin}]}
= \dfrac{\mathrm{kW \cdot h}}{\mathrm{K} / \mathrm{K}} = \mathrm{kWh/an}$ \checkmark
\end{exemple}

\subsection*{Application — Schéma d'installation hydraulique avec symboles ISO GCF}

La figure~\ref{fig:installation-dim-gcf} illustre un circuit de chauffage bitubes
représenté avec les \textbf{symboles normalisés ISO GCF} des composants et annoté
avec les grandeurs physiques mobilisées lors de l'analyse dimensionnelle.

\tikzset{
  pics/pompe/.code={
    \draw[FedOrange, thick] (0,0) circle (0.42);
    \fill[FedOrange!55] (0,0)--(90:0.37)  arc(90:30:0.37) --cycle;
    \fill[FedOrange!55] (0,0)--(210:0.37) arc(210:150:0.37)--cycle;
    \fill[FedOrange!55] (0,0)--(330:0.37) arc(330:270:0.37)--cycle;
    \fill[FedOrange]    (0,0) circle (0.07);
  },
  pics/chaudiere/.code={
    \draw[FedRed, thick] (-0.55,-0.40) rectangle (0.55,0.40);
    \draw[FedRed, line width=1.4pt]
      (-0.10,-0.35)..controls(-0.25,0.05)and(-0.05,0.20)..(0,0.35)
      ..controls(0.05,0.20)and(0.25,0.05)..(0.10,-0.35);
  },
  pics/vase-expansion/.code={
    \draw[FedGreen, thick] (0,0) circle (0.38);
    \draw[FedGreen, thick] (-0.38,0)--(0.38,0);
    \fill[FedGreen!25] (-0.38,0) arc(180:360:0.38)--cycle;
  },
  pics/radiateur/.code={
    \draw[FedRed!75, thick] (-0.45,-0.35) rectangle (0.45,0.35);
    \foreach \x in {-0.22,0.0,0.22}{
      \draw[FedRed!50, thin] (\x,-0.35)--(\x,0.35);
    }
  },
}

\begin{figure}[H]
  \centering
  \begin{tikzpicture}[thick, >=Stealth, font=\small]

    % ---- Équipements — symboles ISO GCF ----
    \pic at (0, 0)     {chaudiere};
    \pic at (0,-2.8)   {pompe};
    \pic at (5.5, 0)   {radiateur};
    \pic at (2.75,1.2) {vase-expansion};

    % ---- Libellés équipements ----
    \node[font=\scriptsize\bfseries, text=FedRed]
      at (0,-0.65)    {Chaudière};
    \node[font=\scriptsize\bfseries, text=FedOrange]
      at (-0.8,-2.8)  {Pompe};
    \node[font=\scriptsize\bfseries, text=FedRed!80]
      at (5.5,-0.65)  {Radiateur};
    \node[font=\scriptsize\bfseries, text=FedGreen!70!black]
      at (2.75,1.80)  {Vase exp.};

    % ---- Tuyauterie départ (rouge) ----
    \draw[->, FedRed, very thick]
      (0.55,0) -- node[above, font=\scriptsize, text=FedRed]
        {$T_{dep}=70^\circ$C \quad $\dot{V}$\,[m³/h]}
      (5.05,0);

    % ---- Raccordement vase d'expansion sur le départ ----
    \draw[FedGreen, thick] (2.75,0)--(2.75,0.82);

    % ---- Tuyauterie retour (bleue) : radiateur → collecteur → pompe ----
    \draw[FedBlue, very thick] (5.5,-0.35)--(5.5,-4.2)--(0,-4.2)--(0,-3.22);

    % ---- Pompe → chaudière (sortie pompe) ----
    \draw[->, FedBlue, very thick]
      (0,-2.38) -- node[right, font=\scriptsize, text=FedBlue]
        {$T_{ret}=55^\circ$C}
      (0,-0.40);

    % ---- Encart formule bilan et grandeurs physiques ----
    \node[draw=FedBlue!35, fill=FedBlue!5, rounded corners=3pt,
          font=\scriptsize, inner sep=6pt, align=left, text width=6.5cm]
      at (3.2,-5.4) {
        \textbf{Bilan thermique :} $\dot{Q}=\dot{V}\cdot\rho\cdot c_p\cdot\Delta T$
        \hfill$[\si{\watt}]$\\[3pt]
        \textbf{Grandeurs :}
        $\dot{V}$\,[m³/h] \quad $\Delta P$\,[\si{\pascal}] \quad
        $\rho$\,[kg/m³] \quad $\dot{m}$\,[kg/s]
      };

  \end{tikzpicture}
  \caption{Circuit de chauffage bitubes — symboles ISO GCF et grandeurs physiques pour
           l'analyse dimensionnelle}
  \label{fig:installation-dim-gcf}
\end{figure}

% ============================================================
\section{Synthèse et bonnes pratiques}
% ============================================================

\begin{tcolorbox}[colback=FedBlue!10, colframe=FedBlue, coltitle=white, title=Checklist — Analyse dimensionnelle et conversions]
  Avant d'appliquer une formule ou un résultat de calcul :

  \begin{itemize}
    \item ✓ Vérifier l'homogénéité dimensionnelle de l'équation
    \item ✓ S'assurer que toutes les variables sont exprimées dans les bonnes unités
    \item ✓ Convertir les données d'entrée si nécessaire
    \item ✓ Effectuer le calcul
    \item ✓ Exprimer le résultat dans l'unité requise
    \item ✓ Valider l'ordre de grandeur du résultat
    \item ✓ Justifier chaque étape de la transformation d'équation pour éviter les erreurs
  \end{itemize}
\end{tcolorbox}

\finchapitre

% % ============================================================
%  Chapitre 01 — Bases de la thermique
%  BTS FED option GCF — Cours complet
%  Généré par agent Claude Code — 2026-04-02
% ============================================================

\chapter{Bases de la thermique}

% ============================================================
\section{Objectifs pédagogiques}
% ============================================================

\subsection*{Compétences GCF mobilisées}

\begin{itemize}
  \item \textbf{C2-3} — Décrire le fonctionnement d'un système : identifier les paramètres
        thermiques (températures, puissances, débits).
  \item \textbf{C3-3} — Dimensionner tout ou partie d'une installation : appliquer les
        méthodes normalisées de calcul thermique.
  \item \textbf{C5-1 / C5-2} — Recueillir et extraire les exigences réglementaires (RE2020,
        DTU, NF EN ISO 6946).
  \item \textbf{C7-4 / C7-6} — Conduire des mesures et interpréter les résultats
        (thermomètres, fluxmètres, sonde de surface).
  \item \textbf{C8-3 / C8-4} — Analyser les performances énergétiques et concevoir des
        actions correctives.
\end{itemize}

\subsection*{Savoirs associés GCF}

\begin{itemize}
  \item \textbf{S4} — Sciences physiques : thermodynamique, états de la matière, échanges
        d'énergie.
  \item \textbf{S5.1} — Réglementation thermique RE2020 : exigences sur les parois opaques
        et le coefficient \(U_{bat}\). (Niveau GCF\,: 3)
  \item \textbf{A1-1} — Thermique des tubes : flux linéique et isolation. (Niveau GCF\,: 2)
  \item \textbf{A1-2} — Hygrothermie : condensation en paroi, dispositions constructives.
        (Niveau GCF\,: 3)
  \item \textbf{A2-3} — Bilan thermique : déperditions, charges thermiques et hydriques.
        (Niveau GCF\,: 3)
  \item \textbf{B1-2} — Émission de chaleur : radiateurs, planchers chauffants,
        ventilo-convecteurs. (Niveau GCF\,: 3)
  \item \textbf{S10.1} — Métrologie appliquée à la thermique. (Niveau GCF\,: 3)
\end{itemize}

% ============================================================
\section{Introduction}
% ============================================================

Le génie climatique et fluidique (GCF) a pour mission de concevoir, installer et maintenir les
installations de chauffage, ventilation, climatisation (CVC) et de production d'eau chaude
sanitaire (ECS). L'ensemble de ces systèmes repose sur un socle commun : la maîtrise des
\textbf{transferts thermiques}.

Un technicien supérieur GCF doit être capable de :
\begin{itemize}
  \item quantifier les pertes ou les apports de chaleur dans un bâtiment ;
  \item dimensionner les équipements de production et d'émission de chaleur ou de froid ;
  \item vérifier la conformité d'une installation vis-à-vis de la réglementation RE2020 ;
  \item diagnostiquer un dysfonctionnement thermique et proposer une action corrective.
\end{itemize}

La thermique du bâtiment est au cœur de la transition énergétique. En France, le secteur du
bâtiment représente environ \SI{45}{\percent} de la consommation finale d'énergie et
\SI{27}{\percent} des émissions de \(\mathrm{CO_2}\). La réglementation RE2020, en vigueur
depuis le 1\textsuperscript{er} janvier 2022, exige une réduction drastique des déperditions
thermiques et un recours accru aux énergies renouvelables.

% ============================================================
\section{Notions fondamentales de thermique}
% ============================================================

\subsection{Chaleur, température et énergie thermique}

\begin{definition}
  \textbf{Température} (\(T\)) : grandeur intensive caractérisant le niveau d'agitation
  thermique des molécules d'un système. Elle s'exprime en kelvin (\si{\kelvin}) dans le
  Système International (SI) ou en degré Celsius (\si{\degreeCelsius}) en pratique courante.
  Relation de conversion : \(T\,[\si{\kelvin}] = \theta\,[\si{\degreeCelsius}] + 273{,}15\).
\end{definition}

\begin{definition}
  \textbf{Chaleur} (\(Q\)) : forme d'énergie transférée entre deux systèmes due à une
  différence de température. S'exprime en joule (\si{\joule}) ou en kilowattheure
  (\si{\kilo\watt\hour}). Rappel : \(\SI{1}{\kilo\watt\hour} = \SI{3600}{\kilo\joule}\).
\end{definition}

\begin{definition}
  \textbf{Puissance thermique} (\(\dot{Q}\) ou \(P_{th}\)) : énergie thermique échangée
  par unité de temps. S'exprime en watt (\si{\watt}) ou en kilowatt (\si{\kilo\watt}).
  \[
    \dot{Q} = \frac{\mathrm{d}Q}{\mathrm{d}t}
  \]
\end{definition}

\begin{definition}
  \textbf{Capacité thermique massique} (\(c_p\)) : quantité d'énergie nécessaire pour élever
  de \SI{1}{\kelvin} la température d'un kilogramme de matière à pression constante.
  Unité : \si{\joule\per\kilogram\per\kelvin}. Valeurs courantes :
  \begin{itemize}
    \item Eau liquide : \(c_p \approx \SI{4186}{\joule\per\kilogram\per\kelvin}\)
    \item Air sec : \(c_p \approx \SI{1006}{\joule\per\kilogram\per\kelvin}\)
    \item Béton : \(c_p \approx \SI{880}{\joule\per\kilogram\per\kelvin}\)
  \end{itemize}
\end{definition}

\subsection{Les trois modes de transfert de chaleur}

La chaleur se propage selon trois mécanismes distincts, souvent présents simultanément dans
une installation CVC :

\begin{itemize}
  \item \textbf{Conduction} : transfert au sein d'un milieu solide (ou fluide au repos) par
        contact direct entre molécules. Prépondérant dans les parois de bâtiment.
  \item \textbf{Convection} : transfert entre un fluide en mouvement et une paroi (ou
        inversement). Prépondérant dans les émetteurs de chaleur et les échangeurs.
  \item \textbf{Rayonnement} : transfert par ondes électromagnétiques, sans support matériel.
        Prépondérant pour les corps à haute température (brûleurs, surfaces chauffantes).
\end{itemize}

% ============================================================
\section{Conduction thermique}
% ============================================================

\subsection{Loi de Fourier}

La conduction thermique est régie par la loi de Fourier. Pour une paroi plane en régime
permanent, le flux thermique surfacique \(\varphi\) (en \si{\watt\per\meter\squared}) est
proportionnel au gradient de température :

\begin{equation}
  \varphi = -\lambda \cdot \frac{\mathrm{d}T}{\mathrm{d}x}
  \label{eq:fourier_diff}
\end{equation}

Pour une paroi homogène d'épaisseur \(e\) (\si{\meter}), de conductivité thermique
\(\lambda\) (\si{\watt\per\meter\per\kelvin}), avec une différence de température
\(\Delta T = T_1 - T_2\) (\si{\kelvin}), la loi de Fourier s'écrit sous forme intégrée :

\begin{equation}
  \varphi = \lambda \cdot \frac{\Delta T}{e}
  \qquad \Rightarrow \qquad
  \dot{Q} = \lambda \cdot \frac{S \cdot \Delta T}{e}
  \label{eq:fourier_integre}
\end{equation}

\textbf{Variables :}
\begin{itemize}
  \item \(\lambda\) : conductivité thermique du matériau [\si{\watt\per\meter\per\kelvin}]
  \item \(e\) : épaisseur de la paroi [\si{\meter}]
  \item \(\Delta T\) : différence de température entre les deux faces [\si{\kelvin}] ou
        [\si{\degreeCelsius}]
  \item \(S\) : surface de la paroi [\si{\meter\squared}]
  \item \(\dot{Q}\) : flux thermique total [\si{\watt}]
\end{itemize}

\textbf{Valeurs de conductivité thermique courantes en GCF :}

\begin{center}
\begin{tabular}{lcc}
  \toprule
  Matériau & \(\lambda\) (\si{\watt\per\meter\per\kelvin}) & Usage typique \\
  \midrule
  Béton lourd & 1,75 & Structure portante \\
  Brique pleine & 0,84 & Maçonnerie \\
  Brique creuse & 0,45 & Paroi légère \\
  Bois massif & 0,15 & Ossature bois \\
  Laine de verre & 0,035–0,040 & Isolation thermique \\
  Laine de roche & 0,035–0,042 & Isolation thermique \\
  Polystyrène expansé (PSE) & 0,030–0,038 & Isolation extérieure \\
  Polyuréthane (PUR) & 0,022–0,028 & Isolation haute performance \\
  Plâtre & 0,35 & Enduit intérieur \\
  Verre simple & 1,00 & Vitrage \\
  \bottomrule
\end{tabular}
\end{center}

\subsection{Résistance thermique}

\begin{definition}
  La \textbf{résistance thermique} d'une couche de matériau (résistance surfacique, notée
  \(R\)) est le rapport de l'épaisseur sur la conductivité thermique :
  \begin{equation}
    R = \frac{e}{\lambda}
    \qquad \left[\si{\meter\squared\kelvin\per\watt}\right]
    \label{eq:resistance_thermique}
  \end{equation}
  Elle représente la capacité d'un matériau à s'opposer au passage de la chaleur : plus
  \(R\) est grand, plus l'isolation est efficace.
\end{definition}

Pour une paroi multicouche composée de \(n\) couches en série, les résistances s'additionnent :

\begin{equation}
  R_{tot} = R_{si} + \sum_{i=1}^{n} \frac{e_i}{\lambda_i} + R_{se}
  \label{eq:resistance_totale}
\end{equation}

avec :
\begin{itemize}
  \item \(R_{si}\) : résistance superficielle intérieure. Valeur standard : \SI{0,13}{\meter\squared\kelvin\per\watt}
  \item \(R_{se}\) : résistance superficielle extérieure. Valeur standard : \SI{0,04}{\meter\squared\kelvin\per\watt}
  \item \(e_i / \lambda_i\) : résistance de chaque couche de matériau
\end{itemize}

\begin{definition}
  Le \textbf{coefficient de transmission thermique} (ou transmittance) \(U\) est l'inverse
  de la résistance thermique totale :
  \begin{equation}
    U = \frac{1}{R_{tot}}
    \qquad \left[\si{\watt\per\meter\squared\per\kelvin}\right]
    \label{eq:coefficient_U}
  \end{equation}
  Il caractérise les déperditions thermiques à travers une paroi. La RE2020 fixe des valeurs
  maximales de \(U\) selon le type de paroi.
\end{definition}

Le flux thermique total traversant une paroi de surface \(S\) est alors :

\begin{equation}
  \dot{Q} = U \cdot S \cdot \Delta T = \frac{\Delta T}{R_{tot}} \cdot S
  \label{eq:flux_total}
\end{equation}

\begin{methode}[title=Calcul de la résistance thermique d'une paroi multicouche]
  \begin{enumerate}
    \item \textbf{Identifier} toutes les couches de la paroi (de l'intérieur vers
          l'extérieur) et noter leur épaisseur \(e_i\) et leur conductivité \(\lambda_i\).
    \item \textbf{Calculer} la résistance de chaque couche :
          \(R_i = e_i / \lambda_i\).
    \item \textbf{Ajouter} les résistances superficielles :
          \(R_{si} = \SI{0,13}{\meter\squared\kelvin\per\watt}\) (intérieur) et
          \(R_{se} = \SI{0,04}{\meter\squared\kelvin\per\watt}\) (extérieur).
    \item \textbf{Sommer} toutes les résistances pour obtenir \(R_{tot}\).
    \item \textbf{Calculer} le coefficient \(U = 1/R_{tot}\).
    \item \textbf{Vérifier} la conformité vis-à-vis de la réglementation (RE2020).
  \end{enumerate}
\end{methode}

% ============================================================
\section{Convection thermique}
% ============================================================

\subsection{Convection naturelle et forcée}

\begin{definition}
  La \textbf{convection naturelle} (ou libre) se produit lorsque les mouvements du fluide
  sont uniquement dus aux différences de masse volumique engendrées par les variations de
  température. Elle est responsable, par exemple, du panache de chaleur au-dessus d'un
  radiateur.
\end{definition}

\begin{definition}
  La \textbf{convection forcée} se produit lorsqu'un dispositif mécanique (pompe, ventilateur)
  impose le mouvement du fluide. Elle est caractéristique des ventilo-convecteurs, des
  échangeurs à plaques et des centrales de traitement d'air (CTA).
\end{definition}

La convection forcée est nettement plus efficace que la convection naturelle : les coefficients
d'échange sont en général \num{5} à \num{10} fois plus élevés.

\subsection{Loi de Newton — coefficient d'échange convectif}

Le flux thermique échangé par convection entre un fluide et une paroi est donné par la
\textbf{loi de Newton} :

\begin{equation}
  \dot{Q} = h \cdot S \cdot (T_{fluide} - T_{paroi})
  \label{eq:newton_convection}
\end{equation}

\textbf{Variables :}
\begin{itemize}
  \item \(h\) : coefficient d'échange convectif [\si{\watt\per\meter\squared\per\kelvin}]
  \item \(S\) : surface d'échange [\si{\meter\squared}]
  \item \(T_{fluide}\) : température du fluide en écoulement [\si{\degreeCelsius}]
  \item \(T_{paroi}\) : température de la surface [\si{\degreeCelsius}]
\end{itemize}

\textbf{Ordres de grandeur du coefficient \(h\) :}

\begin{center}
\begin{tabular}{lc}
  \toprule
  Situation & \(h\) (\si{\watt\per\meter\squared\per\kelvin}) \\
  \midrule
  Convection naturelle dans l'air & 2–10 \\
  Convection forcée dans l'air (ventilation) & 20–100 \\
  Convection forcée dans l'eau (circuit hydraulique) & 500–10\,000 \\
  Ébullition de l'eau & 2\,000–50\,000 \\
  Condensation de vapeur d'eau & 5\,000–100\,000 \\
  \bottomrule
\end{tabular}
\end{center}

Les résistances superficielles \(R_{si}\) et \(R_{se}\) de la norme NF EN ISO 6946 sont
liées aux coefficients d'échange convectif de surface :
\(R_{si} = 1/h_i = 1/7{,}7 \approx \SI{0,13}{\meter\squared\kelvin\per\watt}\) et
\(R_{se} = 1/h_e = 1/25 = \SI{0,04}{\meter\squared\kelvin\per\watt}\).

% ============================================================
\section{Rayonnement thermique}
% ============================================================

\subsection{Loi de Stefan-Boltzmann}

Tout corps à une température supérieure au zéro absolu émet un rayonnement électromagnétique.
Pour un corps noir parfait, la puissance rayonnée par unité de surface est donnée par la
\textbf{loi de Stefan-Boltzmann} :

\begin{equation}
  M^{\circ} = \sigma \cdot T^4
  \qquad \left[\si{\watt\per\meter\squared}\right]
  \label{eq:stefan_boltzmann}
\end{equation}

avec \(\sigma = \SI{5,67e-8}{\watt\per\meter\squared\per\kelvin\tothe{4}}\) la constante
de Stefan-Boltzmann et \(T\) la température en \si{\kelvin}.

\begin{definition}
  \textbf{Corps noir} : corps idéal qui absorbe l'intégralité du rayonnement incident, quelle
  que soit la longueur d'onde et l'angle d'incidence. Il constitue la référence en matière de
  rayonnement thermique.
\end{definition}

\begin{definition}
  \textbf{Émissivité} (\(\varepsilon\)) : rapport de l'émittance d'un corps réel à celle du
  corps noir à la même température. Adimensionnel, compris entre 0 et 1.
  Pour un corps réel (corps gris) : \(M = \varepsilon \cdot \sigma \cdot T^4\).
\end{definition}

\textbf{Valeurs d'émissivité courantes :}
\begin{center}
\begin{tabular}{lc}
  \toprule
  Matériau / Surface & Émissivité \(\varepsilon\) \\
  \midrule
  Corps noir (référence) & 1,00 \\
  Peinture blanche mate & 0,90–0,95 \\
  Béton, plâtre & 0,85–0,95 \\
  Verre & 0,84–0,90 \\
  Bois vernis & 0,80–0,90 \\
  Aluminium anodisé & 0,77–0,84 \\
  Aluminium poli (film de pare-vapeur) & 0,04–0,06 \\
  Acier inoxydable poli & 0,10–0,20 \\
  \bottomrule
\end{tabular}
\end{center}

\begin{attention}
  La loi de Stefan-Boltzmann utilise la \textbf{température absolue en kelvin}, jamais en
  degré Celsius. Une erreur de conversion entraîne des résultats faux d'un facteur très
  important. Vérifier systématiquement : \(T\,[\si{\kelvin}] = \theta\,[\si{\degreeCelsius}]
  + 273{,}15\).

  De plus, le flux net échangé par rayonnement entre deux surfaces dépend de leurs facteurs
  de forme et de leurs émissivités respectives. La formule simplifiée \(M = \varepsilon
  \sigma T^4\) ne donne que l'émittance propre d'une surface, pas le flux net échangé.
\end{attention}

\subsection{Rayonnement en pratique CVC}

En génie climatique, le rayonnement intervient notamment dans :
\begin{itemize}
  \item les \textbf{panneaux rayonnants} et planchers chauffants (part de rayonnement dans
        le confort thermique) ;
  \item les \textbf{apports solaires} à travers les vitrages (pris en compte dans les bilans
        d'été, RE2020) ;
  \item les \textbf{surfaces basse émissivité} des vitrages doubles ou triples, qui réduisent
        le rayonnement vers l'extérieur.
\end{itemize}

% ============================================================
\section{Bilan thermique et enthalpie}
% ============================================================

\subsection{Premier principe de la thermodynamique — système ouvert}

En régime permanent, pour un système ouvert (fluide en circulation continue), le premier
principe s'écrit :

\begin{equation}
  \dot{Q} - \dot{W} = \dot{m} \cdot (h_{sortie} - h_{entrée})
  \label{eq:premier_principe_ouvert}
\end{equation}

En l'absence de travail mécanique (\(\dot{W} = 0\), hypothèse valable pour les circuits
hydrauliques sans pompage à l'intérieur du volume de contrôle) :

\begin{equation}
  \dot{Q} = \dot{m} \cdot (h_{sortie} - h_{entrée})
  \label{eq:bilan_enthalpique}
\end{equation}

\subsection{Enthalpie spécifique}

\begin{definition}
  L'\textbf{enthalpie spécifique} \(h\) est une fonction d'état définie par :
  \begin{equation}
    h = u + P \cdot v \qquad \left[\si{\joule\per\kilogram}\right]
    \label{eq:enthalpie}
  \end{equation}
  où \(u\) est l'énergie interne massique, \(P\) la pression et \(v\) le volume massique.
  Pour un liquide ou un gaz parfait, la variation d'enthalpie à pression constante s'écrit :
  \(\Delta h = c_p \cdot \Delta T\).
\end{definition}

\subsection{Puissance thermique d'un fluide caloporteur}

Pour un circuit de chauffage ou de climatisation où un fluide (eau, air) transporte de
l'énergie thermique, la puissance échangée est :

\begin{equation}
  \boxed{
    \dot{Q} = \dot{m} \cdot c_p \cdot \Delta T
  }
  \label{eq:puissance_fluide}
\end{equation}

\textbf{Variables et unités :}
\begin{itemize}
  \item \(\dot{Q}\) : puissance thermique [\si{\watt}] ou [\si{\kilo\watt}]
  \item \(\dot{m}\) : débit massique du fluide [\si{\kilogram\per\second}] ou
        [\si{\kilogram\per\hour}]
  \item \(c_p\) : capacité thermique massique du fluide [\si{\joule\per\kilogram\per\kelvin}]
        \\ (eau liquide : \SI{4186}{\joule\per\kilogram\per\kelvin} ;
           air : \SI{1006}{\joule\per\kilogram\per\kelvin})
  \item \(\Delta T = T_{sortie} - T_{entrée}\) : écart de température [\si{\kelvin}] ou
        [\si{\degreeCelsius}]
\end{itemize}

\begin{attention}
  Le débit volumique \(\dot{V}\) [\si{\liter\per\hour} ou \si{\meter\cubed\per\hour}] est
  lié au débit massique par : \(\dot{m} = \rho \cdot \dot{V}\), avec \(\rho\) la masse
  volumique du fluide. Pour l'eau à \SI{60}{\degreeCelsius} : \(\rho \approx
  \SI{983}{\kilogram\per\meter\cubed}\). On utilise souvent l'approximation
  \(\rho \approx \SI{1000}{\kilogram\per\meter\cubed}\) pour les calculs courants.
  Conversion pratique : \(\SI{1}{\meter\cubed\per\hour} = \SI{1000}{\liter\per\hour}\).
\end{attention}

La formule pratique avec le débit volumique en \si{\liter\per\hour} et \(c_p\) de l'eau :

\begin{equation}
  \dot{Q}\,[\si{\watt}] = \frac{\dot{V}\,[\si{\liter\per\hour}]}{3600} \times \rho \times c_p \times \Delta T
  \approx \dot{V}\,[\si{\liter\per\hour}] \times 1{,}163 \times \Delta T\,[\si{\kelvin}]
  \label{eq:formule_pratique_eau}
\end{equation}

Le facteur \num{1,163} est le produit \((\rho \cdot c_p) / 3600\) pour l'eau
\(= 1000 \times 4186 / 3\,600\,000 \approx \num{1,163}\,\si{\watt\hour\per\liter\per\kelvin}\).

% ============================================================
\section{Exemples résolus}
% ============================================================

\begin{exemple}[title=Calcul de perte thermique par une paroi multicouche]

\textbf{Énoncé :} Un mur de façade d'un immeuble de bureau est composé (de l'intérieur vers
l'extérieur) des couches suivantes :
\begin{enumerate}
  \item Enduit plâtre intérieur : \(e_1 = \SI{1,5}{\centi\meter}\), \(\lambda_1 = \SI{0,35}{\watt\per\meter\per\kelvin}\)
  \item Béton armé : \(e_2 = \SI{20}{\centi\meter}\), \(\lambda_2 = \SI{1,75}{\watt\per\meter\per\kelvin}\)
  \item Laine de verre (panneaux semi-rigides) : \(e_3 = \SI{10}{\centi\meter}\), \(\lambda_3 = \SI{0,04}{\watt\per\meter\per\kelvin}\)
  \item Enduit ciment extérieur : \(e_4 = \SI{2}{\centi\meter}\), \(\lambda_4 = \SI{1,15}{\watt\per\meter\per\kelvin}\)
\end{enumerate}

\textit{Données complémentaires :} \(R_{si} = \SI{0,13}{\meter\squared\kelvin\per\watt}\),
\(R_{se} = \SI{0,04}{\meter\squared\kelvin\per\watt}\), température intérieure
\(\theta_i = \SI{20}{\degreeCelsius}\), température extérieure
\(\theta_e = \SI{0}{\degreeCelsius}\), surface du mur \(S = \SI{50}{\meter\squared}\).

\textbf{Résolution :}

\textit{Étape 1 — Résistance de chaque couche}

\[
  R_1 = \frac{e_1}{\lambda_1} = \frac{0{,}015}{0{,}35} = \SI{0,043}{\meter\squared\kelvin\per\watt}
\]
\[
  R_2 = \frac{e_2}{\lambda_2} = \frac{0{,}20}{1{,}75} = \SI{0,114}{\meter\squared\kelvin\per\watt}
\]
\[
  R_3 = \frac{e_3}{\lambda_3} = \frac{0{,}10}{0{,}04} = \SI{2,500}{\meter\squared\kelvin\per\watt}
\]
\[
  R_4 = \frac{e_4}{\lambda_4} = \frac{0{,}02}{1{,}15} = \SI{0,017}{\meter\squared\kelvin\per\watt}
\]

\textit{Étape 2 — Résistance thermique totale (selon NF EN ISO 6946)}

\[
  R_{tot} = R_{si} + R_1 + R_2 + R_3 + R_4 + R_{se}
\]
\[
  R_{tot} = 0{,}13 + 0{,}043 + 0{,}114 + 2{,}500 + 0{,}017 + 0{,}04
\]
\[
  \boxed{R_{tot} = \SI{2,844}{\meter\squared\kelvin\per\watt}}
\]

\textit{Étape 3 — Coefficient de transmission thermique}

\[
  U = \frac{1}{R_{tot}} = \frac{1}{2{,}844} = \SI{0,352}{\watt\per\meter\squared\per\kelvin}
\]

\textit{Étape 4 — Flux thermique total pour \(\Delta T = \SI{20}{\kelvin}\)}

\[
  \dot{Q} = U \cdot S \cdot \Delta T = 0{,}352 \times 50 \times 20
\]
\[
  \boxed{\dot{Q} = \SI{352}{\watt}}
\]

\textbf{Conclusion :} Ce mur perd \SI{352}{\watt} pour une différence de température de
\SI{20}{\kelvin}. Avec \(U = \SI{0,352}{\watt\per\meter\squared\per\kelvin}\), ce mur
satisfait l'exigence RE2020 pour les murs donnant sur l'extérieur en zone H1
(\(U_{max} = \SI{0,40}{\watt\per\meter\squared\per\kelvin}\) pour logements neufs).

\end{exemple}

\begin{exemple}[title=Calcul de puissance d'un radiateur eau chaude]

\textbf{Énoncé :} Un radiateur à eau chaude est alimenté avec un débit volumique de
\SI{200}{\liter\per\hour}. La température d'entrée est \(\theta_{entrée} = \SI{70}{\degreeCelsius}\)
et la température de sortie est \(\theta_{sortie} = \SI{55}{\degreeCelsius}\). Calculer la
puissance thermique émise par ce radiateur.

\textbf{Données :}
\begin{itemize}
  \item \(\dot{V} = \SI{200}{\liter\per\hour}\)
  \item \(\Delta T = \theta_{entrée} - \theta_{sortie} = 70 - 55 = \SI{15}{\kelvin}\)
  \item \(\rho_{eau} = \SI{980}{\kilogram\per\meter\cubed}\) (eau à environ \SI{62}{\degreeCelsius})
  \item \(c_p = \SI{4186}{\joule\per\kilogram\per\kelvin}\)
\end{itemize}

\textbf{Résolution — méthode débit massique :}

\textit{Étape 1 — Conversion du débit volumique en débit massique}

\[
  \dot{m} = \rho \cdot \dot{V} = 980 \times \frac{200}{3600} = 980 \times 0{,}0556
  = \SI{54,4}{\kilogram\per\second}
\]

\textit{Étape 2 — Puissance thermique}

\[
  \dot{Q} = \dot{m} \cdot c_p \cdot \Delta T = 54{,}4 \times 4186 \times 15
\]

\begin{attention}
  Erreur courante : oublier la division par 3600 pour convertir le débit de \si{\liter\per\hour}
  en \si{\liter\per\second} avant multiplication. Toujours vérifier la cohérence des unités.
\end{attention}

Correction : \(\dot{m} = 980 \times 200/3600 = \SI{54{,}4}{\kilogram\per\hour}\)
\(\Rightarrow \dot{m} = 54{,}4/3600 = \SI{0{,}0151}{\kilogram\per\second}\)

\[
  \dot{Q} = 0{,}0151 \times 4186 \times 15 = \SI{948}{\watt} \approx \SI{0{,}95}{\kilo\watt}
\]

\textbf{Résolution — méthode pratique (facteur 1,163) :}

\[
  \dot{Q}\,[\si{\watt}] = \dot{V}\,[\si{\liter\per\hour}] \times 1{,}163 \times \Delta T\,[\si{\kelvin}]
\]
\[
  \dot{Q} = 200 \times 1{,}163 \times 15 = \SI{3489}{\watt} \approx \SI{3{,}5}{\kilo\watt}
\]

\textbf{Analyse :} Le facteur \num{1,163} utilise \(\rho = \SI{1000}{\kilogram\per\meter\cubed}\).
Avec la densité réelle de l'eau chaude (\(\rho = \SI{980}{\kilogram\per\meter\cubed}\)), on obtient :

\[
  \dot{Q} = 200 \times \frac{980}{1000} \times 1{,}163 \times 15
  = 200 \times 0{,}980 \times 1{,}163 \times 15 = \SI{3419}{\watt} \approx \SI{3{,}4}{\kilo\watt}
\]

\textbf{Résultat final :} \(\boxed{\dot{Q} \approx \SI{3{,}4}{\kilo\watt}}\)

\textit{Note pratique :} Un débit de \SI{200}{\liter\per\hour} avec un écart de température
de \SI{15}{\kelvin} correspond à une puissance émise de l'ordre de \SI{3,4}{\kilo\watt}, ce
qui est cohérent avec un radiateur de taille standard (\SI{1,5}{\meter} de long, double panneau)
pour une pièce de \num{25} à \SI{30}{\meter\squared}.

\end{exemple}

% ============================================================
\section{Exercices d'application}
% ============================================================

\exercice{Paroi multicouche d'un immeuble collectif}

\textbf{Énoncé :}

La façade d'un immeuble collectif BBC (Bâtiment Basse Consommation) est constituée,
de l'intérieur vers l'extérieur, des couches suivantes :
\begin{itemize}
  \item Plaque de plâtre (BA13) : \(e = \SI{13}{\milli\meter}\), \(\lambda = \SI{0,25}{\watt\per\meter\per\kelvin}\)
  \item Lame d'air ventilée : \textit{résistance fournie} \(R_{air} = \SI{0,18}{\meter\squared\kelvin\per\watt}\)
  \item Laine de roche (panneau rigide) : \(e = \SI{14}{\centi\meter}\), \(\lambda = \SI{0,038}{\watt\per\meter\per\kelvin}\)
  \item Béton cellulaire (bloc YTONG) : \(e = \SI{20}{\centi\meter}\), \(\lambda = \SI{0,11}{\watt\per\meter\per\kelvin}\)
  \item Enduit de façade : \(e = \SI{1,5}{\centi\meter}\), \(\lambda = \SI{0,87}{\watt\per\meter\per\kelvin}\)
\end{itemize}

Résistances superficielles standard : \(R_{si} = \SI{0,13}{\meter\squared\kelvin\per\watt}\),
\(R_{se} = \SI{0,13}{\meter\squared\kelvin\per\watt}\) (paroi verticale, flux horizontal).

\textbf{Questions :}
\begin{enumerate}
  \item Calculer la résistance thermique de chaque couche.
  \item Calculer la résistance thermique totale \(R_{tot}\).
  \item Calculer le coefficient \(U\) de la paroi.
  \item Vérifier la conformité à la RT2012 (\(U_{max} = \SI{0,36}{\watt\per\meter\squared\per\kelvin}\) pour mur en zone H1b).
  \item Calculer la perte thermique pour une surface de \SI{80}{\meter\squared} et une
        différence de température de \SI{25}{\kelvin}\,.
\end{enumerate}

\textbf{Corrigé :}

\textit{1. Résistances des couches}

\[
  R_{pl} = \frac{0{,}013}{0{,}25} = \SI{0,052}{\meter\squared\kelvin\per\watt}
  \qquad
  R_{lr} = \frac{0{,}14}{0{,}038} = \SI{3,684}{\meter\squared\kelvin\per\watt}
\]
\[
  R_{bc} = \frac{0{,}20}{0{,}11} = \SI{1,818}{\meter\squared\kelvin\per\watt}
  \qquad
  R_{end} = \frac{0{,}015}{0{,}87} = \SI{0,017}{\meter\squared\kelvin\per\watt}
\]
\[
  R_{air} = \SI{0,18}{\meter\squared\kelvin\per\watt} \quad \text{(donnée)}
\]

\textit{2. Résistance totale}

\[
  R_{tot} = R_{si} + R_{pl} + R_{air} + R_{lr} + R_{bc} + R_{end} + R_{se}
\]
\[
  R_{tot} = 0{,}13 + 0{,}052 + 0{,}18 + 3{,}684 + 1{,}818 + 0{,}017 + 0{,}13
\]
\[
  \boxed{R_{tot} = \SI{6,011}{\meter\squared\kelvin\per\watt}}
\]

\textit{3. Coefficient U}

\[
  \boxed{U = \frac{1}{6{,}011} = \SI{0,166}{\watt\per\meter\squared\per\kelvin}}
\]

\textit{4. Conformité RT2012}

\(U = \SI{0,166}{\watt\per\meter\squared\per\kelvin} < U_{max} = \SI{0,36}{\watt\per\meter\squared\per\kelvin}\)
\(\Rightarrow\) La paroi est \textbf{conforme} à la RT2012 (et a fortiori à la RE2020 qui est
encore plus exigeante).

\textit{5. Perte thermique}

\[
  \dot{Q} = U \cdot S \cdot \Delta T = 0{,}166 \times 80 \times 25 = \SI{332}{\watt}
\]

Pour une façade de \SI{80}{\meter\squared} avec \(\Delta T = \SI{25}{\kelvin}\), la perte
thermique est seulement \SI{332}{\watt}, ce qui illustre l'excellente performance de cette
composition de paroi.

% ----------------------------------------------------------

\exercice{Dimensionnement d'un plancher chauffant hydraulique}

\textbf{Énoncé :}

Un salon de \SI{30}{\meter\squared} doit être chauffé par un plancher chauffant hydraulique.
La puissance de chauffe nécessaire est estimée à \SI{60}{\watt\per\meter\squared} (déperditions
calculées selon la méthode RE2020). Le plancher chauffant fonctionne avec des températures
d'eau de \SI{35}{\degreeCelsius} (aller) et \SI{28}{\degreeCelsius} (retour).

\textbf{Questions :}
\begin{enumerate}
  \item Calculer la puissance thermique totale nécessaire pour ce salon.
  \item Calculer le débit volumique d'eau chaude nécessaire [\si{\liter\per\hour}] en
        utilisant le facteur pratique \num{1,163}.
  \item Vérifier que la température de surface du plancher reste inférieure à
        \SI{28}{\degreeCelsius} (confort et réglementation). Commentaire.
\end{enumerate}

\textbf{Corrigé :}

\textit{1. Puissance totale nécessaire}

\[
  \dot{Q}_{tot} = \varphi \cdot S = 60 \times 30 = \SI{1800}{\watt} = \SI{1{,}8}{\kilo\watt}
\]

\textit{2. Débit volumique d'eau}

On utilise la formule pratique : \(\dot{Q} = \dot{V} \times 1{,}163 \times \Delta T\)

\[
  \Delta T = 35 - 28 = \SI{7}{\kelvin}
\]
\[
  \dot{V} = \frac{\dot{Q}}{1{,}163 \times \Delta T} = \frac{1800}{1{,}163 \times 7}
  = \frac{1800}{8{,}141} = \SI{221}{\liter\per\hour}
\]

\[
  \boxed{\dot{V} \approx \SI{221}{\liter\per\hour}}
\]

\textit{3. Température de surface et confort}

La réglementation impose que la température de surface d'un plancher chauffant ne dépasse pas
\SI{28}{\degreeCelsius} (norme NF EN 1264-2). La température de surface est liée à la
température moyenne de l'eau : \(T_{moy} = (35+28)/2 = \SI{31{,}5}{\degreeCelsius}\). Avec
les résistances du plancher (chape, revêtement), la surface sera bien inférieure à
\SI{28}{\degreeCelsius}. Ce dimensionnement est conforme.

% ----------------------------------------------------------

\exercice{Bilan thermique d'une pièce de bureau}

\textbf{Énoncé :}

Un bureau de \SI{20}{\meter\squared} est situé dans un immeuble tertiaire. Les déperditions
sont les suivantes :
\begin{itemize}
  \item Par le mur extérieur (\(S = \SI{12}{\meter\squared}\), \(U = \SI{0,30}{\watt\per\meter\squared\per\kelvin}\)) ;
  \item Par la fenêtre double vitrage (\(S = \SI{4}{\meter\squared}\), \(U = \SI{1,40}{\watt\per\meter\squared\per\kelvin}\)) ;
  \item Par le plancher haut non isolé sur combles (\(S = \SI{20}{\meter\squared}\), \(U = \SI{0,25}{\watt\per\meter\squared\per\kelvin}\)) ;
  \item Par ponts thermiques linéiques : \(\Psi \cdot L = \SI{0,45}{\watt\per\kelvin}\) (valeur globale).
\end{itemize}

Les apports internes sont : éclairage \SI{200}{\watt}, occupants (2 personnes × \SI{80}{\watt}),
équipements informatiques \SI{150}{\watt}.

Conditions de calcul : \(\theta_i = \SI{20}{\degreeCelsius}\), \(\theta_e = \SI{-7}{\degreeCelsius}\)
(température extérieure de base en zone H1, Paris, selon EN 12831).

\textbf{Questions :}
\begin{enumerate}
  \item Calculer le coefficient de déperdition total de la pièce (\(H_T\) en \si{\watt\per\kelvin}).
  \item Calculer les déperditions thermiques totales.
  \item Calculer les apports internes totaux.
  \item Déterminer la puissance de chauffage à installer.
\end{enumerate}

\textbf{Corrigé :}

\textit{1. Coefficient de déperdition \(H_T\)}

\[
  H_T = \sum (U_i \cdot S_i) + \Psi \cdot L
\]
\[
  H_T = (0{,}30 \times 12) + (1{,}40 \times 4) + (0{,}25 \times 20) + 0{,}45
\]
\[
  H_T = 3{,}60 + 5{,}60 + 5{,}00 + 0{,}45 = \SI{14{,}65}{\watt\per\kelvin}
\]

\textit{2. Déperditions thermiques}

\[
  \Delta T = \theta_i - \theta_e = 20 - (-7) = \SI{27}{\kelvin}
\]
\[
  \dot{Q}_{dep} = H_T \cdot \Delta T = 14{,}65 \times 27 = \SI{395{,}6}{\watt}
\]

\textit{3. Apports internes}

\[
  \dot{Q}_{int} = 200 + (2 \times 80) + 150 = 200 + 160 + 150 = \SI{510}{\watt}
\]

\textit{4. Puissance de chauffage à installer}

En période de chauffage, les apports internes réduisent le besoin en chauffage :

\[
  \dot{Q}_{chauf} = \dot{Q}_{dep} - \dot{Q}_{int} = 395{,}6 - 510 = \SI{-114{,}4}{\watt}
\]

Le résultat négatif indique que les apports internes dépassent les déperditions dans ces
conditions. \textbf{Aucun chauffage n'est nécessaire} dans cette configuration pour \(\theta_e
= \SI{-7}{\degreeCelsius}\). Cette situation est caractéristique des immeubles tertiaires à
forte densité d'occupation et d'équipements.

\begin{attention}
  En pratique, il faut également prévoir le chauffage pour les périodes sans occupation
  (nuit, week-end) avec apports internes nuls. La puissance à installer devient alors :
  \(\dot{Q}_{chauf} = \SI{395{,}6}{\watt}\), arrondie à \SI{400}{\watt} après marges de sécurité.
\end{attention}

% ============================================================
\section{Éléments de réglementation}
% ============================================================

\subsection{RE2020 — Réglementation Environnementale 2020}

Entrée en vigueur le 1\textsuperscript{er} janvier 2022 pour les maisons individuelles et les
logements collectifs neufs, et progressivement étendue aux bâtiments tertiaires, la RE2020
remplace la RT2012. Ses trois piliers sont :
\begin{enumerate}
  \item \textbf{Efficacité énergétique} : réduction de la consommation d'énergie primaire
        (indicateur \(Ep_{max}\)).
  \item \textbf{Confort d'été} : limitation de l'inconfort en période chaude (indicateur DH,
        degrés-heures d'inconfort).
  \item \textbf{Impact carbone} : prise en compte de l'analyse du cycle de vie (ACV) du
        bâtiment, indicateur \(IC_{bât}\).
\end{enumerate}

\textbf{Exigences sur les parois opaques (résidentiels, zone H1) :}

\begin{center}
\begin{tabular}{lcc}
  \toprule
  Type de paroi & \(U_{max}\) RE2020 & \(U_{max}\) RT2012 (référence) \\
                & [\si{\watt\per\meter\squared\per\kelvin}] & [\si{\watt\per\meter\squared\per\kelvin}] \\
  \midrule
  Murs donnant sur l'extérieur & 0,36 & 0,40 \\
  Plancher bas sur vide sanitaire & 0,27 & 0,36 \\
  Toiture-terrasse & 0,20 & 0,25 \\
  Comble aménagé & 0,20 & 0,22 \\
  \bottomrule
\end{tabular}
\end{center}

\subsection{DTU 45.4 — Isolation thermique des bâtiments}

Le DTU 45.4 définit les règles de mise en œuvre des isolants thermiques dans les bâtiments.
Il précise notamment :
\begin{itemize}
  \item les conditions de mise en œuvre (fixation, joints, continuité de l'isolant) ;
  \item la protection contre l'humidité (pare-vapeur côté chaud en hiver) ;
  \item les incompatibilités entre matériaux ;
  \item les règles de sécurité incendie pour les isolants.
\end{itemize}

\subsection{NF EN ISO 6946 — Résistance thermique des parois}

Cette norme internationale définit la méthode de calcul de la résistance thermique et du
coefficient de transmission thermique des parois de bâtiment composées de couches
thermiquement homogènes. Elle fixe notamment :
\begin{itemize}
  \item les valeurs des résistances superficielles \(R_{si}\) et \(R_{se}\) selon la direction
        du flux et la position de la paroi ;
  \item la méthode de calcul pour les couches non homogènes (ponts thermiques intégrés) ;
  \item les corrections à apporter pour les fixations métalliques et les lames d'air.
\end{itemize}

\textbf{Résistances superficielles selon NF EN ISO 6946 :}

\begin{center}
\begin{tabular}{lccc}
  \toprule
  Direction du flux & Position & \(R_{si}\) [\si{\meter\squared\kelvin\per\watt}]
                    & \(R_{se}\) [\si{\meter\squared\kelvin\per\watt}] \\
  \midrule
  Ascendant (paroi horizontale chauffée par le bas) & Plancher chauffant & 0,10 & 0,13 \\
  Horizontal (paroi verticale) & Mur & 0,13 & 0,04 \\
  Descendant (paroi horizontale chauffée par le haut) & Toiture & 0,17 & 0,04 \\
  \bottomrule
\end{tabular}
\end{center}

% \chapter{Hydraulique des circuits}

% ============================================================
\section{Objectifs pédagogiques}
% ============================================================

\subsection*{Compétences GCF mobilisées}
\begin{itemize}
  \item \textbf{C2-1 à C2-3} — Analyser un système hydraulique : identifier composants, chaînes d'énergie, paramètres de fonctionnement
  \item \textbf{C3-3} — Dimensionner tout ou partie d'un réseau (méthodes normées, notes de calcul justifiées)
  \item \textbf{C7-4 à C7-6} — Conduire des essais de mise en service, mesurer débits et pressions, interpréter les écarts
  \item \textbf{C8-1 à C8-4} — Vérifier les performances, identifier les dérives, concevoir des actions correctives
\end{itemize}

\subsection*{Savoirs associés mobilisés}
\begin{itemize}
  \item \textbf{A3} — Dynamique des fluides : pertes de charge, point de fonctionnement, équilibrage (niveau 3)
  \item \textbf{B4-1} — Réseaux hydrauliques : eau, eau glacée, équilibrage, sécurité, dimensionnement (niveau 3)
  \item \textbf{S5.5} — Réglementations fluides thermiques : chauffage, EFS (niveau 3)
  \item \textbf{S10.3} — Sécurités d'installation, remplissage (niveau 3)
  \item \textbf{S10.5} — Mise en service d'une installation GCF (niveau 3)
  \item \textbf{S10.6} — Critères de bon fonctionnement et optimisation (niveau 3)
\end{itemize}

% ============================================================
\section{Introduction}
% ============================================================

La maîtrise de l'hydraulique est au cœur du métier de technicien supérieur GCF. Un réseau mal dimensionné ou mal équilibré entraîne des déséquilibres de débit, une surconsommation électrique des pompes, un inconfort thermique pour les occupants et une usure prématurée des équipements.

En pratique, un réseau de chauffage ou de climatisation transporte l'énergie thermique depuis la production (chaudière, PAC, groupe froid) jusqu'aux émetteurs (radiateurs, ventilo-convecteurs, planchers chauffants) via un réseau de tuyauteries. Chaque choix — diamètre, vitesse d'eau, type de pompe, organe d'équilibrage — a des conséquences sur l'efficacité énergétique globale de l'installation.

Ce chapitre couvre les bases de la dynamique des fluides appliquées aux réseaux CVC, le calcul des pertes de charge, le dimensionnement des pompes et l'équilibrage des réseaux.

% ============================================================
\section{Rappels de mécanique des fluides}
% ============================================================

\subsection{Grandeurs fondamentales}

\begin{definition}[title=Débit volumique]
Le débit volumique $q_v$ est le volume de fluide traversant une section droite par unité de temps.
\begin{equation}
  q_v = v \cdot S \qquad \left[\si{\cubic\metre\per\second}\right]
\end{equation}
\begin{itemize}
  \item $v$ : vitesse moyenne du fluide (\si{\metre\per\second})
  \item $S$ : section droite de la conduite (\si{\square\metre})
\end{itemize}
En pratique, on utilise souvent le débit en \si{\litre\per\hour} ou \si{\cubic\metre\per\hour}.
\end{definition}

\begin{definition}[title=Débit massique]
\begin{equation}
  \dot{m} = \rho \cdot q_v \qquad \left[\si{\kilogram\per\second}\right]
\end{equation}
\begin{itemize}
  \item $\rho$ : masse volumique du fluide (\si{\kilogram\per\cubic\metre})
\end{itemize}
\end{definition}

\begin{table}[H]
  \centering
  \caption{Masse volumique et viscosité cinématique de l'eau selon la température}
  \begin{tabular}{lcccc}
    \toprule
    $T$ (°C) & $\rho$ (kg/m³) & $c_p$ (J/kg·K) & $\nu \times 10^{6}$ (m²/s) & Facteur $\rho c_p / 3600$ \\
    \midrule
    10 & 999,7 & 4192 & 1,307 & 1,164 \\
    20 & 998,2 & 4182 & 1,004 & 1,163 \\
    40 & 992,2 & 4175 & 0,658 & 1,153 \\
    60 & 983,2 & 4184 & 0,475 & 1,145 \\
    70 & 977,8 & 4190 & 0,415 & 1,139 \\
    80 & 971,8 & 4196 & 0,365 & 1,133 \\
    \bottomrule
  \end{tabular}
\end{table}

\begin{definition}[title=Pression]
La pression $P$ est la force exercée par le fluide par unité de surface. Elle s'exprime en pascal (\si{\pascal}) ou en bar (\SI{1}{bar} = \SI{100000}{\pascal}).

On distingue :
\begin{itemize}
  \item \textbf{Pression absolue} $P_{abs}$ : pression totale par rapport au vide absolu
  \item \textbf{Pression relative} $P_{rel} = P_{abs} - P_{atm}$, référencée à la pression atmosphérique
  \item \textbf{Pression différentielle} $\Delta P$ : différence de pression entre deux points du réseau
\end{itemize}
\end{definition}

\subsection{Régimes d'écoulement — Nombre de Reynolds}

\begin{definition}[title=Nombre de Reynolds]
Le nombre de Reynolds $Re$ caractérise le régime d'écoulement du fluide dans une conduite.
\begin{equation}
  Re = \frac{v \cdot D}{\nu}
\end{equation}
\begin{itemize}
  \item $v$ : vitesse du fluide (\si{\metre\per\second})
  \item $D$ : diamètre intérieur de la conduite (\si{\metre})
  \item $\nu$ : viscosité cinématique (\si{\square\metre\per\second}) — pour l'eau à 50\,°C : $\nu \approx 5{,}5 \times 10^{-7}$\,\si{\square\metre\per\second}
\end{itemize}
\end{definition}

\begin{table}[H]
  \centering
  \caption{Régimes d'écoulement selon le nombre de Reynolds}
  \begin{tabular}{lll}
    \toprule
    Régime & Condition & Caractéristiques \\
    \midrule
    Laminaire & $Re < 2300$ & Filets parallèles, peu de pertes \\
    Transitoire & $2300 < Re < 4000$ & Instable \\
    Turbulent & $Re > 4000$ & Mélange intense, pertes plus élevées \\
    \bottomrule
  \end{tabular}
\end{table}

\begin{attention}
Dans les installations CVC, l'écoulement est presque toujours \textbf{turbulent} ($Re \gg 4000$). Les formules de pertes de charge utilisées dans ce chapitre s'appliquent au régime turbulent.
\end{attention}

\subsection{Équation de Bernoulli généralisée}

L'équation de Bernoulli relie les différentes formes d'énergie du fluide entre deux points 1 et 2 d'un réseau :

\begin{equation}
  \frac{P_1}{\rho g} + \frac{v_1^2}{2g} + z_1 + H_{pompe} = \frac{P_2}{\rho g} + \frac{v_2^2}{2g} + z_2 + \Delta h_{pertes}
  \label{eq:bernoulli}
\end{equation}
\begin{itemize}
  \item $P$ : pression (\si{\pascal})
  \item $v$ : vitesse (\si{\metre\per\second})
  \item $z$ : hauteur géométrique (\si{\metre})
  \item $H_{pompe}$ : hauteur manométrique apportée par la pompe (\si{\metre})
  \item $\Delta h_{pertes}$ : pertes de charge totales (\si{\metre})
  \item $g = \SI{9{,}81}{\metre\per\square\second}$
\end{itemize}

Pour un circuit fermé bouclé (chauffage, climatisation), les termes géométriques et de pression statique se compensent. Seules les pertes de charge et la HMT de la pompe subsistent.

% ============================================================
\section{Pertes de charge}
% ============================================================

Les pertes de charge représentent la chute de pression due aux frottements du fluide dans la conduite et dans les singularités (coudes, vannes, té, etc.). Elles sont exprimées en \si{\pascal} ou en mètre de colonne d'eau (\si{mCE}).

\subsection{Pertes de charge régulières (linéaires)}

Les pertes de charge régulières sont dues aux frottements le long de la conduite. On utilise la formule de \textbf{Darcy-Weisbach} :

\begin{equation}
  \Delta P_r = \lambda \cdot \frac{L}{D} \cdot \frac{\rho v^2}{2} \qquad \left[\si{\pascal}\right]
  \label{eq:darcy}
\end{equation}
\begin{itemize}
  \item $\lambda$ : coefficient de frottement (sans dimension) — dépend de $Re$ et de la rugosité $\varepsilon/D$
  \item $L$ : longueur de la conduite (\si{\metre})
  \item $D$ : diamètre intérieur (\si{\metre})
  \item $\rho$ : masse volumique (\si{\kilogram\per\cubic\metre})
  \item $v$ : vitesse du fluide (\si{\metre\per\second})
\end{itemize}

\begin{definition}[title=Coefficient de frottement $\lambda$ — diagramme de Moody]
En régime turbulent rugueux, $\lambda$ dépend du nombre de Reynolds $Re$ et de la rugosité relative $\varepsilon/D$. Pour les tubes acier courants ($\varepsilon \approx 0{,}046$\,mm) en CVC :
\begin{itemize}
  \item DN 20 à DN 50 : $\lambda \approx 0{,}022$ à $0{,}028$
  \item DN 80 à DN 200 : $\lambda \approx 0{,}018$ à $0{,}022$
\end{itemize}
Pour les tubes en cuivre ou PER (très lisses, $\varepsilon \approx 0{,}001$\,mm) : $\lambda$ légèrement plus faible.
\end{definition}

\begin{methode}[title=Utilisation des abaques de pertes de charge]
En pratique, on utilise des abaques ou des logiciels de calcul qui donnent directement la perte de charge linéique $R$ en \si{\pascal\per\metre} en fonction du débit et du diamètre.

La perte de charge régulière est alors :
\begin{equation}
  \Delta P_r = R \cdot L \qquad \left[\si{\pascal}\right]
\end{equation}

\textbf{Règle de dimensionnement GCF :}
\begin{itemize}
  \item Vitesse recommandée : entre \SI{0{,}5}{} et \SI{1{,}5}{\metre\per\second} (pour limiter bruit et érosion)
  \item Perte de charge linéique cible : entre \SI{100}{} et \SI{300}{\pascal\per\metre}
  \item Pour les colonnes montantes d'immeuble : \SI{100}{\pascal\per\metre} maximum
  \item Pour les réseaux eau glacée (tertiaire) : \SI{200}{} à \SI{400}{\pascal\per\metre}
\end{itemize}
\end{methode}

\subsection{Pertes de charge singulières}

Les pertes de charge singulières sont dues aux perturbations locales de l'écoulement (coudes, vannes, réducteurs, tés, etc.).

\begin{equation}
  \Delta P_s = \zeta \cdot \frac{\rho v^2}{2} \qquad \left[\si{\pascal}\right]
  \label{eq:singuliere}
\end{equation}
\begin{itemize}
  \item $\zeta$ : coefficient de perte de charge singulière (adimensionnel, fourni par les fabricants)
\end{itemize}

\begin{table}[H]
  \centering
  \caption{Valeurs courantes du coefficient $\zeta$ pour les accessoires hydrauliques}
  \begin{tabular}{lll}
    \toprule
    Accessoire & Conditions & $\zeta$ typique \\
    \midrule
    Coude 90° & $R/D = 1$ & 1{,}5 \\
    Coude 90° & $R/D = 2$ & 0{,}9 \\
    Coude 45° & $R/D = 1$ & 0{,}5 \\
    Té passage direct & — & 0{,}5 \\
    Té passage en dérivation & — & 1{,}5 \\
    Vanne à soupape ouverte & — & 10 \\
    Robinet d'arrêt à boisseau sphérique & ouvert & 0{,}1 \\
    Clapet anti-retour & — & 3 à 5 \\
    Filtre en Y & propre & 2 \\
    Réducteur concentrique & — & 0{,}1 à 0{,}3 \\
    Entrée brusque de tuyauterie & — & 0{,}5 \\
    \bottomrule
  \end{tabular}
\end{table}

\subsection{Méthode de la longueur équivalente}

Une méthode simplifiée consiste à remplacer chaque singularité par une longueur équivalente de tuyau droit $L_{eq}$ qui provoquerait la même perte de charge :

\begin{equation}
  \Delta P_{total} = R \cdot (L_{réel} + \sum L_{eq})
\end{equation}

En pratique, on majore souvent la longueur développée d'un forfait de 30 à 50\,\% pour tenir compte des singularités dans les réseaux courants.

\begin{table}[H]
  \centering
  \caption{Longueurs équivalentes typiques (tube acier, eau à 70°C) — extrait abaque}
  \begin{tabular}{lccc}
    \toprule
    Accessoire & DN 20 & DN 32 & DN 50 \\
    \midrule
    Coude 90° ($R/D=1$) & 0,9 m & 1,4 m & 2,2 m \\
    Té dérivation & 0,8 m & 1,3 m & 2,1 m \\
    Robinet d'arrêt & 5,5 m & 9,0 m & 14,0 m \\
    Clapet anti-retour & 1,9 m & 3,0 m & 4,8 m \\
    \bottomrule
  \end{tabular}
\end{table}

% ============================================================
\section{Les pompes de circulation}
% ============================================================

\subsection{Caractéristiques d'une pompe}

\begin{definition}[title=Hauteur manométrique totale (HMT)]
La HMT est la pression que la pompe doit vaincre pour mettre le fluide en mouvement à travers le réseau. Elle s'exprime en mètre de colonne d'eau ou en pascal.
\begin{equation}
  HMT = \Delta h_{pertes\,totales} \qquad \left[\si{mCE}\right]
\end{equation}
Pour un circuit bouclé (chauffage, climatisation), les différences de hauteur géométrique se compensent. Seules les pertes de charge restent.
\end{definition}

\begin{definition}[title=Puissance hydraulique]
\begin{equation}
  P_{hyd} = \rho \cdot g \cdot q_v \cdot HMT \qquad \left[\si{\watt}\right]
\end{equation}
\end{definition}

\begin{definition}[title=Puissance absorbée et rendement]
\begin{equation}
  P_{abs} = \frac{P_{hyd}}{\eta_{pompe}} \qquad \left[\si{\watt}\right]
\end{equation}
Le rendement global d'une pompe en circulateur CVC est typiquement compris entre 30 et 70\,\% selon la qualité de la machine et son point de fonctionnement.
\end{definition}

\begin{definition}[title=Indice d'efficacité énergétique (IEE)]
Le règlement européen 547/2012 impose un $IEE \leq 0{,}23$ pour les circulateurs de chauffage mis sur le marché depuis 2015. Les circulateurs à vitesse variable (moteur EC) atteignent des $IEE \approx 0{,}15$ à $0{,}20$.
\end{definition}

\subsection{Courbe caractéristique de pompe}

La courbe caractéristique d'une pompe donne la HMT en fonction du débit $q_v$. Elle est fournie par le fabricant. La courbe du réseau représente la perte de charge totale en fonction du débit :

\begin{equation}
  \Delta P_{réseau} = k \cdot q_v^2
\end{equation}

\begin{definition}[title=Point de fonctionnement]
Le point de fonctionnement est l'intersection entre la courbe caractéristique de la pompe et la courbe du réseau. C'est le seul point d'équilibre stable du système.
\end{definition}

\begin{attention}
Une pompe trop puissante fonctionne à débit excessif, loin de son rendement optimal. Une pompe sous-dimensionnée ne peut pas vaincre les pertes de charge de la branche défavorisée. Le choix doit être précis, avec une légère marge de sécurité (10\,\% sur HMT et débit).
\end{attention}

\subsection{Régulation de la pompe}

\begin{table}[H]
  \centering
  \caption{Modes de régulation des circulateurs à vitesse variable}
  \begin{tabular}{p{3.5cm}p{4cm}p{4.5cm}}
    \toprule
    Mode & Principe & Usage recommandé \\
    \midrule
    Pression différentielle constante & La HMT est maintenue fixe quelle que soit la demande & Petits réseaux, radiateurs sans robinets thermostatiques \\
    Pression différentielle proportionnelle & La HMT varie proportionnellement au débit ($\Delta P \propto q_v$) & Réseaux avec robinets thermostatiques — économies jusqu'à 50\,\% \\
    Température constante & Le débit est ajusté pour maintenir $\Delta T$ & Grands réseaux, variation de charge importante \\
    \bottomrule
  \end{tabular}
\end{table}

\subsection{Association de pompes}

\begin{table}[H]
  \centering
  \caption{Association de pompes identiques}
  \begin{tabular}{lll}
    \toprule
    Montage & Effet sur le débit & Effet sur la HMT \\
    \midrule
    Série & Identique & Doublée \\
    Parallèle & Doublé & Identique \\
    \bottomrule
  \end{tabular}
\end{table}

\begin{attention}
Le doublement du débit ou de la HMT n'est valable que pour des pompes identiques sur un réseau idéal. En pratique, l'augmentation réelle est moindre à cause de la courbe de réseau. Le montage en parallèle est utilisé pour la redondance (secours) et les périodes de pointe.
\end{attention}

% ============================================================
\section{Équilibrage des réseaux hydrauliques}
% ============================================================

\subsection{Problématique du déséquilibre}

Dans un réseau ramifié, certains émetteurs reçoivent plus d'eau que prévu (les plus proches de la pompe) tandis que d'autres sont sous-alimentés (les plus éloignés). Ce déséquilibre provoque :
\begin{itemize}
  \item Une surchauffe des pièces proches de la production
  \item Un inconfort dans les pièces éloignées
  \item Une consommation électrique excessive (la pompe lutte contre un réseau mal équilibré)
\end{itemize}

\subsection{Organes d'équilibrage}

\begin{table}[H]
  \centering
  \caption{Organes d'équilibrage — caractéristiques et usages}
  \begin{tabular}{p{3.5cm}p{5cm}p{4cm}}
    \toprule
    Organe & Principe & Application \\
    \midrule
    Robinet d'équilibrage manuel & Crée une perte de charge réglable. Nécessite une mesure de débit. & Tout réseau statique \\
    Régulateur de débit automatique (PIC) & Maintient un débit constant indépendamment de la pression différentielle & Réseau à débit variable (robinets thermostatiques) \\
    Vanne de régulation pression différentielle & Maintient une $\Delta P$ constante sur un circuit & Réseaux de distribution \\
    Vanne motorisée 2 voies & Ouverture proportionnelle ou TOR commandée par régulateur & Émetteur avec régulation locale \\
    \bottomrule
  \end{tabular}
\end{table}

\begin{methode}[title=Équilibrage par la méthode de la résistance équivalente (branche index)]
\begin{enumerate}
  \item Calculer les pertes de charge de chaque branche du réseau (du départ au retour)
  \item Identifier la branche la plus défavorisée (pertes de charge les plus élevées) — c'est la branche \textbf{index}
  \item Pour chaque autre branche, calculer l'excédent de pression à absorber :
    \[
      \Delta P_{excédent} = \Delta P_{index} - \Delta P_{branche}
    \]
  \item Régler les organes d'équilibrage pour créer une perte de charge supplémentaire égale à $\Delta P_{excédent}$ sur chaque branche
  \item Vérifier les débits après réglage (débitmètre, mesure différentielle)
\end{enumerate}
\end{methode}

\subsection{Sécurités hydrauliques}

\begin{table}[H]
  \centering
  \caption{Dispositifs de sécurité des installations hydrauliques (DTU 65.11)}
  \begin{tabular}{p{3.5cm}p{4.5cm}p{4.5cm}}
    \toprule
    Dispositif & Rôle & Réglage typique \\
    \midrule
    Vase d'expansion & Absorbe la dilatation de l'eau en chauffe (eau : $+4\,\%$ de 10 à 90°C) & Pression de gonflage : $P_{stat} + 0{,}3$\,bar \\
    Soupape de sécurité & Limite la pression max en cas de surpression & 3 bar (chaudière murale) \\
    Pressostat de sécurité & Coupe le brûleur si pression anormale & Bas : 0,5 bar / Haut : 3 bar \\
    Clapet anti-retour & Empêche la circulation inverse dans les circuits en parallèle & — \\
    Purgeur automatique & Élimine l'air dissous pouvant créer des bouchons & Point haut de chaque circuit \\
    \bottomrule
  \end{tabular}
\end{table}

% ============================================================
\section{Exemples résolus}
% ============================================================

\begin{exemple}[title=Calcul des pertes de charge d'un tronçon]
\textbf{Énoncé :} Un tronçon de réseau de chauffage est constitué d'un tube acier DN 25 (diamètre intérieur \SI{27{,}2}{\milli\metre}) de longueur \SI{15}{\metre}. Il débite \SI{800}{\litre\per\hour} d'eau à 70\,°C. Ce tronçon comporte 4 coudes à 90° ($\zeta = 1{,}5$ chacun) et 1 robinet d'arrêt à boisseau ($\zeta = 0{,}1$). Masse volumique de l'eau à 70\,°C : $\rho = \SI{978}{\kilogram\per\cubic\metre}$.

À partir de l'abaque du fabricant, la perte de charge linéique pour ce diamètre et ce débit est $R = \SI{250}{\pascal\per\metre}$.

\textbf{Étape 1 : Perte de charge régulière}
\[
  \Delta P_r = R \times L = 250 \times 15 = \SI{3750}{\pascal}
\]

\textbf{Étape 2 : Vitesse du fluide}
\[
  S = \frac{\pi \times (0{,}0272)^2}{4} = \SI{5{,}81e-4}{\square\metre}
\]
\[
  q_v = \frac{800}{3600 \times 1000} = \SI{2{,}22e-4}{\cubic\metre\per\second}
\]
\[
  v = \frac{2{,}22 \times 10^{-4}}{5{,}81 \times 10^{-4}} = \SI{0{,}38}{\metre\per\second}
\]

\textbf{Étape 3 : Pression dynamique}
\[
  \frac{\rho v^2}{2} = \frac{978 \times (0{,}38)^2}{2} = \SI{70{,}6}{\pascal}
\]

\textbf{Étape 4 : Pertes de charge singulières}
\[
  \Delta P_s = \left(4 \times 1{,}5 + 1 \times 0{,}1\right) \times 70{,}6 = 6{,}1 \times 70{,}6 = \SI{430}{\pascal}
\]

\textbf{Étape 5 : Perte de charge totale}
\[
  \Delta P_{total} = 3750 + 430 = \SI{4180}{\pascal} \approx \SI{0{,}042}{\bar}
\]

\textbf{Vérification de la vitesse :} $v = 0{,}38$\,m/s — légèrement en dessous de la plage 0,5–1,5\,m/s, acceptable pour une colonne montante (critère acoustique prioritaire).
\end{exemple}

\begin{exemple}[title=Dimensionnement d'une pompe de chauffage]
\textbf{Énoncé :} Un réseau de chauffage délivre une puissance de \SI{45}{\kilo\watt} avec une température de départ à 70\,°C et retour à 55\,°C. La perte de charge totale de la boucle la plus défavorisée est \SI{18000}{\pascal}. Déterminer le débit nécessaire et choisir une pompe.

\textbf{Étape 1 : Calcul du débit}
\[
  \dot{m} = \frac{\dot{Q}}{c_p \cdot \Delta T} = \frac{45\,000}{4186 \times (70 - 55)} = \SI{0{,}717}{\kilogram\per\second}
\]
\[
  q_v = \frac{0{,}717}{978} = \SI{7{,}33e-4}{\cubic\metre\per\second} = \SI{2{,}64}{\cubic\metre\per\hour}
\]

\textbf{Étape 2 : HMT requise}
\[
  HMT = \frac{\Delta P_{réseau}}{\rho \cdot g} = \frac{18\,000}{978 \times 9{,}81} = \SI{1{,}88}{\mCE} \approx \SI{1{,}9}{\mCE}
\]

\textbf{Étape 3 : Sélection pompe (avec marge 10\,\%)}
\begin{itemize}
  \item Débit de sélection : $2{,}64 \times 1{,}10 \approx \SI{2{,}9}{\cubic\metre\per\hour}$
  \item HMT de sélection : $1{,}9 \times 1{,}10 \approx \SI{2{,}1}{\mCE}$
\end{itemize}
Rechercher dans les catalogues fabricants (Grundfos, Wilo, Armstrong) un circulateur délivrant au minimum $q_v = \SI{2{,}9}{\cubic\metre\per\hour}$ à $HMT = \SI{2{,}1}{\mCE}$.

\textbf{Résultat :} Un circulateur de type Grundfos ALPHA2 25-40 ou équivalent convient (plage : 0 à 3\,m³/h, HMT max 4\,mCE), avec régulation de vitesse intégrée pour l'efficacité énergétique ($IEE \leq 0{,}20$).
\end{exemple}

\begin{exemple}[title=Équilibrage d'un réseau bitubes — 3 radiateurs]
\textbf{Énoncé :} Un réseau bitubes alimente 3 radiateurs depuis une nourrice commune. Les pertes de charge calculées de la nourrice jusqu'à chaque radiateur sont :
\begin{itemize}
  \item Radiateur R1 (proche) : $\Delta P_{R1} = \SI{2800}{\pascal}$
  \item Radiateur R2 (intermédiaire) : $\Delta P_{R2} = \SI{5100}{\pascal}$
  \item Radiateur R3 (éloigné, branche index) : $\Delta P_{R3} = \SI{7600}{\pascal}$
\end{itemize}

\textbf{Perte de charge à créer sur chaque robinet d'équilibrage :}
\begin{itemize}
  \item R1 : $\Delta P_{rob,R1} = 7600 - 2800 = \SI{4800}{\pascal}$ → régler le robinet
  \item R2 : $\Delta P_{rob,R2} = 7600 - 5100 = \SI{2500}{\pascal}$ → régler le robinet
  \item R3 : $\Delta P_{rob,R3} = \SI{0}{\pascal}$ → robinet grand ouvert (branche index)
\end{itemize}

\textbf{La pompe doit fournir} au minimum $\SI{7600}{\pascal}$ (pertes de la branche index) plus les pertes du tronçon commun.

\textbf{Vérification :} Après réglage, les 3 radiateurs ont la même résistance hydraulique vue depuis la nourrice → même pression disponible → débits nominaux respectés.
\end{exemple}

\begin{exemple}[title=Dimensionnement d'un vase d'expansion]
\textbf{Énoncé :} Une installation de chauffage contient un volume d'eau total $V_{eau} = \SI{150}{\litre}$ (chaudière + réseau). La pression statique de remplissage est $P_{stat} = \SI{1{,}5}{\bar}$. La pression max de la soupape est $P_{soupape} = \SI{3}{\bar}$. La température passe de 10\,°C (remplissage) à 80\,°C (chauffe). Calculer le volume du vase.

\textbf{Dilatation de l'eau :} $\rho_{10°C} = 999{,}7$\,kg/m³, $\rho_{80°C} = 971{,}8$\,kg/m³.
\[
  \Delta V_{eau} = V_{eau} \times \left(\frac{\rho_{10}}{\rho_{80}} - 1\right) = 150 \times \left(\frac{999{,}7}{971{,}8} - 1\right) = 150 \times 0{,}0287 = \SI{4{,}3}{\litre}
\]

\textbf{Volume utile du vase :}
Le vase à membrane est gonflé à $P_{gonflage} = P_{stat} = \SI{1{,}5}{\bar}$.
\[
  V_{vase} = \Delta V_{eau} \times \frac{P_{soupape} + 1}{P_{soupape} - P_{stat}} = 4{,}3 \times \frac{3 + 1}{3 - 1{,}5} = 4{,}3 \times \frac{4}{1{,}5} = \SI{11{,}5}{\litre}
\]
Choisir un vase de \SI{12}{\litre} (première taille standard supérieure).
\end{exemple}

% ============================================================
\section{Schémas et diagrammes caractéristiques}
% ============================================================

\subsection{Circuit hydraulique bouclé}

\begin{figure}[H]
  \centering
  \begin{tikzpicture}[>=Stealth, font=\small, thick,
      lbl/.style={font=\footnotesize, inner sep=2pt}]

    % ---- CHAUDIÈRE (symbole ISO NF) ----
    \pic[draw=FedBlue] at (0,0) {chaudiere};
    \node[lbl, above=0.65cm] at (0,0) {Chaudière};
    \node[lbl, text=FedBlue, below=0.35cm] at (0,0) {\footnotesize condensation};

    % ---- VASE D'EXPANSION (symbole ISO) ----
    \pic[draw=FedGreen!70!black] at (2.5,1.2) {vase-expansion};
    \node[lbl, above=0.75cm] at (2.5,1.2) {Vase exp.};

    % ---- POMPE DE CIRCULATION (symbole ISO) ----
    \pic[draw=FedOrange] at (0,-2.2) {pompe};
    \node[lbl, right=0.55cm] at (0,-2.2) {Circ.};

    % ---- NOURRICE (rectangle simple — nœud de distribution) ----
    \draw[draw=FedBlue, fill=FedBlue!10, rounded corners=2pt]
      (4.5,-0.25) rectangle (5.5,0.25);
    \node[lbl] at (5.0,0) {Nourrice};

    % ---- RADIATEURS (symboles ISO NF) ----
    \pic[draw=FedRed] at (2.5,-3.5) {radiateur};
    \node[lbl, below=0.55cm] at (2.5,-3.5) {Rad. 1};

    \pic[draw=FedRed] at (5.0,-3.5) {radiateur};
    \node[lbl, below=0.55cm] at (5.0,-3.5) {Rad. 2};

    \pic[draw=FedRed] at (7.5,-3.5) {radiateur};
    \node[lbl, below=0.55cm] at (7.5,-3.5) {Rad. 3};

    % ---- VANNES QUART DE TOUR sur chaque radiateur ----
    \pic[draw=FedBlue!60] at (2.5,-2.70) {vanne-quart};
    \pic[draw=FedBlue!60] at (5.0,-2.70) {vanne-quart};
    \pic[draw=FedBlue!60] at (7.5,-2.70) {vanne-quart};

    % ============ TUYAUTERIE DÉPART (rouge) ============
    % Chaudière (est) → nourrice (ouest)
    \draw[->, draw=FedRed, very thick]
      (0.80,0) -- node[above, lbl, text=FedRed]{Départ 70\,°C} (4.5,0);
    % Vase d'expansion raccordé sur la ligne départ
    \draw[draw=FedGreen!70!black, thick] (2.5,0.70) -- (2.5,0) -- (3.5,0);
    % Nourrice → distribution verticale
    \draw[draw=FedRed, thick] (5.0,-0.25) -- (5.0,-2.35);
    \draw[draw=FedRed, thick] (5.0,-2.35) -- (2.5,-2.35) -- (7.5,-2.35);
    % Départ vers vannes
    \draw[->, draw=FedRed, thick] (2.5,-2.35) -- (2.5,-2.38);
    \draw[->, draw=FedRed, thick] (5.0,-2.35) -- (5.0,-2.38);
    \draw[->, draw=FedRed, thick] (7.5,-2.35) -- (7.5,-2.38);
    % Vannes → radiateurs
    \draw[->, draw=FedRed, thick] (2.5,-3.02) -- (2.5,-3.12);
    \draw[->, draw=FedRed, thick] (5.0,-3.02) -- (5.0,-3.12);
    \draw[->, draw=FedRed, thick] (7.5,-3.02) -- (7.5,-3.12);

    % ============ TUYAUTERIE RETOUR (bleu) ============
    \draw[draw=FedBlue, thick] (2.5,-3.88) -- (2.5,-4.6);
    \draw[draw=FedBlue, thick] (5.0,-3.88) -- (5.0,-4.6);
    \draw[draw=FedBlue, thick] (7.5,-3.88) -- (7.5,-4.6);
    \draw[draw=FedBlue, thick] (2.5,-4.6) -- (7.5,-4.6);
    \draw[draw=FedBlue, thick] (5.0,-4.6) -- (5.0,-5.1) -- (0,-5.1) -- (0,-2.62);
    \draw[->, draw=FedBlue, very thick]
      (0,-2.62) -- node[right, lbl, text=FedBlue]{Retour 55\,°C} (0,-2.62);
    \draw[->, draw=FedBlue, very thick] (0,-2.62) -- (0.0,-2.62);
    % Retour → pompe → chaudière
    \draw[draw=FedBlue, very thick] (0,-5.1) -- (0,-2.62);
    \draw[->, draw=FedBlue, very thick] (0,-1.78) -- (0,-0.42);

    % ============ LÉGENDE ============
    \begin{scope}[yshift=-6.0cm]
      \draw[FedRed,  very thick, ->] (0,0) -- (0.7,0);
      \node[right, lbl, text=FedRed]  at (0.8,0)   {Aller — eau chaude};
      \draw[FedBlue, very thick, ->]  (3.8,0) -- (4.5,0);
      \node[right, lbl, text=FedBlue] at (4.6,0)   {Retour — eau tiède};
      \node[lbl, text=FedGreen!70!black]  at (1.5,-0.55)
        {\tikz\pic[draw=FedGreen!70!black, scale=0.7]{vase-expansion}; \;Vase d'expansion};
      \node[lbl, text=FedOrange]          at (6.5,-0.55)
        {\tikz\pic[draw=FedOrange, scale=0.7]{pompe}; \;Pompe circulatrice};
    \end{scope}

  \end{tikzpicture}
  \caption{Circuit hydraulique bouclé — chauffage bitubes (symboles ISO NF GCF :
    chaudière, pompe circulatrice, vase d'expansion, radiateurs, vannes quart de tour)}
  \label{fig:circuit-boucle}
\end{figure}

\subsection{Courbe caractéristique pompe / réseau}

\begin{figure}[H]
  \centering
  \begin{tikzpicture}[>=Stealth, thick]
    % Axes
    \draw[->] (0,0) -- (7.5,0) node[right, font=\small] {$q_v$ (\si{\cubic\metre\per\hour})};
    \draw[->] (0,0) -- (0,5.5) node[above, font=\small] {$H$ (mCE)};

    % Graduations axe X
    \foreach \x/\lbl in {1/0{,}5, 2/1{,}0, 3/1{,}5, 4/2{,}0, 5/2{,}5, 6/3{,}0} {
      \draw (\x,0.05) -- (\x,-0.05);
      \node[below, font=\tiny] at (\x,0) {\lbl};
    }
    % Graduations axe Y
    \foreach \y/\lbl in {1/1, 2/2, 3/3, 4/4, 5/5} {
      \draw (0.05,\y) -- (-0.05,\y);
      \node[left, font=\tiny] at (0,\y) {\lbl};
    }

    % Courbe pompe : H = 5 - 0.2*qv^2 (parabole descendante)
    \draw[color=FedBlue, very thick, domain=0:6.5, samples=60]
      plot (\x, {5 - 0.1*\x*\x})
      node[right, font=\small, text=FedBlue] {Courbe pompe};

    % Courbe réseau : H = 0.12*qv^2 (parabole ascendante, passe par 0)
    \draw[color=FedOrange, very thick, domain=0:6.2, samples=60]
      plot (\x, {0.13*\x*\x})
      node[above, font=\small, text=FedOrange] {Courbe réseau};

    % Point de fonctionnement : intersection de 5 - 0.1*x^2 = 0.13*x^2
    % 5 = 0.23*x^2 => x^2 = 21.74 => x = 4.66, y = 0.13*21.74 = 2.83
    \draw[dashed, color=FedRed] (4.66,0) -- (4.66,2.83) -- (0,2.83);
    \filldraw[FedRed] (4.66,2.83) circle (3pt);
    \node[above right, font=\small\bfseries, text=FedRed] at (4.66,2.83) {Pt de fonctionnement};
    \node[below, font=\tiny, text=FedRed] at (4.66,0) {$q_{v,0}$};
    \node[left, font=\tiny, text=FedRed] at (0,2.83) {$H_0$};

    % Zone de bon rendement (hachures)
    \draw[color=FedGreen!60, dashed, thick] (3.5,0) -- (3.5,{5 - 0.1*3.5*3.5});
    \draw[color=FedGreen!60, dashed, thick] (5.5,0) -- (5.5,{5 - 0.1*5.5*5.5});
    \node[font=\tiny, text=FedGreen!80!black, rotate=90] at (4.5,1.0) {Zone optimale};
  \end{tikzpicture}
  \caption{Courbe caractéristique pompe et courbe de réseau — détermination du point de fonctionnement}
  \label{fig:courbe-pompe-reseau}
\end{figure}

\begin{methode}[title=Lecture du point de fonctionnement]
Le point de fonctionnement $\left(q_{v,0},\, H_0\right)$ est l'intersection des deux courbes :
\begin{itemize}
  \item \textbf{Courbe pompe} : $H_{pompe}(q_v)$ décroissante — fournie par le fabricant
  \item \textbf{Courbe réseau} : $H_{réseau} = k \cdot q_v^2$ croissante — calculée par l'installateur
\end{itemize}
Si le réseau est modifié (équilibrage, remplacement d'un émetteur), la courbe réseau change et le point de fonctionnement se déplace sur la courbe pompe.
\end{methode}

\subsection{Régimes d'écoulement — Profils de vitesse}

\begin{figure}[H]
  \centering
  \begin{tikzpicture}[>=Stealth, thick, node distance=4.5cm]

    %--- Régime laminaire (gauche) ---
    \begin{scope}[xshift=0cm]
      % Tube
      \draw[draw=FedBlue!50, fill=FedBlue!8] (0,-1.2) rectangle (4,1.2);
      \draw[very thick] (0,-1.2) -- (4,-1.2);
      \draw[very thick] (0, 1.2) -- (4, 1.2);

      % Profil parabolique : v(r) = v_max*(1-(r/R)^2), R=1.2
      % r=0.0 : vlen=1.69; r=±0.3 : vlen=1.55; r=±0.6 : vlen=1.12; r=±0.9 : vlen=0.62; r=±1.1 : vlen=0.23
      \draw[->, color=FedBlue, thick] (0.6, 0.0)  -- ++(1.69,0);
      \draw[->, color=FedBlue, thick] (0.6, 0.3)  -- ++(1.55,0);
      \draw[->, color=FedBlue, thick] (0.6,-0.3)  -- ++(1.55,0);
      \draw[->, color=FedBlue, thick] (0.6, 0.6)  -- ++(1.12,0);
      \draw[->, color=FedBlue, thick] (0.6,-0.6)  -- ++(1.12,0);
      \draw[->, color=FedBlue, thick] (0.6, 0.9)  -- ++(0.62,0);
      \draw[->, color=FedBlue, thick] (0.6,-0.9)  -- ++(0.62,0);
      \draw[->, color=FedBlue, thick] (0.6, 1.1)  -- ++(0.23,0);
      \draw[->, color=FedBlue, thick] (0.6,-1.1)  -- ++(0.23,0);

      % Légende
      \node[below, font=\small\bfseries, text=FedBlue] at (2,-1.5) {Laminaire ($Re < 2300$)};
      \node[below, font=\footnotesize] at (2,-1.9) {Profil parabolique};
      \node[right, font=\tiny] at (2.3,0) {$v_{max}$};

      % Note Reynolds
      \node[font=\footnotesize, text=FedBlue!70!black] at (2,1.6) {Filets parallèles, $\lambda = 64/Re$};
    \end{scope}

    %--- Régime turbulent (droite) ---
    \begin{scope}[xshift=5.5cm]
      % Tube
      \draw[draw=FedOrange!50, fill=FedOrange!8] (0,-1.2) rectangle (4,1.2);
      \draw[very thick] (0,-1.2) -- (4,-1.2);
      \draw[very thick] (0, 1.2) -- (4, 1.2);

      % Profil quasi-uniforme : flèches presque égales, chute rapide en paroi
      \draw[->, color=FedOrange, thick] (0.6, 0.0)  -- ++(1.69,0);
      \draw[->, color=FedOrange, thick] (0.6, 0.3)  -- ++(1.63,0);
      \draw[->, color=FedOrange, thick] (0.6,-0.3)  -- ++(1.63,0);
      \draw[->, color=FedOrange, thick] (0.6, 0.6)  -- ++(1.50,0);
      \draw[->, color=FedOrange, thick] (0.6,-0.6)  -- ++(1.50,0);
      \draw[->, color=FedOrange, thick] (0.6, 0.9)  -- ++(1.20,0);
      \draw[->, color=FedOrange, thick] (0.6,-0.9)  -- ++(1.20,0);
      \draw[->, color=FedOrange, thick] (0.6, 1.1)  -- ++(0.40,0);
      \draw[->, color=FedOrange, thick] (0.6,-1.1)  -- ++(0.40,0);

      % Perturbations schématisées
      \draw[color=FedOrange!60, thin, ->] (1.5, 0.4) -- ++(-0.1,0.25) -- ++(0.15,0.2);
      \draw[color=FedOrange!60, thin, ->] (1.5,-0.4) -- ++(0.1,-0.25) -- ++(-0.15,-0.2);
      \draw[color=FedOrange!60, thin, ->] (2.3, 0.4) -- ++(-0.1,0.25) -- ++(0.15,0.2);
      \draw[color=FedOrange!60, thin, ->] (2.3,-0.4) -- ++(0.1,-0.25) -- ++(-0.15,-0.2);

      % Légende
      \node[below, font=\small\bfseries, text=FedOrange] at (2,-1.5) {Turbulent ($Re > 4000$)};
      \node[below, font=\footnotesize] at (2,-1.9) {Profil quasi-uniforme};
      \node[font=\footnotesize, text=FedOrange!70!black] at (2,1.6) {Mélange transversal, $\lambda = f(Re,\, \varepsilon/D)$};
    \end{scope}

    % Flèche de transition au centre
    \draw[->, very thick, color=FedGreen] (4.2,0) -- node[above, font=\footnotesize, text=FedGreen]{$Re \uparrow$} (5.3,0);

  \end{tikzpicture}
  \caption{Profils de vitesse en régime laminaire et turbulent dans une conduite circulaire}
  \label{fig:profils-ecoulement}
\end{figure}

% ============================================================
\section{Éléments de réglementation}
% ============================================================

\begin{description}
  \item[DTU 65.11 (NF P52-211)] Dispositifs de sécurité des installations de chauffage central. Prescrit les dispositifs de protection (soupapes, vases d'expansion, pressostat) et leurs dimensionnements.
  \item[DTU 60.1 (NF P40-201)] Plomberie sanitaire — règles de calcul et installation. Dimensionnement des réseaux d'eau froide et chaude sanitaire.
  \item[NF EN 12828] Installations de chauffage dans les bâtiments — conception des systèmes de chauffage à eau chaude. Prescrit notamment les pressions de service, températures maximales et dispositifs de sécurité.
  \item[NF EN 14336] Mise en service des installations de chauffage à eau chaude. Procédures d'essais d'étanchéité, rinçage, remplissage et équilibrage.
  \item[NF EN 12831] Méthode de calcul des besoins en chaleur — débit à distribuer par émetteur (base du dimensionnement des tuyauteries).
  \item[RE 2020] Impose une approche de calcul de la consommation d'énergie primaire incluant les auxiliaires (pompes). Favorise les circulateurs à haut rendement (classe ErP 2015 minimum).
  \item[Règlement EU 547/2012] Définit les exigences d'écoconception des pompes à eau. Les circulateurs de chauffage doivent satisfaire l'indice d'efficacité énergétique $IEE \leq 0{,}23$.
\end{description}

\begin{tcolorbox}[
  enhanced, colback=FedBlue!3, colframe=FedBlue!50,
  fonttitle=\bfseries\sffamily, coltitle=FedBlue,
  title={FICHE MÉMO — Hydraulique des circuits}, boxrule=0.8pt, arc=2pt,
  width=\textwidth
]
\begin{multicols}{2}
\textbf{Débit volumique}
\[ q_v = v \cdot S \quad \left[\si{\cubic\metre\per\second}\right] \]
avec $v$ = vitesse [\si{\metre\per\second}], $S$ = section droite [\si{\metre\squared}]

\medskip
\textbf{Débit massique}
\[ \dot{m} = \rho \cdot q_v \quad \left[\si{\kilogram\per\second}\right] \]
avec $\rho$ = masse volumique [\si{\kilogram\per\cubic\metre}]

\medskip
\textbf{Nombre de Reynolds}
\[ Re = \frac{v \cdot D}{\nu} \]
$Re < 2300$ : laminaire \quad $Re > 4000$ : turbulent \\
avec $\nu$ = viscosité cinématique [\si{\metre\squared\per\second}]

\medskip
\textbf{Pertes de charge régulières (Darcy-Weisbach)}
\[ \Delta P_r = \lambda \cdot \frac{L}{D} \cdot \frac{\rho\, v^2}{2} \quad \left[\si{\pascal}\right] \]
avec $\lambda$ = coeff. de frottement, $L$ = longueur [\si{\metre}]

\columnbreak

\textbf{Pertes de charge singulières}
\[ \Delta P_s = \zeta \cdot \frac{\rho\, v^2}{2} \quad \left[\si{\pascal}\right] \]
avec $\zeta$ = coefficient de singularité (coude 90° : $\zeta \approx 1{,}5$)

\medskip
\textbf{Méthode de la longueur équivalente}
\[ \Delta P_{total} = R \cdot \bigl(L_{réel} + \textstyle\sum L_{eq}\bigr) \]
avec $R$ = perte de charge linéique [\si{\pascal\per\metre}] \\
Règle GCF : cible $R$ entre \SI{100}{} et \SI{300}{\pascal\per\metre}

\medskip
\textbf{HMT — Hauteur manométrique totale}
\[ HMT = \Delta h_{pertes\,totales} \quad \left[\si{mCE}\right] \]

\medskip
\textbf{Puissance hydraulique}
\[ P_{hyd} = \rho \cdot g \cdot q_v \cdot HMT \quad \left[\si{\watt}\right] \]

\medskip
\textbf{Puissance absorbée par la pompe}
\[ P_{abs} = \frac{P_{hyd}}{\eta_{pompe}} \quad \left[\si{\watt}\right] \]
\textit{Exigence RE : IEE $\leq 0{,}23$ (circulateurs ErP 2015)}
\end{multicols}
\end{tcolorbox}

\finchapitre

% \chapter{Aéraulique}

% ============================================================
\section{Objectifs pédagogiques}
% ============================================================

\subsection*{Compétences GCF mobilisées}
\begin{itemize}
  \item \textbf{C2-1 à C2-3} — Analyser un réseau aéraulique : identifier composants, chaînes d'énergie, modes de fonctionnement
  \item \textbf{C3-3} — Dimensionner un réseau de gaines et sélectionner les équipements (ventilateurs, diffuseurs)
  \item \textbf{C7-4 à C7-6} — Mesurer débits d'air et pressions, interpréter les résultats en mise en service
  \item \textbf{C8-1 à C8-4} — Vérifier l'équilibrage aéraulique, diagnostiquer les dérives, proposer des actions correctives
\end{itemize}

\subsection*{Savoirs associés mobilisés}
\begin{itemize}
  \item \textbf{A3} — Dynamique des fluides : pertes de charge, point de fonctionnement, équilibrage (niveau 3)
  \item \textbf{A4} — Traitement d'air : diagramme de l'air humide, évolutions, soufflage (niveau 3)
  \item \textbf{B2-1} — Ventilation simple/double flux, free cooling, récupération de chaleur (niveau 3)
  \item \textbf{B4-2} — Réseaux aérauliques : gaines, filtres, diffuseurs, équilibrage, ventilateurs (niveau 3)
  \item \textbf{S5.5} — Réglementations ventilation et climatisation (niveau 3)
  \item \textbf{S10.5} — Mise en service d'une installation GCF (niveau 3)
  \item \textbf{S10.6} — Critères de bon fonctionnement et optimisation (niveau 3)
\end{itemize}

% ============================================================
\section{Introduction}
% ============================================================

L'aéraulique est la discipline qui étudie les écoulements d'air dans les bâtiments et leurs réseaux de distribution. Elle est au cœur des systèmes de ventilation, de climatisation et de traitement d'air (CVC). Un réseau mal dimensionné engendre des nuisances sonores, une mauvaise qualité d'air intérieur (QAI), des surconsommations électriques et un inconfort thermique.

En France, la réglementation impose des débits minimaux de ventilation dans tous les locaux habités (arrêté du 24 mars 1982, NF EN 16798-3). Le technicien supérieur GCF dimensionne les réseaux de gaines, sélectionne les ventilateurs et les diffuseurs, équilibre les débits et vérifie les performances en mise en service.

% ============================================================
\section{Propriétés de l'air humide}
% ============================================================

\subsection{Composition et grandeurs caractéristiques}

L'air atmosphérique est un mélange d'air sec et de vapeur d'eau. Ses propriétés varient avec la température et la teneur en humidité.

\begin{definition}[title=Humidité spécifique (teneur en eau)]
\begin{equation}
  w = \frac{m_{vapeur}}{m_{air\,sec}} \qquad \left[\si{\kilogram\per\kilogram}\right]
\end{equation}
\begin{itemize}
  \item $m_{vapeur}$ : masse de vapeur d'eau (\si{\kilogram})
  \item $m_{air\,sec}$ : masse d'air sec (\si{\kilogram})
\end{itemize}
L'humidité spécifique est aussi notée $x$ ou $W$ selon les sources. En pratique : air extérieur hivernal $w \approx \SI{3}{\gram\per\kilogram}$, air estival $w \approx \SI{10}{\gram\per\kilogram}$.
\end{definition}

\begin{definition}[title=Humidité relative]
\begin{equation}
  HR = \frac{p_{vapeur}}{p_{vapeur\,sat}(T)} \times 100 \qquad [\%]
\end{equation}
\begin{itemize}
  \item $p_{vapeur}$ : pression partielle de la vapeur d'eau dans l'air (\si{\pascal})
  \item $p_{vapeur\,sat}(T)$ : pression de saturation de la vapeur à la température $T$ (\si{\pascal})
\end{itemize}
Confort recommandé : $HR$ entre 40 et 60\,\% (NF EN 15251).
\end{definition}

\begin{definition}[title=Enthalpie de l'air humide]
\begin{equation}
  h = c_{p,a} \cdot T + w \cdot (L_0 + c_{p,v} \cdot T) \qquad \left[\si{\kilo\joule\per\kilogram}\right]
\end{equation}
\begin{itemize}
  \item $c_{p,a} = \SI{1006}{\joule\per\kilogram\per\kelvin}$ : chaleur massique de l'air sec
  \item $c_{p,v} = \SI{1805}{\joule\per\kilogram\per\kelvin}$ : chaleur massique de la vapeur d'eau
  \item $L_0 = \SI{2501}{\kilo\joule\per\kilogram}$ : chaleur latente de vaporisation à 0\,°C
  \item $T$ : température en °C
\end{itemize}
Formule simplifiée pratique : $h \approx 1{,}006 \cdot T + w \cdot (2501 + 1{,}86 \cdot T)$
\end{definition}

\begin{definition}[title=Masse volumique de l'air humide]
À pression atmosphérique standard (\SI{101325}{\pascal}) :
\begin{equation}
  \rho \approx \frac{1{,}293}{1 + 0{,}608 \cdot w} \cdot \frac{273}{273 + T} \qquad \left[\si{\kilogram\per\cubic\metre}\right]
\end{equation}
\textbf{Valeurs usuelles en CVC :}
\begin{itemize}
  \item Air à 20\,°C, $HR = 50\,\%$ : $\rho \approx \SI{1{,}20}{\kilogram\per\cubic\metre}$
  \item Air à 0\,°C : $\rho \approx \SI{1{,}29}{\kilogram\per\cubic\metre}$
\end{itemize}
\end{definition}

\subsection{Le diagramme de l'air humide (diagramme de Mollier)}

Le diagramme de l'air humide représente graphiquement toutes les propriétés de l'air ($T$, $w$, $HR$, $h$, $T_{rosée}$). Il est indispensable pour calculer les évolutions d'air dans une CTA (Centrale de Traitement d'Air).

\begin{table}[H]
  \centering
  \caption{Évolutions de l'air sur le diagramme de Mollier}
  \begin{tabular}{p{3.5cm}p{4cm}p{4cm}}
    \toprule
    Évolution & Paramètre constant & Exemple \\
    \midrule
    Chauffage sensible & $w$ constant & Batterie chaude, résistance \\
    Refroidissement sec & $w$ constant, $T > T_{rosée}$ & Début refroidissement \\
    Refroidissement humide & $w$ diminue & Batterie froide (condensation) \\
    Humidification adiabatique & $h$ constant & Humidificateur à eau \\
    Humidification par vapeur & $w$ augmente, $T \uparrow$ & Injection vapeur \\
    Mélange de deux airs & Point sur le segment & Mélange air neuf / repris \\
    \bottomrule
  \end{tabular}
\end{table}

% ============================================================
\section{Dynamique des fluides gazeux}
% ============================================================

\subsection{Débits d'air}

\begin{definition}[title=Débit volumique d'air]
\begin{equation}
  q_v = v \cdot S \qquad \left[\si{\cubic\metre\per\hour}\right]
\end{equation}
\begin{itemize}
  \item $v$ : vitesse de l'air dans la gaine (\si{\metre\per\second})
  \item $S$ : section droite de la gaine (\si{\square\metre})
\end{itemize}
\end{definition}

\begin{definition}[title=Débit massique d'air]
\begin{equation}
  \dot{m} = \rho \cdot q_v \qquad \left[\si{\kilogram\per\hour}\right]
\end{equation}
\end{definition}

\begin{methode}[title=Vitesses recommandées dans les gaines (NF EN 13779)]
\begin{center}
\begin{tabular}{lll}
  \toprule
  Type de gaine & Vitesse max (m/s) & Application \\
  \midrule
  Soufflage basse vitesse & 2 à 4 & Tertiaire, résidentiel \\
  Soufflage haute vitesse & 6 à 10 & Industrie, grandes surfaces \\
  Gaine principale & 5 à 8 & Transport \\
  Gaine terminale & 2 à 4 & Distribution locale \\
  Prise/rejet air extérieur & 2 à 3 & Façade, toiture \\
  \bottomrule
\end{tabular}
\end{center}
Des vitesses excessives provoquent du bruit et des pertes de charge importantes.
\end{methode}

\subsection{Pertes de charge dans les réseaux de gaines}

\subsubsection{Pertes de charge régulières}

La formule de Darcy-Weisbach s'applique également à l'air :
\begin{equation}
  \Delta P_r = \lambda \cdot \frac{L}{D_h} \cdot \frac{\rho v^2}{2} \qquad \left[\si{\pascal}\right]
\end{equation}

Pour les gaines rectangulaires de dimensions $a \times b$, on utilise le \textbf{diamètre hydraulique équivalent} :
\begin{equation}
  D_h = \frac{2 \cdot a \cdot b}{a + b} \qquad \left[\si{\metre}\right]
\end{equation}

En pratique, on utilise des abaques ou logiciels donnant la perte de charge linéique $R$ (\si{\pascal\per\metre}) en fonction du débit et du diamètre (ou des dimensions).

\subsubsection{Pertes de charge singulières}

\begin{equation}
  \Delta P_s = \zeta \cdot \frac{\rho v^2}{2} \qquad \left[\si{\pascal}\right]
\end{equation}

\begin{table}[H]
  \centering
  \caption{Coefficients $\zeta$ courants pour les accessoires aérauliques}
  \begin{tabular}{lll}
    \toprule
    Accessoire & Conditions & $\zeta$ typique \\
    \midrule
    Coude rectangulaire 90° & sans aube & 1{,}5 à 2{,}5 \\
    Coude rectangulaire 90° & avec aubes & 0{,}3 à 0{,}5 \\
    Piquage en té & dérivation & 0{,}8 à 1{,}5 \\
    Convergent 30° & — & 0{,}05 \\
    Divergent 15° & — & 0{,}10 \\
    Filtre à poches (propre) & — & 20 à 50 Pa (valeur abs.) \\
    Batterie à eau chaude & — & 30 à 80 Pa (valeur abs.) \\
    Registre d'air frais & ouvert & 3 à 5 \\
    \bottomrule
  \end{tabular}
\end{table}

\begin{attention}
Pour les filtres et batteries, les fabricants fournissent directement la perte de charge en \si{\pascal} (valeur absolue, pas un coefficient $\zeta$). Ces valeurs augmentent avec l'encrassement des filtres — prévoir la perte de charge \textbf{filtre colmaté} pour le dimensionnement.
\end{attention}

% ============================================================
\section{Les ventilateurs}
% ============================================================

\subsection{Types de ventilateurs}

\begin{table}[H]
  \centering
  \caption{Types de ventilateurs et domaines d'application en CVC}
  \begin{tabular}{p{3cm}p{4.5cm}p{4.5cm}}
    \toprule
    Type & Principe & Application GCF \\
    \midrule
    Centrifuge à pales en avant & Fort débit, faible pression & Ventilation résidentielle basse vitesse \\
    Centrifuge à pales en arrière & Rendement élevé (55–65\,\%) & CTA tertiaire, industrielle \\
    Hélicoïdal & Faible pression, fort débit & Extraction directe, condenseurs \\
    Tangentiel & Faible débit silencieux & Ventilo-convecteurs, split-systems \\
    EC (moteur à commutation électronique) & Régulation vitesse intégrée, IE4 & Toutes CTA modernes (RE2020) \\
    \bottomrule
  \end{tabular}
\end{table}

\subsection{Caractéristiques d'un ventilateur}

\begin{definition}[title=Pression totale (HMT aéraulique)]
La pression totale $\Delta P_t$ est la différence de pression totale (statique + dynamique) entre le refoulement et l'aspiration du ventilateur.
\begin{equation}
  \Delta P_t = \Delta P_{statique} + \Delta P_{dynamique} = \Delta P_s + \frac{\rho v^2}{2} \qquad \left[\si{\pascal}\right]
\end{equation}
\end{definition}

\begin{definition}[title=Puissance absorbée et rendement]
\begin{equation}
  P_{hyd} = q_v \cdot \Delta P_t \qquad \left[\si{\watt}\right]
\end{equation}
\begin{equation}
  P_{abs} = \frac{P_{hyd}}{\eta_{ventilateur}} \qquad \left[\si{\watt}\right]
\end{equation}
Rendement typique d'un ventilateur centrifuge à pales en arrière : $\eta = 55$ à $65\,\%$.
\end{definition}

\subsection{Lois de similitude des ventilateurs}

Les lois de similitude permettent de prévoir le comportement d'un ventilateur lorsqu'on modifie sa vitesse de rotation. Elles sont fondamentales pour le calcul des économies d'énergie par variation de vitesse.

\begin{methode}[title=Lois de similitude (variation de vitesse)]
Pour un même ventilateur, si la vitesse passe de $n_1$ à $n_2$ :
\begin{align}
  q_{v2} &= q_{v1} \cdot \frac{n_2}{n_1} & \text{(débit)} \label{eq:sim_q}\\
  \Delta P_{t2} &= \Delta P_{t1} \cdot \left(\frac{n_2}{n_1}\right)^2 & \text{(pression)} \label{eq:sim_p}\\
  P_{abs2} &= P_{abs1} \cdot \left(\frac{n_2}{n_1}\right)^3 & \text{(puissance)} \label{eq:sim_p3}
\end{align}

\textbf{Conséquence énergétique majeure :} réduire la vitesse de 20\,\% ($n_2/n_1 = 0{,}8$) divise la puissance absorbée par $0{,}8^3 = 0{,}51$ — soit une économie de 49\,\% sur la consommation électrique du ventilateur.
\end{methode}

% ============================================================
\section{Systèmes de ventilation}
% ============================================================

\subsection{Ventilation mécanique contrôlée (VMC)}

\begin{table}[H]
  \centering
  \caption{Comparatif des systèmes VMC selon la réglementation française}
  \begin{tabular}{p{2.5cm}p{4.5cm}p{4.5cm}}
    \toprule
    Système & Principe & Contexte d'usage \\
    \midrule
    VMC simple flux autoréglable & Extraction mécanique, entrées d'air autoréglables. Débit constant. & Logement neuf standard \\
    VMC simple flux hygroréglable & Débit modulé selon l'humidité (bouches + entrées hygro). & Logement économe en énergie \\
    VMC double flux & Extraction + soufflage avec récupérateur de chaleur (échangeur). Rendement 70–95\,\%. & Logement très performant, BBC, RE2020 \\
    Ventilation naturelle assistée (VNA) & Extraction naturelle + appoint mécanique. & Réhabilitation \\
    \bottomrule
  \end{tabular}
\end{table}

\subsection{Centrale de traitement d'air (CTA)}

Une CTA conditionne l'air neuf (filtration, chauffage, refroidissement, humidification/déshumidification) avant de le distribuer dans le bâtiment.

\textbf{Composants typiques d'une CTA tertiaire (dans l'ordre de l'air) :}
\begin{enumerate}
  \item Filtre grossier (G4) — protection mécanique
  \item Récupérateur de chaleur (échangeur à plaques ou rotatif)
  \item Batterie eau chaude (chauffage hiver)
  \item Filtre fin (F7/F9) — filtration particulaire
  \item Batterie eau glacée (refroidissement été)
  \item Humidificateur (vapeur ou adiabatique)
  \item Ventilateur de soufflage (moteur EC)
  \item Gaine de distribution
  \item Ventilateur d'extraction
\end{enumerate}

\begin{definition}[title=Débit d'air neuf réglementaire]
L'arrêté du 24 mars 1982 (logements) et la NF EN 16798-1 (tertiaire) fixent les débits minimaux d'air neuf selon l'occupation et l'usage.

\textbf{Valeurs courantes tertiaire (NF EN 16798-1, catégorie II) :}
\begin{itemize}
  \item Bureau : 7 à 10 \si{\litre\per\second} par personne + 1 \si{\litre\per\second\per\square\metre}
  \item Salle de réunion : 10 \si{\litre\per\second} par personne
  \item Restauration : 10 à 15 \si{\litre\per\second} par personne
\end{itemize}
\end{definition}

% ============================================================
\section{Équilibrage des réseaux aérauliques}
% ============================================================

L'équilibrage aéraulique garantit que chaque bouche ou diffuseur reçoit exactement le débit nominal.

\begin{methode}[title=Méthode de la branche index (aéraulique)]
Identique à la méthode hydraulique :
\begin{enumerate}
  \item Calculer les pertes de charge de chaque branche (du ventilateur à chaque bouche)
  \item Identifier la branche index (pertes de charge maximales)
  \item Calculer l'excédent de pression sur chaque autre branche
  \item Régler les registres d'équilibrage pour absorber l'excédent
\end{enumerate}

Les registres utilisés sont :
\begin{itemize}
  \item \textbf{Registres à lamelles} — réglage manuel, verrouillage par vis
  \item \textbf{Registres à débit constant} (CAV) — maintien mécanique du débit quelle que soit la pression
  \item \textbf{Boîtes à débit variable} (VAV) — débit modulé selon la demande (occupation, CO\textsubscript{2})
\end{itemize}
\end{methode}

% ============================================================
\section{Exemples résolus}
% ============================================================

\begin{exemple}[title=Dimensionnement d'une gaine de soufflage]
\textbf{Énoncé :} Une gaine rectangulaire de section 400 × 200 mm alimente 4 bureaux identiques. Le débit total est de \SI{2000}{\cubic\metre\per\hour}. L'air est à 18\,°C ($\rho = \SI{1{,}21}{\kilogram\per\cubic\metre}$). La longueur développée de la gaine est de \SI{25}{\metre}. La perte de charge linéique est $R = \SI{0{,}8}{\pascal\per\metre}$. La gaine comporte 2 coudes 90° avec aubes ($\zeta = 0{,}4$) et 4 piquages en té ($\zeta = 1{,}2$ par piquage, vitesse gaine principale).

Calculer la perte de charge totale de cette gaine.

\textbf{Étape 1 : Section et vitesse}
\[
  S = 0{,}4 \times 0{,}2 = \SI{0{,}08}{\square\metre}
\]
\[
  q_v = \frac{2000}{3600} = \SI{0{,}556}{\cubic\metre\per\second}
\]
\[
  v = \frac{0{,}556}{0{,}08} = \SI{6{,}94}{\metre\per\second}
\]
Vitesse acceptable pour une gaine principale tertiaire.

\textbf{Étape 2 : Pression dynamique}
\[
  \frac{\rho v^2}{2} = \frac{1{,}21 \times (6{,}94)^2}{2} = \SI{29{,}1}{\pascal}
\]

\textbf{Étape 3 : Perte de charge régulière}
\[
  \Delta P_r = R \times L = 0{,}8 \times 25 = \SI{20}{\pascal}
\]

\textbf{Étape 4 : Pertes de charge singulières}
\[
  \Delta P_s = \left(2 \times 0{,}4 + 4 \times 1{,}2\right) \times 29{,}1 = 5{,}6 \times 29{,}1 = \SI{163}{\pascal}
\]

\textbf{Étape 5 : Perte de charge totale}
\[
  \Delta P_{total} = 20 + 163 = \SI{183}{\pascal}
\]

\textbf{Note :} Les piquages représentent la majorité des pertes. Des piquages coniques (avec divergent) réduiraient significativement $\zeta$.
\end{exemple}

\begin{exemple}[title=Sélection d'un ventilateur CTA]
\textbf{Énoncé :} Une CTA tertiaire dessert 500\,m² de bureaux. Le débit d'air neuf est de \SI{3600}{\cubic\metre\per\hour}. La perte de charge totale du réseau (réseau soufflage + réseau extraction) est estimée à \SI{450}{\pascal} côté soufflage. Le ventilateur est à moteur EC ($\eta = 60\,\%$).

\textbf{Étape 1 : Puissance hydraulique}
\[
  q_v = \frac{3600}{3600} = \SI{1}{\cubic\metre\per\second}
\]
\[
  P_{hyd} = q_v \times \Delta P_t = 1 \times 450 = \SI{450}{\watt}
\]

\textbf{Étape 2 : Puissance absorbée}
\[
  P_{abs} = \frac{450}{0{,}60} = \SI{750}{\watt}
\]

\textbf{Étape 3 : Économie par variation de vitesse}
En période d'occupation partielle (débit réduit à 60\,\%) :
\[
  P_{abs,60\%} = 750 \times (0{,}6)^3 = 750 \times 0{,}216 = \SI{162}{\watt}
\]
Économie : $750 - 162 = \SI{588}{\watt}$ soit 78\,\% d'économie par rapport au plein débit.
\end{exemple}

% ============================================================
\section{Exercices d'application}
% ============================================================

\exercice{Calcul de débit d'air neuf réglementaire}

Un plateau de bureaux de \SI{800}{\square\metre} accueille 60 personnes en occupation nominale. En dehors des heures de bureau, le plateau est inoccupé. Utiliser la catégorie II de la NF EN 16798-1 : 10 \si{\litre\per\second} par personne + 1 \si{\litre\per\second\per\square\metre}.

\begin{enumerate}
  \item Calculer le débit d'air neuf minimal en occupation nominale (en m³/h).
  \item Calculer le débit réduit hors occupation (10\,\% du nominal, exigence minimale réglementaire de balayage).
  \item Un système VAV (débit variable) réduit le débit à 30\,\% en heures creuses (20\,h/sem) par rapport aux heures pleines (40\,h/sem). Le ventilateur absorbe \SI{2{,}2}{\kilo\watt} en plein débit. Estimer la consommation hebdomadaire en kWh avec et sans variation de vitesse.
\end{enumerate}

\textbf{Corrigé :}

\textit{Question 1 — Débit nominal :}
\[
  q_{v,pers} = 60 \times 10 = \SI{600}{\litre\per\second}
\]
\[
  q_{v,surf} = 800 \times 1 = \SI{800}{\litre\per\second}
\]
\[
  q_{v,total} = (600 + 800) \times 3{,}6 = 1400 \times 3{,}6 = \SI{5040}{\cubic\metre\per\hour}
\]

\textit{Question 2 — Débit hors occupation :}
\[
  q_{v,vide} = 0{,}10 \times 5040 = \SI{504}{\cubic\metre\per\hour}
\]

\textit{Question 3 — Consommation avec/sans variation de vitesse :}

\textit{Sans variation (plein débit permanent) :}
\[
  W = 2{,}2 \times (40 + 20) = \SI{132}{\kilo\watt\heure\per\semaine}
\]

\textit{Avec variation (30\,\% en heures creuses) :}
\[
  P_{hc} = 2{,}2 \times (0{,}30)^3 = 2{,}2 \times 0{,}027 = \SI{0{,}059}{\kilo\watt}
\]
\[
  W = 2{,}2 \times 40 + 0{,}059 \times 20 = 88 + 1{,}18 = \SI{89{,}2}{\kilo\watt\heure\per\semaine}
\]
Économie : $132 - 89{,}2 = \SI{42{,}8}{\kilo\watt\heure\per\semaine}$ soit 32\,\%.

\exercice{Équilibrage d'un réseau de soufflage}

Un réseau de soufflage dessert 5 salles. Après calcul des pertes de charge de chaque branche (du ventilateur à la bouche), on obtient :

\begin{center}
\begin{tabular}{lcc}
  \toprule
  Salle & $\Delta P$ branche (Pa) & Débit nominal (m³/h) \\
  \midrule
  S1 (proche) & 180 & 300 \\
  S2 & 220 & 250 \\
  S3 & 310 & 200 \\
  S4 & 410 & 180 \\
  S5 (index) & 520 & 160 \\
  \bottomrule
\end{tabular}
\end{center}

\begin{enumerate}
  \item Identifier la branche index et la pression totale que doit fournir le ventilateur.
  \item Calculer la perte de charge supplémentaire à créer par le registre d'équilibrage de chaque salle.
  \item Quel type d'organe d'équilibrage recommander pour une installation tertiaire avec variation d'occupation ?
\end{enumerate}

\textbf{Corrigé :}

\textit{Question 1 :} La branche index est S5 avec $\Delta P = \SI{520}{\pascal}$. Le ventilateur doit fournir au minimum \SI{520}{\pascal}.

\textit{Question 2 :}
\begin{itemize}
  \item S1 : registre $= 520 - 180 = \SI{340}{\pascal}$
  \item S2 : registre $= 520 - 220 = \SI{300}{\pascal}$
  \item S3 : registre $= 520 - 310 = \SI{210}{\pascal}$
  \item S4 : registre $= 520 - 410 = \SI{110}{\pascal}$
  \item S5 : registre $= \SI{0}{\pascal}$ (branche index, pas de réglage)
\end{itemize}

\textit{Question 3 :} Pour un tertiaire avec variation d'occupation, recommander des \textbf{boîtes à débit variable (VAV)} pilotées par sonde CO\textsubscript{2} ou capteur de présence, couplées à un ventilateur à vitesse variable. Le ventilateur adapte sa pression différentielle à la demande (régulation de pression constante sur la gaine principale).

\exercice{Évolution de l'air dans une CTA — diagramme de Mollier}

Une CTA traite de l'air extérieur dans les conditions suivantes :
\begin{itemize}
  \item Air extérieur hiver : $T_1 = \SI{2}{\celsius}$, $HR_1 = 70\,\%$, $w_1 = \SI{3{,}5}{\gram\per\kilogram}$, $h_1 = \SI{10{,}8}{\kilo\joule\per\kilogram}$
  \item Conditions de soufflage souhaitées : $T_s = \SI{22}{\celsius}$, $HR_s = 45\,\%$, $w_s = \SI{7{,}4}{\gram\per\kilogram}$, $h_s = \SI{41{,}0}{\kilo\joule\per\kilogram}$
  \item Débit massique : $\dot{m} = \SI{2500}{\kilogram\per\hour}$
\end{itemize}

\begin{enumerate}
  \item Identifier les évolutions successives de l'air dans la CTA.
  \item Calculer la puissance de la batterie chaude (chauffage sensible de 2\,°C à 15\,°C, $w$ constant).
  \item Calculer la puissance de l'humidificateur (de 15\,°C à 22\,°C avec humidification, $h_{s,humid} = \SI{41{,}0}{\kilo\joule\per\kilogram}$, $h_{avant\_humid} = \SI{26{,}2}{\kilo\joule\per\kilogram}$).
\end{enumerate}

\textbf{Corrigé :}

\textit{Question 1 — Évolutions :}
\begin{enumerate}
  \item Préchauffage par récupérateur (2\,°C → 8\,°C, $w$ constant)
  \item Chauffage batterie chaude (8\,°C → 15\,°C, $w$ constant)
  \item Humidification adiabatique ou vapeur (15\,°C → 22\,°C, $w$ augmente)
\end{enumerate}

\textit{Question 2 — Puissance batterie chaude :}
\[
  \dot{Q}_{BC} = \dot{m} \cdot c_{p,a} \cdot \Delta T = \frac{2500}{3600} \times 1006 \times (15 - 8) = \SI{4902}{\watt} \approx \SI{4{,}9}{\kilo\watt}
\]

\textit{Question 3 — Puissance humidificateur :}
\[
  \dot{Q}_{hum} = \dot{m} \cdot (h_s - h_{av}) = \frac{2500}{3600} \times (41{,}0 - 26{,}2) \times 1000 = \SI{10\,278}{\watt} \approx \SI{10{,}3}{\kilo\watt}
\]

% ============================================================
\section{Éléments de réglementation}
% ============================================================

\begin{description}
  \item[Arrêté du 24 mars 1982 (modifié)] Aération des logements neufs — fixe les débits minimaux d'extraction par pièce (cuisine : 75 ou 135\,m³/h selon cuisson, WC : 15\,m³/h, bain/douche : 15 ou 30\,m³/h).

  \item[NF EN 16798-3] Ventilation des bâtiments non résidentiels — critères de performance et classification des systèmes. Remplace EN 13779.

  \item[NF EN 15251] Critères d'ambiance intérieure pour la conception et l'évaluation de la performance énergétique des bâtiments. Définit 4 catégories de qualité d'air intérieur (I à IV).

  \item[DTU 68.3 (NF P50-410)] Installations de VMC. Règles de conception et de mise en œuvre des réseaux aérauliques dans les logements.

  \item[RE 2020] Impose une ventilation double flux avec récupération pour atteindre les niveaux de performance énergétique exigés dans les logements neufs. Le coefficient $b_{vent}$ pénalise les systèmes sans récupération.

  \item[Règlement EU 1253/2014 (écoconception)] Définit les exigences de rendement minimal des unités de ventilation résidentielles et non résidentielles. Impose une efficacité de récupération de chaleur minimale pour les systèmes double flux.

  \item[NF EN 12599] Procédures d'essais et méthodes de mesure pour la réception des installations de ventilation et de climatisation — mesure des débits, pressions, températures et humidité.
\end{description}

% \chapter{Production de chaleur}

% ============================================================
\section{Objectifs pédagogiques}
% ============================================================

\subsection*{Compétences GCF mobilisées}
\begin{itemize}
  \item \textbf{C2-1 à C2-3} — Analyser les systèmes de production de chaleur : identifier composants, chaînes d'énergie, paramètres de fonctionnement
  \item \textbf{C3-2 et C3-3} — Comparer et sélectionner un générateur de chaleur selon critères technico-économiques et environnementaux ; dimensionner la puissance de production
  \item \textbf{C5-1 à C5-3} — Appliquer les réglementations thermiques et environnementales liées aux générateurs
  \item \textbf{C7-4 à C7-6} — Mesurer les performances en mise en service (rendement, COP), interpréter les écarts
  \item \textbf{C8-1 à C8-4} — Vérifier les performances saisonnières, proposer des optimisations
\end{itemize}

\subsection*{Savoirs associés mobilisés}
\begin{itemize}
  \item \textbf{A8} — Combustion appliquée : excès d'air, fumées, condensation, rendement (niveau 3)
  \item \textbf{A5} — Thermodynamique appliquée — machine frigorifique côté chaud (niveau 2–3)
  \item \textbf{B1-1} — Équipements de chauffage — production : chaudières, PAC, réseau urbain, échangeurs (niveau 3)
  \item \textbf{B6-3} — Cogénération / micro-cogénération (niveau 2)
  \item \textbf{S5.5} — Réglementations fluides thermiques : chauffage, combustibles (niveau 3)
  \item \textbf{S10.5 et S10.6} — Mise en service et critères de bon fonctionnement (niveau 3)
\end{itemize}

% ============================================================
\section{Introduction}
% ============================================================

La production de chaleur est le point de départ de tout système CVC. Le technicien supérieur GCF doit être capable de dimensionner la puissance nécessaire, de sélectionner la technologie la plus adaptée (chaudière, PAC, solaire, réseau urbain) et de vérifier les performances en exploitation.

Le choix du générateur engage l'installation pour 15 à 25 ans. Il dépend des besoins thermiques du bâtiment, de l'énergie disponible (gaz, électricité, bois, réseau chaleur), des contraintes réglementaires (RE2020, labels énergétiques) et de la rentabilité économique. En 2024, la RE2020 oriente fortement vers les pompes à chaleur et les énergies renouvelables, au détriment des chaudières à gaz dans les logements neufs.

% ============================================================
\section{Combustion et chaudières à gaz}
% ============================================================

\subsection{Principes de la combustion}

La combustion du gaz naturel (méthane : CH\textsubscript{4}) est une réaction d'oxydation exothermique :

\begin{equation}
  \text{CH}_4 + 2\,\text{O}_2 \rightarrow \text{CO}_2 + 2\,\text{H}_2\text{O} + \text{chaleur}
\end{equation}

\textbf{Volume d'air stœchiométrique} pour brûler 1\,m³ de gaz naturel : environ 9,5\,m³ d'air (valeur standard réseau GDF-Suez).

\begin{definition}[title=Pouvoir calorifique]
\begin{itemize}
  \item \textbf{PCI} (Pouvoir Calorifique Inférieur) : énergie dégagée en supposant que la vapeur d'eau des fumées reste à l'état gazeux. Pour le gaz naturel : PCI $\approx \SI{10{,}55}{\kilo\watt\heure\per\cubic\metre}$.
  \item \textbf{PCS} (Pouvoir Calorifique Supérieur) : énergie dégagée en récupérant également la chaleur de condensation de la vapeur d'eau. PCS $\approx \SI{11{,}7}{\kilo\watt\heure\per\cubic\metre}$.
  \item Rapport : $\text{PCS} / \text{PCI} \approx 1{,}11$ pour le gaz naturel (chaleur latente de condensation de l'eau).
\end{itemize}
\end{definition}

\begin{definition}[title=Excès d'air]
Pour assurer une combustion complète, on fournit plus d'air que la quantité stœchiométrique théorique. L'excès d'air $e$ est :
\begin{equation}
  e = \frac{V_{air\,réel} - V_{air\,stœchio}}{V_{air\,stœchio}} \times 100 \qquad [\%]
\end{equation}
Un excès d'air trop faible → combustion incomplète (CO, suies, risque intoxication). Trop élevé → refroidissement excessif des fumées → pertes par les fumées chaudes → rendement dégradé.

\textbf{Valeurs typiques :} brûleur atmosphérique 20–40\,\%, brûleur à prémélange modulant 5–15\,\%.
\end{definition}

\begin{table}[H]
  \centering
  \caption{Composition des fumées de gaz naturel selon l'excès d'air (valeurs indicatives)}
  \begin{tabular}{lccccc}
    \toprule
    Excès d'air & CO\textsubscript{2} (\%) & O\textsubscript{2} (\%) & CO (ppm) & $T_{point\,rosée}$ & Indice Bacharach \\
    \midrule
    0\,\% (stœchio) & 11,7 & 0 & élevé & 57\,°C & — \\
    10\,\% & 10,7 & 1,6 & traces & 57\,°C & 0 \\
    20\,\% & 9,8 & 3,2 & 0 & 57\,°C & 0 \\
    40\,\% & 8,3 & 6,0 & 0 & 55\,°C & 0 \\
    \bottomrule
  \end{tabular}
\end{table}

\subsection{Rendement des chaudières}

\begin{definition}[title=Rendement de combustion]
\begin{equation}
  \eta_{comb} = 1 - \frac{\text{pertes par les fumées} + \text{pertes par incombus\-tion}}{\text{PCI}} \qquad [\%]
\end{equation}
\end{definition}

\begin{definition}[title=Rendement saisonnier (SCOP thermique)]
Le rendement saisonnier intègre les pertes en veille, les cycles de démarrage et les pertes en eau chaude sanitaire. Il est obligatoirement affiché sur l'étiquette énergie EU 811/2013.
\end{definition}

\begin{table}[H]
  \centering
  \caption{Comparatif des technologies de chaudières gaz}
  \begin{tabular}{p{3.5cm}p{3cm}p{3cm}p{2.5cm}}
    \toprule
    Technologie & Rendement PCI & Principe & RE2020 logement \\
    \midrule
    Chaudière basse température & 88–92\,\% & Eau retour $> 45$\,°C & Interdit neuf \\
    Chaudière à condensation & 95–109\,\% & Récupère chaleur latente fumées & Interdit neuf depuis 2022 \\
    Micro-cogénération gaz & 85–92\,\% therm. & Produit électricité + chaleur & Cas par cas \\
    \bottomrule
  \end{tabular}
\end{table}

\begin{attention}
Depuis le 1\textsuperscript{er} janvier 2022, la RE2020 interdit les chaudières gaz comme système principal de chauffage dans les \textbf{logements neufs}. Elles restent autorisées en rénovation et dans le tertiaire. Le marché de la rénovation reste donc très important pour le technicien GCF.
\end{attention}

\subsection{Chaudières à condensation — fonctionnement}

La chaudière à condensation récupère la chaleur latente de condensation de la vapeur d'eau contenue dans les fumées. Cette condensation se produit lorsque la température des fumées descend en dessous du \textbf{point de rosée} (environ 57\,°C pour le gaz naturel).

\begin{methode}[title=Conditions de fonctionnement en condensation]
La condensation est effective si la température de retour eau est inférieure à 57\,°C. Pour maximiser le rendement :
\begin{itemize}
  \item Température départ/retour : 60/45\,°C ou 50/30\,°C (plancher chauffant)
  \item Modulation de puissance : le brûleur à prémélange module de 20 à 100\,\%
  \item Neutralisation des condensats : pH 3 à 5 → kit de neutralisation obligatoire (Loi sur l'eau)
  \item Évacuation des fumées : conduit plastique PVC ou inox (fumées condensées corrosives)
\end{itemize}
\end{methode}

% ---- Diagramme TikZ : Schéma fonctionnel chaudière à condensation ----
\begin{figure}[H]
  \centering
  \begin{tikzpicture}[
      node distance=1.4cm and 2.0cm,
      bloc/.style={draw, rectangle, rounded corners=3pt, minimum width=2.4cm,
                   minimum height=0.9cm, align=center, font=\small\bfseries,
                   fill=#1!12, draw=#1, thick},
      arr/.style={-{Stealth[length=6pt]}, thick},
      lbl/.style={font=\footnotesize, midway}
    ]

    % Composants chaudière (colonne centrale)
    \node[bloc=FedBlue]  (bruleur) {Brûleur\\prémélange};
    \node[bloc=FedOrange, below=of bruleur] (echangeur) {Échangeur\\primaire};
    \node[bloc=FedGreen,  below=of echangeur] (echange2) {Échangeur\\condenseur};
    \node[bloc=FedRed,    below=of echange2]  (siphon)   {Siphon +\\neutralisateur};

    % Circuits extérieurs
    \node[font=\small, above=0.6cm of bruleur, text=FedBlue] (gaz) {Gaz naturel};
    \node[font=\small, right=2.5cm of echangeur] (dep) {Circuit\\chauffage\\(départ)};
    \node[font=\small, right=2.5cm of echange2]  (ret) {Circuit\\chauffage\\(retour $<57$\,°C)};
    \node[font=\small, below=0.5cm of siphon,  text=FedGreen] (cond) {Condensats neutralisés\\(pH 6,5–7) → évacuation};
    \node[font=\small, left=2.2cm of echange2, text=FedOrange] (fum) {Fumées condensées\\→ conduit PVC/inox};

    % Flèches internes
    \draw[arr, FedBlue]   (gaz)      -- (bruleur);
    \draw[arr, FedOrange] (bruleur)  -- node[lbl, right]{gaz chauds}  (echangeur);
    \draw[arr, FedOrange] (echangeur)-- node[lbl, right]{fumées $\downarrow T$} (echange2);
    \draw[arr, FedGreen]  (echange2) -- node[lbl, right]{condensats}  (siphon);
    \draw[arr, FedGreen]  (siphon)   -- (cond);

    % Circuit eau
    \draw[arr, FedOrange] (echangeur.east) -- (dep);
    \draw[arr, FedBlue]   (ret) -- (echange2.east);

    % Sortie fumées
    \draw[arr, FedOrange, dashed] (echange2.west) -- (fum);

    % Annotation point de rosée
    \draw[FedGreen, dashed, thick] (-1.5,-3.0) -- (3.5,-3.0)
      node[right, font=\footnotesize, text=FedGreen]{Point de rosée $\approx 57$\,°C};
  \end{tikzpicture}
  \caption{Schéma fonctionnel d'une chaudière à condensation (récupération de la chaleur latente des fumées)}
  \label{fig:chaudiere-condensation}
\end{figure}

\subsubsection{Chaudières biomasse (bois)}

\begin{table}[H]
  \centering
  \caption{Comparatif des chaudières bois — caractéristiques}
  \begin{tabular}{p{3cm}p{3.5cm}p{3.5cm}p{2.5cm}}
    \toprule
    Type & Combustible & Rendement & Automatisation \\
    \midrule
    Bûches & Bois bûches (H $\leq 20\,\%$) & 75–85\,\% & Manuelle \\
    Plaquettes & Plaquettes forestières (H 25–35\,\%) & 80–90\,\% & Automatique \\
    Granulés (pellets) & Granulés normés ENplus & 85–95\,\% & Entièrement automatique \\
    \bottomrule
  \end{tabular}
\end{table}

% ============================================================
\section{Pompes à chaleur (PAC)}
% ============================================================

\subsection{Principe thermodynamique}

Une pompe à chaleur transfère de la chaleur d'une source froide (air extérieur, sol, nappe phréatique) vers une source chaude (circuit de chauffage), en utilisant un apport d'énergie électrique.

\begin{definition}[title=Coefficient de performance (COP)]
\begin{equation}
  COP = \frac{\dot{Q}_{chaud}}{P_{élec}} = \frac{\dot{Q}_{froid} + P_{élec}}{P_{élec}}
  \label{eq:cop}
\end{equation}
\begin{itemize}
  \item $\dot{Q}_{chaud}$ : puissance thermique fournie au circuit de chauffage (\si{\kilo\watt})
  \item $P_{élec}$ : puissance électrique consommée (\si{\kilo\watt})
  \item $\dot{Q}_{froid}$ : puissance prélevée sur la source froide (\si{\kilo\watt})
\end{itemize}
Un COP = 3 signifie que pour \SI{1}{\kilo\watt} d'électricité consommée, la PAC produit \SI{3}{\kilo\watt} de chaleur.
\end{definition}

\begin{definition}[title=COP de Carnot (COP maximal théorique)]
\begin{equation}
  COP_{Carnot} = \frac{T_{chaud}}{T_{chaud} - T_{froid}} \qquad \text{(températures en Kelvin)}
  \label{eq:cop_carnot}
\end{equation}
Le COP réel est toujours inférieur au COP de Carnot (irréversibilités). Rapport typique : $COP_{réel} / COP_{Carnot} \approx 0{,}4$ à $0{,}6$.
\end{definition}

\subsection{Cycle frigorifique côté chaud}

Le cycle à compression de vapeur comporte 4 étapes :

\begin{enumerate}
  \item \textbf{Évaporation (source froide)} : le fluide frigorigène s'évapore en prélevant de la chaleur à la source froide. Pression basse, température basse.
  \item \textbf{Compression} : le compresseur élève la pression et la température du fluide (apport d'énergie électrique).
  \item \textbf{Condensation (source chaude)} : le fluide cède sa chaleur au circuit de chauffage en se condensant. \textbf{C'est ici que la chaleur utile est produite.}
  \item \textbf{Détente} : le fluide se détend dans le détendeur (électronique, thermostatique ou à tube capillaire), sa pression et sa température chutent. Retour à l'état initial.
\end{enumerate}

\begin{table}[H]
  \centering
  \caption{Fluides frigorigènes courants pour PAC résidentielle}
  \begin{tabular}{lcccl}
    \toprule
    Fluide & GWP & Pression basse (bar) & Pression haute (bar) & Statut \\
    \midrule
    R410A & 2088 & 4,9 & 24,3 & En cours d'élimination (F-Gas) \\
    R32 & 675 & 5,8 & 26,1 & Très répandu (PAC air/eau) \\
    R290 (propane) & 3 & 4,7 & 18,5 & GWP très faible, inflammable \\
    R454B & 466 & 5,4 & 24,8 & Successeur R410A \\
    \bottomrule
  \end{tabular}
\end{table}

\subsection{Types de PAC et performances}

\begin{table}[H]
  \centering
  \caption{Comparatif des types de PAC — performances et contextes d'usage}
  \begin{tabular}{p{2.5cm}p{3cm}p{2.5cm}p{4cm}}
    \toprule
    Type & Source froide & COP typique (A7/W35) & Contexte \\
    \midrule
    Air/Air & Air extérieur & 2,5–4,5 & Climatisation réversible \\
    Air/Eau & Air extérieur & 2,8–4,5 & Chauffage + ECS maison/collectif \\
    Eau glycolée/Eau & Sol horizontal & 3,5–5,5 & Maison individuelle, terrain disponible \\
    Eau/Eau & Nappe phréatique & 4,0–6,0 & Performances élevées, contraintes hydrogéol. \\
    PAC sur eaux grises & Eau usées bâtiment & 3,0–4,5 & Immeuble collectif, récupération énergie \\
    \bottomrule
  \end{tabular}
\end{table}

\begin{definition}[title=SCOP (Seasonal COP)]
Le SCOP est le COP moyen calculé sur l'ensemble de la saison de chauffage, selon la norme EN 14825. Il intègre les variations de température extérieure, les cycles de dégivrage (PAC air/eau) et les modes de relève.

\textbf{Valeurs indicatives (zone H2, émetteurs basse température) :}
\begin{itemize}
  \item PAC air/eau haute performance : SCOP 3,5–4,5
  \item PAC géothermique eau glycolée : SCOP 4,0–5,5
\end{itemize}

\textbf{Conditions normalisées d'essai EN 14825 :}
\begin{itemize}
  \item A7/W35 : air à 7\,°C, eau départ à 35\,°C (référence basse température)
  \item A2/W35 : air à 2\,°C, eau départ à 35\,°C (point intermédiaire)
  \item A-7/W35 : air à -7\,°C, eau départ à 35\,°C (froid hivernal)
\end{itemize}
\end{definition}

% ---- Diagramme TikZ : Cycle frigorifique PAC ----
\begin{figure}[H]
  \centering
  \begin{tikzpicture}[
      node distance=2.8cm,
      box/.style={draw, rectangle, rounded corners=4pt, minimum width=2.6cm,
                  minimum height=1.1cm, align=center, font=\small\bfseries,
                  fill=#1!15, draw=#1, thick},
      arr/.style={-{Stealth[length=7pt]}, thick, draw=#1}
    ]

    % Nœuds principaux
    \node[box=FedBlue]   (comp) {Compresseur\\{\footnotesize(élec. $P_{él}$)}};
    \node[box=FedOrange, right=of comp] (cond) {Condenseur\\{\footnotesize(source chaude)}};
    \node[box=FedGreen,  below=of cond] (dete) {Détendeur\\{\footnotesize(valve / cap.)}};
    \node[box=FedBlue,   below=of comp] (evap) {Évaporateur\\{\footnotesize(source froide)}};

    % Flèches fluide frigorigène
    \draw[arr=FedBlue]   (comp) -- node[above, font=\footnotesize]{HP vapeur} (cond);
    \draw[arr=FedOrange] (cond) -- node[right, font=\footnotesize]{HP liquide} (dete);
    \draw[arr=FedGreen]  (dete) -- node[below, font=\footnotesize]{BP liquide} (evap);
    \draw[arr=FedBlue]   (evap) -- node[left,  font=\footnotesize]{BP vapeur}  (comp);

    % Échanges énergétiques
    \draw[arr=FedOrange, dashed] (cond.north) -- ++(0, 0.9)
          node[above, font=\small, text=FedOrange]
          {$\dot{Q}_{chaud} = \dot{Q}_{froid} + P_{él}$};
    \draw[arr=FedBlue, dashed]   (evap.south) -- ++(0,-0.9)
          node[below, font=\small, text=FedBlue]
          {$\dot{Q}_{froid}$ prélevé sur source froide};
    \draw[arr=FedBlue, dashed]   (comp.west) -- ++(-1.1, 0)
          node[left, font=\small, text=FedBlue]{$P_{él}$};

    % Légende COP
    \node[draw=FedOrange, fill=FedOrange!10, rounded corners, font=\small,
          below=0.5cm of evap, xshift=1.4cm, text width=4.8cm, align=center]
      {$\displaystyle COP = \frac{\dot{Q}_{chaud}}{P_{él}} > 1$};
  \end{tikzpicture}
  \caption{Cycle frigorifique d'une pompe à chaleur (cycle à compression de vapeur)}
  \label{fig:cycle-pac}
\end{figure}

\subsection{Point de bivalence}

La puissance des PAC air/eau diminue lorsque la température extérieure baisse. En dessous d'une certaine température (point de bivalence), un générateur d'appoint (résistance électrique, chaudière) prend le relais.

\begin{methode}[title=Détermination du point de bivalence]
\begin{enumerate}
  \item Tracer la courbe de puissance de la PAC en fonction de $T_{ext}$ (fournie par le fabricant)
  \item Tracer la droite des déperditions du bâtiment : $\dot{Q}_{dep} = H_T \cdot (T_{int} - T_{ext})$
  \item L'intersection des deux courbes donne le \textbf{point de bivalence} ($T_{biv}$, $\dot{Q}_{biv}$)
  \item En dessous de $T_{biv}$, l'appoint assure la différence de puissance
\end{enumerate}
\textbf{Valeur courante :} $T_{biv}$ entre $-5$\,°C et $+5$\,°C selon le dimensionnement choisi.

Un point de bivalence bas (−5\,°C) signifie une PAC surdimensionnée qui couvre presque tous les besoins sans appoint : confort maximal, coût d'investissement élevé.
Un point de bivalence haut (+5\,°C) : PAC plus petite, appoint plus sollicité.
\end{methode}

% ============================================================
\section{Solaire thermique}
% ============================================================

\subsection{Principe et composants}

Le solaire thermique capte le rayonnement solaire pour produire de la chaleur, principalement pour l'eau chaude sanitaire (ECS) et en appoint de chauffage.

\begin{table}[H]
  \centering
  \caption{Types de capteurs solaires thermiques}
  \begin{tabular}{p{3cm}p{4cm}p{4.5cm}}
    \toprule
    Type & Principe & Performance / Température \\
    \midrule
    Plan vitré & Absorbeur plan sous vitrage & $\eta_0 \approx 0{,}75$, jusqu'à 80\,°C \\
    Tube sous vide & Tube en verre sous vide, pertes minimales & $\eta_0 \approx 0{,}70$, jusqu'à 120\,°C \\
    Concentrateur & Miroir parabolique, rayonnement focalisé & Températures très élevées, usage industriel \\
    \bottomrule
  \end{tabular}
\end{table}

\begin{definition}[title=Rendement d'un capteur solaire plan]
\begin{equation}
  \eta = \eta_0 - a_1 \cdot \frac{T_m - T_{ext}}{G} - a_2 \cdot \frac{(T_m - T_{ext})^2}{G}
\end{equation}
\begin{itemize}
  \item $\eta_0$ : rendement optique (sans pertes thermiques)
  \item $a_1$, $a_2$ : coefficients de pertes thermiques (fournis par le fabricant, certifiés Solar Keymark)
  \item $T_m$ : température moyenne du fluide dans le capteur (\si{\celsius})
  \item $G$ : irradiance solaire (\si{\watt\per\square\metre}) — en France : 150 à 900\,W/m²
\end{itemize}
\end{definition}

\begin{definition}[title=Taux de couverture solaire]
\begin{equation}
  f_{sol} = \frac{\text{énergie solaire utile annuelle}}{\text{besoins thermiques annuels}} \times 100 \qquad [\%]
\end{equation}
Pour une installation ECS résidentielle bien dimensionnée : $f_{sol} = 50$ à 70\,\%.
\end{definition}

\subsection{Systèmes solaires combinés (SSC)}

Un SSC couvre à la fois les besoins ECS et une partie du chauffage. La surface de capteurs est plus importante (20 à 40\,m² pour une maison). Le taux de couverture annuel global est typiquement de 25 à 40\,\% pour l'ensemble chauffage + ECS.

\begin{table}[H]
  \centering
  \caption{Surfaces de capteurs recommandées selon l'usage (France, zone H2)}
  \begin{tabular}{lccc}
    \toprule
    Usage & Personnes / Surface & Surface capteurs & Volume stockage \\
    \midrule
    ECS seule & 4 personnes & 3–4 m² & 200–300 L \\
    ECS seule & 6 personnes & 5–6 m² & 300–400 L \\
    SSC (ECS + chauffage) & Maison 100 m² RT2012 & 15–20 m² & 1000–2000 L \\
    SSC (ECS + chauffage) & Maison 150 m² RE2020 & 20–30 m² & 2000–3000 L \\
    \bottomrule
  \end{tabular}
\end{table}

% ============================================================
\section{Réseaux de chaleur urbains}
% ============================================================

\begin{definition}[title=Réseau de chaleur urbain]
Un réseau de chaleur distribue de l'eau chaude ou de la vapeur produite par une ou plusieurs chaufferies centralisées (cogénération, géothermie, biomasse, récupération chaleur industrielle) vers des sous-stations d'abonnés via un réseau de canalisations enterrées.

\textbf{Avantages :}
\begin{itemize}
  \item Mutualisation de la production, meilleur rendement global
  \item Facilite l'intégration des énergies renouvelables et de récupération (EnR\&R)
  \item Comptabilisé comme EnR dans la RE2020 si taux EnR\&R $\geq 50\,\%$
\end{itemize}
\end{definition}

\begin{definition}[title=Sous-station d'abonné]
La sous-station est le point de livraison de l'énergie thermique. Elle comprend un échangeur de chaleur qui sépare le réseau primaire (réseau urbain) du réseau secondaire (installation de l'abonné). Le comptage de l'énergie livrée est assuré par un compteur d'énergie thermique (calorimètre).
\end{definition}

% ============================================================
\section{Dimensionnement de la puissance de production}
% ============================================================

\begin{methode}[title=Méthode de calcul de la puissance installée]
\begin{enumerate}
  \item \textbf{Calculer les déperditions thermiques} du bâtiment à la température de base (selon la méthode NF EN 12831)
  \item \textbf{Ajouter les besoins ECS} : $\dot{Q}_{ECS} = \dot{m}_{ECS} \cdot c_p \cdot (T_{ECS} - T_{froide})$
  \item \textbf{Ajouter les besoins de ventilation} (préchauffage de l'air neuf en hiver)
  \item \textbf{Appliquer un coefficient de foisonnement} (0,7 à 0,9 pour le collectif)
  \item \textbf{Puissance installée} = somme des besoins / coefficient foisonnement
\end{enumerate}

\textbf{Règle de majoration courante :} majorer de 10 à 15\,\% pour tenir compte des pertes de distribution et des pointes.
\end{methode}

\begin{table}[H]
  \centering
  \caption{Coefficients de foisonnement ECS selon le type de bâtiment}
  \begin{tabular}{lcc}
    \toprule
    Type de bâtiment & Coefficient de foisonnement & Commentaire \\
    \midrule
    Maison individuelle & 1,0 & Pas de foisonnement \\
    Petit collectif ($\leq$ 10 logements) & 0,9 & Foisonnement faible \\
    Grand collectif (10–50 logements) & 0,7–0,8 & Foisonnement significatif \\
    Tertiaire (bureaux) & 0,6–0,7 & Pointes ECS faibles \\
    Hôtel & 0,5–0,7 & Forte diversité des besoins \\
    \bottomrule
  \end{tabular}
\end{table}

% ============================================================
\section{Exemples résolus}
% ============================================================

\begin{exemple}[title=Rendement d'une chaudière à condensation — PCI et PCS]
\textbf{Énoncé :} Une chaudière à condensation brûle \SI{2{,}5}{\cubic\metre\per\hour} de gaz naturel (PCI = \SI{10{,}55}{\kilo\watt\heure\per\cubic\metre}, PCS = \SI{11{,}7}{\kilo\watt\heure\per\cubic\metre}). Elle délivre une puissance de \SI{24}{\kilo\watt} au circuit de chauffage. Calculer le rendement PCI et le rendement PCS.

\textbf{Puissance chimique PCI :}
\[
  \dot{Q}_{PCI} = 2{,}5 \times 10{,}55 = \SI{26{,}375}{\kilo\watt}
\]

\textbf{Rendement PCI :}
\[
  \eta_{PCI} = \frac{24}{26{,}375} = 0{,}910 = \mathbf{91{,}0\,\%}
\]

\textbf{Puissance chimique PCS :}
\[
  \dot{Q}_{PCS} = 2{,}5 \times 11{,}7 = \SI{29{,}25}{\kilo\watt}
\]

\textbf{Rendement PCS :}
\[
  \eta_{PCS} = \frac{24}{29{,}25} = 0{,}820 = \mathbf{82{,}0\,\%}
\]

\textbf{Interprétation :} Le rendement PCS est inférieur au rendement PCI car le PCS intègre la chaleur latente de condensation (la chaudière la récupère mais le calcul PCS comptabilise plus d'énergie en entrée). Pour comparer les chaudières entre elles en Europe, on utilise le rendement PCS (directive Rendement). Un rendement PCS > 100\,\% est possible pour les chaudières à condensation car la récupération de la chaleur latente est comptabilisée en sortie.
\end{exemple}

\begin{exemple}[title=COP d'une PAC air/eau et comparaison économique]
\textbf{Énoncé :} Une PAC air/eau fonctionne avec : $T_{ext} = \SI{7}{\celsius}$ (A7), $T_{dep} = \SI{35}{\celsius}$ (W35), COP fabricant = 4,1. Puissance électrique absorbée : \SI{4}{\kilo\watt}.

\textit{1 — Puissance thermique produite :}
\[
  \dot{Q}_{chaud} = COP \times P_{élec} = 4{,}1 \times 4 = \SI{16{,}4}{\kilo\watt}
\]
\[
  \dot{Q}_{froid} = \dot{Q}_{chaud} - P_{élec} = 16{,}4 - 4 = \SI{12{,}4}{\kilo\watt} \text{ (prélevé sur l'air ext.)}
\]

\textit{2 — COP de Carnot et rendement exergétique :}
\[
  T_{chaud} = 35 + 273 = \SI{308}{\kelvin}, \quad T_{froid} = 7 + 273 = \SI{280}{\kelvin}
\]
\[
  COP_{Carnot} = \frac{308}{308 - 280} = \frac{308}{28} = 11{,}0
\]
\[
  \eta_{ex} = \frac{COP_{réel}}{COP_{Carnot}} = \frac{4{,}1}{11{,}0} = 0{,}373 = \mathbf{37{,}3\,\%}
\]

\textit{3 — Coût de production par kWh chaleur (prix gaz 0,12\,€/kWh, élec 0,22\,€/kWh) :}
\begin{itemize}
  \item \textbf{PAC :} $c_{PAC} = \frac{0{,}22}{4{,}1} = \SI{0{,}054}{\euro\per\kilo\watt\heure}_{th}$
  \item \textbf{Chaudière gaz ($\eta_{PCI} = 95\,\%$) :} $c_{gaz} = \frac{0{,}12}{0{,}95} = \SI{0{,}126}{\euro\per\kilo\watt\heure}_{th}$
\end{itemize}
La PAC est \textbf{2,3 fois moins chère} à l'usage dans ces conditions. L'écart se réduit si le prix de l'électricité augmente plus vite que le gaz.
\end{exemple}

\begin{exemple}[title=Dimensionnement de la puissance d'une chaudière collective]
\textbf{Énoncé :} Un immeuble collectif de 20 logements de 70\,m² chacun est situé en zone climatique H2 ($T_{base} = -7$\,°C, $T_{int} = 19$\,°C). Le coefficient de déperdition moyen est $G = \SI{1{,}2}{\watt\per\square\metre\per\kelvin}$. Les besoins ECS représentent \SI{8}{\kilo\watt} en pointe. Le coefficient de foisonnement est 0,8.

\textit{1 — Déperditions totales :}
\[
  S_{totale} = 20 \times 70 = \SI{1400}{\square\metre}
\]
\[
  \dot{Q}_{dep} = G \times S \times (T_{int} - T_{base}) = 1{,}2 \times 1400 \times (19 - (-7)) = \SI{43\,680}{\watt} \approx \SI{43{,}7}{\kilo\watt}
\]

\textit{2 — Puissance nécessaire (chauffage + ECS, avec foisonnement) :}
\[
  \dot{Q}_{total} = \frac{\dot{Q}_{dep} + \dot{Q}_{ECS}}{foisonnement} = \frac{43{,}7 + 8}{0{,}8} = \frac{51{,}7}{0{,}8} = \SI{64{,}6}{\kilo\watt}
\]

\textit{3 — Choix :} Le modèle A (\SI{60}{\kilo\watt}) est insuffisant (60 < 64,6). Le \textbf{modèle B (\SI{80}{\kilo\watt})} est retenu — marge de 24\,\%, suffisante sans surdimensionnement excessif.
\end{exemple}

\begin{exemple}[title=Analyse des performances d'une PAC en exploitation]
\textbf{Énoncé :} Une PAC air/eau est mesurée lors de sa mise en service :
\begin{itemize}
  \item Puissance électrique mesurée : \SI{3{,}8}{\kilo\watt}
  \item Température départ eau : \SI{45}{\celsius}, retour eau : \SI{40}{\celsius}
  \item Débit eau : \SI{1200}{\litre\per\hour}, $\rho_{eau} = \SI{992}{\kilogram\per\cubic\metre}$, $c_p = \SI{4179}{\joule\per\kilogram\per\kelvin}$
  \item Température extérieure : \SI{2}{\celsius}
\end{itemize}
Le fabricant annonce COP = 3,8 en A2/W45.

\textit{1 — Puissance thermique produite :}
\[
  \dot{m} = \frac{1200 \times 992}{3600 \times 1000} = \SI{0{,}331}{\kilogram\per\second}
\]
\[
  \dot{Q} = 0{,}331 \times 4179 \times 5 = \SI{6{,}92}{\kilo\watt}
\]

\textit{2 — COP mesuré :}
\[
  COP_{mes} = \frac{6{,}92}{3{,}8} = \mathbf{1{,}82}
\]

\textit{3 — Analyse :} Le COP mesuré (1,82) est très inférieur au COP annoncé (3,8 en A2/W45). Pistes de diagnostic :
\begin{itemize}
  \item Vérifier l'état du circuit frigorifique (charge en fluide frigorigène — sous-charge typique)
  \item Vérifier l'encrassement ou le givrage de l'échangeur air/fluide
  \item Vérifier si la PAC est en mode dégivrage lors de la mesure
  \item Vérifier le débit d'air sur l'évaporateur (filtre encrassé ou obstrué)
  \item Vérifier l'absence de recyclage d'air chaud vers l'aspiration
\end{itemize}
\end{exemple}

% ---- Diagramme TikZ : Comparaison COP PAC vs rendement chaudière ----
\begin{figure}[H]
  \centering
  \begin{tikzpicture}[x=1.6cm, y=0.9cm]
    % Axes
    \draw[-{Stealth}, thick] (0,0) -- (6.5,0)
      node[right, font=\small]{$T_{ext}$ (\si{\celsius})};
    \draw[-{Stealth}, thick] (0,0) -- (0,6.0)
      node[above, font=\small]{Performance (\%)};

    % Graduations axe x : -10, -5, 0, 5, 7, 10, 15
    \foreach \x/\lbl in {0.5/-10, 1.5/-5, 2.5/0, 3.5/5, 4.0/7, 4.5/10, 5.5/15}{
      \draw (\x,0.05) -- (\x,-0.05);
      \node[below, font=\footnotesize] at (\x,0) {\lbl};
    }

    % Graduations axe y : 100, 200, 300, 400, 500
    \foreach \y/\lbl in {1/100, 2/200, 3/300, 4/400, 5/500}{
      \draw (0.05,\y) -- (-0.05,\y);
      \node[left, font=\footnotesize] at (0,\y) {\lbl\,\%};
    }

    % Courbe PAC air/eau — COP×100 (linearly increases with T_ext)
    % T_ext: -10→COP 2.2, -5→2.6, 0→3.1, 5→3.6, 7→4.1, 10→4.4, 15→5.0
    \draw[FedBlue, very thick]
      (0.5,2.2) -- (1.5,2.6) -- (2.5,3.1) -- (3.5,3.6) -- (4.0,4.1) -- (4.5,4.4) -- (5.5,5.0);
    \node[FedBlue, font=\small\bfseries, right] at (5.5,5.0) {PAC air/eau (COP×100)};

    % Ligne chaudière condensation ≈ 105 % PCI = constant
    \draw[FedOrange, very thick, dashed] (0.5,1.05) -- (5.5,1.05);
    \node[FedOrange, font=\small\bfseries, right] at (5.5,1.05)
      {Chaudière condensation ($\eta_{PCI}\approx105\,\%$)};

    % Ligne chaudière basse température ≈ 90 %
    \draw[FedRed, thick, dotted] (0.5,0.90) -- (5.5,0.90);
    \node[FedRed, font=\small, right] at (5.5,0.90)
      {Chaudière basse T° ($\eta\approx90\,\%$)};

    % Annotation point de référence A7/W35
    \draw[FedBlue, dashed] (4.0,0) -- (4.0,4.1);
    \node[FedBlue, font=\footnotesize, above] at (4.0,4.1) {A7/W35};

    % Titre
    \node[font=\small\itshape, above] at (3.0,5.8)
      {Performances comparées en fonction de la température extérieure};
  \end{tikzpicture}
  \caption{Comparaison COP d'une PAC air/eau et rendement d'une chaudière en fonction de la température extérieure}
  \label{fig:comparaison-cop-rendement}
\end{figure}

% ============================================================
\section{Éléments de réglementation}
% ============================================================

\begin{description}
  \item[RE2020 (décret du 29 juillet 2021)] Interdit les chaudières à gaz comme système principal dans les logements neufs depuis le 1\textsuperscript{er} janvier 2022. Favorise les PAC (coefficient $b_{EF} = 2{,}3$ pour l'électricité) et les réseaux de chaleur à EnR\&R. Fixe un seuil d'émissions de GES (Ic\textsubscript{énergie}).

  \item[NF EN 12831] Méthode de calcul des déperditions thermiques de bâtiment — base du dimensionnement des générateurs de chaleur.

  \item[NF EN 14825] Méthode de calcul du SCOP des pompes à chaleur — conditions d'essai normalisées (A7/W35, A2/W35, etc.) pour la comparaison des équipements.

  \item[Règlement EU 813/2013 (écoconception chaudières)] Impose un rendement énergétique saisonnier minimal pour les chaudières de chauffage (classe ErP). Depuis 2015, seules les chaudières à condensation ou haute performance sont autorisées dans l'UE.

  \item[Règlement EU 811/2013 (étiquetage énergie)] Étiquette énergie de A+++ à G pour les générateurs de chaleur — obligatoire sur tout équipement vendu.

  \item[Arrêté du 23 juin 1978 (modifié)] Installations fixes destinées au chauffage et à l'alimentation en eau chaude sanitaire. Règles de sécurité : pression max, température max, soupapes, vases d'expansion.

  \item[DTU 61.1 — NF P45-204] Installations de gaz dans les bâtiments — règles de sécurité, raccordements, ventilation des locaux contenant des appareils à gaz.

  \item[Règlement (UE) 517/2014 (F-Gas)] Réglementation sur les gaz fluorés à fort effet de serre. Interdit progressivement les fluides frigorigènes à GWP élevé (R410A interdit dans les nouveaux équipements à partir de 2025). Impose une attestation de capacité (certificat fluides frigorigènes) pour la manipulation.

  \item[MaPrimeRénov' et CEE] Dispositifs d'aide à la rénovation énergétique : bonus pour les PAC (SCOP $\geq 3{,}5$), bonus pour les chaudières biomasse. Les CEE (Certificats d'Économies d'Énergie) financent la rénovation via les obligations des fournisseurs d'énergie.
\end{description}

\begin{tcolorbox}[
  enhanced, colback=FedBlue!3, colframe=FedBlue!50,
  fonttitle=\bfseries\sffamily, coltitle=FedBlue,
  title={FICHE MÉMO — Production de chaleur}, boxrule=0.8pt, arc=2pt,
  width=\textwidth
]
\textbf{Combustion — Chaudières gaz}
\[
  e = \frac{V_{air\,réel} - V_{air\,stœchio}}{V_{air\,stœchio}} \times 100 \quad [\%]
  \qquad \text{excès d'air (brûleur atm. : 20–40\,\%, prémélange : 5–15\,\%)}
\]
\[
  \eta_{comb} = \frac{\dot{Q}_{utile}}{\dot{Q}_{PCI}} \quad [\%]
  \qquad \text{condensation effective si } T_{retour} < 57\,\text{°C (point de rosée gaz)}
\]

\textbf{Pompe à chaleur (PAC) — Cycle à compression}
\[
  COP = \frac{\dot{Q}_{chaud}}{P_{élec}} \qquad \text{avec } \dot{Q}_{chaud} = \dot{Q}_{froid} + P_{élec} \quad [\si{\kilo\watt}]
\]
\[
  COP_{Carnot} = \frac{T_{chaud}}{T_{chaud} - T_{froid}} \quad \text{(températures en \si{\kelvin})}
  \qquad COP_{réel} \approx 0{,}4 \text{ à } 0{,}6 \times COP_{Carnot}
\]
\[
  SCOP = \frac{E_{thermique\,annuelle}}{E_{électrique\,annuelle}}
  \qquad \text{(NF EN 14825 — intègre variations saisonnières et dégivrage)}
\]

\textbf{Solaire thermique}
\[
  \eta = \eta_0 - a_1 \cdot \frac{T_m - T_{ext}}{G} - a_2 \cdot \frac{(T_m - T_{ext})^2}{G}
  \qquad \text{avec } G = \text{irradiance} [\si{\watt\per\square\metre}]
\]
\[
  f_{sol} = \frac{E_{solaire\,utile}}{E_{besoin\,total}} \times 100 \quad [\%]
  \qquad \text{ECS résidentielle bien dimensionnée : } f_{sol} = 50 \text{ à } 70\,\%
\]

\textbf{Dimensionnement de la puissance}
\[
  \dot{Q}_{ECS} = \dot{m}_{ECS} \cdot c_p \cdot (T_{ECS} - T_{froide}) \quad [\si{\kilo\watt}]
  \qquad \text{déperditions : } \dot{Q}_{dep} = H_T \cdot (T_{int} - T_{ext})
\]
\end{tcolorbox}

\finchapitre

% \chapter{Distribution et émission de chaleur}

% ============================================================
\section{Objectifs pédagogiques}
% ============================================================

\subsection*{Compétences GCF mobilisées}
\begin{itemize}
  \item \textbf{C2-1 à C2-3} — Analyser un réseau de distribution et ses émetteurs : topologie, chaîne d'énergie, paramètres de fonctionnement
  \item \textbf{C3-3} — Dimensionner les émetteurs (puissance, sélection sur catalogue) et les réseaux de distribution (bitubes, monotube)
  \item \textbf{C4-1 et C4-2} — Élaborer les schémas hydrauliques et les plans d'implantation des émetteurs
  \item \textbf{C7-4 à C7-6} — Mesurer les températures, débits et puissances en mise en service ; interpréter les écarts
  \item \textbf{C8-1 à C8-4} — Vérifier le confort thermique, diagnostiquer les déséquilibres, proposer des optimisations
\end{itemize}

\subsection*{Savoirs associés mobilisés}
\begin{itemize}
  \item \textbf{B1-2} — Émission de chaleur : radiateurs, planchers chauffants, ventilo-convecteurs, aérothermes (niveau 3)
  \item \textbf{B4-1} — Réseaux hydrauliques : eau chaude, équilibrage, sécurité, dimensionnement (niveau 3)
  \item \textbf{A2-3} — Bilan thermique : déperditions, charges thermiques (niveau 3)
  \item \textbf{S5.5} — Réglementations fluides thermiques : chauffage (niveau 3)
  \item \textbf{S10.3 et S10.5} — Sécurités d'installation et mise en service (niveau 3)
\end{itemize}

% ============================================================
\section{Introduction}
% ============================================================

La distribution et l'émission de chaleur constituent le maillon final entre la production (chaudière, PAC) et le confort des occupants. Un émetteur bien dimensionné et correctement alimenté garantit la température de consigne dans chaque local, avec un minimum de consommation énergétique.

En GCF, le technicien supérieur est amené à concevoir les réseaux de distribution (tracé, diamètres, topologie), à sélectionner les émetteurs sur catalogue, à vérifier l'équilibrage hydraulique et à intervenir en cas de dysfonctionnement. Ce chapitre couvre les émetteurs statiques (radiateurs, planchers chauffants) et dynamiques (ventilo-convecteurs), ainsi que les architectures de réseaux bitubes et monotubes.

% ============================================================
\section{Topologies des réseaux de distribution}
% ============================================================

\subsection{Réseau bitube}

\begin{definition}[title=Réseau bitube]
Dans un réseau bitube, chaque émetteur est alimenté par une canalisation de départ et une canalisation de retour \textbf{indépendantes}. L'eau chaude arrive à la même température à tous les émetteurs.

\textbf{Avantages :}
\begin{itemize}
  \item Température d'alimentation identique pour tous les émetteurs
  \item Facilité d'équilibrage (robinets thermostatiques + robinets d'équilibrage)
  \item Standard en construction neuve tertiaire et résidentiel collectif
\end{itemize}

\textbf{Inconvénients :} Coût de tuyauterie plus élevé qu'un réseau monotube.
\end{definition}

\subsection{Réseau monotube}

\begin{definition}[title=Réseau monotube]
Dans un réseau monotube, les émetteurs sont raccordés en série sur une boucle unique. L'eau se refroidit progressivement en traversant chaque émetteur.

\textbf{Conséquences :}
\begin{itemize}
  \item La température d'alimentation \textbf{diminue} d'un émetteur au suivant
  \item Les émetteurs en aval doivent être surdimensionnés pour compenser la perte de température
  \item Utilisé principalement en réhabilitation ou petits bâtiments
\end{itemize}
\end{definition}

\begin{table}[H]
  \centering
  \caption{Comparatif réseau bitube vs monotube}
  \begin{tabular}{lll}
    \toprule
    Critère & Bitube & Monotube \\
    \midrule
    Température d'alimentation émetteurs & Identique pour tous & Décroissante en série \\
    Dimensionnement émetteurs & Homogène & Variable (croissant en aval) \\
    Coût tuyauterie & Plus élevé & Plus faible \\
    Facilité d'équilibrage & Bonne & Difficile \\
    Régulation locale & Robinet thermostatique & Bypass nécessaire \\
    Usage typique & Neuf, tertiaire, collectif & Réhabilitation, petit pavillon \\
    \bottomrule
  \end{tabular}
\end{table}

\subsection{Réseau en antenne vs réseau en anneau}

\begin{itemize}
  \item \textbf{Réseau en antenne (arborescent)} : la plus répandue. Distribution depuis une nourrice centrale vers des branches. Simple à équilibrer avec la méthode de la branche index.
  \item \textbf{Réseau en anneau (boucle)} : tolérance aux pannes, utilisé dans les grands réseaux collectifs. Équilibrage plus complexe.
  \item \textbf{Réseau en nourrice} : toutes les branches partent du même collecteur, longueurs quasi-égales. Équilibrage facilité. Standard pour les planchers chauffants.
\end{itemize}

% ============================================================
\section{Radiateurs à eau chaude}
% ============================================================

\subsection{Principe et types}

Un radiateur transfère la chaleur de l'eau chaude vers l'air ambiant par convection (majoritaire, 70--80\,\%) et rayonnement (20--30\,\%). Les types courants sont :

\begin{table}[H]
  \centering
  \caption{Types de radiateurs à eau chaude}
  \begin{tabular}{p{3cm}p{4.5cm}p{4cm}}
    \toprule
    Type & Construction & Usage typique \\
    \midrule
    Panneau acier & Tôles embouties, ailettes. 1, 2 ou 3 panneaux + ailettes. & Résidentiel, tertiaire courant \\
    Fonte & Éléments assemblés. Grande inertie thermique. & Réhabilitation, bâtiments anciens \\
    Aluminium & Éléments extrudés. Faible inertie, montée en T rapide. & Résidentiel neuf \\
    Sèche-serviettes & Radiateur tubulaire + résistance électrique d'appoint. & Salles de bain \\
    \bottomrule
  \end{tabular}
\end{table}

\subsection{Puissance nominale et correction}

La puissance des radiateurs est donnée dans les catalogues fabricants pour des conditions normalisées :
\begin{itemize}
  \item Température d'eau : départ 75\,°C / retour 65\,°C (moyenne 70\,°C)
  \item Température ambiante : 20\,°C
  \item Excès de température : $\Delta T_{nom} = 70 - 20 = \SI{50}{\kelvin}$
\end{itemize}

En conditions réelles (températures différentes), il faut corriger la puissance :

\begin{equation}
  \Phi_{réel} = \Phi_{nom} \cdot \left(\frac{\Delta T_{réel}}{\Delta T_{nom}}\right)^n
  \label{eq:radiateur_correction}
\end{equation}
\begin{itemize}
  \item $\Delta T_{réel} = T_{moy,eau} - T_{ambiante}$ : écart réel (\si{\kelvin})
  \item $\Delta T_{nom} = \SI{50}{\kelvin}$ : écart nominal
  \item $n$ : exposant du radiateur — typiquement $n = 1{,}3$ pour un panneau acier, $n = 1{,}25$ pour un radiateur fonte
\end{itemize}

\begin{attention}
Avec des émetteurs basse température (50/40\,°C, compatible PAC), $\Delta T_{réel} = 45 - 20 = \SI{25}{\kelvin}$. Le radiateur ne délivre que :
\[
  \Phi_{50/40} = \Phi_{nom} \cdot \left(\frac{25}{50}\right)^{1{,}3} = \Phi_{nom} \times 0{,}406
\]
Un radiateur prévu pour 75/65\,°C devra être surdimensionné d'un facteur $\approx 2{,}5$ pour fonctionner à basse température. C'est un point critique lors du passage à une PAC.
\end{attention}

\subsection{Sélection d'un radiateur}

\begin{methode}[title=Sélection d'un radiateur sur catalogue]
\begin{enumerate}
  \item Calculer la puissance de déperdition du local à couvrir $\dot{Q}_{local}$ (W)
  \item Déterminer les températures de fonctionnement ($T_{dep}$, $T_{ret}$, $T_{amb}$)
  \item Calculer $\Delta T_{réel}$
  \item Calculer la puissance nominale équivalente nécessaire :
  \[
    \Phi_{nom,néc} = \frac{\dot{Q}_{local}}{\left(\dfrac{\Delta T_{réel}}{50}\right)^n}
  \]
  \item Sélectionner dans le catalogue un radiateur de puissance nominale $\geq \Phi_{nom,néc}$
  \item Vérifier l'encombrement et la position (sous fenêtre recommandé pour le confort)
\end{enumerate}
\end{methode}

\subsection{Robinets thermostatiques}

Le robinet thermostatique régule le débit d'eau dans le radiateur en fonction de la température ambiante. Il comporte une tête thermostatique à dilatation de cire ou de liquide, réglable de 1 à 5 (correspondant à environ 10--30\,°C).

\begin{definition}[title=Kvs du robinet thermostatique]
Le $K_{vs}$ (ou $K_v$ à ouverture maximale) est le débit en m³/h d'eau à travers le robinet pour une perte de charge de \SI{1}{\bar}. Il caractérise la capacité de débit du robinet.
\begin{equation}
  q_v = K_{vs} \cdot \sqrt{\Delta P} \qquad \left[q_v \text{ en m}^3/\text{h}, \Delta P \text{ en bar}\right]
\end{equation}
\end{definition}

% ============================================================
\section{Planchers chauffants hydrauliques}
% ============================================================

\subsection{Principe}

Le plancher chauffant hydraulique (PCH) est un émetteur à très basse température (35/28\,°C typique) intégré dans la dalle. Il offre un confort thermique excellent (chaleur rayonnante au sol) et est parfaitement adapté aux PAC.

\begin{definition}[title=Température de surface maximale réglementaire]
La NF EN 1264 limite la température de surface du plancher à :
\begin{itemize}
  \item \textbf{30\,°C} dans les zones habitées (pièces principales)
  \item \textbf{35\,°C} dans les zones de bordure (bande périmétrique de 1\,m)
  \item \textbf{33\,°C} dans les salles de bain
\end{itemize}
Au-delà, le confort est dégradé et le risque d'endommagement du revêtement de sol augmente.
\end{definition}

\subsection{Dimensionnement d'un plancher chauffant}

\begin{methode}[title=Calcul du pas de tube (NF EN 1264)]
La puissance surfacique $q$ d'un plancher chauffant dépend de la température moyenne de l'eau ($T_{moy}$), de la température ambiante ($T_{amb}$) et du pas $p$ entre les tubes :

\begin{equation}
  q = B \cdot \frac{(T_{moy} - T_{amb})}{p} \qquad \left[\si{\watt\per\square\metre}\right]
\end{equation}

Où $B$ est un coefficient dépendant du revêtement de sol (résistance thermique $R_{revêt}$) :
\begin{equation}
  B = \frac{1}{0{,}173 + R_{revêt}} \qquad [\si{\watt\per\square\metre\per\kelvin}]
\end{equation}

\textbf{Étapes de dimensionnement :}
\begin{enumerate}
  \item Calculer la puissance surfacique nécessaire : $q_{néc} = \dot{Q}_{local} / S_{local}$
  \item Choisir le revêtement de sol (carrelage : $R = 0{,}02$\,m²K/W ; parquet : $R = 0{,}10$\,m²K/W)
  \item Fixer les températures ($T_{dep}$, $T_{ret}$, $T_{moy}$, $T_{amb}$)
  \item Calculer $B$ puis le pas $p = B \cdot (T_{moy} - T_{amb}) / q_{néc}$
  \item Arrondir à un pas normalisé : 75, 100, 150, 200, 250, 300\,mm
  \item Vérifier que la température de surface ne dépasse pas la limite
\end{enumerate}
\end{methode}

\begin{attention}
Sous un parquet massif ($R_{revêt} > 0{,}15$\,m²K/W), la puissance surfacique maximale chute fortement. Vérifier la compatibilité du revêtement avant de valider le dimensionnement (label « Compatible plancher chauffant »).
\end{attention}

\subsection{Collecteur de distribution}

Le collecteur distribue l'eau chaude vers les différentes boucles du plancher. Il comporte :
\begin{itemize}
  \item Des débitmètres par boucle (en verre ou numériques) pour l'équilibrage
  \item Des têtes thermoélectriques (actionneurs) commandées par les thermostats de pièce
  \item Un organe de sécurité (soupape) et un purgeur
\end{itemize}

\textbf{Longueur maximale d'une boucle :} 100 à 120\,m pour limiter les pertes de charge (pression différentielle au collecteur typique : 15 à 30\,kPa).

% ============================================================
\section{Ventilo-convecteurs}
% ============================================================

\begin{definition}[title=Ventilo-convecteur (fan-coil)]
Un ventilo-convecteur est un émetteur dynamique : il comporte un échangeur à eau (batterie chaude ou froide) et un ventilateur qui force la circulation d'air à travers la batterie. Il peut fonctionner en chaud et en froid (2 tubes) ou alternativement chaud et froid (4 tubes).
\end{definition}

\begin{table}[H]
  \centering
  \caption{Comparatif ventilo-convecteurs 2 tubes vs 4 tubes}
  \begin{tabular}{lll}
    \toprule
    Critère & 2 tubes & 4 tubes \\
    \midrule
    Nombre de circuits hydrauliques & 1 (chaud \textbf{ou} froid) & 2 (chaud \textbf{et} froid) \\
    Chauffage et froid simultanés & Non & Oui \\
    Coût d'installation & Faible & Élevé \\
    Flexibilité & Limitée (saison) & Totale \\
    Usage typique & Logement collectif & Hôtel, bureaux, tertiaire \\
    \bottomrule
  \end{tabular}
\end{table}

\textbf{Puissances typiques :} de 0,5 à 5\,kW en chauffage selon le modèle. Vitesse du ventilateur réglable (silencieux / moyen / fort).

% ============================================================
\section{Aérothermes}
% ============================================================

\begin{definition}[title=Aérotherme]
Un aérotherme est un grand ventilo-convecteur à soufflage d'air chaud, utilisé pour le chauffage de grands volumes (ateliers, entrepôts, halls). Il est alimenté en eau chaude et dispose d'un ventilateur hélicoïdal de grande taille.
\end{definition}

\textbf{Puissances typiques :} 10 à 100\,kW. Température de soufflage : 40 à 60\,°C. Portée du jet d'air : 10 à 30\,m.

\begin{attention}
L'aérotherme crée des mouvements d'air importants, ce qui peut causer de la poussière en suspension. À éviter dans les locaux propres (agro-alimentaire, médical) et à combiner avec une étude acoustique dans les locaux de travail (directive bruit au travail).
\end{attention}

% ============================================================
\section{Exemples résolus}
% ============================================================

\begin{exemple}[title=Sélection d'un radiateur pour fonctionnement PAC]
\textbf{Énoncé :} Un séjour de 30\,m² présente des déperditions de \SI{1400}{\watt} pour $T_{ext} = -7$\,°C, $T_{amb} = 20$\,°C. La PAC fonctionne en 45/38\,°C. Sélectionner le radiateur (exposant $n = 1{,}3$).

\textbf{Étape 1 : Écart de température réel}
\[
  T_{moy,eau} = \frac{45 + 38}{2} = \SI{41{,}5}{\celsius}
\]
\[
  \Delta T_{réel} = 41{,}5 - 20 = \SI{21{,}5}{\kelvin}
\]

\textbf{Étape 2 : Puissance nominale équivalente nécessaire}
\[
  \Phi_{nom,néc} = \frac{1400}{\left(\dfrac{21{,}5}{50}\right)^{1{,}3}} = \frac{1400}{(0{,}430)^{1{,}3}} = \frac{1400}{0{,}371} = \SI{3774}{\watt}
\]

\textbf{Étape 3 : Sélection catalogue}
Retenir un radiateur panneau acier type 22 (2 panneaux + 2 rangées d'ailettes) de puissance nominale $\geq$ \SI{3800}{\watt}.

\textbf{Exemple :} Radiateur 22/600/1600 → puissance nominale 3936\,W (à vérifier catalogue fabricant). ✅

\textbf{À noter :} Le même local avec des radiateurs 75/65\,°C n'aurait nécessité qu'un radiateur de \SI{1400}{\watt} nominal — le passage PAC multiplie par 2,5 la puissance nominale requise.
\end{exemple}

\begin{exemple}[title=Dimensionnement d'une boucle de plancher chauffant]
\textbf{Énoncé :} Une chambre de 12\,m² présente des déperditions de \SI{480}{\watt}. Le revêtement est du carrelage ($R_{revêt} = \SI{0{,}01}{\square\metre\kelvin\per\watt}$). Le PCH fonctionne en 35/28\,°C, $T_{amb} = 20$\,°C.

\textbf{Étape 1 : Puissance surfacique nécessaire}
\[
  q_{néc} = \frac{480}{12} = \SI{40}{\watt\per\square\metre}
\]

\textbf{Étape 2 : Coefficient B}
\[
  B = \frac{1}{0{,}173 + 0{,}01} = \frac{1}{0{,}183} = \SI{5{,}46}{\watt\per\square\metre\per\kelvin}
\]

\textbf{Étape 3 : Température moyenne eau}
\[
  T_{moy} = \frac{35 + 28}{2} = \SI{31{,}5}{\celsius}
\]

\textbf{Étape 4 : Pas de tube}
\[
  p = \frac{B \cdot (T_{moy} - T_{amb})}{q_{néc}} = \frac{5{,}46 \times (31{,}5 - 20)}{40} = \frac{5{,}46 \times 11{,}5}{40} = \frac{62{,}8}{40} = \SI{1{,}57}{\metre}
\]

\begin{attention}
Un pas de 1,57\,m est irréaliste ! Recalculons : $B \cdot (T_{moy} - T_{amb}) = 5{,}46 \times 11{,}5 = 62{,}8$\,W/m². La puissance surfacique disponible (62,8\,W/m²) est déjà supérieure à la puissance nécessaire (40\,W/m²). Choisir le pas maximal : \textbf{300\,mm}.
\end{attention}

\textbf{Vérification avec $p = 0{,}30$\,m :}
\[
  q_{réel} = B \cdot \frac{(T_{moy} - T_{amb})}{p} = 5{,}46 \times \frac{11{,}5}{0{,}30} = \SI{209}{\watt\per\square\metre}
\]
Mais cette formule simplifiée n'est valable que pour $q < q_{max}$ (environ 100\,W/m²). En pratique, avec $p = 300$\,mm et ces températures, la puissance surfacique est bien suffisante (40\,W/m² << 100\,W/m² maxi).

\textbf{Longueur de tube pour la chambre :}
\[
  L_{tube} = \frac{S_{chambre}}{p} = \frac{12}{0{,}30} = \SI{40}{\metre}
\]
Acceptable (< 100\,m). Une seule boucle suffit.
\end{exemple}

% ============================================================
\section{Exercices d'application}
% ============================================================

\exercice{Réseau bitube — sélection des radiateurs d'un appartement}

Un appartement de 4 pièces doit être équipé de radiateurs à eau chaude 70/55\,°C ($T_{amb} = 20$\,°C, exposant $n = 1{,}3$). Les déperditions par pièce sont :

\begin{center}
\begin{tabular}{lcc}
  \toprule
  Pièce & Déperditions (W) & Surface (m²) \\
  \midrule
  Séjour & 1600 & 25 \\
  Chambre 1 & 800 & 12 \\
  Chambre 2 & 700 & 10 \\
  Cuisine & 500 & 8 \\
  \bottomrule
\end{tabular}
\end{center}

\begin{enumerate}
  \item Calculer la puissance nominale nécessaire pour chaque pièce.
  \item Calculer la puissance totale à fournir. Vérifier la cohérence avec le dimensionnement de la chaudière.
\end{enumerate}

\textbf{Corrigé :}

\textit{Question 1 — Écart de température réel :}
\[
  T_{moy} = \frac{70 + 55}{2} = \SI{62{,}5}{\celsius}, \quad \Delta T_{réel} = 62{,}5 - 20 = \SI{42{,}5}{\kelvin}
\]
\[
  \text{Facteur correction} = \left(\frac{42{,}5}{50}\right)^{1{,}3} = (0{,}85)^{1{,}3} = 0{,}826
\]

\begin{center}
\begin{tabular}{lcc}
  \toprule
  Pièce & Déperditions (W) & $\Phi_{nom,néc}$ = Dép. / 0,826 (W) \\
  \midrule
  Séjour & 1600 & 1937 \\
  Chambre 1 & 800 & 969 \\
  Chambre 2 & 700 & 848 \\
  Cuisine & 500 & 605 \\
  \midrule
  \textbf{Total} & \textbf{3600} & \textbf{4359} \\
  \bottomrule
\end{tabular}
\end{center}

\textit{Question 2 :} La puissance totale de déperditions est 3600\,W. La chaudière devra fournir au minimum 3,6\,kW, majoré des pertes de distribution (5–10\,\%) → chaudière $\geq$ 4\,kW.

\exercice{Plancher chauffant — dimensionnement complet d'une pièce}

Un salon de 20\,m² présente des déperditions de \SI{900}{\watt}. Revêtement : parquet flottant compatible PCH ($R_{revêt} = \SI{0{,}10}{\square\metre\kelvin\per\watt}$). Le PCH fonctionne en 40/32\,°C, $T_{amb} = 20$\,°C.

\begin{enumerate}
  \item Calculer la puissance surfacique nécessaire.
  \item Calculer le coefficient $B$ et la puissance surfacique disponible avec $p = 150$\,mm.
  \item Calculer la longueur de tube nécessaire.
  \item Vérifier si une seule boucle suffit.
\end{enumerate}

\textbf{Corrigé :}

\textit{Question 1 :}
\[
  q_{néc} = \frac{900}{20} = \SI{45}{\watt\per\square\metre}
\]

\textit{Question 2 :}
\[
  B = \frac{1}{0{,}173 + 0{,}10} = \frac{1}{0{,}273} = \SI{3{,}66}{\watt\per\square\metre\per\kelvin}
\]
\[
  T_{moy} = \frac{40 + 32}{2} = \SI{36}{\celsius}, \quad \Delta T = 36 - 20 = \SI{16}{\kelvin}
\]
\[
  q_{disp} = B \times \frac{\Delta T}{p} = 3{,}66 \times \frac{16}{0{,}15} = \SI{390}{\watt\per\square\metre}
\]
La puissance disponible (390\,W/m²) est très largement supérieure à la puissance nécessaire (45\,W/m²). Même avec $p = 300$\,mm, $q = 195$\,W/m² >> 45\,W/m². Choisir $p = \SI{300}{\milli\metre}$.

\textit{Question 3 — Longueur de tube :}
\[
  L_{tube} = \frac{S}{p} = \frac{20}{0{,}30} = \SI{67}{\metre}
\]

\textit{Question 4 :} 67\,m < 100\,m → une seule boucle suffit. ✅

\exercice{Diagnostic d'un radiateur sous-performant}

Un technicien mesure sur un radiateur de bureau les valeurs suivantes :
\begin{itemize}
  \item Température entrée eau : 68\,°C
  \item Température sortie eau : 58\,°C
  \item Débit eau mesuré : 40\,L/h
  \item Température ambiante : 18\,°C
\end{itemize}
Le local présente des déperditions de 1200\,W. Le radiateur a une puissance nominale de 1800\,W (en 75/65/20\,°C, $n = 1{,}3$).

\begin{enumerate}
  \item Calculer la puissance thermique réellement délivrée.
  \item Calculer la puissance théorique attendue dans ces conditions.
  \item Analyser l'écart et proposer des pistes de diagnostic.
\end{enumerate}

\textbf{Corrigé :}

\textit{Question 1 — Puissance réelle mesurée :}
\[
  \dot{m} = \frac{40}{3600} \times 1000 \times \frac{1}{1000} = \SI{0{,}0111}{\kilogram\per\second}
\]
\[
  \dot{Q}_{réel} = \dot{m} \times c_p \times \Delta T = 0{,}0111 \times 4186 \times (68 - 58) = \SI{465}{\watt}
\]

\textit{Question 2 — Puissance théorique attendue :}
\[
  T_{moy} = \frac{68 + 58}{2} = \SI{63}{\celsius}, \quad \Delta T_{réel} = 63 - 18 = \SI{45}{\kelvin}
\]
\[
  \dot{Q}_{théo} = 1800 \times \left(\frac{45}{50}\right)^{1{,}3} = 1800 \times (0{,}90)^{1{,}3} = 1800 \times 0{,}870 = \SI{1566}{\watt}
\]

\textit{Question 3 — Analyse :} Écart massif : 465\,W mesuré vs 1566\,W théorique (−70\,\%). Pistes de diagnostic :
\begin{itemize}
  \item Débit très faible (40\,L/h vs débit nominal probablement 100–150\,L/h) → robinet thermostatique bloqué ou partiellement fermé
  \item Présence d'air dans le radiateur → purge nécessaire
  \item Encrassement intérieur (boues) réduisant le débit et l'échange thermique
  \item Vanne d'équilibrage trop fermée → mesurer la pression différentielle
\end{itemize}

% ============================================================
\section{Éléments de réglementation}
% ============================================================

\begin{description}
  \item[NF EN 1264] Systèmes de chauffage et de rafraîchissement intégrés dans les planchers — exigences de calcul, dimensionnement et essais. Norme de référence pour les planchers chauffants hydrauliques. Fixe les températures de surface maximales.

  \item[NF EN 442] Radiateurs et convecteurs — spécifications et essais techniques. Définit les conditions normalisées de mesure des puissances (75/65/20\,°C) et les méthodes d'essai.

  \item[NF EN 15316] Systèmes de chauffage dans les bâtiments — méthode de calcul des besoins en énergie et des rendements du système (distribution, émission, régulation). Utilisée pour le calcul réglementaire RE2020.

  \item[RE 2020] Favorise les émetteurs basse température (planchers chauffants, radiateurs basse T) pour optimiser le SCOP des PAC. Les logements neufs sont dimensionnés pour fonctionner à 45/40\,°C maximum.

  \item[DTU 65.14 — NF P52-307] Exécution des planchers chauffants à eau chaude. Prescrit les règles de pose, les distances minimales aux murs, les joints de dilatation et les conditions de mise en eau.

  \item[DTU 65.11 — NF P52-211] Dispositifs de sécurité pour installations de chauffage central. S'applique aux réseaux de distribution : soupapes de sécurité, vases d'expansion, pression maximale de service.
\end{description}

% \chapter{Climatisation et traitement d'air}

% ============================================================
\section{Objectifs pédagogiques}
% ============================================================

\subsection*{Compétences GCF mobilisées}
\begin{itemize}
  \item \textbf{C2-1 à C2-3} — Analyser un système de climatisation : cycle frigorifique, architecture, paramètres de fonctionnement
  \item \textbf{C3-2 et C3-3} — Comparer et sélectionner un système de climatisation selon les critères technico-économiques et environnementaux ; dimensionner la puissance frigorifique
  \item \textbf{C5-1 à C5-3} — Appliquer les réglementations fluides frigorigènes et les normes de sécurité
  \item \textbf{C7-4 à C7-6} — Mesurer les performances en mise en service (EER, COP), interpréter les écarts
  \item \textbf{C8-1 à C8-4} — Vérifier les performances saisonnières, diagnostiquer les dérives
\end{itemize}

\subsection*{Savoirs associés mobilisés}
\begin{itemize}
  \item \textbf{A5} — Thermodynamique appliquée — machine frigorifique : fonctionnel, COP, impact environnemental (niveau 2–3)
  \item \textbf{A4} — Traitement d'air : diagramme air humide, évolutions (niveau 3)
  \item \textbf{B2-2} — Climatisation : systèmes thermodynamiques, eau, air — CTA, VRV (niveau 3)
  \item \textbf{B5-5} — Systèmes de rejet de chaleur : condenseur air/eau, tour aéroréfrigérante (niveau 3)
  \item \textbf{S5.5} — Réglementations fluides frigorigènes (niveau 2)
  \item \textbf{S10.5 et S10.6} — Mise en service et critères de bon fonctionnement (niveau 3)
\end{itemize}

% ============================================================
\section{Introduction}
% ============================================================

La climatisation désigne l'ensemble des techniques permettant de maintenir les conditions de température et d'humidité dans un local, quelle que soit la saison. En été, il s'agit de refroidir et souvent de déshumidifier ; en hiver, dans les bâtiments à fort apport interne (bureaux, commerces, data centers), de refroidir malgré les basses températures extérieures.

En France, la consommation liée à la climatisation représente environ 10\,\% de la consommation électrique des bâtiments tertiaires et croît régulièrement avec le réchauffement climatique. Le technicien supérieur GCF dimensionne les systèmes, sélectionne les fluides frigorigènes compatibles avec la réglementation F-Gaz, et vérifie les performances à la mise en service.

% ============================================================
\section{Cycle frigorifique à compression de vapeur}
% ============================================================

\subsection{Les quatre composants du cycle}

Le cycle frigorifique à compression de vapeur est le fondement de tous les systèmes de climatisation. Il repose sur quatre composants principaux :

\begin{table}[H]
  \centering
  \caption{Composants du cycle frigorifique et leur rôle}
  \begin{tabular}{p{3cm}p{4.5cm}p{4cm}}
    \toprule
    Composant & Transformation & Effet sur le fluide \\
    \midrule
    Évaporateur & Échange à basse pression & Vaporisation : absorbe chaleur (effet froid) \\
    Compresseur & Compression adiabatique & Pression et température augmentent \\
    Condenseur & Échange à haute pression & Condensation : cède chaleur (rejet) \\
    Détendeur & Détente isenthalpique & Pression et température chutent \\
    \bottomrule
  \end{tabular}
\end{table}

\subsection{Diagramme enthalpique (diagramme de Mollier frigorifique)}

Le diagramme pression-enthalpie ($\log P$ / $h$) représente le cycle frigorifique. Il permet de calculer tous les bilans énergétiques du cycle.

\begin{definition}[title=Effet frigorifique]
\begin{equation}
  q_0 = h_1 - h_4 \qquad \left[\si{\kilo\joule\per\kilogram}\right]
  \label{eq:effet_frigo}
\end{equation}
\begin{itemize}
  \item $h_1$ : enthalpie à la sortie de l'évaporateur (vapeur saturée sèche)
  \item $h_4$ : enthalpie à l'entrée de l'évaporateur (liquide sous-refroidi après détente)
\end{itemize}
\end{definition}

\begin{definition}[title=Travail de compression]
\begin{equation}
  w_c = h_2 - h_1 \qquad \left[\si{\kilo\joule\per\kilogram}\right]
  \label{eq:travail_compresseur}
\end{equation}
\begin{itemize}
  \item $h_2$ : enthalpie à la sortie du compresseur (vapeur surchauffée)
\end{itemize}
\end{definition}

\begin{definition}[title=Effet de condensation]
\begin{equation}
  q_c = h_2 - h_3 \qquad \left[\si{\kilo\joule\per\kilogram}\right]
  \label{eq:effet_condenseur}
\end{equation}
\begin{itemize}
  \item $h_3$ : enthalpie à la sortie du condenseur (liquide sous-refroidi)
\end{itemize}
\end{definition}

\begin{definition}[title=Bilan énergétique du cycle]
Le premier principe appliqué au cycle donne :
\begin{equation}
  q_c = q_0 + w_c
  \label{eq:bilan_cycle}
\end{equation}
Le condenseur rejette la somme de l'effet frigorifique et du travail de compression.
\end{definition}

\subsection{Coefficient de performance en froid (EER)}

\begin{definition}[title=EER — Energy Efficiency Ratio]
\begin{equation}
  EER = \frac{\dot{Q}_0}{P_{élec}} = \frac{q_0}{w_c} \qquad \left[\text{sans dimension}\right]
  \label{eq:eer}
\end{equation}
\begin{itemize}
  \item $\dot{Q}_0$ : puissance frigorifique (\si{\kilo\watt})
  \item $P_{élec}$ : puissance électrique absorbée par le compresseur (\si{\kilo\watt})
\end{itemize}

\textbf{Valeurs typiques :}
\begin{itemize}
  \item Split-system résidentiel : EER = 2,5 à 4,5
  \item Groupe froid à eau glacée haute efficacité : EER = 4 à 7
  \item EER de Carnot (limite théorique) : $EER_{Carnot} = T_{froid}/(T_{chaud}-T_{froid})$
\end{itemize}
\end{definition}

\begin{definition}[title=SEER — Seasonal EER]
Le SEER est l'EER moyen calculé sur toute la saison de refroidissement (norme EN 14825). Il est obligatoirement affiché sur l'étiquette énergie EU 206/2012.

\textbf{Valeurs de référence RE2020 :} SEER $\geq$ 3,5 recommandé pour les bâtiments tertiaires.
\end{definition}

\subsection{Les fluides frigorigènes}

\begin{table}[H]
  \centering
  \caption{Fluides frigorigènes courants en climatisation}
  \begin{tabular}{p{2cm}p{2.5cm}p{1.5cm}p{1.5cm}p{4cm}}
    \toprule
    Fluide & Famille & GWP & ODP & Statut réglementaire \\
    \midrule
    R410A & HFC & 2088 & 0 & Phase-out F-Gaz 2025 (neuf) \\
    R32 & HFC & 675 & 0 & Autorisé, transition \\
    R290 (propane) & HC & 3 & 0 & Autorisé, charges limitées \\
    R744 (CO\textsubscript{2}) & naturel & 1 & 0 & Autorisé, haute pression \\
    R1234yf & HFO & 4 & 0 & Remplaçant R134a \\
    R454B & HFO/HFC & 466 & 0 & Remplaçant R410A \\
    \bottomrule
  \end{tabular}
\end{table}

\begin{attention}
Le règlement F-Gaz EU 517/2014 (révisé 2024) prévoit un calendrier de réduction des HFC à fort GWP. Depuis 2025, les nouveaux équipements de climatisation résidentielle ne peuvent utiliser de fluides avec GWP $>$ 750 dans l'UE. Le technicien GCF doit impérativement vérifier la conformité du fluide avant toute installation neuve.

\textbf{GWP} = Global Warming Potential (potentiel de réchauffement global, référence CO\textsubscript{2} = 1).\\
\textbf{ODP} = Ozone Depleting Potential (potentiel de destruction de l'ozone, référence R11 = 1).
\end{attention}

% ============================================================
\section{Systèmes de climatisation}
% ============================================================

\subsection{Systèmes split et multi-split}

\begin{definition}[title=Système split]
Un système split comporte une unité extérieure (condenseur + compresseur) et une unité intérieure (évaporateur + ventilateur), reliées par des liaisons frigorifiques (cuivre). Il est destiné au confort d'un seul local.
\end{definition}

\begin{table}[H]
  \centering
  \caption{Types d'unités intérieures split}
  \begin{tabular}{p{3cm}p{5cm}p{4cm}}
    \toprule
    Type & Installation & Usage \\
    \midrule
    Mural (high-wall) & Fixé en hauteur sur un mur & Résidentiel, petit bureau \\
    Cassette plafonnière & Encastrée dans le faux-plafond, soufflage 4 voies & Tertiaire \\
    Console (floor) & Posée au sol ou en allège & Salles de réunion, hôtels \\
    Gainable & Dans le plénum, réseau de gaines & Grande surface, discrétion \\
    \bottomrule
  \end{tabular}
\end{table}

\begin{definition}[title=Système multi-split]
Un multi-split connecte plusieurs unités intérieures à une seule unité extérieure. Les unités fonctionnent indépendamment (régulation individuelle) mais partagent le même circuit frigorifique.

\textbf{Limite :} toutes les unités fonctionnent en chaud ou en froid simultanément (pas de simultanéité chaud/froid). Pour la simultanéité, il faut un système VRV/VRF.
\end{definition}

\subsection{Systèmes VRV / VRF}

\begin{definition}[title=VRV / VRF (Volume de Réfrigérant Variable)]
Un système VRV (marque Daikin) ou VRF est une évolution du multi-split qui permet de faire varier le débit de fluide frigorigène vers chaque unité intérieure, et d'assurer du \textbf{chauffage et du refroidissement simultanés} grâce à un module de récupération de chaleur (Heat Recovery).
\end{definition}

\textbf{Architecture typique :}
\begin{itemize}
  \item 1 à 3 unités extérieures (compresseurs à vitesse variable INVERTER)
  \item Jusqu'à 64 unités intérieures par système
  \item Module BS (Branch Selector) pour la récupération de chaleur simultanée
  \item Réseau de cuivre de petit diamètre (pas besoin de réseau eau)
\end{itemize}

\textbf{Avantages :} flexibilité, SEER élevé (4 à 8), pas de réseau eau, contrôle individuel par local.\\
\textbf{Inconvénients :} maintenance frigoriste qualifiée obligatoire, longueur des liaisons limitée (100–200\,m), charge en fluide frigorigène élevée (enjeu F-Gaz).

\subsection{Climatisation à eau glacée}

\begin{definition}[title=Système à eau glacée (Chilled Water)]
Dans un système à eau glacée, un groupe froid produit de l'eau refroidie (typiquement 6–7\,°C, retour 12\,°C) distribuée par un réseau hydraulique vers des unités terminales (ventilo-convecteurs, CTA à batterie froide, poutres froides).
\end{definition}

\textbf{Composants principaux :}
\begin{itemize}
  \item \textbf{Groupe froid (chiller)} : compresseur(s) + évaporateur (côté eau) + condenseur (air ou eau)
  \item \textbf{Réseau de distribution eau glacée} : pompes, équilibrage, isolation thermique
  \item \textbf{Unités terminales} : ventilo-convecteurs 4 tubes, CTA, poutres froides actives
  \item \textbf{Tour aéroréfrigérante ou condenseur à air} : rejet de chaleur vers l'extérieur
\end{itemize}

\begin{table}[H]
  \centering
  \caption{Comparatif groupes froids selon type de condenseur}
  \begin{tabular}{p{3.5cm}p{4cm}p{4cm}}
    \toprule
    Type condenseur & EER typique & Contraintes \\
    \midrule
    Air (aircooled) & 2,8 à 4,5 & Simple, pas d'eau, bruit \\
    Eau (watercooled) & 4,5 à 7,0 & Meilleur EER, tour aéroréfrig. ou puits \\
    Eau à condensation géothermique & 5,0 à 8,0 & Très haut EER, contrainte hydrogéologique \\
    \bottomrule
  \end{tabular}
\end{table}

\subsection{Rejet de chaleur}

\begin{definition}[title=Tour aéroréfrigérante (TAR)]
La tour aéroréfrigérante refroidit l'eau de condensation par évaporation partielle. L'eau chaude (condenseur) ruisselle sur un garnissage, un ventilateur force le passage d'air. Une partie de l'eau s'évapore, refroidissant le reste.

\textbf{Avantage :} température de condensation proche de la température humide extérieure (plus basse que la température sèche) → meilleur EER.\\
\textbf{Contrainte majeure :} risque Légionella — entretien et traitement de l'eau réglementés (arrêté du 14 décembre 2013, plan de gestion du risque légionelle).
\end{definition}

% ============================================================
\section{Calcul des charges thermiques de refroidissement}
% ============================================================

\begin{definition}[title=Charge thermique de refroidissement]
La charge thermique de refroidissement $\dot{Q}_{froid}$ est la puissance que le système de climatisation doit extraire pour maintenir la température de consigne. Elle comprend :
\begin{itemize}
  \item \textbf{Apports solaires} par les vitrages ($Q_{sol}$)
  \item \textbf{Apports internes} : éclairage, équipements, personnes ($Q_{int}$)
  \item \textbf{Apports par transmission} des parois extérieures ($Q_{trans}$)
  \item \textbf{Apports par renouvellement d'air} (air neuf chaud en été) ($Q_{air}$)
\end{itemize}
\end{definition}

\begin{methode}[title=Calcul simplifié des apports internes]
\begin{align}
  Q_{personnes} &= n_{pers} \times q_{pers} & q_{pers} &\approx \SI{70}{\watt} \text{ (bureau)} \\
  Q_{éclairage} &= P_{inst} \times S_{local} & P_{inst} &\approx \SI{10}{\watt\per\square\metre} \text{ (LED)} \\
  Q_{informatique} &= P_{PC} \times n_{PC} & P_{PC} &\approx \SI{150}{\watt\per\poste}
\end{align}
\end{methode}

\begin{definition}[title=Apport solaire par vitrage]
\begin{equation}
  Q_{sol} = S_{vit} \times G_{sol} \times g \times F_{ombr}
\end{equation}
\begin{itemize}
  \item $S_{vit}$ : surface vitrée (\si{\square\metre})
  \item $G_{sol}$ : irradiance solaire sur la paroi (\si{\watt\per\square\metre}) — dépend de l'orientation et de la saison
  \item $g$ : facteur solaire du vitrage (sans dimension) — typiquement 0,2 à 0,6
  \item $F_{ombr}$ : facteur d'ombrage (stores, débords, etc.) — 0,3 à 1,0
\end{itemize}
\end{definition}

% ============================================================
\section{Exemples résolus}
% ============================================================

\begin{exemple}[title=Calcul de l'EER d'un groupe froid sur diagramme enthalpique]
\textbf{Énoncé :} Un groupe froid au R32 fonctionne dans les conditions suivantes (lues sur le diagramme enthalpique R32) :
\begin{itemize}
  \item $h_1 = \SI{503}{\kilo\joule\per\kilogram}$ (sortie évaporateur)
  \item $h_2 = \SI{550}{\kilo\joule\per\kilogram}$ (sortie compresseur)
  \item $h_3 = \SI{280}{\kilo\joule\per\kilogram}$ (sortie condenseur)
  \item $h_4 = h_3 = \SI{280}{\kilo\joule\per\kilogram}$ (sortie détendeur, isenthalpique)
\end{itemize}
Le débit massique de fluide est $\dot{m} = \SI{0{,}12}{\kilogram\per\second}$.

\textbf{Effet frigorifique massique :}
\[
  q_0 = h_1 - h_4 = 503 - 280 = \SI{223}{\kilo\joule\per\kilogram}
\]

\textbf{Travail de compression massique :}
\[
  w_c = h_2 - h_1 = 550 - 503 = \SI{47}{\kilo\joule\per\kilogram}
\]

\textbf{Vérification du bilan :}
\[
  q_c = h_2 - h_3 = 550 - 280 = \SI{270}{\kilo\joule\per\kilogram} = q_0 + w_c = 223 + 47 ✓
\]

\textbf{Puissances :}
\[
  \dot{Q}_0 = \dot{m} \times q_0 = 0{,}12 \times 223 = \SI{26{,}76}{\kilo\watt}
\]
\[
  P_{élec} = \dot{m} \times w_c = 0{,}12 \times 47 = \SI{5{,}64}{\kilo\watt}
\]
\[
  \dot{Q}_c = \dot{m} \times q_c = 0{,}12 \times 270 = \SI{32{,}4}{\kilo\watt}
\]

\textbf{EER :}
\[
  EER = \frac{\dot{Q}_0}{P_{élec}} = \frac{26{,}76}{5{,}64} = \mathbf{4{,}75}
\]

\textbf{EER de Carnot} (évaporation à 5\,°C, condensation à 40\,°C) :
\[
  EER_{Carnot} = \frac{278}{313 - 278} = \frac{278}{35} = 7{,}94
\]
\[
  \eta_{ex} = \frac{4{,}75}{7{,}94} = 0{,}598 = \mathbf{59{,}8\,\%}
\]
\end{exemple}

\begin{exemple}[title=Dimensionnement d'une climatisation de salle de réunion]
\textbf{Énoncé :} Une salle de réunion de 40\,m² (hauteur 2,7\,m) est orientée sud. Elle accueille 20 personnes. Données :
\begin{itemize}
  \item Vitrage sud : 8\,m², $g = 0{,}35$, $F_{ombr} = 0{,}5$, $G_{sol} = \SI{500}{\watt\per\square\metre}$
  \item Éclairage LED : 10\,W/m²
  \item Ordinateurs : 10 postes à 150\,W chacun
  \item Personnes : 80\,W/pers (réunion, activité modérée)
  \item Apports par transmission et ventilation : \SI{1200}{\watt} (estimé)
\end{itemize}
Calculer la puissance frigorifique nécessaire et sélectionner un système.

\textbf{Apports solaires :}
\[
  Q_{sol} = 8 \times 500 \times 0{,}35 \times 0{,}5 = \SI{700}{\watt}
\]

\textbf{Apports internes :}
\[
  Q_{pers} = 20 \times 80 = \SI{1600}{\watt}
\]
\[
  Q_{écl} = 10 \times 40 = \SI{400}{\watt}
\]
\[
  Q_{info} = 10 \times 150 = \SI{1500}{\watt}
\]

\textbf{Total des charges :}
\[
  \dot{Q}_{froid} = 700 + 1600 + 400 + 1500 + 1200 = \SI{5400}{\watt} = \SI{5{,}4}{\kilo\watt}
\]

\textbf{Sélection :} Système split mural ou cassette 4 voies de 6\,kW nominal (arrondi au calibre supérieur disponible, ex. : Daikin FTXM60R ou équivalent). Vérifier le SEER $\geq$ 3,5 pour conformité RE2020 tertiaire.
\end{exemple}

% ============================================================
\section{Exercices d'application}
% ============================================================

\exercice{Analyse du cycle frigorifique R454B}

Un climatiseur au R454B présente les enthalpies suivantes mesurées au banc d'essai :
$h_1 = \SI{490}{\kilo\joule\per\kilogram}$, $h_2 = \SI{535}{\kilo\joule\per\kilogram}$, $h_3 = \SI{270}{\kilo\joule\per\kilogram}$, $h_4 = h_3$.

Le débit massique mesuré est $\dot{m} = \SI{0{,}08}{\kilogram\per\second}$.

\begin{enumerate}
  \item Calculer les puissances frigorifique, de compression et de condensation.
  \item Vérifier le bilan énergétique.
  \item Calculer l'EER. Comparer avec l'EER de Carnot ($T_{0} = 7$\,°C, $T_{c} = 45$\,°C).
  \item Le fabricant annonce EER = 4,1 dans ces conditions. Analyser l'écart.
\end{enumerate}

\textbf{Corrigé :}

\textit{Question 1 — Puissances :}
\[
  q_0 = 490 - 270 = \SI{220}{\kilo\joule\per\kilogram}
\]
\[
  w_c = 535 - 490 = \SI{45}{\kilo\joule\per\kilogram}
\]
\[
  q_c = 535 - 270 = \SI{265}{\kilo\joule\per\kilogram}
\]
\[
  \dot{Q}_0 = 0{,}08 \times 220 = \SI{17{,}6}{\kilo\watt}, \quad P_{élec} = 0{,}08 \times 45 = \SI{3{,}6}{\kilo\watt}, \quad \dot{Q}_c = 0{,}08 \times 265 = \SI{21{,}2}{\kilo\watt}
\]

\textit{Question 2 — Bilan :} $q_0 + w_c = 220 + 45 = 265 = q_c$ ✓

\textit{Question 3 — EER :}
\[
  EER_{mes} = \frac{17{,}6}{3{,}6} = 4{,}89
\]
\[
  EER_{Carnot} = \frac{280}{318 - 280} = \frac{280}{38} = 7{,}37
\]
\[
  \eta_{ex} = \frac{4{,}89}{7{,}37} = 66{,}4\,\%
\]

\textit{Question 4 — Écart EER :} EER mesuré (4,89) > EER annoncé (4,1). L'écart peut s'expliquer par : conditions d'essai différentes des conditions normalisées EN 14825, instrumentation de mesure du débit massique, incertitudes de mesure des enthalpies. L'EER annoncé intègre les pertes auxiliaires (ventilateurs, régulation) que le calcul enthalpique n'inclut pas.

\exercice{Charges thermiques d'un open-space et sélection d'un système}

Un open-space de 200\,m² est situé en façade ouest. Données estivales :
\begin{itemize}
  \item Vitrage ouest : 30\,m², $g = 0{,}25$ (vitrage solaire), $F_{ombr} = 0{,}6$, $G_{sol,ouest} = \SI{450}{\watt\per\square\metre}$
  \item 40 personnes, 80\,W/pers
  \item Éclairage : 8\,W/m²
  \item Informatique : 40 postes à 120\,W
  \item Transmission et ventilation : \SI{3000}{\watt}
\end{itemize}

\begin{enumerate}
  \item Calculer chaque type d'apport et la puissance frigorifique totale.
  \item Proposer deux solutions techniques (une à fluide frigorigène, une à eau glacée) en justifiant le choix.
  \item Calculer la consommation électrique mensuelle du système à eau glacée si EER = 4,0 et si le système fonctionne 8\,h/jour, 22\,jours/mois à plein régime.
\end{enumerate}

\textbf{Corrigé :}

\textit{Question 1 — Apports :}
\[
  Q_{sol} = 30 \times 450 \times 0{,}25 \times 0{,}6 = \SI{2025}{\watt}
\]
\[
  Q_{pers} = 40 \times 80 = \SI{3200}{\watt}
\]
\[
  Q_{écl} = 8 \times 200 = \SI{1600}{\watt}
\]
\[
  Q_{info} = 40 \times 120 = \SI{4800}{\watt}
\]
\[
  \dot{Q}_{froid} = 2025 + 3200 + 1600 + 4800 + 3000 = \SI{14\,625}{\watt} \approx \SI{14{,}6}{\kilo\watt}
\]

\textit{Question 2 — Solutions :}
\begin{itemize}
  \item \textbf{VRV/VRF (fluide frigorigène)} : adapté si bâtiment multi-zones, flexibilité élevée, pas besoin de réseau eau glacée. Recommandé pour réhabilitation ou bâtiment de taille moyenne.
  \item \textbf{Eau glacée + ventilo-convecteurs 4 tubes} : adapté si bâtiment neuf de grande taille avec besoins chaud et froid simultanés, meilleur EER, maintenance centralisée. Recommandé pour immeuble tertiaire neuf.
\end{itemize}

\textit{Question 3 — Consommation électrique mensuelle :}
\[
  P_{élec} = \frac{\dot{Q}_{froid}}{EER} = \frac{14{,}6}{4{,}0} = \SI{3{,}65}{\kilo\watt}
\]
\[
  E_{mensuelle} = 3{,}65 \times 8 \times 22 = \SI{642}{\kilo\watt\heure}
\]

\exercice{Diagnostic d'un groupe froid sous-performant}

Un groupe froid à eau glacée de 80\,kW nominal présente les mesures suivantes lors de la mise en service :
\begin{itemize}
  \item Puissance électrique mesurée : \SI{28}{\kilo\watt}
  \item Température eau glacée départ : 9\,°C (consigne : 7\,°C)
  \item Température eau glacée retour : 13\,°C
  \item Débit eau glacée : \SI{30}{\cubic\metre\per\hour}, $\rho = \SI{999}{\kilogram\per\cubic\metre}$, $c_p = \SI{4182}{\joule\per\kilogram\per\kelvin}$
  \item Température extérieure : 28\,°C
\end{itemize}

\begin{enumerate}
  \item Calculer la puissance frigorifique produite.
  \item Calculer l'EER mesuré.
  \item Identifier les anomalies et les causes probables.
\end{enumerate}

\textbf{Corrigé :}

\textit{Question 1 — Puissance frigorifique :}
\[
  \dot{m} = \frac{30 \times 999}{3600} = \SI{8{,}33}{\kilogram\per\second}
\]
\[
  \dot{Q}_0 = 8{,}33 \times 4182 \times (13 - 9) = \SI{139\,333}{\watt} \approx \SI{139}{\kilo\watt}
\]

\textit{Question 2 — EER mesuré :}
\[
  EER = \frac{139}{28} = \mathbf{4{,}96}
\]

\textit{Question 3 — Anomalies :}
\begin{itemize}
  \item La puissance produite (139\,kW) est \textbf{supérieure} à la puissance nominale (80\,kW) → le groupe est en surcharge ou les conditions sont plus favorables que prévues (température extérieure plus basse, ou nominal mesuré à des conditions défavorables)
  \item La température de départ (9\,°C) est \textbf{au-dessus} de la consigne (7\,°C) → vérifier la régulation, le sous-refroidissement insuffisant, ou un débit trop élevé
  \item Piste : recalibrer la vanne de régulation du débit ; vérifier la charge en fluide frigorigène ; vérifier la consigne du thermostat de départ
\end{itemize}

% ============================================================
\section{Éléments de réglementation}
% ============================================================

\begin{description}
  \item[Règlement F-Gaz EU 517/2014 (révisé 2024)] Calendrier de réduction des HFC. Depuis 2025 : interdiction des fluides GWP $>$ 750 dans les nouvelles unités de climatisation résidentielle divisée. Obligations de contrôle d'étanchéité selon la charge en CO\textsubscript{2} équivalent.

  \item[Arrêté du 14 décembre 2013 (tours aéroréfrigérantes)] Prescriptions techniques relatives aux systèmes de refroidissement par dispersion d'eau dans un flux d'air. Plan de gestion du risque légionelles, analyses bactériologiques périodiques, traitement chimique de l'eau.

  \item[Attestation de capacité (décret 2015-1790)] Toute entreprise qui installe, entretient ou répare des équipements contenant des fluides frigorigènes HFC doit disposer d'une attestation de capacité délivrée par un organisme agréé.

  \item[NF EN 14825] Méthode de calcul du SEER et SCOP des climatiseurs et pompes à chaleur — conditions d'essai normalisées pour la comparaison des équipements.

  \item[Règlement EU 206/2012 (écoconception et étiquetage)] Exigences de rendement minimal et étiquetage énergie pour les climatiseurs résidentiels. Classes de A+++ à D.

  \item[RE 2020] Fixe un seuil de consommation en énergie primaire incluant la climatisation. Impose une réflexion sur les protections solaires passives (casquettes, stores) avant de dimensionner la climatisation.

  \item[NF EN 15251] Critères d'ambiance intérieure été : température opérative maximale selon catégorie (I : 25\,°C, II : 26\,°C, III : 27\,°C) pour les bâtiments non refroidis mécaniquement.
\end{description}

% % ============================================================
%  Chapitre 07 — Régulation et GTB
%  BTS FED — Option GCF
%  Savoirs mobilisés : B11 (niveau 3), B8-1 (niveau 3),
%                      D1-2 (niveau 2), C6, C2, C8
% ============================================================

\chapter{Régulation et GTB}

% ------------------------------------------------------------
\section{Objectifs pédagogiques}
% ------------------------------------------------------------

À l'issue de ce chapitre, l'étudiant sera capable de :

\begin{itemize}
  \item Expliquer le principe d'une boucle de régulation en boucle fermée (savoir B11, niveau 3).
  \item Distinguer les actions TOR, proportionnelle, intégrale et dérivée (PID) et choisir la loi adaptée à une application CVC.
  \item Analyser et paramétrer une régulation par loi d'eau ou par cascade sur une installation de chauffage.
  \item Décrire l'architecture d'une GTB/GTC, identifier les niveaux (terrain, automatisme, supervision).
  \item Identifier les protocoles de communication BACnet/IP et KNX utilisés dans le bâtiment tertiaire (savoir D1-2, niveau 2).
  \item Mettre en œuvre les outils de pilotage d'un système CVC (compétence C6).
  \item Vérifier et adapter les performances d'un système régulé (compétence C8).
\end{itemize}

\textbf{Savoirs associés du référentiel :}
\begin{itemize}
  \item B11 — Régulation (boucles, TOR, proportionnel, optimisation) — niveau 3
  \item B8-1 — Architecture systèmes centralisés (capteurs, actionneurs, supervision, protocoles) — niveau 3
  \item D1-2 — Standards de communication (protocoles, compatibilités) — niveau 2
  \item C2-2 — Identifier les chaînes d'énergie et d'information — niveau 3
\end{itemize}

% ------------------------------------------------------------
\section{Introduction}
% ------------------------------------------------------------

Dans une installation CVC (Chauffage, Ventilation, Climatisation), maintenir le confort des occupants tout en minimisant la consommation d'énergie exige un pilotage dynamique et précis des équipements. La \textbf{régulation} assure en temps réel l'ajustement automatique des grandeurs physiques (température, débit, pression, hygrométrie) autour de valeurs consignes, en réponse aux perturbations extérieures (température extérieure, apports solaires, occupation variable). La \textbf{gestion technique du bâtiment (GTB)}, également appelée GTC (gestion technique centralisée), fédère l'ensemble des régulations locales, des comptages énergétiques et des alarmes au sein d'une infrastructure numérique unifiée. Elle constitue aujourd'hui un levier incontournable pour atteindre les objectifs de performance énergétique fixés par la RE2020, et pour satisfaire les exigences de reporting et de commissionnement des bâtiments tertiaires soumis au décret BACS (Building Automation and Control Systems, 2020).

% ------------------------------------------------------------
\section{Bases de la régulation automatique}
% ------------------------------------------------------------

\subsection{Grandeurs fondamentales d'une boucle}

\begin{definition}
  Une \textbf{boucle de régulation} est un système automatique qui compare en permanence la valeur mesurée d'une grandeur physique (\textbf{mesure} $y$) à une valeur de référence (\textbf{consigne} $w$), puis agit sur un \textbf{actionneur} pour réduire l'\textbf{écart} (erreur) $e = w - y$.
\end{definition}

Les composants types d'une boucle fermée sont :

\begin{center}
\begin{tabular}{ll}
  \toprule
  Composant & Rôle \\
  \midrule
  Capteur / sonde & Mesure la grandeur réglée ($y$) \\
  Régulateur (ou automate) & Calcule la commande à partir de $e$ \\
  Actionneur & Agit sur le processus (vanne, variateur\ldots) \\
  Processus & Installation physique à réguler \\
  Perturbations & Variations non maîtrisées (température ext., occupation\ldots) \\
  \bottomrule
\end{tabular}
\end{center}

% --- Schéma A — Boucle de régulation fermée ---
\begin{figure}[H]
\centering
\begin{tikzpicture}[
  >=Stealth,
  bloc/.style={
    draw=btsprimary, fill=btssecondary,
    rounded corners=4pt,
    minimum width=2.6cm, minimum height=0.9cm,
    text centered, font=\small
  },
  comparateur/.style={
    draw=btsprimary, fill=white,
    circle, minimum size=0.8cm,
    font=\small\bfseries
  },
  fl/.style={->, thick},
  lbl/.style={font=\footnotesize\itshape}
]

% --- Noeuds chaîne directe (ligne du haut) ---
\node[comparateur]           (comp) at (0,0)    {$\Sigma$};
\node[bloc] (pid)  at (3.2,0)   {Correcteur\\(PID)};
\node[bloc] (act)  at (6.8,0)   {Actionneur\\(vanne/VFD)};
\node[bloc] (proc) at (10.8,0)  {Processus\\(installation)};

% Sortie
\node[font=\small] (sortie) at (13.6,0) {$y$};

% Consigne
\node[font=\small] (cons) at (-2.0,0) {Consigne $w$};

% Capteur (chaîne de retour)
\node[bloc] (capt) at (5.4,-2.2) {Capteur\\(sonde T°)};

% Perturbation
\node[font=\footnotesize\itshape, color=red!70!black] (pert) at (10.8,1.8) {Perturbation};

% --- Flèches chaîne directe ---
\draw[fl] (cons)  -- (comp);
\draw[fl] (comp)  -- node[above, lbl]{Erreur $e$} (pid);
\draw[fl] (pid)   -- node[above, lbl]{Signal $u$} (act);
\draw[fl] (act)   -- (proc);
\draw[fl] (proc)  -- node[above, lbl]{$y$} (sortie);

% --- Retour via capteur ---
% Depuis la sortie vers le bas, puis vers la gauche via capteur, puis vers le comp
\draw[fl] (sortie.south) -- ++(0,-2.2) -| (capt.east);
\draw[fl] (capt.west) -| (comp.south);

% --- Perturbation sur le processus ---
\draw[fl, dashed, color=red!70!black] (pert) -- (proc.north);

% --- Labels +/- sur comparateur ---
\node[lbl, above left=1pt and 1pt of comp] {$+$};
\node[lbl, below=2pt of comp] {$-$};

\end{tikzpicture}
\caption{Schéma A — Boucle de régulation en boucle fermée (principe général)}
\label{fig:boucle-fermee}
\end{figure}

\subsection{Régulation Tout-ou-Rien (TOR)}

\begin{definition}
  La \textbf{régulation TOR} (Tout-Ou-Rien) ne connaît que deux états : \textit{marche} ou \textit{arrêt}. Un \textbf{hystérésis} $\Delta w$ est introduit pour éviter les commutations trop fréquentes (pompage).
\end{definition}

Si la consigne est $w = \SI{20}{\celsius}$ avec une bande d'hystérésis $\Delta w = \SI{1}{\kelvin}$ :
\begin{itemize}
  \item Le chauffage s'enclenche quand $y < \SI{19.5}{\celsius}$.
  \item Le chauffage s'arrête quand $y > \SI{20.5}{\celsius}$.
\end{itemize}

\begin{attention}
  La régulation TOR génère des oscillations permanentes autour de la consigne. Elle convient aux applications peu exigeantes (résistances électriques, petits chaudières domestiques) mais est inadaptée aux processus thermiques à forte inertie ou aux exigences de précision élevées.
\end{attention}

\subsection{Régulation proportionnelle (P)}

L'action proportionnelle produit une commande proportionnelle à l'erreur :

\begin{equation}
  u(t) = K_P \cdot e(t) + u_0
\end{equation}

où :
\begin{itemize}
  \item $u(t)$ : signal de commande [$\%$ ou signal 0–10 V]
  \item $K_P$ : gain proportionnel [sans dimension ou \%/K]
  \item $e(t) = w - y(t)$ : erreur à l'instant $t$ [\si{\kelvin}]
  \item $u_0$ : point de fonctionnement nominal [\%]
\end{itemize}

\begin{definition}
  La \textbf{bande proportionnelle} $X_P$ (ou \textit{proportional band} en anglais) est l'écart sur la grandeur mesurée qui correspond à une variation de $100\,\%$ du signal de sortie. On a :
  \begin{equation}
    X_P = \frac{1}{K_P} \times 100 \quad [\si{\kelvin}]
  \end{equation}
  Une bande proportionnelle étroite ($X_P$ faible, $K_P$ élevé) améliore la rapidité mais peut provoquer des instabilités.
\end{definition}

L'action P seule laisse subsister un \textbf{écart statique résiduel} en régime permanent.

\subsection{Régulation proportionnelle-intégrale (PI)}

L'ajout d'une action intégrale annule l'écart statique :

\begin{equation}
  u(t) = K_P \left[ e(t) + \frac{1}{T_i} \int_0^t e(\tau)\,\mathrm{d}\tau \right]
\end{equation}

où $T_i$ est le \textbf{temps d'intégration} [\si{\second}] (ou constante d'action intégrale).

\begin{itemize}
  \item $T_i$ petit $\Rightarrow$ action intégrale rapide, risque de dépassement et d'oscillations.
  \item $T_i$ grand $\Rightarrow$ action intégrale lente, retour lent à la consigne.
\end{itemize}

Le régulateur PI est le \textbf{standard de l'industrie CVC} : il assure précision et stabilité sur la majorité des boucles thermiques et hydrauliques.

\subsection{Régulation proportionnelle-intégrale-dérivée (PID)}

\begin{equation}
  u(t) = K_P \left[ e(t) + \frac{1}{T_i} \int_0^t e(\tau)\,\mathrm{d}\tau + T_d \frac{\mathrm{d}e}{\mathrm{d}t} \right]
\end{equation}

où $T_d$ est le \textbf{temps dérivé} [\si{\second}] (anticipation sur la vitesse de variation de l'erreur).

\begin{tabular}{lll}
  \toprule
  Action & Effet & Inconvénient \\
  \midrule
  P & Rapidité proportionnelle à $K_P$ & Écart statique résiduel \\
  I & Annule l'écart statique & Risque de dépassement \\
  D & Amortit les oscillations, anticipe & Amplifie le bruit de mesure \\
  \bottomrule
\end{tabular}

\vspace{0.5em}
\begin{attention}
  En CVC, l'action dérivée est rarement utilisée sur les boucles de température ambiante (signal bruité, inertie thermique importante). Elle peut être utile sur les boucles de pression d'eau rapides.
\end{attention}

\subsection{Réglage des paramètres PID — méthode pratique}

\begin{methode}
  \textbf{Méthode par essais successifs (méthode de Ziegler-Nichols simplifiée en CVC) :}
  \begin{enumerate}
    \item Mettre l'installation en régime stable (charge thermique constante).
    \item Partir avec $T_i = \infty$ (intégrale désactivée) et $T_d = 0$.
    \item Augmenter $K_P$ jusqu'à obtenir des oscillations entretenues : noter le gain critique $K_{Pc}$ et la période d'oscillation $T_{osc}$.
    \item Appliquer les valeurs recommandées pour un régulateur PI :
      $K_P \approx 0{,}45 \cdot K_{Pc}$ et $T_i \approx 0{,}83 \cdot T_{osc}$.
    \item Affiner empiriquement sur site (réponse à un échelon de consigne).
  \end{enumerate}
\end{methode}

% ------------------------------------------------------------
\section{Applications en installations CVC}
% ------------------------------------------------------------

\subsection{Régulation de la température ambiante}

La boucle de base d'une zone chauffée met en jeu :
\begin{itemize}
  \item \textbf{Capteur :} sonde de température ambiante (NTC ou PT1000), précision typique $\pm\SI{0.3}{\kelvin}$.
  \item \textbf{Consigne :} température de confort (ex. $w = \SI{20}{\celsius}$) ou de réduit ($w = \SI{16}{\celsius}$).
  \item \textbf{Actionneur :} vanne 3 voies motorisée ou vanne 2 voies sur émetteur, variateur de fréquence sur pompe.
  \item \textbf{Perturbation principale :} température extérieure $T_{ext}$, apports internes.
\end{itemize}

\subsection{Régulation par loi d'eau (régulation climatique)}

\begin{definition}
  La \textbf{loi d'eau} (ou courbe de chauffe) est une relation linéaire entre la température de départ du réseau de chauffage $T_{dep}$ et la température extérieure $T_{ext}$. Elle anticipe les besoins thermiques sans attendre qu'une erreur de confort se manifeste — c'est une \textbf{régulation en boucle ouverte} sur $T_{ext}$, corrigée en boucle fermée sur $T_{dep}$.
\end{definition}

La loi d'eau s'exprime par :
\begin{equation}
  T_{dep} = T_{dep,base} + K_{loi} \cdot (T_{ext,base} - T_{ext})
\end{equation}

où :
\begin{itemize}
  \item $T_{dep,base}$ : température de départ à la température extérieure de base [\si{\celsius}]
  \item $T_{ext,base}$ : température extérieure de base de dimensionnement [\si{\celsius}]
  \item $K_{loi}$ : pente de la loi d'eau [K/K, sans dimension]
\end{itemize}

\begin{exemple}
  \textbf{Paramétrage d'une loi d'eau pour un réseau de radiateurs :}

  Données du projet :
  \begin{itemize}
    \item Température extérieure de base : $T_{ext,base} = \SI{-7}{\celsius}$ (Paris, zone H1b)
    \item Température de départ à la base : $T_{dep,base} = \SI{75}{\celsius}$ (radiateurs acier)
    \item Température de départ minimale souhaitée (à $T_{ext} = \SI{15}{\celsius}$) : $T_{dep,min} = \SI{30}{\celsius}$
  \end{itemize}

  Calcul de la pente :
  \begin{equation}
    K_{loi} = \frac{T_{dep,base} - T_{dep,min}}{T_{ext,min} - T_{ext,base}} = \frac{75 - 30}{15 - (-7)} = \frac{45}{22} \approx 2{,}05
  \end{equation}

  À $T_{ext} = \SI{0}{\celsius}$ :
  \begin{equation}
    T_{dep} = 30 + 2{,}05 \times (15 - 0) = 30 + 30{,}75 \approx \SI{60{,}8}{\celsius}
  \end{equation}

  Le régulateur climatique fixe la consigne de départ à $\SI{60{,}8}{\celsius}$ lorsque la température extérieure est nulle.
\end{exemple}

\begin{attention}
  La loi d'eau doit être \textbf{adaptée à chaque émetteur} : les planchers chauffants (émetteurs basse température, $T_{dep} \leq \SI{45}{\celsius}$) nécessitent une pente bien plus faible que des radiateurs haute température.
\end{attention}

% --- Schéma B — Courbe de chauffe (loi d'eau) — TikZ pur ---
% Axe X : T_ext de -10 à +18°C  → plage = 28 K, échelle : 0.35 cm/K → 9.8 cm
% Axe Y : T_dep de 25 à 80°C    → plage = 55 K, échelle : 0.13 cm/K → 7.15 cm
% Origine repère TikZ : (0,0) = (T_ext=-10, T_dep=25)
% Transformation : x_tikz = (T_ext - (-10)) * 0.35   y_tikz = (T_dep - 25) * 0.13
%   Radiateurs   : (-10,70)→(0,5.85)   (18,30)→(9.8,0.65)
%   PC           : (-10,45)→(0,2.60)   (18,28)→(9.8,0.39)
%   Bivalence    : T_ext=-2 → x=(8)*0.35=2.80  T_dep=50 → y=(25)*0.13=3.25
\begin{figure}[H]
\centering
\begin{tikzpicture}[x=0.35cm, y=0.13cm]
  % --- Repère ---
  % Axes
  \draw[->, thick] (0,0) -- (30,0) node[right, font=\small] {$T_{ext}$ (°C)};
  \draw[->, thick] (0,0) -- (0,60) node[above, font=\small] {$T_{dep}$ (°C)};

  % Graduations X : T_ext de -10 à 18 par pas de 4
  % x_tikz = (T_ext - (-10)) * 1 (unité x=0.35cm → 1 unité = 1K depuis -10)
  \foreach \Text/\Tpos in {-10/0,-6/4,-2/8,2/12,6/16,10/20,14/24,18/28}{
    \draw (\Tpos, -1.5) -- (\Tpos, 0);
    \node[below, font=\footnotesize] at (\Tpos, -2) {\Text};
  }

  % Graduations Y : T_dep de 25 à 80 par pas de 5
  % y_tikz = (T_dep - 25) * 1 (unité y=0.13cm → 1 unité = 1K depuis 25)
  \foreach \Tdep/\Ypos in {25/0,30/5,35/10,40/15,45/20,50/25,55/30,60/35,65/40,70/45,75/50,80/55}{
    \draw (-0.8, \Ypos) -- (0, \Ypos);
    \node[left, font=\footnotesize] at (-1.2, \Ypos) {\Tdep};
  }

  % --- Grille légère ---
  \foreach \Tpos in {0,4,8,12,16,20,24,28}{
    \draw[gray!25, thin] (\Tpos, 0) -- (\Tpos, 55);
  }
  \foreach \Ypos in {0,5,10,15,20,25,30,35,40,45,50,55}{
    \draw[gray!25, thin] (0, \Ypos) -- (28, \Ypos);
  }

  % --- Courbe 1 : Radiateurs (pente forte) ---
  % Points : T_ext=-10,T_dep=70 → (0,45) ; T_ext=18,T_dep=30 → (28,5)
  \draw[btsprimary, line width=1.8pt] (0,45) -- (28,5);
  \node[btsprimary, font=\footnotesize\bfseries, anchor=south west] at (0,45)
    {Radiateurs};

  % --- Courbe 2 : Plancher chauffant (pente douce) ---
  % Points : T_ext=-10,T_dep=45 → (0,20) ; T_ext=18,T_dep=28 → (28,3)
  \draw[btsaccent, line width=1.8pt, dashed] (0,20) -- (28,3);
  \node[btsaccent, font=\footnotesize\bfseries, anchor=north west] at (0,20)
    {Plancher chauffant};

  % --- Point de bivalence PAC ---
  % T_ext=-2 → x=8 ; T_dep=50 → y=25
  \filldraw[red!70!black] (8,25) circle (0.6ex);
  \draw[red!70!black, thin, ->, shorten >=2pt]
    (8,25) -- ++(5,6)
    node[anchor=south west, font=\footnotesize, color=red!70!black,
         align=left]
      {Point de bivalence PAC\\$T_{ext}=-2$°C, $T_{dep}=50$°C};

  % --- Ligne de coupure chauffage T_ext=18°C ---
  \draw[gray!60, dotted, thin] (28,0) -- (28,55);
  \node[gray, font=\footnotesize, rotate=90, anchor=south] at (28.8,10)
    {Coupure (18°C)};

  % --- Légende ---
  \draw[btsprimary, line width=1.5pt] (16,50) -- (20,50);
  \node[right, font=\footnotesize] at (20.3,50) {Radiateurs};
  \draw[btsaccent, line width=1.5pt, dashed] (16,46) -- (20,46);
  \node[right, font=\footnotesize] at (20.3,46) {Plancher chauffant};

\end{tikzpicture}
\caption{Schéma B — Courbe de chauffe (loi d'eau) : radiateurs vs plancher chauffant,
         avec point de bivalence PAC ($T_{biv}=-2$°C)}
\label{fig:loi-eau}
\end{figure}

\subsection{Régulation en cascade (maître-esclave)}

\begin{definition}
  La \textbf{régulation en cascade} utilise deux boucles imbriquées : une boucle \textbf{maître} lente calcule la consigne de la boucle \textbf{esclave} rapide. La boucle esclave corrige les perturbations rapides avant qu'elles n'atteignent la grandeur principale.
\end{definition}

Exemple typique : régulation de la température ambiante (maître) $\rightarrow$ régulation de la température de départ réseau (esclave).

\begin{itemize}
  \item \textbf{Boucle maître :} mesure $T_{ambiance}$, compare à $w_{ambiance} = \SI{20}{\celsius}$, corrige la consigne de départ $w_{dep}$.
  \item \textbf{Boucle esclave :} mesure $T_{dep}$, compare à $w_{dep}$, agit sur la vanne mélangeuse 3 voies.
\end{itemize}

La cascade offre une meilleure immunité aux perturbations hydrauliques (variations de débit) que la boucle simple.

\subsection{Régulation de la pression différentielle hydraulique}

Sur un réseau hydraulique à débit variable (vannes 2 voies sur les émetteurs), la pression différentielle aux bornes des émetteurs doit rester dans une plage acceptable pour maintenir le débit nominal de chaque zone.

\begin{equation}
  \Delta P_{regulé} = \text{constante} \quad \Rightarrow \quad \text{variateur de fréquence sur pompe}
\end{equation}

La consigne de pression différentielle est typiquement fixée à \SI{30}{\percent} à \SI{50}{\percent} de la pression différentielle maximale. Le variateur adapte la vitesse de la pompe par une action PI.

\subsection{Régulation hygrométrique}

En traitement d'air, la boucle de régulation de l'humidité relative agit sur :
\begin{itemize}
  \item \textbf{Humidification :} vapeur d'eau électrique ou à gaz, buse haute pression.
  \item \textbf{Déshumidification :} batterie froide (en dessous du point de rosée).
\end{itemize}

Le régulateur compare la consigne d'humidité relative $HR_c$ à la mesure $HR_m$ et pilote le modulant du système correspondant.

% ------------------------------------------------------------
\section{Optimisation énergétique par la régulation}
% ------------------------------------------------------------

\subsection{Optimisation du démarrage (relance optimale)}

\begin{definition}
  La \textbf{relance optimale} (ou \textit{optimal start}) est un algorithme de régulation qui calcule dynamiquement l'heure de démarrage du chauffage afin d'atteindre la consigne de confort exactement au début de l'occupation, sans chauffer inutilement pendant la nuit.
\end{definition}

L'heure de démarrage optimale $t_{start}$ est calculée par :
\begin{equation}
  t_{start} = t_{occupation} - \frac{w - T_{amb,nuit}}{K_{bâtiment} \times (T_{dep,max} - w)}
\end{equation}

où $K_{bâtiment}$ est une constante d'apprentissage ajustée par l'automate en fonction des relances précédentes.

\textbf{Économie typique :} 5 à 15\,\% de la consommation de chauffage sur un bâtiment de bureaux (EN 15232, passage de classe C à classe B).

\subsection{Discontinuité de chauffage et hors-gel}

\begin{table}[H]
  \centering
  \caption{Modes de fonctionnement types d'une installation CVC}
  \begin{tabular}{p{3cm}p{3cm}p{4.5cm}}
    \toprule
    Mode & Consigne typique & Déclenchement \\
    \midrule
    Confort & $w = 20$--21\,°C & Heures d'occupation \\
    Réduit & $w = 16$--17\,°C & Nuit, week-end \\
    Hors-gel & $w = 7$--8\,°C & Arrêt total (vacances) \\
    Arrêt & — & Moteur arrêté, chaufferie en veille \\
    \bottomrule
  \end{tabular}
\end{table}

\begin{attention}
  Le mode hors-gel ne doit jamais être désactivé sur les réseaux hydrauliques exposés au gel (chaufferies mal isolées, réseaux extérieurs). La gelée des canalisations est responsable de sinistres importants sur les installations CVC non gardiennées.
\end{attention}

\subsection{Régulation en cascade PAC + chaufferie appoint}

Dans les installations combinant une PAC et une chaudière d'appoint :
\begin{enumerate}
  \item La PAC est prioritaire : elle couvre la charge jusqu'à sa puissance maximale.
  \item La chaudière prend le relais (\textbf{appoint}) quand la PAC ne suffit plus ($T_{ext}$ très basse) ou quand la puissance demandée dépasse la PAC.
  \item Le \textbf{point de bivalence} est la température extérieure en dessous de laquelle l'appoint s'enclenche (typiquement $T_{biv} = -3$\,à\,$+2$\,°C selon le dimensionnement).
\end{enumerate}

\begin{equation}
  \dot{Q}_{PAC}(T_{ext}) = \dot{Q}_{PAC,max} \times \frac{T_{ext} - T_{min}}{T_{max} - T_{min}}
\end{equation}

Le dimensionnement de la PAC au point de bivalence (couverture 70--80\,\% des besoins) minimise l'investissement tout en maintenant un SCOP élevé.

% ------------------------------------------------------------
\section{Gestion Technique du Bâtiment (GTB)}
% ------------------------------------------------------------

\subsection{Définitions et périmètre}

\begin{definition}
  La \textbf{Gestion Technique du Bâtiment (GTB)}, ou \textit{Building Management System (BMS)} en anglais, est un système informatique centralisé qui supervise, contrôle et optimise l'ensemble des équipements techniques d'un bâtiment : CVC, éclairage, sécurité, comptage énergétique, contrôle d'accès.

  La \textbf{Gestion Technique Centralisée (GTC)} est une appellation synonyme, historiquement utilisée en France pour les bâtiments industriels.
\end{definition}

Le \textbf{décret BACS} (2020, transposition de la directive EPBD 2018) rend obligatoire l'installation d'une GTB pour les bâtiments tertiaires neufs dès \SI{290}{\kilo\watt} de puissance nominale CVC, et pour les bâtiments existants au-delà de \SI{290}{\kilo\watt} lors de rénovations, à partir de 2025.

\subsection{Architecture en trois niveaux}

La GTB est structurée selon trois niveaux fonctionnels :

\begin{center}
\begin{tabular}{lll}
  \toprule
  Niveau & Nom & Fonctions \\
  \midrule
  1 & \textbf{Terrain} & Capteurs, actionneurs, équipements CVC \\
  2 & \textbf{Automatisme} & Automates programmables (UTA, régulateurs) \\
  3 & \textbf{Supervision} & Interface homme-machine (IHM), serveur, reporting \\
  \bottomrule
\end{tabular}
\end{center}

\begin{description}
  \item[Niveau 1 — Terrain :] sonde de température, hygromètre, débitmètre, vanne motorisée, variateur de fréquence, chaudière, groupe froid. Les signaux sont analogiques (0–10 V, 4–20 mA, PT100, NTC) ou numériques (TOR, bus de terrain).
  \item[Niveau 2 — Automatisme :] régulateur de zone, automate programmable industriel (API), unité de traitement d'air (UTA) intégrée. Ce niveau implémente les boucles de régulation et les séquences de démarrage/arrêt. Il communique avec le niveau 1 (câblage direct ou bus) et avec le niveau 3 (réseau IP).
  \item[Niveau 3 — Supervision :] logiciel SCADA (Supervisory Control And Data Acquisition) ou interface web. Fonctions : visualisation synoptique, historisation des données, gestion des alarmes, édition des rapports, optimisation des programmes horaires.
\end{description}

\subsection{Protocoles de communication}

\subsubsection{BACnet (Building Automation and Control Networks)}

\begin{definition}
  \textbf{BACnet} (norme ASHRAE/ISO 16484-5) est le protocole ouvert de référence pour la communication entre équipements GTB de fabricants différents. Il définit des \textbf{objets} (points de données) standard : Analog Input, Binary Output, Schedule, Trend Log, etc.
\end{definition}

\begin{itemize}
  \item \textbf{BACnet/IP :} transport sur réseau Ethernet/TCP-IP — utilisé entre régulateurs et serveur de supervision.
  \item \textbf{BACnet MS/TP :} transport sur liaison série RS-485 — utilisé au niveau terrain (coût réduit).
  \item \textbf{BACnet/MSTP} permet d'atteindre des débits jusqu'à \SI{76.8}{\kilo bit\per\second}.
  \item Interopérabilité garantie par les \textbf{profils de conformité (BTL)} délivrés par BACnet Testing Laboratories.
\end{itemize}

\subsubsection{KNX}

\begin{definition}
  \textbf{KNX} (norme EN 50090 / ISO 14543) est un protocole de bus décentralisé, conçu initialement pour la domotique et les bâtiments tertiaires. Il couvre l'éclairage, les stores, le CVC et la sécurité.
\end{definition}

\begin{itemize}
  \item \textbf{Support physique :} câble bus KNX 2 fils (alimentation + données à \SI{9.6}{\kilo bit\per\second}), RF, IP.
  \item \textbf{Topologie :} lignes, zones, domaines — jusqu'à $57\,375$ adresses individuelles.
  \item \textbf{Programmation :} logiciel ETS (Engineering Tool Software), paramétrage de chaque composant.
  \item \textbf{Avantage CVC :} pilotage des thermostats de zone, actionneurs de vanne, centrales de mesure, avec priorité de sécurité intégrée.
\end{itemize}

\subsubsection{Autres protocoles rencontrés en CVC}

\begin{center}
\begin{tabular}{lll}
  \toprule
  Protocole & Domaine & Caractéristiques \\
  \midrule
  Modbus RTU/TCP & CVC industriel & Simple, robuste, RS-485 ou Ethernet \\
  LonWorks & Bâtiment tertiaire & Bus neurone, interopérable \\
  M-Bus & Comptage énergie & EN 13757, relevé à distance \\
  DALI & Éclairage & Bus numérique 2 fils, 64 adresses \\
  \bottomrule
\end{tabular}
\end{center}

\subsection{Fonctions de la supervision}

Le logiciel de supervision offre, à minima, les fonctions suivantes :

\begin{enumerate}
  \item \textbf{Visualisation synoptique :} schémas animés des installations (état des équipements, valeurs mesurées en temps réel).
  \item \textbf{Gestion des alarmes :} détection des dépassements de seuil, des défauts équipements — journalisation et envoi de notifications (email, SMS).
  \item \textbf{Historisation (trending) :} enregistrement des variables à pas de temps paramétrable (ex. toutes les 15 min) pour analyse des dérives et suivi des consommations.
  \item \textbf{Programmation horaire :} calendriers de fonctionnement (jours ouvrés, week-end, jours fériés), confort/réduit/hors gel.
  \item \textbf{Optimisation du démarrage :} calcul de l'heure de relance optimale en fonction de la température extérieure et de la température ambiante — économie d'énergie en évitant le démarrage trop précoce.
  \item \textbf{Suivi énergétique :} index de compteurs, calcul des indices de performance (IPE, kWh/m²), tableaux de bord conformes au décret tertiaire.
\end{enumerate}

\subsection{Classes de performance GTB selon EN 15232}

La norme NF EN 15232 (2012, remplacée en partie par EN 52120-1) définit quatre classes d'efficacité de la GTB :

\begin{center}
\begin{tabular}{cll}
  \toprule
  Classe & Niveau & Description \\
  \midrule
  D & Non automatisé & Installations sans régulation automatique \\
  C & Automatisation standard & Régulation de base (TOR, P) \\
  B & Automatisation avancée & Régulation PI, loi d'eau, optimisation démarrage \\
  A & Haute performance & GTB intégrée, optimisation multi-énergie, reporting \\
  \bottomrule
\end{tabular}
\end{center}

Le passage de la classe D à la classe B permet de réduire la consommation d'énergie de \textbf{30 à 40\,\%} pour un bâtiment de bureaux (données EN 15232, facteur de pondération).

% ------------------------------------------------------------
\section{Exemples résolus}
% ------------------------------------------------------------

\begin{exemple}[title={Calcul de la loi d'eau pour plancher chauffant}]
  \textbf{Données :}
  \begin{itemize}
    \item Bâtiment de bureaux à Strasbourg — zone H1c, $T_{ext,base} = \SI{-12}{\celsius}$
    \item Plancher chauffant : température de départ maximale $T_{dep,max} = \SI{45}{\celsius}$, température de départ minimale $T_{dep,min} = \SI{28}{\celsius}$ (limite de confort)
    \item Coupure chauffage à $T_{ext} = \SI{18}{\celsius}$
  \end{itemize}

  \textbf{Calcul de la pente :}
  \begin{equation}
    K_{loi} = \frac{T_{dep,max} - T_{dep,min}}{T_{ext,coupure} - T_{ext,base}} = \frac{45 - 28}{18 - (-12)} = \frac{17}{30} \approx 0{,}57
  \end{equation}

  \textbf{Loi d'eau :}
  \begin{equation}
    T_{dep}(T_{ext}) = 28 + 0{,}57 \times (18 - T_{ext})
  \end{equation}

  \textbf{Application numérique} à $T_{ext} = \SI{-5}{\celsius}$ :
  \begin{equation}
    T_{dep} = 28 + 0{,}57 \times (18 - (-5)) = 28 + 0{,}57 \times 23 = 28 + 13{,}1 = \SI{41{,}1}{\celsius}
  \end{equation}

  \textbf{Conclusion :} La loi d'eau fixe $T_{dep} \approx \SI{41}{\celsius}$ pour une température extérieure de $\SI{-5}{\celsius}$. Ce réglage respecte les limites du plancher chauffant et optimise le COP de la pompe à chaleur associée (plus $T_{dep}$ est basse, meilleur est le COP).
\end{exemple}

\begin{exemple}[title={Dimensionnement d'un réseau GTB BACnet/IP}]
  \textbf{Contexte :} bâtiment tertiaire de $\SI{5000}{\metre\squared}$ comportant :
  \begin{itemize}
    \item 8 régulateurs de zone (1 par plateau, régulation température + CO$_2$)
    \item 2 unités de traitement d'air (UTA) avec régulation CVC intégrée
    \item 1 groupe froid + 1 chaudière à condensation — gestion centralisée
    \item 1 serveur de supervision BACnet/IP
  \end{itemize}

  \textbf{Architecture retenue :}
  \begin{itemize}
    \item \textbf{Niveau terrain} : sondes PT1000 câblées en direct sur les régulateurs, vannes motorisées avec retour de position 0–10 V.
    \item \textbf{Niveau automatisme} : 8 régulateurs de zone + 2 automates UTA + 2 automates production — tous adressés en BACnet/IP sur un réseau Ethernet dédié.
    \item \textbf{Niveau supervision} : serveur Windows avec logiciel SCADA BACnet/IP. Interface web accessible depuis smartphone pour l'exploitant.
  \end{itemize}

  \textbf{Points de données estimés :}
  \begin{center}
  \begin{tabular}{lr}
    \toprule
    Sous-système & Nombre de points \\
    \midrule
    Régulateurs de zone ($8 \times 15$ points) & 120 \\
    UTA ($2 \times 40$ points) & 80 \\
    Production ($2 \times 25$ points) & 50 \\
    Comptage énergie ($\times 6$) & 6 \\
    \textbf{Total} & \textbf{256 points} \\
    \bottomrule
  \end{tabular}
  \end{center}

  \textbf{Conclusion :} Le réseau BACnet/IP avec switch dédié offre l'interopérabilité nécessaire entre les équipements de fabricants différents. La classe GTB visée est la \textbf{classe B} (EN 15232), permettant une économie d'énergie estimée à \textbf{25\,\%} par rapport à une installation sans GTB (classe C).
\end{exemple}

\begin{exemple}[title={Paramétrage d'un régulateur PI sur boucle de pression différentielle}]
  \textbf{Contexte :} Un réseau hydraulique à débit variable (vannes 2 voies sur 8 zones) est équipé d'une pompe à vitesse variable. On souhaite réguler la pression différentielle $\Delta P$ à la valeur consigne $w_{\Delta P} = \SI{80}{\kilo\pascal}$.

  \textbf{Données mesurées lors du réglage :}
  \begin{itemize}
    \item Gain critique $K_{Pc} = 0{,}6\,\%/\text{kPa}$ (gain de bande proportionnelle à l'apparition d'oscillations).
    \item Période d'oscillation critique : $T_{osc} = \SI{12}{\second}$.
  \end{itemize}

  \textbf{Application Ziegler-Nichols pour un régulateur PI :}
  \begin{align}
    K_P &= 0{,}45 \times K_{Pc} = 0{,}45 \times 0{,}6 = 0{,}27\,\%/\text{kPa}\\
    T_i  &= 0{,}83 \times T_{osc} = 0{,}83 \times 12 = \SI{10}{\second}
  \end{align}

  \textbf{Vérification du comportement en régime :} Si toutes les vannes sont ouvertes (pleine charge), $\Delta P$ chute à $\SI{60}{\kilo\pascal}$. L'erreur est $e = 80 - 60 = \SI{20}{\kilo\pascal}$.
  \begin{itemize}
    \item Action P : $\Delta u_P = 0{,}27 \times 20 = 5{,}4\,\%$ (augmentation de consigne VFD).
    \item Après stabilisation PI : erreur nulle, $\Delta P = \SI{80}{\kilo\pascal}$.
  \end{itemize}

  \textbf{Économie d'énergie en mi-charge :} Avec 4 zones sur 8 ouvertes, la pompe opère à $n_2 = 0{,}75 \times n_1$. Par les lois de similitude :
  \[
    P_2 = (0{,}75)^3 \times P_1 = 0{,}422 \times P_1
  \]
  La puissance est réduite de \textbf{58\,\%} par rapport au fonctionnement à pleine charge. La régulation de pression différentielle par VFD permet donc des économies d'énergie considérables en bâtiment tertiaire à occupation variable.
\end{exemple}

% ------------------------------------------------------------
\section{Éléments de réglementation}
% ------------------------------------------------------------

\subsection{Décret BACS (2020)}

Le \textbf{décret n°2020-658 du 29 mai 2020} (transposition de la directive européenne EPBD 2018/844) impose l'installation de systèmes d'automatisation et de contrôle des bâtiments (BACS) dans les bâtiments tertiaires :

\begin{itemize}
  \item \textbf{Bâtiments neufs} : dès \SI{290}{\kilo\watt} de puissance nominale de chauffage ou de climatisation — obligation depuis le 1\textsuperscript{er} janvier 2021.
  \item \textbf{Bâtiments existants} : au-delà de \SI{290}{\kilo\watt}, lors du renouvellement du système de chauffage ou de climatisation — obligation depuis le 1\textsuperscript{er} janvier 2025.
  \item Le BACS doit permettre la surveillance, la journalisation et le contrôle de l'efficacité énergétique de l'installation.
\end{itemize}

\subsection{RE2020}

La \textbf{Réglementation Environnementale 2020} (RE2020), entrée en vigueur le 1\textsuperscript{er} janvier 2022 pour les logements neufs, intègre la GTB comme levier d'atteinte des objectifs de performance :
\begin{itemize}
  \item La GTB contribue à la réduction du Bbio (besoin bioclimatique) et du Cep (consommation d'énergie primaire) par l'optimisation des programmes horaires et la récupération de chaleur.
  \item Les calculs réglementaires (moteur de calcul Th-BCE) intègrent les classes GTB selon EN 52120.
\end{itemize}

\subsection{Norme NF EN 15232-1 / EN 52120-1}

\begin{itemize}
  \item \textbf{NF EN 15232-1 (2017)} : Efficacité énergétique des bâtiments — Impact de l'automatisation, de la régulation et de la gestion technique. Définit les classes A à D et les facteurs de pondération par type de bâtiment.
  \item \textbf{EN 52120-1 (2021)} : Révision harmonisée à l'échelle européenne, intégrant les nouvelles exigences BACS.
\end{itemize}

\subsection{Normes de protocoles et d'installation}

\begin{itemize}
  \item \textbf{ASHRAE 135 / ISO 16484-5} : Norme BACnet — spécification complète du protocole, des objets et des services.
  \item \textbf{EN 50090 / ISO 14543} : Norme KNX — systèmes électroniques pour les habitations et les bâtiments (HBES).
  \item \textbf{NF EN 13757} : Communication des compteurs d'eau, de gaz et d'énergie thermique — inclut le M-Bus.
  \item \textbf{DTU 65.11} (NF P52-305) : Règles d'installation des régulations pour les systèmes de chauffage. Prescrit les emplacements des sondes, les sections de câblage et les distances minimales par rapport aux sources de perturbations électromagnétiques.
  \item \textbf{NF C 15-100} : Installations électriques basse tension — applicable au câblage des automates et des actionneurs.
\end{itemize}

\subsection{Bonnes pratiques de mise en service}

\begin{methode}
  \textbf{Procédure de commissionnement d'une GTB (selon FG BIM / COSTIC) :}
  \begin{enumerate}
    \item Vérifier le câblage de chaque point de mesure et d'action (test E/S).
    \item Calibrer les capteurs (comparaison à un instrument de référence étalonné).
    \item Vérifier le sens d'action des actionneurs (ouverture/fermeture des vannes).
    \item Paramétrer les boucles de régulation (consignes, gains, plages horaires).
    \item Réaliser des tests de réponse à un échelon sur chaque boucle.
    \item Valider les alarmes : simuler un défaut et vérifier la remontée au superviseur.
    \item Rédiger le procès-verbal de mise en service (PV de recette GTB).
  \end{enumerate}
\end{methode}

% \chapter{Efficacité énergétique et RE2020}

% ============================================================
\section{Objectifs pédagogiques}
% ============================================================

À l'issue de ce chapitre, l'étudiant sera capable de :

\begin{itemize}
  \item Identifier et calculer les indicateurs réglementaires de la RE2020 (Bbio, Cep, Cep,nr, DH, IC\textsubscript{énergie}, IC\textsubscript{construction}).
  \item Lire et interpréter un Diagnostic de Performance Énergétique (DPE).
  \item Comparer les systèmes CVC à haute efficacité : pompes à chaleur, CTA double flux, récupération de chaleur.
  \item Calculer les consommations conventionnelles d'une installation CVC selon la méthode réglementaire.
  \item Situer les labels et certifications énergétiques (BBC, BEPOS, HQE, E+C-) dans leur contexte professionnel.
\end{itemize}

\medskip
\noindent\textbf{Compétences GCF mobilisées :}
\begin{itemize}
  \item C1-2 — Exprimer les contraintes réglementaires, environnementales, économiques.
  \item C3-2 — Comparer et proposer des solutions (critères technique, économique, environnemental, normatif).
  \item C3-3 — Dimensionner tout ou partie selon les normes (notes de calcul justifiées).
  \item C5 — Appliquer les réglementations en vigueur (S5.1 — Réglementation thermique, niveau 3).
  \item C8 — Vérifier, adapter les performances d'un système (S10.6 — Critères de bon fonctionnement).
\end{itemize}

\medskip
\noindent\textbf{Savoirs associés principaux :} S5.1, S5.10, S6.6, A2-3, A2-4, B1-1, B2-1, B6-3, B9.

% ============================================================
\section{Introduction}
% ============================================================

La transition énergétique du secteur du bâtiment constitue un enjeu majeur en France : les bâtiments représentent environ 45~\% de la consommation d'énergie finale nationale et 25~\% des émissions de gaz à effet de serre. Pour un technicien supérieur en génie climatique et fluidique, maîtriser les exigences réglementaires et les indicateurs de performance énergétique est indispensable, qu'il s'agisse de concevoir une installation neuve conforme à la RE2020, d'optimiser un système existant en vue d'un audit énergétique, ou de justifier le choix d'équipements à haute efficacité face à un maître d'ouvrage. Ce chapitre présente les fondements réglementaires actuels, les outils de calcul associés, et les systèmes techniques permettant d'atteindre — voire de dépasser — les exigences en vigueur.

% ============================================================
\section{La réglementation thermique RE2020}
% ============================================================

\subsection{Historique et champ d'application}

La Réglementation Environnementale 2020 (RE2020) est entrée en vigueur pour les logements neufs au \textbf{1\textsuperscript{er} janvier 2022}, puis pour les bâtiments de bureaux et d'enseignement au 1\textsuperscript{er} juillet 2022. Elle succède à la RT2012 en renforçant trois axes :

\begin{enumerate}
  \item La \textbf{sobriété énergétique} : réduire les besoins en énergie du bâti.
  \item La \textbf{décarbonation} : réduire les émissions de CO\textsubscript{2} liées à l'énergie et à la construction.
  \item Le \textbf{confort d'été} : limiter l'inconfort thermique sans recourir à la climatisation.
\end{enumerate}

\begin{definition}[title=Champ d'application de la RE2020]
La RE2020 s'applique aux bâtiments neufs à usage de logement, bureaux et enseignement dont le permis de construire est déposé après les dates d'entrée en vigueur. Les bâtiments existants restent soumis à des réglementations spécifiques (décret tertiaire, obligation de rénovation énergétique).
\end{definition}

\subsection{Les indicateurs réglementaires RE2020}

La RE2020 introduit \textbf{six indicateurs} principaux. Trois existaient sous forme modifiée en RT2012 ; trois sont nouveaux.

\subsubsection{Bbio — Besoin bioclimatique}

\begin{definition}[title=Bbio — Besoin bioclimatique conventionnel]
Le Bbio (sans unité) quantifie les besoins en énergie du bâti pour le chauffage, le refroidissement et l'éclairage, \textbf{indépendamment des systèmes énergétiques}. Il caractérise la qualité de l'enveloppe et de la conception architecturale.
\end{definition}

\begin{equation}
  \text{Bbio} = \text{Bbio}_{\text{ch}} + \text{Bbio}_{\text{fr}} + \text{Bbio}_{\text{éclairage}}
\end{equation}

La RE2020 impose :
\begin{equation}
  \text{Bbio} \leq \text{Bbio}_{\text{max}} \times M_{\text{zone}} \times M_{\text{alt}} \times M_{\text{surface}}
\end{equation}

où $M_{\text{zone}}$, $M_{\text{alt}}$, $M_{\text{surface}}$ sont des modulateurs géographiques et dimensionnels.

\begin{attention}
Le Bbio est calculé par un logiciel réglementaire (COMETH ou équivalent) à partir des données architecturales. Son abaissement nécessite d'agir sur l'isolation, l'inertie, les apports solaires et l'étanchéité à l'air — \textbf{avant} de dimensionner les équipements CVC.
\end{attention}

\subsubsection{Cep — Consommation d'énergie primaire conventionnelle}

\begin{definition}[title=Cep — Consommation en énergie primaire]
Le Cep (en \si{\kWh_{ep} \per \m^2 \per \year}) est la consommation conventionnelle annuelle en énergie primaire pour le chauffage, le refroidissement, l'eau chaude sanitaire (ECS), l'éclairage et les auxiliaires (pompes, ventilateurs).
\end{definition}

\begin{equation}
  \text{Cep} = \sum_{i} \frac{C_{\text{ef},i} \times \text{fep}_i}{S_{\text{ref}}}
  \label{eq:cep}
\end{equation}

où :
\begin{itemize}
  \item $C_{\text{ef},i}$ — consommation en énergie finale pour l'usage $i$ [\si{\kWh \per \year}]
  \item $\text{fep}_i$ — coefficient de conversion énergie finale $\rightarrow$ énergie primaire pour l'énergie $i$
  \item $S_{\text{ref}}$ — surface de référence du bâtiment [\si{\m^2}]
\end{itemize}

\medskip
\noindent\textbf{Coefficients fep réglementaires (RE2020) :}

\begin{table}[H]
  \centering
  \begin{tabular}{lc}
    \toprule
    Énergie & $\text{fep}$ \\
    \midrule
    Électricité & 2,3 \\
    Gaz naturel & 1,0 \\
    Fioul domestique & 1,0 \\
    Réseau de chaleur & variable (arrêté) \\
    Bois — biomasse & 1,0 \\
    \bottomrule
  \end{tabular}
  \caption{Coefficients de conversion énergie finale / énergie primaire (RE2020)}
\end{table}

\subsubsection{Cep,nr — Consommation d'énergie primaire non renouvelable}

\begin{definition}[title=Cep,nr — Indicateur clé RE2020]
Le Cep,nr (en \si{\kWh_{ep,nr} \per \m^2 \per \year}) est le \textbf{nouvel indicateur central} de la RE2020 : il ne comptabilise que la part \textbf{non renouvelable} de l'énergie primaire consommée. L'énergie solaire, la pompe à chaleur (fraction renouvelable) et la biomasse sont valorisées.
\end{definition}

\begin{equation}
  \text{Cep,nr} = \sum_{i} \frac{C_{\text{ef},i} \times \text{fep,nr}_i}{S_{\text{ref}}}
\end{equation}

Le seuil Cep,nr\textsubscript{max} varie de \SI{35}{\kWh_{ep,nr} \per \m^2 \per \year} (zones climatiques clémentes) à \SI{70}{\kWh_{ep,nr} \per \m^2 \per \year} (zones froides, petites surfaces).

\subsubsection{DH — Degrés-heures d'inconfort thermique estival}

\begin{definition}[title=DH — Indicateur de confort d'été]
Le DH (en \si{\degreeCelsius \cdot h \per \year}) mesure l'inconfort thermique estival. Il comptabilise la somme des écarts positifs entre la température opérative intérieure et un seuil de confort, sur l'ensemble de la période estivale.
\end{definition}

\begin{equation}
  \text{DH} = \sum_{\text{été}} \max\!\left(0,\; T_{\text{op}} - T_{\text{confort}}\right) \times \Delta t
\end{equation}

La RE2020 fixe :
\begin{itemize}
  \item DH $\leq$ 350 \si{\degreeCelsius \cdot h} : aucune pénalité
  \item 350 $<$ DH $\leq$ 1\,250 \si{\degreeCelsius \cdot h} : exigence maximale tolérée
  \item DH $>$ 1\,250 \si{\degreeCelsius \cdot h} : non-conformité
\end{itemize}

\subsubsection{IC\textsubscript{énergie} et IC\textsubscript{construction} — Impact carbone}

\begin{definition}[title=IC — Indicateurs carbone (nouveautés RE2020)]
\begin{itemize}
  \item \textbf{IC\textsubscript{énergie}} [\si{\kgCO_2eq \per \m^2}] : émissions de GES liées aux consommations d'énergie sur la durée de vie du bâtiment (50 ans conventionnels).
  \item \textbf{IC\textsubscript{construction}} [\si{\kgCO_2eq \per \m^2}] : émissions liées aux matériaux et procédés de construction (énergie grise). Calculé via les FDES (Fiches de Déclaration Environnementale et Sanitaire).
\end{itemize}
Ces indicateurs sont calculés par Analyse de Cycle de Vie (ACV) dynamique.
\end{definition}

Les seuils IC\textsubscript{construction} se renforcent par paliers : 2025, 2028, 2031, conduisant progressivement vers des constructions biosourcées ou à faible empreinte carbone.

% ============================================================
\section{Diagnostic de Performance Énergétique (DPE)}
% ============================================================

\subsection{Définition et cadre réglementaire}

Le DPE est un document réglementaire (loi Grenelle, réforme 2021) obligatoire pour la vente et la location de tout bâtiment à usage d'habitation. Depuis le \textbf{1\textsuperscript{er} juillet 2021}, le DPE est opposable juridiquement et sa méthode de calcul est unifiée (3CL-DPE 2021).

\begin{definition}[title=Étiquettes DPE]
Le DPE classe le logement sur deux échelles de A à G :
\begin{itemize}
  \item \textbf{Énergie} : consommation en énergie primaire [\si{\kWh_{ep} \per \m^2 \per \year}]
  \item \textbf{Climat} : émissions de GES [\si{\kgCO_2eq \per \m^2 \per \year}]
\end{itemize}
La classe retenue est la moins favorable des deux.
\end{definition}

\begin{table}[H]
  \centering
  \begin{tabular}{ccc}
    \toprule
    Classe & Énergie primaire [\si{\kWh_{ep} \per \m^2 \per \year}] & GES [\si{\kgCO_2eq \per \m^2 \per \year}] \\
    \midrule
    A & $\leq 70$  & $\leq 6$   \\
    B & $\leq 110$ & $\leq 11$  \\
    C & $\leq 180$ & $\leq 30$  \\
    D & $\leq 250$ & $\leq 50$  \\
    E & $\leq 330$ & $\leq 70$  \\
    F & $\leq 420$ & $\leq 100$ \\
    G & $> 420$    & $> 100$    \\
    \bottomrule
  \end{tabular}
  \caption{Seuils DPE 2021 — logements résidentiels}
\end{table}

\subsection{Audit énergétique réglementaire}

Depuis le 1\textsuperscript{er} avril 2023, les logements classés F ou G mis en vente doivent être accompagnés d'un \textbf{audit énergétique} (décret n°2022-780). Cet audit inclut :
\begin{itemize}
  \item L'analyse de l'enveloppe et des systèmes existants.
  \item Au moins deux scénarios de rénovation par étapes menant à la classe B.
  \item L'estimation des coûts et des aides financières (MaPrimeRénov', CEE).
\end{itemize}

Le technicien GCF intervient en proposant les solutions CVC adaptées : PAC, ventilation double flux, régulation, isolation des réseaux.

% ============================================================
\section{Systèmes CVC à haute efficacité énergétique}
% ============================================================

\subsection{Pompes à chaleur (PAC)}

\subsubsection{Principe et indicateurs de performance}

\begin{definition}[title=Coefficient de Performance — COP et SCOP]
\begin{itemize}
  \item \textbf{COP} (Coefficient Of Performance) : rapport de la puissance thermique utile sur la puissance électrique absorbée en régime permanent, à des conditions normalisées (EN 14511).
  \item \textbf{SCOP} (Seasonal COP) : COP saisonnier, intégré sur toute la saison de chauffage selon EN 14825. C'est l'indicateur réglementaire pour la RE2020.
\end{itemize}
\end{definition}

\begin{equation}
  \text{COP} = \frac{P_{\text{ch}}}{P_{\text{élec}}}
  \label{eq:cop}
\end{equation}

\begin{equation}
  \text{SCOP} = \frac{Q_{\text{ch,saison}}}{W_{\text{élec,saison}}}
  \label{eq:scop}
\end{equation}

\medskip
\noindent\textbf{Valeurs typiques pour les PAC air/eau :}

\begin{table}[H]
  \centering
  \begin{tabular}{lcc}
    \toprule
    Type de PAC & COP (A7/W35) & SCOP moyen \\
    \midrule
    Air/eau monobloc & 3,2 -- 4,5 & 2,8 -- 3,8 \\
    Air/eau split & 3,5 -- 5,0 & 3,0 -- 4,2 \\
    Eau glycolée/eau (géothermique) & 4,5 -- 6,5 & 4,0 -- 5,5 \\
    \bottomrule
  \end{tabular}
  \caption{COP et SCOP typiques des PAC (conditions EN 14511)}
\end{table}

\begin{attention}
Le COP dépend fortement des températures de source froide et de condensation. Une PAC air/eau dimensionnée pour des planchers chauffants (35°C) aura un COP nettement supérieur à une PAC alimentant des radiateurs haute température (70°C).
\end{attention}

\subsubsection{Fraction renouvelable et Cep,nr}

La part renouvelable de l'énergie produite par une PAC est :

\begin{equation}
  Q_{\text{renouvelable}} = Q_{\text{ch}} - W_{\text{élec}} = Q_{\text{ch}} \left(1 - \frac{1}{\text{COP}}\right)
\end{equation}

Dans le calcul RE2020, la PAC électrique bénéficie d'un coefficient fep,nr réduit car une fraction de sa production est d'origine renouvelable, ce qui améliore directement le Cep,nr.

\subsection{Centrale de traitement d'air double flux avec récupération}

\subsubsection{Rendement de récupération thermique}

\begin{definition}[title=Rendement de récupération thermique $\eta_r$]
Le rendement de récupération (ou efficacité thermique) d'un échangeur air/air exprime la fraction de l'énergie de l'air extrait récupérée pour préchauffer (ou pré-refroidir) l'air neuf.
\end{definition}

\begin{equation}
  \eta_r = \frac{T_{\text{soufflage}} - T_{\text{extérieur}}}{T_{\text{extrait}} - T_{\text{extérieur}}}
  \label{eq:eta_recup}
\end{equation}

où :
\begin{itemize}
  \item $T_{\text{soufflage}}$ — température de l'air neuf après l'échangeur [\si{\degreeCelsius}]
  \item $T_{\text{extérieur}}$ — température de l'air extérieur [\si{\degreeCelsius}]
  \item $T_{\text{extrait}}$ — température de l'air extrait [\si{\degreeCelsius}]
\end{itemize}

Les échangeurs à contre-courant atteignent $\eta_r = $ 75 à 95~\%. La norme EN 308 définit les conditions d'essai. La RE2020 exige $\eta_r \geq $ 75~\% pour les systèmes double flux collectifs.

\subsubsection{Économie d'énergie par rapport à un système simple flux}

L'économie annuelle d'énergie thermique apportée par la récupération est :

\begin{equation}
  \Delta Q_{\text{récup}} = \dot{m}_{\text{air}} \times c_p \times \eta_r \times \int_{\text{hiver}} \left(T_{\text{intérieur}} - T_{\text{ext}}(t)\right) \, dt
  \label{eq:recup_annuel}
\end{equation}

En pratique, pour les calculs réglementaires, on utilise les DJU (Degrés-Jours Unifiés) de la zone climatique.

\subsection{Récupération de chaleur sur réseaux d'eaux usées}

Les systèmes de récupération de chaleur sur eaux grises (RECUEIL) permettent de préchauffer l'eau froide sanitaire avec la chaleur des eaux de douche. Le rendement typique est de 40 à 60~\%, ce qui réduit la consommation d'ECS de 15 à 25~\%.

\begin{equation}
  Q_{\text{ECS, économisée}} = \dot{m}_{\text{EFS}} \times c_p \times \eta \times (T_{\text{ECS}} - T_{\text{EFS,ambiante}})
\end{equation}

\subsection{Récupération de chaleur sur condenseur (free-cooling et desurchauffe)}

En installation de climatisation ou de froid commercial, la chaleur rejetée au condenseur peut être valorisée pour la production d'ECS ou de chauffage. Le condenseur à récupération totale peut fournir de l'eau chaude à 55--65°C.

\begin{equation}
  Q_{\text{récup,condenseur}} = Q_{\text{évap}} + W_{\text{comp}} = Q_{\text{évap}} \times \left(1 + \frac{1}{\text{COP}_{\text{froid}}}\right)
\end{equation}

% ============================================================
\section{Calcul des consommations conventionnelles CVC}
% ============================================================

\subsection{Méthode de calcul réglementaire}

Le calcul réglementaire RE2020 s'appuie sur la méthode \textbf{Th-BCE 2020} (Thermique Bâtiment — Calcul Énergétique), décrite dans l'arrêté du 4 août 2021. Les logiciels homologués (PLEIADES, COMETH, etc.) implémentent cette méthode.

\begin{methode}[title=Démarche de calcul des consommations CVC]
\begin{enumerate}
  \item Définir les \textbf{données d'entrée} : géographie (zone H1a, H2a…), orientation, surface, composition des parois, systèmes.
  \item Calculer les \textbf{déperditions} par transmission et renouvellement d'air (méthode NF EN 12831).
  \item Déterminer les \textbf{besoins} en chauffage, refroidissement et ECS.
  \item Calculer les \textbf{consommations en énergie finale} en tenant compte des rendements système (génération, distribution, émission, régulation).
  \item Convertir en \textbf{énergie primaire} via les coefficients fep.
  \item Vérifier les \textbf{exigences} Cep,nr, Bbio, DH, IC.
\end{enumerate}
\end{methode}

\subsection{Rendements de système selon NF EN 15316}

La consommation en énergie finale d'un système de chauffage se décompose selon :

\begin{equation}
  C_{\text{ef,ch}} = \frac{Q_{\text{besoin,ch}}}{\eta_{\text{gén}} \times \eta_{\text{dist}} \times \eta_{\text{ém}} \times \eta_{\text{rég}}}
  \label{eq:cef_chauffage}
\end{equation}

où :
\begin{itemize}
  \item $Q_{\text{besoin,ch}}$ — besoin en chaleur du bâtiment [\si{\kWh}]
  \item $\eta_{\text{gén}}$ — rendement de génération (chaudière, PAC…)
  \item $\eta_{\text{dist}}$ — rendement de distribution (pertes réseau)
  \item $\eta_{\text{ém}}$ — rendement d'émission (corps de chauffe)
  \item $\eta_{\text{rég}}$ — rendement de régulation
\end{itemize}

\begin{table}[H]
  \centering
  \begin{tabular}{lcccc}
    \toprule
    Système & $\eta_{\text{gén}}$ & $\eta_{\text{dist}}$ & $\eta_{\text{ém}}$ & $\eta_{\text{rég}}$ \\
    \midrule
    Chaudière gaz condensation & 1,05 & 0,95 & 0,97 & 0,97 \\
    PAC air/eau (SCOP 3,5) & 3,50 & 0,97 & 0,98 & 0,97 \\
    Réseau de chaleur urbain & 0,90 & 0,95 & 0,97 & 0,97 \\
    Électrique direct & 1,00 & 1,00 & 0,99 & 0,95 \\
    \bottomrule
  \end{tabular}
  \caption{Rendements de système de chauffage typiques (source : Th-BCE 2020)}
\end{table}

\subsection{DJU et besoins annuels en chauffage}

Les Degrés-Jours Unifiés (DJU) permettent d'estimer le besoin annuel en chauffage :

\begin{equation}
  Q_{\text{besoin,ch}} = \frac{(\text{Ubât} \times S_{\text{dép}} + 0{,}34 \times \dot{V}_{\text{inf}}) \times \text{DJU} \times 24}{1\,000}
  \label{eq:besoin_dju}
\end{equation}

où :
\begin{itemize}
  \item $\text{Ubât}$ — coefficient de déperdition global du bâti [\si{\W \per \m^2 \per \kelvin}]
  \item $S_{\text{dép}}$ — surface déperditive [\si{\m^2}]
  \item $\dot{V}_{\text{inf}}$ — débit d'air par infiltration [\si{\m^3 \per \hour}]
  \item $\text{DJU}$ — degrés-jours unifiés annuels (base 18°C) [K.j]
\end{itemize}

\medskip
\noindent\textbf{DJU de référence (base 18°C) pour quelques villes françaises :}

\begin{table}[H]
  \centering
  \begin{tabular}{lcc}
    \toprule
    Ville & Zone RE2020 & DJU chauffage \\
    \midrule
    Bordeaux & H2a & 1\,540 \\
    Lyon & H1b & 2\,130 \\
    Marseille & H3 & 960 \\
    Paris & H1a & 2\,400 \\
    Strasbourg & H1c & 3\,050 \\
    \bottomrule
  \end{tabular}
  \caption{DJU chauffage de référence RE2020 (base 18°C)}
\end{table}

% ============================================================
\section{Décret tertiaire et audit énergétique}
% ============================================================

\subsection{Décret tertiaire (décret ELAN — Dispositif Éco-Énergie Tertiaire)}

\begin{definition}[title=Décret tertiaire]
Le \textbf{décret n°2019-771 du 23 juillet 2019} (dit décret tertiaire) impose aux bâtiments tertiaires de surface $\geq \SI{1000}{\metre\squared}$ une réduction progressive de leur consommation d'énergie finale par rapport à une année de référence.
\end{definition}

\textbf{Objectifs de réduction :}
\begin{center}
\begin{tabular}{cc}
  \toprule
  Échéance & Réduction minimale \\
  \midrule
  2030 & $-40\,\%$ \\
  2040 & $-50\,\%$ \\
  2050 & $-60\,\%$ \\
  \bottomrule
\end{tabular}
\end{center}

Ces objectifs peuvent être atteints par rapport à une consommation de référence (année 2010 recommandée) ou par rapport à un seuil de consommation en valeur absolue (exprimé en \si{\kWh_{ef} \per \metre\squared \per \year} selon l'usage).

\textbf{Plateforme OPERAT :} les assujettis déclarent annuellement leurs consommations sur la plateforme numérique de l'ADEME (Observatoire de la Performance Énergétique, de la Rénovation et des Actions du Tertiaire).

\subsection{Aides financières à la rénovation énergétique}

\begin{table}[H]
  \centering
  \caption{Principales aides financières pour la rénovation CVC (2024)}
  \begin{tabular}{p{3.5cm}p{4cm}p{4cm}}
    \toprule
    Aide & Bénéficiaire & Dispositif \\
    \midrule
    MaPrimeRénov' & Ménages, copropriétés & PAC, VMC, isolation \\
    CEE (Certificats d'Économies d'Énergie) & Entreprises, collectivités & Tous travaux économies \\
    Éco-PTZ & Propriétaires occupants & Prêt à taux zéro rénovation \\
    TVA 5,5\,\% & Travaux sur logements $>$ 2 ans & Rénovation énergétique \\
    DEEE (Décret Éco-Énergie Tertiaire) & Gestionnaires tertiaires & Accompagnement ADEME \\
    \bottomrule
  \end{tabular}
\end{table}

\subsection{Méthode de calcul du temps de retour sur investissement}

\begin{methode}[title=Temps de retour sur investissement (TRI simplifié)]
\begin{equation}
  TRI = \frac{C_{\text{invest}}}{\Delta C_{\text{énergie,annuel}}}
\end{equation}
avec :
\begin{itemize}
  \item $C_{\text{invest}}$ : coût d'investissement après déduction des aides [\si{\euro}]
  \item $\Delta C_{\text{énergie,annuel}}$ : économie annuelle sur la facture énergétique [\si{\euro \per \year}]
\end{itemize}
\textbf{Un TRI $\leq$ 10 ans} est généralement accepté pour les projets de rénovation énergétique en tertiaire.
\end{methode}

% ============================================================
\section{Labels et certifications énergétiques}
% ============================================================

\subsection{BBC — Bâtiment Basse Consommation}

Le label BBC (arrêté du 3 mai 2007, géré par Effinergie) a anticipé la RT2012. Il impose Cep $\leq$ 50 \si{\kWh_{ep} \per \m^2 \per \year} (modulé par zone et altitude). Bien que la RT2012 ait rendu le BBC obligatoire, ce label reste une référence de qualité pour les rénovations.

\subsection{BEPOS — Bâtiment à Énergie POSitive}

\begin{definition}[title=BEPOS — Bâtiment à Énergie Positive]
Un bâtiment BEPOS produit, sur site, davantage d'énergie renouvelable qu'il n'en consomme. La production excédentaire est injectée sur le réseau ou stockée. Le label Effinergie+ définit Cep,nr $\leq$ 0 \si{\kWh_{ep,nr} \per \m^2 \per \year}.
\end{definition}

\subsection{Label E+C- — Énergie Positive et Réduction Carbone}

Ce label expérimental (2016--2022) préfigurait la RE2020. Il comportait quatre niveaux d'énergie (E1 à E4) et deux niveaux carbone (C1 et C2). Les projets labellisés E3C2 ou E4C2 ont servi de base à l'étalonnage des seuils RE2020.

\subsection{HQE — Haute Qualité Environnementale}

\begin{definition}[title=Certification HQE]
La HQE (Cerqual/Certivéa) est une certification globale de la qualité environnementale du bâtiment couvrant 14 cibles organisées en 4 domaines : Éco-construction, Éco-gestion, Confort, Santé. Elle intègre mais dépasse le seul volet énergétique, en abordant les nuisances acoustiques, la qualité de l'air intérieur et l'entretien-maintenance.
\end{definition}

\subsection{Tableau comparatif des labels}

\begin{table}[H]
  \centering
  \begin{tabularx}{\linewidth}{lXXc}
    \toprule
    Label & Organisme & Critères principaux & Statut \\
    \midrule
    BBC & Effinergie & Cep $\leq$ 50 kWh$_{ep}$/m².an & Actif (réno) \\
    BBC Effinergie & Effinergie & Cep conforme RT2012 & Actif (neuf) \\
    Effinergie+ & Effinergie & Cep,nr $\leq$ 0 (BEPOS) & Actif \\
    E+C- & Min. Logement & Énergie + Carbone & Arrêté (2022) \\
    HQE & Cerqual/Certivéa & 14 cibles qualité globale & Actif \\
    BREEAM & BRE Global & Méthode UK, notes A--Outstanding & Actif \\
    LEED & US GBC & Points, niveaux Bronze--Platinum & Actif \\
    \bottomrule
  \end{tabularx}
  \caption{Principaux labels et certifications énergétiques en France}
\end{table}

% ============================================================
\section{Exemples résolus}
% ============================================================

\begin{exemple}[title=Calcul du Cep,nr d'un logement neuf — PAC vs chaudière gaz]

\textbf{Données du projet :} Maison individuelle, $S_{\text{ref}} = \SI{120}{\m^2}$, zone H1a (Paris).
Besoin en chauffage : $Q_{\text{ch}} = \SI{12\,000}{\kWh \per \year}$.
Besoin ECS : $Q_{\text{ECS}} = \SI{2\,400}{\kWh \per \year}$.

\textbf{Scénario A — Chaudière gaz condensation} ($\eta_{\text{syst}} = 1{,}05 \times 0{,}95 \times 0{,}97 \times 0{,}97 = 0{,}940$)

Consommation en énergie finale chauffage + ECS :
\[
C_{\text{ef,A}} = \frac{12\,000}{0{,}940} + \frac{2\,400}{0{,}940} = \SI{15\,319}{\kWh \per \year}
\]
Avec $\text{fep,nr}(\text{gaz}) = 1{,}0$ :
\[
\text{Cep,nr}_A = \frac{15\,319}{120} = \SI{127{,}7}{\kWh_{ep,nr} \per \m^2 \per \year}
\]

\textbf{Scénario B — PAC air/eau} (SCOP chauffage = 3,5 ; SCOP ECS = 2,8 ; $\eta_{\text{dist}} = 0{,}97$)

Consommation électrique chauffage :
\[
C_{\text{ef,ch,B}} = \frac{12\,000}{3{,}5 \times 0{,}97} = \SI{3\,534}{\kWh \per \year}
\]
Consommation électrique ECS :
\[
C_{\text{ef,ECS,B}} = \frac{2\,400}{2{,}8} = \SI{857}{\kWh \per \year}
\]
Avec $\text{fep,nr}(\text{élec}) = 2{,}3$ et déduction de la part renouvelable :

Dans le cadre simplifié RE2020, le calcul Cep,nr intègre la pondération réglementaire :
\[
\text{Cep,nr}_B = \frac{(3\,534 + 857) \times 2{,}3 \times f_{\text{nr,PAC}}}{120}
\]
La fraction non renouvelable de l'électricité en France est fixée à $0{,}60$ par la RE2020 (arrêté du 4 août 2021) : cela signifie que 60\,\% de l'énergie électrique consommée est comptabilisée comme énergie primaire non renouvelable, les 40\,\% restants étant considérés comme renouvelables (part du nucléaire valorisée partiellement et des EnR dans le mix électrique national). Ainsi $f_{\text{nr,PAC}} \approx 0{,}60$ (fraction non renouvelable réglementaire PAC air/eau) :
\[
\text{Cep,nr}_B = \frac{4\,391 \times 2{,}3 \times 0{,}60}{120} = \frac{6\,060}{120} = \SI{50{,}5}{\kWh_{ep,nr} \per \m^2 \per \year}
\]

\textbf{Conclusion :} La PAC permet de diviser le Cep,nr par 2,5 par rapport à la chaudière gaz, et est conforme à la RE2020 (seuil $\approx$ \SI{65}{\kWh_{ep,nr} \per \m^2 \per \year} en H1a pour cette surface). La chaudière gaz seule est non conforme.

\end{exemple}

\begin{exemple}[title=Rendement de récupération d'une CTA double flux]

\textbf{Données :}
Débit d'air : $\dot{V} = \SI{2\,400}{\m^3 \per \hour}$, soit $\dot{m} = 2\,400 / 3\,600 \times 1{,}2 = \SI{0{,}800}{\kg \per \s}$.
Température extérieure : $T_{\text{ext}} = \SI{2}{\degreeCelsius}$.
Température de l'air extrait (local) : $T_{\text{ext,local}} = \SI{22}{\degreeCelsius}$.
Température soufflage après échangeur : $T_{\text{soufflage}} = \SI{17}{\degreeCelsius}$.

\textbf{Calcul du rendement de récupération :}
\[
\eta_r = \frac{17 - 2}{22 - 2} = \frac{15}{20} = 0{,}75 \quad \text{soit } 75\,\%
\]

\textbf{Puissance récupérée :}
\[
P_{\text{récup}} = \dot{m} \times c_p \times (T_{\text{soufflage}} - T_{\text{ext}}) = 0{,}800 \times 1\,006 \times 15 = \SI{12\,072}{\W} \approx \SI{12{,}1}{\kW}
\]

\textbf{Économie annuelle} (2\,400 h de fonctionnement hivernal) :
\[
\Delta Q = 12{,}1 \times 2\,400 = \SI{29\,040}{\kWh \per \year}
\]

À un coût de l'énergie de \SI{0{,}12}{\euro \per \kWh} (gaz), l'économie annuelle est de \textbf{\SI{3\,485}{\euro}}, justifiant l'investissement dans un échangeur performant.

\end{exemple}

\begin{exemple}[title=Calcul du besoin de chauffage et impact du double flux]

\textbf{Données :} Bâtiment résidentiel collectif de \SI{800}{\m^2} (zone H1b, Lyon). Ventilation simple flux hygro B avec débit $\dot{V} = \SI{3\,200}{\m^3 \per \hour}$. On envisage de le remplacer par un système double flux ($\eta_r = 85\,\%$). Coût du kWh gaz : \SI{0{,}10}{\euro}. Surcoût investissement double flux : \SI{18\,000}{\euro}.

\textbf{Besoin de chauffage annuel — ventilation simple flux :}
\[
\dot{m}_{\text{air}} = \frac{3\,200}{3\,600} \times 1{,}2 = \SI{1{,}067}{\kg \per \s}
\]
\[
Q_{\text{vent,SF}} = \dot{m}_{\text{air}} \times c_p \times \text{DJU} \times 24 \times 3{,}6^{-1}
= 1{,}067 \times 1\,006 \times (2\,130 \times 24) / 1\,000 = \SI{54\,852}{\kWh \per \year}
\]

\textbf{Économie apportée par le double flux :}
\[
\Delta Q = 0{,}85 \times 54\,852 = \SI{46\,624}{\kWh \per \year}
\]
Économie financière avec $\eta_{\text{syst,gaz}} = 0{,}95$ :
\[
\Delta E = \frac{46\,624}{0{,}95} \times 0{,}10 = \SI{4\,908}{\euro \per \year}
\]

\textbf{Temps de retour :}
\[
T_{\text{retour}} = \frac{18\,000}{4\,908} \approx \mathbf{3{,}7\,\text{ans}}
\]

\textbf{Conclusion :} Le retour sur investissement en moins de 4 ans justifie pleinement le passage au double flux. De plus, la RE2020 exige $\eta_r \geq 75\,\%$ pour les systèmes double flux collectifs, ce qui est satisfait ($85\,\% > 75\,\%$).
\end{exemple}

% ============================================================
\section{Éléments de réglementation}
% ============================================================

\subsection{Textes réglementaires applicables}

\begin{itemize}
  \item \textbf{Arrêté du 4 août 2021} — Réglementation Environnementale 2020 (RE2020) : exigences de performance énergétique et environnementale des bâtiments neufs.
  \item \textbf{Arrêté du 13 avril 2022} — Méthode de calcul Th-BCE 2020 (logiciels homologués).
  \item \textbf{Loi Énergie-Climat du 8 novembre 2019} — Décret tertiaire, objectifs de réduction des consommations des bâtiments tertiaires existants (-40\,\% en 2030, -50\,\% en 2040, -60\,\% en 2050).
  \item \textbf{Décret n°2022-780 du 4 mai 2022} — Obligation d'audit énergétique pour la vente des logements classés F ou G.
  \item \textbf{Loi ELAN (2018)} — Réforme du DPE, opposabilité depuis le 1\textsuperscript{er} juillet 2021.
  \item \textbf{Arrêté du 31 mars 2021} — Nouvelle méthode DPE 3CL 2021.
\end{itemize}

\subsection{Normes NF EN applicables}

\begin{table}[H]
  \centering
  \begin{tabularx}{\linewidth}{lX}
    \toprule
    Norme & Objet \\
    \midrule
    NF EN 12831 & Calcul des déperditions thermiques (charge de dimensionnement) \\
    NF EN 13370 & Transmission thermique par le sol \\
    NF EN 14511 & PAC — Conditions d'essai et de mesure du COP \\
    NF EN 14825 & PAC — Mesure à charge partielle et calcul du SCOP/SEER \\
    NF EN 15316 & Systèmes de chauffage et ECS — Méthode de calcul des performances \\
    NF EN 13779 & Ventilation des bâtiments non résidentiels — Exigences de performance \\
    NF EN 308 & Échangeurs de chaleur air/air — Procédures d'essai \\
    NF EN ISO 13789 & Performance thermique — Coefficients de transmission \\
    ISO 50001 & Systèmes de management de l'énergie \\
    \bottomrule
  \end{tabularx}
  \caption{Normes NF EN applicables à l'efficacité énergétique des bâtiments}
\end{table}

\subsection{DTU et documents techniques}

\begin{itemize}
  \item \textbf{DTU 65.11} — Dispositifs d'évacuation des produits de combustion des chaudières (impact sur l'étanchéité à l'air).
  \item \textbf{DTU 68.3} — Travaux de bâtiment — Installations de ventilation mécanique (qualité de l'installation et des liaisons).
  \item \textbf{Guide RAGE} — Règles de l'Art Grenelle Environnement : fiches thématiques sur l'isolation, la ventilation, les PAC (disponibles sur le site Règles de l'Art).
  \item \textbf{Guide ADEME} — Méthodes de calcul des économies d'énergie, aide au dimensionnement des systèmes de récupération.
  \item \textbf{Référentiel QUALIBAT / QualiPAC} — Qualification des installateurs de pompes à chaleur (obligatoire pour certaines aides financières).
\end{itemize}

\begin{attention}
La RE2020 évolue par paliers successifs (2025, 2028, 2031) avec des seuils de Cep,nr et d'IC\textsubscript{construction} de plus en plus contraignants. Le technicien GCF doit veiller à se référer aux seuils en vigueur à la date de dépôt du permis de construire, disponibles sur le site officiel \texttt{rt-re2020.fr}.
\end{attention}

% ============================================================
\section{Diagrammes et représentations graphiques}
% ============================================================

\subsection{Diagramme de Sankey simplifié d'un bâtiment résidentiel}

Le diagramme de Sankey ci-dessous illustre les flux d'énergie primaire d'un logement type (chaudière gaz) : de l'énergie primaire consommée aux usages finaux, en passant par les pertes des systèmes et de l'enveloppe.

\begin{figure}[H]
  \centering
  \begin{tikzpicture}[
      >=Stealth,
      node distance=0pt,
      every node/.style={font=\small},
      flux/.style={draw, fill=#1, minimum height=#2, minimum width=2.2cm, inner sep=3pt, align=center},
    ]

    % --- Colonne 1 : énergie primaire ---
    \node[flux=FedBlue!80!black, minimum height=4cm, minimum width=2.5cm] (EP) at (0,0)
      {\textcolor{white}{\textbf{Énergie\\primaire\\100~\%\\(gaz)}}};

    % --- Colonne 2 : énergie finale après génération ---
    \node[flux=FedBlue!50, minimum height=3.2cm, minimum width=2.2cm, right=2.5cm of EP.north east, anchor=north west] (EF)
      {\textbf{Énergie\\finale\\chaudière\\$\eta=0{,}94$}};
    % Perte génération
    \node[flux=FedRed!70, minimum height=0.7cm, minimum width=2.2cm, below=0.05cm of EF.south west, anchor=north west] (PG)
      {\textcolor{white}{Pertes chaudière 6~\%}};

    % --- Colonne 3 : usages finaux ---
    \node[flux=FedGreen!70, minimum height=1.6cm, minimum width=2.2cm, right=2.5cm of EF.north east, anchor=north west] (CH)
      {\textbf{Chauffage\\~65~\%}};
    \node[flux=FedGreen!40, minimum height=0.8cm, minimum width=2.2cm, below=0.05cm of CH.south west, anchor=north west] (ECS)
      {ECS\\20~\%};
    \node[flux=FedOrange!60, minimum height=0.5cm, minimum width=2.2cm, below=0.05cm of ECS.south west, anchor=north west] (AUX)
      {Auxiliaires\\5~\%};
    \node[flux=FedRed!50, minimum height=0.4cm, minimum width=2.2cm, below=0.05cm of AUX.south west, anchor=north west] (PD)
      {\footnotesize Pertes dist.~4~\%};

    % --- Flèches ---
    \draw[->, line width=3pt, FedBlue!60]
      (EP.east) -- node[above, midway, font=\footnotesize]{100~\%} ++(2.5,0);
    \draw[->, line width=2.5pt, FedBlue!50]
      (EF.east) -- ++(2.5,0);
    \draw[->, line width=1pt, FedRed!80]
      (PG.east) -- ++(1.5,0) node[right]{\footnotesize $\rightarrow$ fumées};

    % --- Légende ---
    \node[below=1.2cm of PD.south, anchor=north, font=\footnotesize, align=center]
      {Logement 120~m², zone H1a — chaudière gaz condensation ($\eta_{\text{syst}}=0{,}94$)};

  \end{tikzpicture}
  \caption{Diagramme de Sankey simplifié — flux d'énergie primaire d'un logement résidentiel (chaudière gaz)}
  \label{fig:sankey-batiment}
\end{figure}

\subsection{Pyramide des indicateurs RE2020}

La figure suivante représente la hiérarchie des indicateurs RE2020 : le Bbio caractérise la qualité du bâti (base), le Cep,nr quantifie les consommations d'énergie non renouvelable, et l'IC\textsubscript{énergie} traduit l'empreinte carbone globale sur 50 ans.

\begin{figure}[H]
  \centering
  \begin{tikzpicture}[>=Stealth, every node/.style={font=\small}]

    % Niveau 1 (base) — Bbio
    \fill[FedBlue!70]
      (-4.5,-0.5) -- (4.5,-0.5) -- (3.2,0.8) -- (-3.2,0.8) -- cycle;
    \node[text=white, font=\bfseries\small] at (0, 0.15)
      {Bbio — Besoin bioclimatique (sans unité)};
    \node[font=\footnotesize, align=center] at (0,-1.0)
      {Qualité de l'enveloppe, inertie, apports solaires\\
       \textbf{Bbio $\leq$ Bbio\textsubscript{max} $\times$ M\textsubscript{zone} $\times$ M\textsubscript{alt} $\times$ M\textsubscript{surface}}};

    % Niveau 2 — Cep,nr
    \fill[FedOrange!80]
      (-3.2,0.9) -- (3.2,0.9) -- (1.9,2.2) -- (-1.9,2.2) -- cycle;
    \node[text=white, font=\bfseries\small] at (0, 1.55)
      {Cep,nr [\si{\kWh_{ep,nr}/m^2/an}]};
    \node[font=\footnotesize, align=center] at (0, 3.0)
      {Consommation énergie primaire non renouvelable\\
       \textbf{Cep,nr $\leq$ 35 à 70~\si{\kWh_{ep,nr}/m^2/an}} selon zone};

    % Niveau 3 (sommet) — IC_énergie
    \fill[FedGreen!70]
      (-1.9,2.3) -- (1.9,2.3) -- (0.6,3.4) -- (-0.6,3.4) -- cycle;
    \node[text=white, font=\bfseries\footnotesize, align=center] at (0, 2.85)
      {IC\textsubscript{én.}\\[\si{kgCO_2eq}]};
    \node[font=\footnotesize, align=center] at (0, 4.3)
      {IC\textsubscript{énergie} [\si{\kgCO_2eq/m^2}]\\
       Émissions GES énergie sur 50 ans};

    % Flèche DH (latérale)
    \draw[->, thick, FedRed!80] (4.8, 0.15) -- (3.3, 0.15)
      node[left, font=\footnotesize]{\textcolor{FedRed!80}{}};
    \node[right, font=\footnotesize, FedRed!80, align=left] at (4.85, 0.15)
      {DH $\leq$ 1\,250\,\si{\degreeCelsius\cdot h}\\(confort été)};

    % Flèche IC_construction (latérale)
    \draw[->, thick, FedGreen!60] (-4.8, 1.55) -- (-3.3, 1.55);
    \node[left, font=\footnotesize, FedGreen!60, align=right] at (-4.85, 1.55)
      {IC\textsubscript{construction}\\[\si{kgCO_2eq/m^2}]\\(matériaux)};

    % Titre
    \node[font=\bfseries, align=center] at (0, 5.1)
      {Hiérarchie des indicateurs RE2020};

  \end{tikzpicture}
  \caption{Pyramide des indicateurs réglementaires RE2020 — du bâti (Bbio) aux consommations (Cep,nr) et à l'empreinte carbone (IC\textsubscript{énergie})}
  \label{fig:pyramide-re2020}
\end{figure}

\subsection{Comparaison BBC RT2012 / RE2020 — niveaux Cep}

Ce diagramme en barres illustre l'évolution des exigences réglementaires en énergie primaire : du label BBC (2007) jusqu'aux seuils RE2020 par paliers.

\begin{figure}[H]
  \centering
  \begin{tikzpicture}[>=Stealth, every node/.style={font=\small}]

    % Axes
    \draw[->] (-0.3, 0) -- (10.8, 0) node[right]{\footnotesize Réglementation / Label};
    \draw[->] (0, -0.3) -- (0, 6.8) node[above, align=center]{\footnotesize Cep ou Cep,nr\\\footnotesize [\si{\kWh_{ep}/m^2/an}]};

    % Graduations axe Y (chaque 0.5 cm = 10 kWh)
    \foreach \y/\val in {0/0, 1/20, 2/40, 3/60, 4/80, 5/100, 6/120}{
      \draw (-0.15, \y) -- (0.15, \y);
      \node[left, font=\footnotesize] at (-0.2, \y) {\val};
    }

    % Barres (largeur 1.2cm, positions en x)
    % BBC 2007 : Cep <= 50 kWh_ep -> 2.5 cm
    \fill[FedBlue!80] (1.0, 0) rectangle (2.2, 2.5)
      node[above, font=\footnotesize\bfseries, midway, xshift=0pt, yshift=1.5cm]{\textcolor{FedBlue!80}{$\leq 50$}};
    \node[below, font=\footnotesize, align=center] at (1.6, -0.3) {BBC\\2007};

    % RT2012 : Cep <= 50 kWh_ep -> 2.5 cm (même exigence de base)
    \fill[FedBlue!60] (3.0, 0) rectangle (4.2, 2.5)
      node[above, font=\footnotesize\bfseries, midway, xshift=0pt, yshift=1.5cm]{\textcolor{FedBlue!60}{$\leq 50$}};
    \node[below, font=\footnotesize, align=center] at (3.6, -0.3) {RT2012\\(BBC obligatoire)};

    % RE2020 Cep,nr zone H1a : <= 65 kWh_ep,nr -> 3.25 cm
    \fill[FedOrange!80] (5.0, 0) rectangle (6.2, 3.25)
      node[above, font=\footnotesize\bfseries, midway, xshift=0pt, yshift=1.5cm]{\textcolor{FedOrange!80}{$\leq 65$}};
    \node[below, font=\footnotesize, align=center] at (5.6, -0.3) {RE2020\\H1a (2022)};

    % RE2020 palier 2025 : Cep,nr renforcé ~ 50
    \fill[FedOrange!60] (7.0, 0) rectangle (8.2, 2.5)
      node[above, font=\footnotesize\bfseries, midway, xshift=0pt, yshift=1.5cm]{\textcolor{FedOrange!60}{$\leq 50$}};
    \node[below, font=\footnotesize, align=center] at (7.6, -0.3) {RE2020\\palier 2025};

    % BEPOS / Effinergie+ : Cep,nr <= 0 -> barre symbolique à 0.15 cm
    \fill[FedGreen!70] (9.0, 0) rectangle (10.2, 0.15);
    \node[above, font=\footnotesize\bfseries] at (9.6, 0.2){\textcolor{FedGreen!70}{$\approx 0$}};
    \node[below, font=\footnotesize, align=center] at (9.6, -0.3) {BEPOS\\Effinergie+};

    % Ligne de référence RE2020 seuil maxi
    \draw[dashed, FedRed!80, thick] (-0.1, 3.5) -- (10.5, 3.5)
      node[right, font=\footnotesize, FedRed!80]{Seuil RT2012 non labellisé ($\leq 70$)};

    % Note
    \node[font=\footnotesize, align=left] at (5.0, 6.5)
      {\textit{Cep en \si{\kWh_{ep}/m^2/an} — logement zone H1a, surface $\geq 100$~m²}};

  \end{tikzpicture}
  \caption{Évolution des exigences réglementaires en énergie primaire : BBC (2007), RT2012, RE2020 et BEPOS}
  \label{fig:comparatif-bbc-re2020}
\end{figure}

% ============================================================
\section{Exercices}
% ============================================================

\begin{exercice}[title=Exercice 1 — Indicateurs RE2020 et conformité]

Un bureau de 350~m² (zone H1b, Lyon) est équipé d'une PAC air/eau (SCOP = 3,8) pour le chauffage et d'une chaudière gaz condensation pour l'ECS ($\eta_{\text{syst}} = 0{,}90$).

\textbf{Données :}
\begin{itemize}
  \item Besoin chauffage : $Q_{\text{ch}} = \SI{28\,000}{\kWh \per \year}$
  \item Besoin ECS : $Q_{\text{ECS}} = \SI{5\,600}{\kWh \per \year}$
  \item $\text{fep,nr}(\text{élec}) = 2{,}3$ ; fraction non renouvelable PAC air/eau : $f_{\text{nr}} = 0{,}60$
  \item $\text{fep,nr}(\text{gaz}) = 1{,}0$
  \item Seuil Cep,nr en zone H1b pour bureaux : $\leq \SI{60}{\kWh_{ep,nr} \per \m^2 \per \year}$
\end{itemize}

\textbf{Questions :}
\begin{enumerate}
  \item Calculer la consommation en énergie finale de chaque système.
  \item Calculer le Cep,nr du bâtiment.
  \item Le bâtiment est-il conforme à la RE2020 ?
  \item Quelle modification du système permettrait d'améliorer le Cep,nr ?
\end{enumerate}

\tcblower
\textbf{Corrigé :}

\textbf{1. Consommations en énergie finale :}
\[
C_{\text{ef,ch}} = \frac{Q_{\text{ch}}}{\text{SCOP}} = \frac{28\,000}{3{,}8} = \SI{7\,368}{\kWh \per \year}
\]
\[
C_{\text{ef,ECS}} = \frac{Q_{\text{ECS}}}{\eta_{\text{syst,gaz}}} = \frac{5\,600}{0{,}90} = \SI{6\,222}{\kWh \per \year}
\]

\textbf{2. Calcul du Cep,nr :}

Contribution PAC (électricité) :
\[
\text{Cep,nr}_{\text{PAC}} = \frac{C_{\text{ef,ch}} \times \text{fep,nr} \times f_{\text{nr}}}{S_{\text{ref}}}
= \frac{7\,368 \times 2{,}3 \times 0{,}60}{350} = \frac{10\,169}{350} = \SI{29{,}1}{\kWh_{ep,nr} \per \m^2 \per \year}
\]

Contribution chaudière gaz (ECS) :
\[
\text{Cep,nr}_{\text{gaz}} = \frac{C_{\text{ef,ECS}} \times \text{fep,nr}(\text{gaz})}{S_{\text{ref}}}
= \frac{6\,222 \times 1{,}0}{350} = \SI{17{,}8}{\kWh_{ep,nr} \per \m^2 \per \year}
\]

Total :
\[
\text{Cep,nr} = 29{,}1 + 17{,}8 = \SI{46{,}9}{\kWh_{ep,nr} \per \m^2 \per \year}
\]

\textbf{3. Conformité RE2020 :}
$46{,}9 < 60 \; \si{\kWh_{ep,nr} \per \m^2 \per \year}$ \quad \checkmark \quad Le bâtiment est \textbf{conforme} à la RE2020.

\textbf{4. Amélioration possible :}
Remplacer la chaudière gaz ECS par une PAC thermodynamique dédiée ECS (COP $\approx 2{,}5$) réduirait la contribution ECS au Cep,nr d'environ 30~\%, en supprimant la consommation de gaz fossile.

\end{exercice}

\begin{exercice}[title=Exercice 2 — DPE et classe énergétique]

Un appartement de 75~m² présente les caractéristiques suivantes :
\begin{itemize}
  \item Consommation énergie primaire : \SI{195}{\kWh_{ep} \per \m^2 \per \year}
  \item Émissions de GES : \SI{45}{\kgCO_2eq \per \m^2 \per \year}
\end{itemize}

\textbf{Questions :}
\begin{enumerate}
  \item Quelle est la classe DPE énergie de cet appartement ?
  \item Quelle est la classe DPE GES ?
  \item Quelle est la classe DPE retenue ? Justifier.
  \item Quels travaux prioritaires recommanderiez-vous pour atteindre la classe C ?
\end{enumerate}

\tcblower
\textbf{Corrigé :}

\textbf{1. Classe DPE énergie :}

D'après le tableau des seuils DPE 2021 : $180 < 195 \leq 250$ $\Rightarrow$ \textbf{Classe D}.

\textbf{2. Classe DPE GES :}

$30 < 45 \leq 50$ $\Rightarrow$ \textbf{Classe D}.

\textbf{3. Classe DPE retenue :}

La classe retenue est la \textbf{moins favorable} des deux. Les deux classes sont D, donc la classe DPE est \textbf{D}.

\textbf{4. Travaux pour atteindre la classe C :}

Pour passer de D à C, il faut atteindre $\leq 180\;\si{\kWh_{ep} \per \m^2 \per \year}$ et $\leq 30\;\si{\kgCO_2eq \per \m^2 \per \year}$. Les actions prioritaires sont :
\begin{itemize}
  \item Isolation des combles et des parois (gain estimé : $-20$ à $-30$~\% sur les besoins).
  \item Remplacement du système de chauffage (chaudière ancienne $\rightarrow$ PAC ou chaudière gaz condensation).
  \item Installation d'une VMC double flux pour récupérer la chaleur de l'air extrait.
  \item Remplacement des fenêtres simple vitrage par double vitrage faible émissivité.
\end{itemize}

\end{exercice}

\begin{exercice}[title=Exercice 3 — Temps de retour sur investissement PAC]

Un gestionnaire de bâtiment tertiaire (bureaux, 500~m², zone H2a — Bordeaux) souhaite remplacer sa chaudière gaz fioul (rendement 78~\%) par une PAC air/eau (SCOP = 3,2). Le besoin en chauffage annuel est de $\SI{40\,000}{\kWh \per \year}$.

\textbf{Données économiques :}
\begin{itemize}
  \item Prix du fioul : \SI{0{,}12}{\euro \per \kWh}
  \item Prix de l'électricité : \SI{0{,}18}{\euro \per \kWh}
  \item Coût d'installation PAC (après MaPrimeRénov' et CEE) : \SI{14\,000}{\euro}
\end{itemize}

\textbf{Questions :}
\begin{enumerate}
  \item Calculer la consommation annuelle en énergie finale de la chaudière fioul.
  \item Calculer la consommation annuelle en électricité de la PAC.
  \item Calculer les coûts annuels de chaque solution.
  \item Calculer le temps de retour sur investissement.
\end{enumerate}

\tcblower
\textbf{Corrigé :}

\textbf{1. Consommation fioul (chaudière existante) :}
\[
C_{\text{ef,fioul}} = \frac{Q_{\text{ch}}}{\eta} = \frac{40\,000}{0{,}78} = \SI{51\,282}{\kWh \per \year}
\]

\textbf{2. Consommation électrique PAC :}
\[
C_{\text{ef,élec,PAC}} = \frac{Q_{\text{ch}}}{\text{SCOP}} = \frac{40\,000}{3{,}2} = \SI{12\,500}{\kWh \per \year}
\]

\textbf{3. Coûts annuels :}
\[
\text{Coût fioul} = 51\,282 \times 0{,}12 = \SI{6\,154}{\euro \per \year}
\]
\[
\text{Coût PAC} = 12\,500 \times 0{,}18 = \SI{2\,250}{\euro \per \year}
\]

\textbf{4. Temps de retour sur investissement :}
\[
\Delta C_{\text{annuel}} = 6\,154 - 2\,250 = \SI{3\,904}{\euro \per \year}
\]
\[
TRI = \frac{C_{\text{invest}}}{\Delta C_{\text{annuel}}} = \frac{14\,000}{3\,904} \approx \mathbf{3{,}6\,\text{ans}}
\]

\textbf{Conclusion :} Avec un TRI de moins de 4 ans, l'investissement dans la PAC est très rentable. De plus, le bâtiment bénéficiera d'une réduction du Cep,nr, améliorant sa conformité réglementaire et la valeur verte du bâtiment.

\end{exercice}

\finchapitre

% % ============================================================
%  Chapitre 09 — Acoustique du bâtiment
%  BTS FED — Option GCF
%  Savoir associé : S5.2 (niveau 3) et A7 (niveau 3)
% ============================================================

\chapter{Acoustique du bâtiment}

% ------------------------------------------------------------
\section{Objectifs pédagogiques}
% ------------------------------------------------------------

À l'issue de ce chapitre, l'étudiant sera capable de :

\begin{itemize}
  \item Maîtriser les grandeurs acoustiques fondamentales : niveau de pression sonore $L_p$, niveau de puissance sonore $L_w$, pondération A.
  \item Calculer la propagation et l'affaiblissement du son dans un local et en champ libre.
  \item Appliquer la loi de masse pour l'isolement acoustique des parois et interpréter les indices $R_w$ et $D_{nT,w}$.
  \item Identifier les sources de bruit des installations CVC (ventilateurs, gaines, terminaux, pompes) et quantifier leurs niveaux.
  \item Appliquer la réglementation acoustique en vigueur (arrêté du 30/06/1999, NF EN ISO 10140, RE2020 / NRA).
  \item Dimensionner des solutions de traitement acoustique : silencieux à baffles, plénums, caissons insonorisés.
\end{itemize}

\textbf{Compétences mobilisées :} C3-3 (dimensionner), C5-1 à C5-3 (réglementation), C7-4 à C7-6 (mesures), C8-4 (action corrective).

% ------------------------------------------------------------
\section{Introduction}
% ------------------------------------------------------------

Le bruit constitue l'une des principales sources de nuisances dans les bâtiments. Les installations de génie climatique et fluidique (CVC) — ventilateurs, pompes, gaines, diffuseurs, groupes frigorifiques — sont des sources de bruit à part entière. Le technicien supérieur GCF doit intégrer l'acoustique dès la phase de conception pour garantir le confort des occupants et respecter les exigences réglementaires.

\textbf{Domaines d'application en GCF :}
\begin{itemize}
  \item Choix et implantation des équipements (niveaux de puissance sonore $L_w$ des constructeurs).
  \item Dimensionnement des réseaux aérauliques pour limiter les vitesses génératrices de bruit.
  \item Traitement acoustique des locaux techniques et des gaines.
  \item Vérification réglementaire lors de la réception.
\end{itemize}

% ------------------------------------------------------------
\section{Grandeurs acoustiques fondamentales}
% ------------------------------------------------------------

\subsection{Le son — nature physique}

\begin{definition}
Le son est une onde de pression mécanique qui se propage dans un milieu élastique (air, eau, solide) sous forme de variations de pression autour de la pression atmosphérique. Sa vitesse de propagation dans l'air est :
\[
  c = 331 + 0{,}6 \cdot \theta \quad [\si{\metre\per\second}]
\]
où $\theta$ est la température de l'air en \si{\degreeCelsius}. À 20~\si{\degreeCelsius}, $c \approx 343$~\si{\metre\per\second}.
\end{definition}

\textbf{Domaine fréquentiel audible :} de 20~\si{\hertz} à 20~000~\si{\hertz}. Les tiers d'octave normalisés vont de 100~\si{\hertz} à 3\,150~\si{\hertz} pour les applications courantes du bâtiment.

\subsection{Niveau de pression sonore \texorpdfstring{$L_p$}{Lp}}

\begin{definition}
Le \textbf{niveau de pression sonore} $L_p$ est défini par :
\[
  L_p = 10 \cdot \log_{10}\!\left(\frac{p^2}{p_0^2}\right) = 20 \cdot \log_{10}\!\left(\frac{p}{p_0}\right) \quad [\si{\decibel}]
\]
avec :
\begin{itemize}
  \item $p$ : pression acoustique efficace [\si{\pascal}]
  \item $p_0 = 20~\si{\micro\pascal}$ : pression de référence (seuil d'audibilité à 1~\si{\kilo\hertz})
\end{itemize}
\end{definition}

\begin{attention}
L'échelle en décibels est \textbf{logarithmique}. Doubler la pression acoustique n'augmente $L_p$ que de 6~\si{\decibel}. Additionner deux sources identiques n'augmente $L_p$ que de 3~\si{\decibel}.
\end{attention}

\textbf{Ordres de grandeur courants :}

\begin{center}
\begin{tabular}{lcc}
\toprule
Situation & $L_p$ (dB) & Ressenti \\
\midrule
Seuil d'audibilité & 0 & Silence total \\
Chambre calme & 30 & Très calme \\
Bureau ouvert & 50--60 & Normal \\
Salle de classe & 40--50 & Acceptable \\
Local technique CVC & 70--90 & Bruyant \\
Ventilateur industriel & 85--100 & Très bruyant \\
Seuil de douleur & 120 & Dangereux \\
\bottomrule
\end{tabular}
\end{center}

\subsection{Addition de sources sonores}

\begin{methode}
Pour additionner $n$ sources de niveaux $L_1, L_2, \ldots, L_n$ (en \si{\decibel}), on utilise :
\[
  L_{\text{total}} = 10 \cdot \log_{10}\!\left(\sum_{i=1}^{n} 10^{L_i/10}\right) \quad [\si{\decibel}]
\]
\textbf{Cas particulier :} $n$ sources identiques de niveau $L_0$ :
\[
  L_{\text{total}} = L_0 + 10 \cdot \log_{10}(n) \quad [\si{\decibel}]
\]
\end{methode}

\begin{exemple}
Deux ventilateurs installés dans un local technique émettent chacun $L = 75$~\si{\decibel}. Quel est le niveau résultant ?
\[
  L_{\text{total}} = 10 \cdot \log_{10}(10^{75/10} + 10^{75/10}) = 10 \cdot \log_{10}(2 \times 10^{7{,}5}) = 75 + 10 \cdot \log_{10}(2) \approx 75 + 3 = 78~\si{\decibel}
\]
\end{exemple}

\subsection{Niveau de puissance sonore \texorpdfstring{$L_w$}{Lw}}

\begin{definition}
Le \textbf{niveau de puissance sonore} $L_w$ caractérise la source de bruit indépendamment de son environnement :
\[
  L_w = 10 \cdot \log_{10}\!\left(\frac{W}{W_0}\right) \quad [\si{\decibel}]
\]
avec :
\begin{itemize}
  \item $W$ : puissance acoustique émise par la source [\si{\watt}]
  \item $W_0 = 10^{-12}$~\si{\watt} : puissance de référence
\end{itemize}
$L_w$ est une caractéristique intrinsèque de l'équipement, fournie par le constructeur (données de la fiche technique).
\end{definition}

\begin{attention}
$L_w$ (puissance sonore) $\neq$ $L_p$ (pression sonore). La pression dépend de la distance à la source et de l'acoustique du local. $L_w$ est toujours $>$ $L_p$ mesuré à une distance usuelle.
\end{attention}

\subsection{Pondération A — décibels A \texorpdfstring{dB(A)}{dB(A)}}

\begin{definition}
La \textbf{pondération A} est un filtre fréquentiel qui reproduit la sensibilité de l'oreille humaine. Elle atténue fortement les basses fréquences et les très hautes fréquences.

La correction A à appliquer à chaque bande de tiers d'octave est (valeurs normalisées) :

\begin{center}
\begin{tabular}{lcccccc}
\toprule
Fréquence (Hz) & 125 & 250 & 500 & 1\,000 & 2\,000 & 4\,000 \\
\midrule
Correction A (dB) & $-16{,}1$ & $-8{,}6$ & $-3{,}2$ & $0$ & $+1{,}2$ & $+1{,}0$ \\
\bottomrule
\end{tabular}
\end{center}

Le niveau global pondéré A est ensuite calculé par addition logarithmique des bandes corrigées.
\end{definition}

La réglementation acoustique du bâtiment utilise systématiquement le dB(A) pour les exigences de niveaux de bruit résiduel.

\subsection{Propagation en champ libre — loi en distance}

En espace libre (source ponctuelle), le niveau de pression sonore décroît avec la distance $r$ :
\[
  L_p = L_w - 20 \cdot \log_{10}(r) - 11 \quad [\si{\decibel}]
\]

Pour une source rayonnant en demi-espace (sur un plan réfléchissant) :
\[
  L_p = L_w - 20 \cdot \log_{10}(r) - 8 \quad [\si{\decibel}]
\]

\begin{methode}
\textbf{Règle pratique :} chaque doublement de la distance réduit $L_p$ de 6~\si{\decibel} pour une source ponctuelle.
\end{methode}

\subsection{Niveau de pression dans un local — formule de Sabine}

Dans un local, la pression sonore résultante combine le champ direct et le champ réverbéré :
\[
  L_p = L_w + 10 \cdot \log_{10}\!\left(\frac{Q}{4\pi r^2} + \frac{4}{R}\right) \quad [\si{\decibel}]
\]
avec :
\begin{itemize}
  \item $Q$ : facteur de directivité de la source (sphère $= 1$, demi-sphère $= 2$, coin $= 8$)
  \item $r$ : distance source--récepteur [\si{\metre}]
  \item $R = \dfrac{S \cdot \bar{\alpha}}{1 - \bar{\alpha}}$ : constante du local [\si{\metre\squared}]
  \item $S$ : surface totale des parois du local [\si{\metre\squared}]
  \item $\bar{\alpha}$ : coefficient d'absorption moyen (sans unité, $0 < \bar{\alpha} < 1$)
\end{itemize}

\textbf{Absorption acoustique :} $A = S \cdot \bar{\alpha}$ exprimée en sabins [\si{\metre\squared}].

% ------------------------------------------------------------
\section{Isolement acoustique des parois}
% ------------------------------------------------------------

\subsection{Transmission du son à travers une paroi}

Lorsqu'une onde sonore frappe une paroi, une partie est réfléchie, une partie est absorbée et une partie est transmise. L'isolement caractérise la résistance de la paroi à cette transmission.

\begin{definition}
\textbf{Indice d'affaiblissement acoustique $R$} (ou \textit{Schalldämm-Maß} en allemand, mesuré en laboratoire) :
\[
  R = L_1 - L_2 + 10 \cdot \log_{10}\!\left(\frac{S}{A}\right) \quad [\si{\decibel}]
\]
avec :
\begin{itemize}
  \item $L_1$ : niveau de pression en salle d'émission [\si{\decibel}]
  \item $L_2$ : niveau de pression en salle de réception [\si{\decibel}]
  \item $S$ : surface de la paroi [\si{\metre\squared}]
  \item $A$ : absorption de la salle de réception [sabins]
\end{itemize}
\end{definition}

\subsection{Loi de masse}

\begin{definition}
La \textbf{loi de masse} établit que l'affaiblissement acoustique d'une paroi homogène croît avec sa masse surfacique $m_s$ [\si{\kilogram\per\metre\squared}] et avec la fréquence :
\[
  R \approx 20 \cdot \log_{10}(m_s \cdot f) - 48 \quad [\si{\decibel}]
\]
avec $f$ la fréquence en \si{\hertz}.
\end{definition}

\textbf{Conséquences pratiques :}
\begin{itemize}
  \item Doubler la masse surfacique apporte environ $+6$~\si{\decibel} d'isolement.
  \item Doubler la fréquence apporte également $+6$~\si{\decibel}.
  \item Les basses fréquences (100--250~\si{\hertz}), typiques des ventilateurs et des vibreurs, sont les plus difficiles à traiter.
\end{itemize}

\subsection{Indice d'affaiblissement pondéré \texorpdfstring{$R_w$}{Rw}}

L'indice $R_w$ est la valeur unique (en dB) représentant la performance globale d'une paroi, déterminée en laboratoire selon NF EN ISO 10140. Il est accompagné de deux termes d'adaptation spectraux :
\begin{itemize}
  \item $C$ : terme pour le bruit rose et le bruit de circulation (fréquences moyennes).
  \item $C_{tr}$ : terme pour le bruit de trafic routier et ferroviaire (basses fréquences dominantes).
\end{itemize}

\begin{center}
\begin{tabular}{lc}
\toprule
Type de paroi & $R_w$ (dB) \\
\midrule
Cloison légère plaque de plâtre 72 mm & 38 \\
Mur béton 15 cm & 48 \\
Mur béton 20 cm & 52 \\
Porte acoustique renforcée & 42--47 \\
Vitrage simple 6 mm & 30 \\
Double vitrage 4/16/4 mm & 32 \\
\bottomrule
\end{tabular}
\end{center}

\subsection{Isolement normalisé \texorpdfstring{$D_{nT,w}$}{DnTw}}

\begin{definition}
$D_{nT,w}$ est l'isolement acoustique mesuré \textit{in situ} entre deux locaux, normalisé par rapport au temps de réverbération de référence $T_0 = 0{,}5$~\si{\second} :
\[
  D_{nT,w} = R_w - 10 \cdot \log_{10}\!\left(\frac{S}{0{,}16 \cdot V / T_0}\right) \quad [\si{\decibel}]
\]
C'est cet indice qui est directement comparé aux exigences de l'arrêté du 30/06/1999.
\end{definition}

\begin{attention}
$D_{nT,w}$ peut être inférieur à $R_w$ si le local de réception est grand et peu absorbant, ou supérieur si le local est petit et très absorbant. La mesure in situ est toujours requise pour la réception réglementaire.
\end{attention}

% ------------------------------------------------------------
\section{Bruit des installations CVC}
% ------------------------------------------------------------

\subsection{Sources de bruit en CVC}

Les installations de chauffage, ventilation et climatisation génèrent plusieurs types de bruit :

\begin{itemize}
  \item \textbf{Bruit aéraulique :} turbulence dans les ventilateurs, les gaines, les singularités (coudes, piquages, registres), les diffuseurs.
  \item \textbf{Bruit mécanique :} vibrations des moteurs, des compresseurs, des pompes — transmis par rayonnement direct et par voie solidienne.
  \item \textbf{Bruit de régulation :} vannes motorisées en chasse, robinets thermostatiques, détendeurs.
\end{itemize}

\subsection{Bruit des ventilateurs}

Le niveau de puissance sonore d'un ventilateur est lié à son débit $q_v$ [\si{\metre\cubed\per\second}] et à sa pression totale $\Delta P_t$ [\si{\pascal}] :

\begin{methode}
\textbf{Formule empirique de Regenscheit (indicative) :}
\[
  L_w \approx K_v + 10 \cdot \log_{10}(q_v) + 20 \cdot \log_{10}(\Delta P_t) \quad [\si{\decibel}]
\]
avec $K_v \approx 31$ pour un ventilateur à aubes rétroverses de bonne qualité (valeur constructeur à préférer).

\textbf{Pratiquement :} toujours utiliser la valeur $L_w$ par bande d'octave fournie par le constructeur.
\end{methode}

\textbf{Niveaux de puissance sonore typiques (ventilateurs CVC) :}

\begin{center}
\begin{tabular}{lcc}
\toprule
Type de ventilateur & $L_w$ global (dB) & Application \\
\midrule
CTA résidentielle (200 m³/h, 80 Pa) & 60--70 & Logement VMC \\
CTA tertiaire (3\,000 m³/h, 200 Pa) & 75--85 & Bureau, commerce \\
Extracteur de toiture & 65--80 & Parking, cuisine \\
Ventilateur axial (grand débit) & 80--95 & Local technique \\
\bottomrule
\end{tabular}
\end{center}

\subsection{Bruit aéraulique dans les gaines}

La vitesse de l'air dans les gaines génère un bruit auto-induit proportionnel à la puissance 5 à 6 de la vitesse.

\begin{methode}
\textbf{Vitesses recommandées pour limiter le bruit :}

\begin{center}
\begin{tabular}{lcc}
\toprule
Zone du réseau & Vitesse max. (m/s) & Commentaire \\
\midrule
Gaine principale (tertiaire) & 6--8 & \\
Gaine principale (logement) & 3--5 & VMC double flux \\
Gaine secondaire & 4--6 & \\
Terminal (bouche, grille) & 2--3 & Critère confort \\
Diffuseur basse vitesse & 1,5--2,5 & Locaux calmes \\
\bottomrule
\end{tabular}
\end{center}
\end{methode}

\begin{attention}
Chaque réduction de vitesse de 20~\% dans une gaine réduit le bruit auto-induit de environ 6~dB. Le surdimensionnement des gaines est la première mesure acoustique à adopter.
\end{attention}

\subsection{Bruit des terminaux et diffuseurs}

Les terminaux (bouches, grilles, diffuseurs, inducteurs) ont une courbe NC (\textit{Noise Criteria}) ou NR (\textit{Noise Rating}) fournie par le constructeur en fonction du débit. Le dimensionnement acoustique consiste à choisir le terminal pour que son niveau NC soit inférieur à l'exigence du local :

\begin{center}
\begin{tabular}{lc}
\toprule
Type de local & Niveau NC recommandé \\
\midrule
Studio d'enregistrement & NC 15--20 \\
Salle de réunion & NC 25--30 \\
Bureau individuel & NC 30--35 \\
Open space, commerce & NC 35--40 \\
Cuisine industrielle & NC 45--50 \\
\bottomrule
\end{tabular}
\end{center}

\subsection{Bruit des pompes et équipements hydrauliques}

Les pompes centrifuges transmettent leurs vibrations aux canalisations par voie solidienne. Les niveaux de puissance sonore des pompes CVC sont généralement de 50 à 75~dB(A) selon la puissance. Les mesures constructives à adopter :

\begin{itemize}
  \item Supports anti-vibratiles (ressorts ou plots silentblocs) entre la pompe et sa semelle.
  \item Manchons souples sur les raccordements hydrauliques.
  \item Passage des canalisations dans des fourreaux avec garnissage souple au franchissement des parois.
\end{itemize}

% ------------------------------------------------------------
\section{Chaîne acoustique en installation CVC}
% ------------------------------------------------------------

\subsection{Bilan acoustique sur réseau aéraulique}

Le niveau sonore en bout de réseau résulte d'une somme algébrique de contributions et d'atténuations successives :

\begin{equation}
  L_p(\text{local}) = L_w(\text{ventilateur}) - \text{Att}_{\text{réseau}} - \text{Att}_{\text{silencieux}} + DI(\text{diffuseur})
\end{equation}

où :
\begin{itemize}
  \item $\text{Att}_{\text{réseau}}$ : atténuation naturelle du réseau (gaines droites, coudes, piquages, expansions) [\si{\decibel}]
  \item $\text{Att}_{\text{silencieux}}$ : atténuation du (des) silencieux [\si{\decibel}]
  \item $DI(\text{diffuseur})$ : indice directivité du diffuseur terminal [\si{\decibel}]
\end{itemize}

\begin{table}[H]
  \centering
  \caption{Atténuations naturelles indicatives des éléments aérauliques (par bande d'octave)}
  \begin{tabular}{lcccc}
    \toprule
    Élément & 125\,Hz & 500\,Hz & 1\,000\,Hz & 2\,000\,Hz \\
    \midrule
    Gaine rectiligne (acier, 1\,m) & 0,1 & 0,1 & 0,1 & 0,1 \\
    Gaine flex isolée (1\,m) & 0,5 & 1,0 & 1,5 & 2,0 \\
    Coude 90° sans directeur & 0 & 3 & 4 & 4 \\
    Piquage droit ($75\,\%$ / $25\,\%$) & — & 6 & 6 & 6 \\
    Réduction de section ($\times$ 2) & 3 & 3 & 3 & 3 \\
    \bottomrule
  \end{tabular}
\end{table}

\begin{methode}
\textbf{Démarche de dimensionnement acoustique d'un réseau CVC :}
\begin{enumerate}
  \item Relever $L_w$ du ventilateur par bande d'octave (données constructeur).
  \item Appliquer les atténuations naturelles du réseau pour chaque bande.
  \item Calculer l'atténuation résiduelle nécessaire pour satisfaire l'exigence du local (niveau NC ou dB(A) cible).
  \item Dimensionner le(s) silencieux en conséquence (longueur, espacement des baffles).
  \item Vérifier que la perte de charge du silencieux est compatible avec le ventilateur.
  \item Vérifier que la vitesse dans le silencieux est $\leq 6$\,m/s (risque de bruit auto-induit).
\end{enumerate}
\end{methode}

\subsection{Niveaux admissibles par type de local (référence NC/NR)}

\begin{table}[H]
  \centering
  \caption{Niveaux de bruit de fond recommandés selon l'usage du local}
  \begin{tabular}{lcc}
    \toprule
    Type de local & Niveau NC recommandé & Niveau $L_p$ dB(A) approx. \\
    \midrule
    Studio radio, salle de concert & NC 15--20 & 25--30 \\
    Salle de réunion, conférence & NC 25--30 & 33--38 \\
    Bureau individuel calme & NC 30--35 & 38--43 \\
    Salle de classe & NC 30--35 & 38--43 \\
    Open space, bureau collectif & NC 35--40 & 43--48 \\
    Hôtel — chambre & NC 30--35 & 38--43 \\
    Restaurant & NC 40--45 & 48--53 \\
    Cuisine professionnelle & NC 45--50 & 53--58 \\
    Local technique CVC & NC 50--60 & 58--68 \\
    \bottomrule
  \end{tabular}
\end{table}

% ------------------------------------------------------------
\section{Traitement acoustique}
% ------------------------------------------------------------

\subsection{Absorption acoustique des matériaux}

\begin{definition}
Le \textbf{coefficient d'absorption} $\alpha$ d'un matériau représente la fraction d'énergie acoustique absorbée ($\alpha = 1$ : absorption totale, $\alpha = 0$ : réflexion totale).

\begin{center}
\begin{tabular}{lcccc}
\toprule
Matériau & \multicolumn{4}{c}{Coefficient $\alpha$ par fréquence} \\
 & 250 Hz & 500 Hz & 1\,000 Hz & 2\,000 Hz \\
\midrule
Béton nu & 0,02 & 0,03 & 0,04 & 0,05 \\
Plaque de plâtre & 0,05 & 0,06 & 0,07 & 0,09 \\
Laine minérale 50 mm & 0,30 & 0,70 & 0,90 & 0,95 \\
Mousse mélamine 50 mm & 0,20 & 0,55 & 0,85 & 0,95 \\
Dalle de faux plafond & 0,40 & 0,65 & 0,75 & 0,80 \\
Moquette épaisse & 0,10 & 0,25 & 0,45 & 0,60 \\
\bottomrule
\end{tabular}
\end{center}
\end{definition}

\subsection{Silencieux à baffles}

\begin{definition}
Un \textbf{silencieux à baffles} est un dispositif intercalé dans la gaine aéraulique, constitué de lames absorbantes parallèles au sens d'écoulement. Il atténue le bruit en forçant l'onde sonore à traverser un chemin tortueux garni de matériaux absorbants (laine de verre, mousse).
\end{definition}

\textbf{Paramètres de dimensionnement :}
\begin{itemize}
  \item \textbf{Longueur active} $L$ : plus elle est grande, plus l'atténuation est élevée.
  \item \textbf{Largeur des canaux libres} $d$ : plus $d$ est faible, plus l'atténuation en basses fréquences est efficace.
  \item \textbf{Perte de charge} : toujours vérifier que la perte de charge du silencieux est compatible avec le réseau ($\Delta P$ typique : 20--80~Pa).
\end{itemize}

\textbf{Atténuations typiques (silencieux 1 m, canaux de 100 mm) :}

\begin{center}
\begin{tabular}{lcccccc}
\toprule
Fréquence (Hz) & 125 & 250 & 500 & 1\,000 & 2\,000 & 4\,000 \\
\midrule
Atténuation (dB) & 5 & 12 & 20 & 25 & 22 & 18 \\
\bottomrule
\end{tabular}
\end{center}

\subsection{Plénum acoustique}

\begin{definition}
Un \textbf{plénum acoustique} est une boîte de connexion entre le ventilateur et la gaine principale, dont les parois intérieures sont recouvertes de matériaux absorbants. Il offre une atténuation importante sur toutes les fréquences grâce à l'expansion du flux et à l'absorption.
\end{definition}

Son atténuation peut être estimée par :
\[
  \text{Att} = 10 \cdot \log_{10}\!\left(\frac{S_{\text{sortie}}}{S_{\text{plénum}} \cdot \bar{\alpha}}\right) \quad [\si{\decibel}]
\]
Un plénum bien garni (laine de verre 50 mm) peut atténuer de 10 à 20~dB selon la fréquence.

\subsection{Caissons insonorisés}

Les groupes de ventilation, compresseurs et pompes peuvent être placés dans des \textbf{caissons insonorisés} (ou capots acoustiques) constitués :
\begin{itemize}
  \item D'une paroi extérieure lourde (acier galvanisé, béton projeté) assurant l'isolement ($R_w \geq 25$~dB).
  \item D'un garnissage intérieur absorbant limitant la réverbération.
  \item De bouches de ventilation du caisson équipées de silencieux.
  \item De passages de gaines avec joints acoustiques.
\end{itemize}

\subsection{Isolement solidien — découplage vibratoire}

\begin{methode}
\textbf{Découplage des équipements vibrants :}
\begin{enumerate}
  \item Interposer des \textbf{plots anti-vibratiles} ou \textbf{ressorts} entre l'équipement et son support (choisir la raideur selon la fréquence d'excitation).
  \item Monter le groupe sur une \textbf{dalle flottante} (dalle béton sur plots silentblocs).
  \item Raccorder les canalisations par des \textbf{manchons souples} (caoutchouc ou EPDM).
  \item Passer les canalisations dans des fourreaux avec garnissage souple.
\end{enumerate}
La fréquence propre du système masse-ressort doit être inférieure à $0{,}5 \times f_{\text{excitation}}$ pour une isolation efficace.
\end{methode}

% ------------------------------------------------------------
\section{Exemples résolus}
% ------------------------------------------------------------

\begin{exemple}[title={Calcul du niveau de pression dans un local technique}]
\textbf{Données :}
\begin{itemize}
  \item CTA double flux : $L_w = 82$~dB (valeur globale constructeur)
  \item Local technique : $V = 80$~\si{\metre\cubed}, $S = 100$~\si{\metre\squared}, $\bar{\alpha} = 0{,}10$
  \item Distance d'évaluation : $r = 2$~\si{\metre}, source au sol ($Q = 2$)
\end{itemize}

\textbf{Calcul de la constante du local :}
\[
  R = \frac{S \cdot \bar{\alpha}}{1 - \bar{\alpha}} = \frac{100 \times 0{,}10}{1 - 0{,}10} = \frac{10}{0{,}90} \approx 11{,}1~\si{\metre\squared}
\]

\textbf{Niveau de pression à 2 m :}
\[
  L_p = 82 + 10 \cdot \log_{10}\!\left(\frac{2}{4\pi \times 4} + \frac{4}{11{,}1}\right) = 82 + 10 \cdot \log_{10}(0{,}0398 + 0{,}360)
\]
\[
  L_p = 82 + 10 \cdot \log_{10}(0{,}400) = 82 - 4{,}0 \approx 78~\si{\decibel}
\]
\textbf{Conclusion :} le niveau de 78~dB est acceptable pour un local technique non occupé en permanence.
\end{exemple}

\begin{exemple}[title={Dimensionnement d'un silencieux}]
\textbf{Situation :} une CTA (bureau open space, exigence NC 35) émet un bruit de 72~dB(A) en sortie de ventilateur. L'atténuation naturelle du réseau est estimée à 12~dB(A). Quelle atténuation doit apporter le silencieux ?

\textbf{Niveau NC 35 $\approx$ 42~dB(A) en réception.}

Atténuation nécessaire du silencieux :
\[
  \text{Att}_{\text{silencieux}} = (72 - 12) - 42 = 18~\si{\decibel(A)}
\]

Un silencieux à baffles de 1~m avec canaux de 100~mm (atténuation 20~dB(A) en moyenne) convient. On vérifie que la perte de charge ajoutée ($\approx 35$~Pa) est compatible avec le ventilateur choisi.
\end{exemple}

\begin{exemple}[title={Vérification réglementaire isolement logement}]
\textbf{Données :}
\begin{itemize}
  \item Paroi séparative entre deux logements : mur béton 20~cm, $R_w = 52$~dB.
  \item Logement réception : $V = 60$~\si{\metre\cubed}, $T = 0{,}6$~\si{\second}, $S_{\text{paroi}} = 12$~\si{\metre\squared}.
\end{itemize}

\textbf{Isolement normalisé :}
\[
  D_{nT,w} = R_w + 10 \cdot \log_{10}(T) - 10 \cdot \log_{10}(T_0) - 10 \cdot \log_{10}\!\left(\frac{S}{0{,}16 \cdot V / T_0}\right)
\]

Formule simplifiée :
\[
  D_{nT,w} = R_w + 10 \cdot \log_{10}\!\left(\frac{T}{T_0}\right) - 10 \cdot \log_{10}\!\left(\frac{S \cdot T_0}{0{,}16 \cdot V}\right)
\]
\[
  D_{nT,w} = 52 + 10 \cdot \log_{10}\!\left(\frac{0{,}6}{0{,}5}\right) - 10 \cdot \log_{10}\!\left(\frac{12 \times 0{,}5}{0{,}16 \times 60}\right)
  = 52 + 0{,}8 - 10 \cdot \log_{10}(0{,}625)
  \approx 52 + 0{,}8 + 2{,}0 = 54{,}8~\si{\decibel}
\]

\textbf{Exigence arrêté 1999 :} $D_{nT,w} \geq 53$~dB pour les bruits aériens entre logements.

\textbf{Conclusion :} la paroi est conforme ($54{,}8 > 53$~dB).
\end{exemple}

\begin{exemple}[title={Dimensionnement acoustique d'une installation VRV — local technique}]
\textbf{Situation :} Un groupe condenseur VRV extérieur ($L_w = 90$~dB) est installé en terrasse, au-dessus d'un bureau dont la façade est à $r = 5$~m du condenseur. La source est posée sur un plan réfléchissant (dalle terrasse), donc $Q = 2$.

\textbf{Niveau en façade bureau (champ libre, demi-espace) :}
\[
  L_p = L_w - 20 \cdot \log_{10}(r) - 8 = 90 - 20 \cdot \log_{10}(5) - 8 = 90 - 14{,}0 - 8 = 68~\text{dB}
\]

\textbf{Exigence réglementaire (logement ou bureau) :} $L_{nA} \leq 30$~dB(A) à l'intérieur — soit, compte tenu d'une façade avec $R_w = 30$~dB (vitrage double standard) :
\[
  L_{intérieur} \approx L_p - R_{façade} = 68 - 30 = 38~\text{dB(A)}
\]

\textbf{L'exigence de 30~dB(A) n'est pas atteinte} — écart de 8~dB(A).

\textbf{Solutions envisagées :}
\begin{enumerate}
  \item Augmenter la distance : déplacer le groupe à $r = 12$~m → $L_p = 90 - 21{,}6 - 8 = 60{,}4$~dB → $L_{int} = 30{,}4$~dB(A) (limite acceptable).
  \item Capot acoustique sur le groupe : atténuation de 8~dB(A) → $L_p = 60$~dB → $L_{int} = 30$~dB(A).
  \item Améliorer la façade : passer à $R_w = 38$~dB (double vitrage renforcé 6/16/6) → $L_{int} = 68 - 38 = 30$~dB(A).
\end{enumerate}

\textbf{Conclusion :} On préférera la solution 3 (amélioration de la façade) si le déplacement de l'unité extérieure est impossible. Elle est souvent moins coûteuse qu'un capot acoustique sur le groupe VRV.
\end{exemple}

% ------------------------------------------------------------
\section{Éléments de réglementation}
% ------------------------------------------------------------

\subsection{Arrêté du 30 juin 1999 — Bâtiments d'habitation}

L'arrêté du 30 juin 1999 fixe les exigences acoustiques minimales pour les logements neufs :

\begin{center}
\begin{tabular}{lcc}
\toprule
Exigence & Indicateur & Valeur minimale \\
\midrule
Isolement bruit aérien entre logements & $D_{nT,w}$ & $\geq 53$~dB \\
Isolement bruit aérien façade & $D_{nT,w}$ & $\geq 30$ à 45~dB (selon zone) \\
Bruit de choc entre logements & $L'_{nT,w}$ & $\leq 58$~dB \\
Bruit d'équipements individuels & $L_{\text{nA}}$ & $\leq 35$~dB(A) \\
Bruit d'équipements collectifs & $L_{\text{nA}}$ & $\leq 30$~dB(A) \\
\bottomrule
\end{tabular}
\end{center}

\begin{attention}
Les équipements CVC (ventilateurs, pompes, extracteurs, groupes de production) sont classés parmi les \textbf{équipements collectifs}. L'exigence de 30~dB(A) en pièces principales de nuit est contraignante et nécessite un soin particulier dans le choix et l'installation des équipements.
\end{attention}

\subsection{NF EN ISO 10140 — Mesures en laboratoire}

La norme NF EN ISO 10140 (série de 5 parties) définit les méthodes de mesure en laboratoire des performances acoustiques des éléments de construction :
\begin{itemize}
  \item \textbf{Partie 1 :} Règles d'application pour des produits spécifiques.
  \item \textbf{Partie 2 :} Mesure de l'isolation aux bruits aériens.
  \item \textbf{Partie 3 :} Mesure de l'isolation aux bruits de choc.
  \item \textbf{Partie 4 :} Procédures et exigences de mesure.
  \item \textbf{Partie 5 :} Exigences pour les installations d'essai et d'équipement.
\end{itemize}

Les valeurs $R_w$ et $L_{n,w}$ sont obtenues selon cette norme et reportées sur les fiches techniques des produits.

\subsection{NF EN ISO 16032 — Bruit des équipements dans les bâtiments}

Cette norme définit la méthode de mesure du niveau de bruit des équipements techniques dans les bâtiments (pompes, ventilateurs, ascenseurs). Elle est complémentaire de l'arrêté du 30/06/1999.

\subsection{NRA — Nouvelle Réglementation Acoustique (2000)}

La NRA est issue de l'arrêté du 30/06/1999 et de ses textes d'application. Elle définit :
\begin{itemize}
  \item Les niveaux d'isolement exigés selon la zone de bruit (I à V) pour la façade.
  \item Les procédures de réception acoustique (contrôle par mesures).
  \item Les dispositions particulières pour les bâtiments autres que les logements (bureaux, établissements d'enseignement, hôtels).
\end{itemize}

Pour les \textbf{établissements d'enseignement} (arrêté du 25/04/2003) :
\begin{itemize}
  \item Isolement entre locaux d'enseignement : $D_{nT,w} \geq 43$~dB.
  \item Durée de réverbération en salles de classe : $T \leq 0{,}7$~\si{\second} (volume $< 250$~\si{\metre\cubed}).
\end{itemize}

\subsection{RE2020 et acoustique}

La Réglementation Environnementale 2020 (RE2020), en vigueur depuis le 1\textsuperscript{er} janvier 2022 pour les logements neufs, traite principalement de la performance énergétique et de l'empreinte carbone. Elle n'introduit pas d'exigences acoustiques nouvelles par rapport à l'arrêté de 1999, mais :

\begin{itemize}
  \item La RE2020 encourage le recours à la ventilation naturelle et double flux. Ces systèmes doivent être compatibles avec les exigences acoustiques existantes.
  \item Les objectifs de compacité et d'isolation thermique renforcée (triple vitrage, murs à forte inertie) bénéficient souvent aussi aux performances acoustiques.
  \item L'analyse de cycle de vie (ACV) intégrée à la RE2020 peut prendre en compte les matériaux d'isolation acoustique et leur empreinte carbone.
\end{itemize}

\begin{methode}
\textbf{Démarche de vérification acoustique en projet GCF :}
\begin{enumerate}
  \item Identifier le type de bâtiment et la réglementation applicable (arrêté 1999, arrêté 2003, etc.).
  \item Déterminer les exigences ($D_{nT,w}$, $L'_{nT,w}$, niveaux d'équipements).
  \item Choisir les parois séparatives ($R_w$ catalogue) et vérifier $D_{nT,w}$ calculé.
  \item Sélectionner les équipements CVC sur critères acoustiques ($L_w$, courbe NC).
  \item Dimensionner les silencieux, plénums et solutions de découplage.
  \item Prévoir la réception acoustique par mesures (contrôle $D_{nT,w}$ in situ).
\end{enumerate}
\end{methode}

% \chapter{Électrotechnique appliquée CVC}

% ============================================================
\section{Objectifs pédagogiques}
% ============================================================

À l'issue de ce chapitre, l'étudiant sera capable de :
\begin{itemize}
  \item Maîtriser les grandeurs électriques fondamentales (tension, courant, puissance active, réactive, apparente, facteur de puissance) et les calculer dans des circuits CVC monophasés et triphasés.
  \item Identifier les moteurs asynchrones triphasés utilisés en CVC (pompes, ventilateurs, compresseurs), lire leur plaque signalétique et choisir leur mode de démarrage.
  \item Dimensionner et justifier le choix d'un variateur de vitesse (VFD), en appliquant les lois de similitude (débit, HMT, puissance).
  \item Sélectionner les organes de protection adaptés (disjoncteurs, fusibles, relais thermiques, différentiels) selon la norme NF C 15-100.
  \item Lire et produire des schémas électriques CVC (puissance et commande) conformes à la norme NF EN 60617.
  \item Établir un bilan de puissance d'une installation CVC complète.
\end{itemize}

\textbf{Savoirs référentiel mobilisés :} S4 (électricité), S5.4 (réglementation électrique), S8-C1 (distribution BT), S8-C2 (moteurs, variateurs).

% ============================================================
\section{Introduction}
% ============================================================

Les installations de chauffage, ventilation et climatisation (CVC) sont fortement consommatrices d'énergie électrique : pompes de circulation, ventilateurs, compresseurs, résistances électriques, régulateurs, actionneurs\dots{} Selon l'Agence Internationale de l'Énergie, les moteurs électriques représentent environ 45\,\% de la consommation mondiale d'électricité, dont une grande part est imputable aux systèmes CVC.

La maîtrise de l'électrotechnique appliquée est donc indispensable pour le technicien GCF, tant pour le dimensionnement correct des équipements que pour la sécurité des personnes et des biens, et pour l'optimisation des consommations énergétiques.

Ce chapitre couvre les fondamentaux électriques, les moteurs, les variateurs de vitesse, la protection des circuits, la lecture des schémas, et se conclut par un bilan de puissance global d'installation CVC.

% ============================================================
\section{Rappels d'électrotechnique}
% ============================================================

\subsection{Grandeurs électriques fondamentales}

\begin{definition}
  \textbf{Tension} $U$ (en volts, V) : différence de potentiel électrique entre deux points d'un circuit. En France, la distribution BT résidentielle est \textbf{230\,V monophasé} (phase--neutre) et \textbf{400\,V triphasé} (entre phases). Fréquence : $f = 50\,\text{Hz}$.

  \textbf{Courant} $I$ (en ampères, A) : débit de charges électriques dans un conducteur.

  \textbf{Résistance} $R$ (en ohms, $\Omega$) : opposition d'un conducteur au passage du courant.
\end{definition}

\subsubsection{Loi d'Ohm}

\begin{equation}
  U = R \cdot I
\end{equation}
avec $U$ en V, $R$ en $\Omega$, $I$ en A.

\subsection{Puissances en courant alternatif}

En courant alternatif, on distingue trois puissances :

\begin{definition}
  \begin{itemize}
    \item \textbf{Puissance active} $P$ (watts, W) : énergie effectivement convertie en travail mécanique ou thermique. C'est la puissance « utile ».
    \item \textbf{Puissance réactive} $Q$ (volt-ampères réactifs, var) : énergie échangée entre le réseau et les éléments inductifs (bobines de moteurs) ou capacitifs. Elle ne produit pas de travail mais circule dans les câbles.
    \item \textbf{Puissance apparente} $S$ (volt-ampères, VA) : puissance totale fournie par le réseau, combinaison de $P$ et $Q$.
  \end{itemize}
\end{definition}

\subsubsection{Circuit monophasé}

\begin{equation}
  S = U \cdot I \quad [\text{VA}]
\end{equation}
\begin{equation}
  P = U \cdot I \cdot \cos\varphi \quad [\text{W}]
\end{equation}
\begin{equation}
  Q = U \cdot I \cdot \sin\varphi \quad [\text{var}]
\end{equation}
\begin{equation}
  S^2 = P^2 + Q^2
\end{equation}

avec $\varphi$ : déphasage entre la tension et le courant (rad ou degrés).

\subsubsection{Circuit triphasé équilibré}

\begin{equation}
  S = \sqrt{3} \cdot U \cdot I \quad [\text{VA}]
\end{equation}
\begin{equation}
  P = \sqrt{3} \cdot U \cdot I \cdot \cos\varphi \quad [\text{W}]
\end{equation}
\begin{equation}
  Q = \sqrt{3} \cdot U \cdot I \cdot \sin\varphi \quad [\text{var}]
\end{equation}

où $U$ est la tension entre phases (V) et $I$ le courant de ligne (A).

\subsection{Facteur de puissance}

\begin{definition}
  Le \textbf{facteur de puissance} (ou $\cos\varphi$) est le rapport entre la puissance active et la puissance apparente :
  \begin{equation}
    \cos\varphi = \frac{P}{S}
  \end{equation}
  Il est sans dimension, compris entre 0 et 1. Un $\cos\varphi$ proche de 1 indique que l'énergie est quasi intégralement utilisée. Un $\cos\varphi$ faible (moteur en charge partielle, éclairage fluo non compensé) signifie un transit inutile de courant réactif qui chauffe les câbles et dégrade la tension.
\end{definition}

\begin{attention}
  Les fournisseurs d'énergie pénalisent les installations dont le $\cos\varphi < 0{,}93$ (tarif HTA). En BTS CVC, on visera $\cos\varphi \geq 0{,}85$ en conception. La compensation est réalisée par des batteries de condensateurs.
\end{attention}

\begin{exemple}[title={Calcul de puissance — pompe de chauffage triphasée}]
  Une pompe de circuit primaire est alimentée en triphasé 400\,V. Son courant absorbé est $I = 8{,}5\,\text{A}$ avec $\cos\varphi = 0{,}82$.

  \textbf{Puissance apparente :}
  \[
    S = \sqrt{3} \times 400 \times 8{,}5 = 5{,}887\,\text{kVA}
  \]
  \textbf{Puissance active :}
  \[
    P = S \times \cos\varphi = 5{,}887 \times 0{,}82 = 4{,}83\,\text{kW}
  \]
  \textbf{Puissance réactive :}
  \[
    Q = \sqrt{S^2 - P^2} = \sqrt{5{,}887^2 - 4{,}83^2} = 3{,}35\,\text{kvar}
  \]
\end{exemple}

% ============================================================
\section{Moteurs électriques en CVC}
% ============================================================

\subsection{Généralités — le moteur asynchrone triphasé}

Le \textbf{moteur asynchrone triphasé} (MAT) est le moteur le plus utilisé en CVC pour entraîner pompes, ventilateurs et compresseurs.

\begin{definition}
  \textbf{Principe :} Le stator (partie fixe) crée un champ magnétique tournant à la \textbf{vitesse de synchronisme} $n_s$. Le rotor (partie tournante) tourne à une vitesse légèrement inférieure $n$ : c'est le \textbf{glissement} $g$.

  \begin{equation}
    n_s = \frac{60 \cdot f}{p} \quad [\text{tr/min}]
  \end{equation}
  avec $f$ la fréquence (Hz) et $p$ le nombre de paires de pôles.

  \begin{equation}
    g = \frac{n_s - n}{n_s} \quad [\text{sans unité, typiquement } 2\,\text{à}\,5\,\%]
  \end{equation}
\end{definition}

\subsubsection{Vitesses de synchronisme à 50\,Hz}

\begin{center}
  \begin{tabular}{cccc}
    \toprule
    Paires de pôles $p$ & Pôles totaux & $n_s$ (tr/min) & $n$ typique (tr/min) \\
    \midrule
    1 & 2 & 3000 & 2900 \\
    2 & 4 & 1500 & 1450 \\
    3 & 6 & 1000 & 960 \\
    4 & 8 & 750  & 720 \\
    \bottomrule
  \end{tabular}
\end{center}

\subsubsection{Plaque signalétique — données essentielles}

\begin{center}
  \begin{tabular}{ll}
    \toprule
    Grandeur & Exemple \\
    \midrule
    Puissance nominale & $P_n = 4\,\text{kW}$ \\
    Tension d'alimentation & 230/400\,V (couplage $\triangle$/$Y$) \\
    Courant nominal & $I_n = 9{,}7\,\text{A}$ (400\,V $Y$) \\
    Fréquence & 50\,Hz \\
    Vitesse nominale & 1450\,tr/min \\
    Facteur de puissance & $\cos\varphi = 0{,}82$ \\
    Rendement & $\eta = 87\,\%$ (classe IE3) \\
    Classe d'isolation & F (155\,$^\circ$C) \\
    Degré de protection & IP55 \\
    \bottomrule
  \end{tabular}
\end{center}

\subsection{Classes d'efficacité énergétique (IE)}

Depuis le règlement UE 640/2009 et ses révisions, les moteurs électriques sont classés selon leur rendement :

\begin{center}
  \begin{tabular}{clll}
    \toprule
    Classe & Dénomination & Rendement typique (4\,kW) & Obligation \\
    \midrule
    IE1 & Standard           & $\approx 84\,\%$ & Interdit depuis 2011 \\
    IE2 & Haut rendement     & $\approx 87\,\%$ & Minimum depuis 2011 \\
    IE3 & Premium            & $\approx 89\,\%$ & Obligatoire $\geq 0{,}75\,\text{kW}$ depuis 2015 \\
    IE4 & Super premium      & $\approx 91\,\%$ & Recommandé, moteurs EC \\
    \bottomrule
  \end{tabular}
\end{center}

\textbf{Référence normative :} NF EN 60034-30-1 (classes IE) et règlement UE 2019/1781.

\begin{attention}
  Depuis le 1\textsuperscript{er} juillet 2021, les moteurs de 0,75 à 1000\,kW doivent atteindre la classe IE3. Les moteurs avec variateur de vitesse associé doivent satisfaire IE2 au minimum (règlement UE 2019/1781).
\end{attention}

\subsection{Démarrage des moteurs}

Le courant de démarrage d'un moteur asynchrone peut atteindre \textbf{5 à 8 fois} le courant nominal, ce qui provoque des chutes de tension sur le réseau et sollicite mécaniquement les organes entraînés.

\subsubsection{Démarrage direct (DOL — Direct On Line)}

\begin{definition}
  Le moteur est connecté directement au réseau. Courant de démarrage : $I_d = 5$ à $8 \times I_n$. Simple et économique, mais choc mécanique élevé et perturbations réseau importantes. \textbf{Réservé aux moteurs de faible puissance ($\leq 5{,}5\,\text{kW}$) ou aux applications non sensibles.}
\end{definition}

\subsubsection{Démarrage étoile-triangle (Y/$\triangle$)}

\begin{definition}
  Au démarrage, les enroulements sont couplés en \textbf{étoile} (tension par enroulement réduite d'un facteur $\sqrt{3}$). Après quelques secondes ($\approx 5$ à 10\,s), la commutation en \textbf{triangle} s'effectue pour le régime normal.

  \begin{itemize}
    \item Courant de démarrage : $I_{d,Y} = \dfrac{I_{d,\triangle}}{3}$ (divisé par 3)
    \item Couple de démarrage : $C_{d,Y} = \dfrac{C_{d,\triangle}}{3}$ (divisé par 3)
  \end{itemize}

  \textbf{Condition d'emploi :} le moteur doit être câblé pour 230/400\,V ($\triangle$/$Y$) et la tension de ligne est 400\,V. La charge doit accepter un faible couple au démarrage (pompes centrifuges, ventilateurs).
\end{definition}

\begin{methode}
  \textbf{Séquence de commande étoile-triangle :}
  \begin{enumerate}
    \item Fermeture du contacteur principal $Q1$ + contacteur étoile $Q3$ : moteur alimenté en Y.
    \item Temporisation (5 à 10\,s selon application).
    \item Ouverture de $Q3$ (étoile), puis fermeture du contacteur triangle $Q2$ avec interverrouillage.
    \item Moteur en fonctionnement normal en $\triangle$.
  \end{enumerate}
  \textbf{Interverrouillage électrique et mécanique entre Q2 et Q3 est obligatoire.}
\end{methode}

\subsubsection{Démarreur progressif (soft-starter)}

\begin{definition}
  Un démarreur progressif pilote la tension appliquée au moteur par des thyristors, réduisant progressivement la tension de 0 à $U_n$ sur une rampe réglable (2 à 30\,s). Le courant de démarrage est limité à \textbf{2 à 4 $\times$ $I_n$} selon le réglage.

  \textbf{Avantages :} démarrage très doux, réduction des à-coups hydrauliques (coup de bélier), peu d'usure mécanique. Idéal pour les pompes de réseau hydraulique.

  \textbf{Inconvénient :} génère des harmoniques ; n'économise pas l'énergie en fonctionnement.
\end{definition}

\subsubsection{Variateur de vitesse (VFD / VSV)}

Le variateur de fréquence (Variable Frequency Drive, ou Variateur de Vitesse) est traité en détail à la section suivante.

% ============================================================
\section{Variation de vitesse et économies d'énergie}
% ============================================================

\subsection{Principe du variateur de fréquence (VFD)}

\begin{definition}
  Un \textbf{variateur de fréquence} (VFD, également appelé variateur de vitesse ou VSV) convertit la tension alternative du réseau (400\,V, 50\,Hz) en une tension alternative de fréquence et d'amplitude variables, permettant de faire varier la vitesse de rotation du moteur.

  \textbf{Composition :}
  \begin{enumerate}
    \item \textbf{Redresseur} : convertit l'AC réseau en DC.
    \item \textbf{Bus DC} (condensateurs de filtrage) : stocke l'énergie.
    \item \textbf{Onduleur MLI} (Modulation de Largeur d'Impulsion) : reconstitue un AC de fréquence variable.
  \end{enumerate}
  La vitesse du moteur est proportionnelle à la fréquence de sortie : $n \propto f_{sortie}$.
\end{definition}

\begin{attention}
  Un VFD génère des \textbf{harmoniques} courants qui perturbent le réseau et peuvent échauffer les câbles et les transformateurs. Des filtres d'entrée (inductances de ligne) sont souvent nécessaires, ainsi qu'un câblage blindé entre VFD et moteur pour limiter les perturbations CEM.
\end{attention}

\subsection{Lois de similitude (affinité des turbomachines)}

Les lois de similitude (ou lois de Rateau) relient les performances d'une pompe ou d'un ventilateur centrifuge à sa vitesse de rotation.

\begin{definition}[\textbf{Lois de similitude}]
  Pour une turbomachine (pompe, ventilateur) fonctionnant à deux vitesses $n_1$ et $n_2$ :
  \begin{align}
    \frac{Q_2}{Q_1} &= \frac{n_2}{n_1} \label{eq:loi1} \tag{Débit} \\[6pt]
    \frac{H_2}{H_1} &= \left(\frac{n_2}{n_1}\right)^2 \label{eq:loi2} \tag{HMT / pression} \\[6pt]
    \frac{P_2}{P_1} &= \left(\frac{n_2}{n_1}\right)^3 \label{eq:loi3} \tag{Puissance absorbée}
  \end{align}
  avec :
  \begin{itemize}
    \item $Q$ : débit volumique ($\text{m}^3/\text{h}$ ou $\text{m}^3/\text{s}$)
    \item $H$ : hauteur manométrique totale (HMT en m) ou pression ($\Delta p$ en Pa)
    \item $P$ : puissance absorbée (W)
    \item $n$ : vitesse de rotation (tr/min)
  \end{itemize}
\end{definition}

\subsection{Économies d'énergie par variation de vitesse}

La loi en \textbf{cube} sur la puissance est la conséquence majeure pour l'économie d'énergie :

\begin{center}
  \begin{tabular}{cccc}
    \toprule
    Vitesse relative $n_2/n_1$ & Débit relatif $Q_2/Q_1$ & Pression relative $H_2/H_1$ & Puissance relative $P_2/P_1$ \\
    \midrule
    100\,\% & 100\,\% & 100\,\% & 100\,\% \\
    80\,\%  & 80\,\%  & 64\,\%  & 51\,\% \\
    60\,\%  & 60\,\%  & 36\,\%  & 22\,\% \\
    50\,\%  & 50\,\%  & 25\,\%  & 12{,}5\,\% \\
    \bottomrule
  \end{tabular}
\end{center}

\textbf{Conclusion :} Une réduction de 20\,\% de la vitesse divise la puissance consommée par près de 2. C'est l'intérêt majeur des VFD sur les ventilateurs de CTA et les pompes à débit variable.

\begin{exemple}[title={VFD sur ventilateur de centrale de traitement d'air (CTA)}]
  Un ventilateur de CTA fonctionne initialement à pleine vitesse :
  \begin{itemize}
    \item Débit nominal : $Q_1 = 10\,000\,\text{m}^3/\text{h}$
    \item Pression nominale : $\Delta p_1 = 800\,\text{Pa}$
    \item Puissance absorbée : $P_1 = 4{,}0\,\text{kW}$
    \item Vitesse : $n_1 = 1450\,\text{tr/min}$
  \end{itemize}
  En régulation de débit (nuit), la consigne est réduite à 60\,\% du débit. La vitesse est abaissée par le VFD à $n_2 = 0{,}6 \times 1450 = 870\,\text{tr/min}$.

  \textbf{Application des lois de similitude :}
  \[
    Q_2 = 0{,}6 \times 10\,000 = 6\,000\,\text{m}^3/\text{h}
  \]
  \[
    \Delta p_2 = 0{,}6^2 \times 800 = 0{,}36 \times 800 = 288\,\text{Pa}
  \]
  \[
    P_2 = 0{,}6^3 \times 4{,}0 = 0{,}216 \times 4{,}0 = 0{,}864\,\text{kW}
  \]
  \textbf{Économie :} $\Delta P = 4{,}0 - 0{,}864 = 3{,}136\,\text{kW}$, soit une réduction de \textbf{78\,\%} de la consommation pour 40\,\% de débit en moins.
\end{exemple}

\subsection{Paramétrage d'un VFD en CVC}

Les paramètres essentiels à configurer lors de la mise en service d'un VFD :

\begin{methode}
  \begin{enumerate}
    \item \textbf{Données moteur :} fréquence nominale (50\,Hz), tension nominale (400\,V), courant nominal ($I_n$), puissance nominale.
    \item \textbf{Rampes d'accélération/décélération :} typiquement 10 à 30\,s pour les pompes (éviter les coups de bélier).
    \item \textbf{Fréquence minimale :} généralement 15 à 25\,Hz (en dessous, refroidissement insuffisant des moteurs à ventilation propre).
    \item \textbf{Fréquence maximale :} 50\,Hz (ou 60\,Hz si surboosting ponctuel autorisé).
    \item \textbf{Source de consigne :} 0–10\,V ou 4–20\,mA depuis le régulateur GTB.
    \item \textbf{Protection moteur intégrée :} seuil de courant maximum, protection thermique électronique.
  \end{enumerate}
\end{methode}

% ============================================================
\section{Protection et appareillage électrique}
% ============================================================

\subsection{Principes de protection}

\begin{definition}
  La protection électrique des circuits CVC vise à protéger :
  \begin{itemize}
    \item Les \textbf{câbles et canalisations} : contre les courts-circuits et les surcharges (échauffement, incendie).
    \item Les \textbf{personnes} : contre les contacts directs et indirects (norme NF C 15-100, IEC 60364).
    \item Les \textbf{matériels} : contre les surtensions, défauts de phase, déséquilibres.
  \end{itemize}
\end{definition}

\subsection{Fusibles}

\begin{definition}
  Les \textbf{fusibles} protègent uniquement contre les \textbf{courts-circuits}. Leur calibre ($I_n$) est choisi selon le courant admissible de la canalisation. Ils sont classés selon leur pouvoir de coupure et leur type (gG pour usage général, aM pour protection moteur).

  \textbf{Type gG} (usage général) : protège câbles et circuits.\\
  \textbf{Type aM} (accompagnement moteur) : laisse passer la pointe de courant au démarrage.

  \textbf{Référence normative :} IEC 60269 / NF EN 60269.
\end{definition}

\subsection{Disjoncteurs}

\begin{definition}
  Un \textbf{disjoncteur} assure la protection contre les \textbf{surcharges} (déclencheur thermique) et les \textbf{courts-circuits} (déclencheur magnétique). Il est réarmable manuellement après déclenchement.

  \begin{itemize}
    \item \textbf{Déclencheur thermique :} bimétallique, protège en cas de surcharge durable ($I > I_n$).
    \item \textbf{Déclencheur magnétique :} à action instantanée, déclenche pour les forts courts-circuits ($I \gg I_n$).
  \end{itemize}

  En CVC, on utilise les \textbf{disjoncteurs moteurs} (type GV de Schneider ou MS de ABB) réglables selon le courant nominal du moteur.

  \textbf{Référence normative :} IEC 60947-2 (disjoncteurs industriels), IEC 60947-4-1 (démarreurs moteur).
\end{definition}

\subsection{Relais thermiques}

\begin{definition}
  Le \textbf{relais thermique} est associé à un contacteur pour protéger un moteur contre les \textbf{surcharges prolongées}. Il mesure le courant par effet thermique (bimétallique) et ouvre un contact auxiliaire qui coupera la bobine du contacteur.

  Son réglage est effectué à $I_n$ du moteur (courant nominal de la plaque signalétique).

  \textbf{Attention :} le relais thermique seul ne protège pas contre les courts-circuits — il doit être associé à des fusibles ou un disjoncteur.
\end{definition}

\subsection{Dispositifs différentiels}

\begin{definition}
  Un \textbf{dispositif différentiel} (DDR — Dispositif Différentiel Résiduel) détecte le courant de fuite vers la terre (contact avec une masse ou un conducteur actif). Il déclenche si le courant de fuite dépasse le seuil $I_{\Delta n}$.

  \begin{itemize}
    \item \textbf{30\,mA} : protection des personnes (obligatoire pour les circuits prises et circuits alimentant des équipements en zone humide — locaux CVC souvent concernés).
    \item \textbf{300\,mA ou 500\,mA} : protection incendie.
    \item \textbf{1\,A et au-delà} : protection de tête de tableau (sélectivité).
  \end{itemize}

  \textbf{Référence normative :} IEC 60755, NF EN 61008 (interrupteurs différentiels), NF C 15-100.
\end{definition}

\subsection{Choix des sections de câbles (NF C 15-100)}

\begin{methode}
  \textbf{Procédure simplifiée de choix de section :}
  \begin{enumerate}
    \item Calculer le courant d'emploi $I_B$ (courant en service normal).
    \item Choisir le courant nominal du dispositif de protection $I_n \geq I_B$.
    \item Choisir une section de câble telle que le courant admissible $I_z \geq I_n$.
    \item Appliquer les \textbf{facteurs de correction} $k$ selon le mode de pose (chemins de câbles, fourreaux, etc.) et la température ambiante.
    \item Vérifier la \textbf{chute de tension} : $\leq 3\,\%$ pour les circuits terminaux (NF C 15-100).
  \end{enumerate}

  \textbf{Formule de chute de tension (circuit triphasé) :}
  \begin{equation}
    \Delta U = \frac{\sqrt{3} \cdot \rho \cdot L \cdot I}{S} \quad [\text{V}]
  \end{equation}
  avec $\rho$ : résistivité du cuivre ($\approx 0{,}0225\,\Omega\cdot\text{mm}^2/\text{m}$ à 20\,$^\circ$C), $L$ : longueur du câble (m), $I$ : courant (A), $S$ : section (mm$^2$).
\end{methode}

\begin{exemple}[title={Choix de la section de câble — moteur pompe 4\,kW}]
  Données : moteur 4\,kW, 400\,V triphasé, $\cos\varphi = 0{,}82$, $\eta = 87\,\%$, longueur de câble $L = 25\,\text{m}$, pose sur chemin de câbles (facteur de correction $k = 1$).

  \textbf{Courant d'emploi :}
  \[
    I_B = \frac{P}{\sqrt{3} \cdot U \cdot \cos\varphi \cdot \eta} = \frac{4\,000}{\sqrt{3} \times 400 \times 0{,}82 \times 0{,}87} = \frac{4\,000}{494{,}5} \approx 8{,}1\,\text{A}
  \]

  \textbf{Choix du disjoncteur moteur :} calibré à $I_n = 10\,\text{A}$ (cran normalisé au-dessus de $I_B$).

  \textbf{Section minimale :} $I_z \geq 10\,\text{A}$ sur câble 3G2,5\,mm$^2$ Cu ($I_z = 21\,\text{A}$ — suffisant).

  \textbf{Vérification chute de tension :}
  \[
    \Delta U = \frac{\sqrt{3} \times 0{,}0225 \times 25 \times 8{,}1}{2{,}5} = \frac{7{,}90}{2{,}5} = 3{,}16\,\text{V}
  \]
  \[
    \Delta U\,\% = \frac{3{,}16}{400} \times 100 = 0{,}79\,\% \leq 3\,\% \quad \checkmark
  \]
\end{exemple}

% ============================================================
\section{Schémas électriques CVC}
% ============================================================

\subsection{Normalisation des schémas — NF EN 60617}

\begin{definition}
  Les schémas électriques industriels et CVC sont réalisés selon la norme \textbf{NF EN 60617} (symboles graphiques pour schémas électriques). Les principaux symboles utilisés en CVC sont :

  \begin{center}
    \begin{tabular}{ll}
      \toprule
      Symbole & Désignation \\
      \midrule
      Rectangle (3 traits verticaux) & Moteur triphasé asynchrone \\
      Contact fermé en repos & Contact normalement fermé (NF) \\
      Contact ouvert en repos & Contact normalement ouvert (NO) \\
      Cercle avec A & Bobine de contacteur \\
      Zigzag & Résistance / relais thermique \\
      Carré barré en diagonale & Fusible \\
      $\Delta U$ avec trait & Disjoncteur \\
      \bottomrule
    \end{tabular}
  \end{center}
\end{definition}

\subsection{Schéma de puissance}

Le \textbf{schéma de puissance} représente le circuit principal d'alimentation du moteur : réseau triphasé, disjoncteur/fusibles, contacteur(s), relais thermique, bornes moteur.

\begin{definition}
  Composants types d'un schéma de puissance pompe CVC (démarrage direct) :
  \begin{enumerate}
    \item Jeu de barres triphasé L1, L2, L3
    \item Disjoncteur moteur $Q1$ (réglé à $I_n$ moteur)
    \item Contacteur $KM1$ (3 pôles principaux)
    \item Relais thermique $F1$ (réglé à $I_n$ moteur)
    \item Bornes moteur U1, V1, W1
  \end{enumerate}
\end{definition}

\subsection{Schéma de commande}

Le \textbf{schéma de commande} représente le circuit basse tension (souvent 24\,V AC ou DC, ou 230\,V AC) qui pilote les bobines des contacteurs et les signalisations.

\begin{definition}
  Éléments types d'un schéma de commande (marche/arrêt pompe) :
  \begin{enumerate}
    \item Alimentation commande (ex. 24\,V DC depuis alimentation TBTS)
    \item Contact NF du relais thermique $F1$ (protection thermique)
    \item Contact NF de l'arrêt d'urgence $S0$ (bouton coup de poing)
    \item Bouton-poussoir \textbf{MARCHE} $S1$ (contact NO)
    \item Bouton-poussoir \textbf{ARRÊT} $S2$ (contact NF)
    \item Contact d'auto-maintien du contacteur $KM1$ (contact NO en parallèle de $S1$)
    \item Bobine du contacteur $KM1$
    \item Voyant de signalisation \textbf{EN MARCHE}
  \end{enumerate}
\end{definition}

\begin{methode}
  \textbf{Logique de commande — séquence de marche/arrêt :}
  \begin{enumerate}
    \item Appui sur $S1$ (MARCHE) : courant dans la bobine $KM1$ si $F1$ NF et $S0$ NF.
    \item $KM1$ s'excite : contacts principaux ferment (puissance) + contact d'auto-maintien ferme.
    \item Relâcher $S1$ : le courant continue par l'auto-maintien.
    \item Appui sur $S2$ (ARRÊT) : coupure de l'auto-maintien, $KM1$ retombe.
    \item En cas de surcharge : $F1$ déclenche (contact NF s'ouvre), $KM1$ retombe automatiquement.
  \end{enumerate}
\end{methode}

\subsection{Commande par GTB/GTC}

En CVC moderne, la commande locale (boutons-poussoirs) est souvent remplacée ou supervisée par un automate ou une GTB. Le contacteur est alors commandé par un \textbf{contact sec} issu d'un automate (sortie TOR 24\,V DC) ou directement par un \textbf{variateur de vitesse} qui intègre les fonctions de protection.

% ============================================================
\section{Exemples résolus}
% ============================================================

\begin{exemple}[title={Bilan de puissance électrique d'une installation CVC complète}]
  \textbf{Données de l'installation :} immeuble de bureaux, 2000\,m$^2$, système CVC comprenant :

  \begin{center}
    \begin{tabular}{lcccc}
      \toprule
      Équipement & Nb. & $P_n$ unit. (kW) & $\cos\varphi$ & $P_{abs}$ totale (kW) \\
      \midrule
      Pompe circuit primaire & 2 & 4{,}0 & 0{,}82 & 8{,}0 \\
      Pompe circuit secondaire & 4 & 1{,}5 & 0{,}80 & 6{,}0 \\
      Ventilateur CTA soufflage & 2 & 7{,}5 & 0{,}84 & 15{,}0 \\
      Ventilateur CTA reprise & 2 & 5{,}5 & 0{,}83 & 11{,}0 \\
      Extracteur WC & 6 & 0{,}18 & 0{,}75 & 1{,}08 \\
      Régulateurs/GTB & — & — & — & 0{,}5 \\
      \midrule
      \textbf{Total} & & & & \textbf{41{,}58} \\
      \bottomrule
    \end{tabular}
  \end{center}

  \textbf{Note :} cette puissance est la somme des puissances absorbées nominales. En pratique, un coefficient de foisonnement de 0,75 à 0,85 est appliqué (tous les équipements ne fonctionnent pas simultanément à pleine charge).

  \textbf{Puissance de bilan avec foisonnement ($k_f = 0{,}80$) :}
  \[
    P_{bilan} = 41{,}58 \times 0{,}80 = 33{,}3\,\text{kW}
  \]

  \textbf{Puissance réactive totale estimée} (avec $\cos\varphi$ moyen $\approx 0{,}82$) :
  \[
    Q = P \cdot \tan\varphi = 33{,}3 \times \tan(\arccos 0{,}82) = 33{,}3 \times 0{,}698 = 23{,}2\,\text{kvar}
  \]

  \textbf{Puissance apparente totale :}
  \[
    S = \frac{P}{\cos\varphi} = \frac{33{,}3}{0{,}82} = 40{,}6\,\text{kVA}
  \]

  \textbf{Courant total absorbé (triphasé 400\,V) :}
  \[
    I_{total} = \frac{S}{\sqrt{3} \cdot U} = \frac{40\,600}{\sqrt{3} \times 400} = \frac{40\,600}{693} = 58{,}6\,\text{A}
  \]

  Cela permet de dimensionner le disjoncteur général de l'armoire CVC et la section du câble d'alimentation depuis le TGBT.
\end{exemple}

\begin{exemple}[title={Choix et économie d'un VFD sur pompe secondaire}]
  Une pompe de circuit secondaire de chauffage est équipée d'un VFD. Elle fonctionne :
  \begin{itemize}
    \item 8\,h/j à 100\,\% (pleine charge hivernale) : $P_1 = 1{,}5\,\text{kW}$
    \item 10\,h/j à 70\,\% de vitesse (mi-saison) : $P_2 = 1{,}5 \times (0{,}7)^3 = 0{,}515\,\text{kW}$
    \item 6\,h/j à 50\,\% de vitesse (nuit) : $P_3 = 1{,}5 \times (0{,}5)^3 = 0{,}188\,\text{kW}$
  \end{itemize}

  \textbf{Consommation journalière avec VFD :}
  \[
    W_{jour} = (1{,}5 \times 8) + (0{,}515 \times 10) + (0{,}188 \times 6) = 12 + 5{,}15 + 1{,}13 = 18{,}28\,\text{kWh}
  \]

  \textbf{Sans VFD (vitesse fixe toute la journée) :}
  \[
    W_{jour,ref} = 1{,}5 \times (8 + 10 + 6) = 1{,}5 \times 24 = 36\,\text{kWh}
  \]

  \textbf{Économie journalière :}
  \[
    \Delta W = 36 - 18{,}28 = 17{,}72\,\text{kWh/jour}
  \]
  soit une réduction de \textbf{49\,\%} de la consommation de la pompe, sur 7 mois de chauffe ($\approx 210$ jours) :
  \[
    \Delta W_{annuel} = 17{,}72 \times 210 \approx 3\,721\,\text{kWh/an}
  \]
  Au prix de 0,18\,€/kWh : économie $\approx \textbf{670\,€/an}$ par pompe.
\end{exemple}

% ============================================================
\section{Exercices d'application}
% ============================================================

\exercice{Calcul de puissances — moteur de compresseur VRV}
Un compresseur de système VRV est alimenté en triphasé 400\,V. Sa plaque indique : $P_n = 11\,\text{kW}$, $I_n = 23\,\text{A}$, $\cos\varphi = 0{,}88$, rendement $\eta = 91\,\%$.
\begin{enumerate}
  \item Calculer la puissance apparente $S$ et la puissance réactive $Q$ absorbées par le moteur.
  \item Vérifier la cohérence avec la puissance active $P_n$ (tenir compte du rendement).
  \item Quelle valeur de condensateur (en kvar) permettrait de relever le $\cos\varphi$ à 0,95\,?
\end{enumerate}

\exercice{Démarrage étoile-triangle — ventilateur de désenfumage}
Un ventilateur de désenfumage est équipé d'un moteur 15\,kW, 400\,V, $I_n = 29\,\text{A}$, courant de démarrage en triangle $I_{d,\triangle} = 7 \times I_n$.
\begin{enumerate}
  \item Calculer le courant de démarrage direct et le courant de démarrage en étoile.
  \item Justifier l'intérêt du démarrage étoile-triangle pour ce type d'application.
  \item Quel dispositif de protection choisir (disjoncteur ou fusible aM) et quel calibre\,?
\end{enumerate}

\exercice{Lois de similitude — économies par variateur}
Un ventilateur de CTA a les caractéristiques nominales : $Q_1 = 15\,000\,\text{m}^3/\text{h}$, $\Delta p_1 = 1\,200\,\text{Pa}$, $P_1 = 8\,\text{kW}$ à $n_1 = 1\,450\,\text{tr/min}$.
\begin{enumerate}
  \item Calculer $Q_2$, $\Delta p_2$ et $P_2$ si la vitesse est réduite à 75\,\%.
  \item Exprimer l'économie en \% sur la puissance absorbée.
  \item Sur une année de fonctionnement (4\,000\,h à 100\,\% + 4\,000\,h à 75\,\%), calculer l'économie d'énergie annuelle en kWh et en euros (tarif 0,18\,€/kWh).
\end{enumerate}

\exercice{Dimensionnement d'armoire CVC}
Une armoire électrique CVC doit alimenter les équipements suivants (triphasé 400\,V) :
\begin{itemize}
  \item 2 pompes : $P = 2{,}2\,\text{kW}$ chacune, $\cos\varphi = 0{,}80$
  \item 1 ventilateur : $P = 5{,}5\,\text{kW}$, $\cos\varphi = 0{,}83$
  \item 1 groupe de froid : $P = 18\,\text{kW}$, $\cos\varphi = 0{,}86$
\end{itemize}
\begin{enumerate}
  \item Calculer la puissance active totale et la puissance apparente totale (coefficient de foisonnement : 0,85).
  \item Déduire le courant total absorbé et calibrer le disjoncteur général.
  \item Vérifier si une batterie de condensateurs de 10\,kvar améliore significativement le $\cos\varphi$ global.
\end{enumerate}

% ============================================================
\section{Éléments de réglementation}
% ============================================================

\subsection{NF C 15-100 — Installations électriques BT}

La norme \textbf{NF C 15-100} (transposition française de la norme IEC 60364) fixe les règles de conception, réalisation et vérification des installations électriques basse tension. Elle s'applique à toutes les installations CVC raccordées au réseau BT.

Points clés pour le technicien CVC :
\begin{itemize}
  \item \textbf{Schémas de liaison à la terre (SLT) :} TT (résidentiel, neutre mis à la terre — disjoncteur + DDR 30\,mA), TN-S (industriel, conducteur PE séparé), IT (hôpitaux, salles propres).
  \item \textbf{Sections minimales} des conducteurs selon le courant admissible et le mode de pose.
  \item \textbf{Chute de tension maximale} : 3\,\% pour les circuits terminaux d'éclairage et de force motrice.
  \item \textbf{Degrés de protection IP} des matériels selon l'environnement (locaux techniques CVC : IP55 minimum pour les matériels exposés).
  \item \textbf{Protection différentielle} : DDR 30\,mA obligatoire pour tous les circuits alimentant des équipements en locaux humides (chaufferies, locaux de traitement d'air avec buées, etc.).
\end{itemize}

\subsection{IEC 60947 — Appareillage industriel}

La famille de normes \textbf{IEC 60947} couvre les dispositifs de protection et de commande industriels :
\begin{itemize}
  \item \textbf{IEC 60947-1} : règles générales.
  \item \textbf{IEC 60947-2} : disjoncteurs industriels.
  \item \textbf{IEC 60947-4-1} : contacteurs et démarreurs moteur (y compris relais thermiques).
  \item \textbf{IEC 60947-4-2} : démarreurs progressifs (soft-starters).
  \item \textbf{IEC 60947-7} : variateurs de vitesse (VFD).
\end{itemize}

\subsection{NF EN 60034 — Machines électriques tournantes}

La norme \textbf{NF EN 60034} (IEC 60034) définit les caractéristiques des moteurs électriques :
\begin{itemize}
  \item \textbf{NF EN 60034-1} : valeurs nominales et caractéristiques.
  \item \textbf{NF EN 60034-5} : degrés de protection (IP).
  \item \textbf{NF EN 60034-30-1} : classes d'efficacité (IE1 à IE4).
\end{itemize}

\subsection{Règlement UE 2019/1781 — Écoconception des moteurs}

Ce règlement européen d'application directe en France impose :
\begin{itemize}
  \item Depuis juillet 2021 : moteurs 0,75–1000\,kW $\rightarrow$ classe IE3 minimum (ou IE2 avec VFD obligatoirement associé).
  \item Depuis juillet 2023 : extension à 0,12–0,75\,kW.
  \item Les moteurs à vitesse variable (avec VFD intégré) doivent satisfaire IE2 au minimum.
\end{itemize}

\begin{attention}
  \textbf{Habilitation électrique (NF C 18-510) :} tout travail ou intervention sur une installation électrique nécessite une habilitation adaptée. En CVC, le technicien effectuant des travaux sur les armoires électriques doit au minimum disposer d'une habilitation \textbf{B1V} (travaux hors tension, voisinage de pièces nues sous tension). Les habilitations sont délivrées par l'employeur après formation.
\end{attention}

\subsection{Marquage CE et conformité}

Tout équipement électrique mis sur le marché européen doit porter le \textbf{marquage CE}, attestant la conformité aux directives applicables :
\begin{itemize}
  \item \textbf{Directive Basse Tension} (2014/35/UE) : sécurité électrique des matériels.
  \item \textbf{Directive CEM} (2014/30/UE) : compatibilité électromagnétique (particulièrement critique pour les VFD).
  \item \textbf{Directive Écoconception} (2009/125/CE) et règlements associés pour les moteurs.
\end{itemize}

% \chapter{Études de cas et dimensionnement}

% ============================================================
\section{Objectifs pédagogiques}
% ============================================================

À l'issue de ce chapitre, l'étudiant sera capable de :
\begin{itemize}
  \item Conduire une démarche complète de dimensionnement CVC en partant du cahier des charges ;
  \item Calculer les déperditions thermiques d'un bâtiment selon la méthode réglementaire (RE2020 / NF EN 12831) ;
  \item Sélectionner et dimensionner les équipements de production, distribution et émission de chaleur ;
  \item Dimensionner un réseau aéraulique et une centrale de traitement d'air (CTA) ;
  \item Appliquer les critères d'efficacité énergétique et vérifier la conformité réglementaire ;
  \item Intégrer la régulation et la GTB dans une étude globale ;
  \item Rédiger une note de calcul justifiée à présenter en situation professionnelle.
\end{itemize}

\textbf{Compétences mobilisées :} C1, C2, C3, C5, C7, C8, C9 \quad
\textbf{Savoirs associés :} S5.1, S8-A2, S8-A3, S8-A4, S8-A5, S8-B1, S8-B2, S8-B4, S8-B11

% ============================================================
\section{Introduction}
% ============================================================

Le technicien supérieur GCF est régulièrement amené à produire des études techniques complètes : appels d'offres, avant-projets sommaires (APS), avant-projets définitifs (APD), dossiers de consultation des entreprises (DCE). Ces études intègrent plusieurs disciplines enseignées séparément dans les chapitres précédents.

Ce chapitre propose trois études de cas représentatives :
\begin{enumerate}
  \item Un \textbf{immeuble résidentiel RE2020} — déperditions, production chaleur par PAC, distribution hydraulique, régulation ;
  \item Un \textbf{bâtiment tertiaire climatisé} — CTA double flux, réseau aéraulique, production froid, GTB ;
  \item Une \textbf{installation industrielle} — process thermique, sécurités, efficacité énergétique.
\end{enumerate}

\begin{definition}[title={Note de calcul}]
  La note de calcul est le document technique central d'une étude CVC. Elle regroupe :
  les données d'entrée (géométrie, occupation, conditions climatiques), les hypothèses
  retenues, les formules appliquées avec leurs références normatives, les résultats
  numériques et la sélection du matériel. Elle doit être \textbf{traçable, vérifiable et justifiée}.
\end{definition}

% ============================================================
\section{Méthodologie de dimensionnement CVC}
% ============================================================

\subsection{Démarche générale}

\begin{methode}[title={Étapes du dimensionnement CVC}]
  \begin{enumerate}
    \item \textbf{Collecte des données} : plans architectural, orientation, localisation climatique, destination des locaux, réglementation applicable.
    \item \textbf{Bilan thermique d'hiver} : calcul des déperditions par transmission et renouvellement d'air (NF EN 12831).
    \item \textbf{Bilan thermique d'été} : calcul des charges de refroidissement (apports solaires, internes, renouvellement).
    \item \textbf{Sélection des générateurs} : chaudière, PAC, groupe froid — selon besoins de pointe et modulabilité.
    \item \textbf{Distribution} : réseau hydraulique (dimensionnement des tuyauteries, pompes, équilibrage) et/ou aéraulique (gaines, ventilateurs, diffuseurs).
    \item \textbf{Émission} : radiateurs, planchers chauffants, ventilo-convecteurs, soufflage.
    \item \textbf{Régulation et GTB} : stratégies de commande, sondes, actionneurs.
    \item \textbf{Vérification RE2020} : calcul Bbio, Cep, DH.
    \item \textbf{Rédaction de la note de calcul} et sélection du matériel sur catalogue.
  \end{enumerate}
\end{methode}

\subsection{Données climatiques de référence}

Les températures extérieures de base $T_{ext,base}$ sont définies par zone climatique (NF EN 12831 et données Météo-France) :

\begin{center}
\begin{tabular}{llcc}
  \toprule
  Zone & Exemples de villes & $T_{ext,base}$ (\si{\celsius}) & DJU (base 18) \\
  \midrule
  H1a & Paris, Lille         & $-7$  & 2\,400 \\
  H1b & Reims, Nancy         & $-9$  & 2\,800 \\
  H1c & Strasbourg           & $-11$ & 3\,100 \\
  H2a & Nantes, La Rochelle  & $-4$  & 1\,800 \\
  H2b & Lyon, Grenoble       & $-8$  & 2\,400 \\
  H2c & Bordeaux, Toulouse   & $-3$  & 1\,600 \\
  H2d & Nice, Marseille      & 0     & 1\,200 \\
  H3   & Guadeloupe, Réunion  & +18  & 0 \\
  \bottomrule
\end{tabular}
\end{center}

\subsection{Rappel des formules fondamentales}

\subsubsection{Déperditions par transmission}

\begin{equation}
  \Phi_{T} = H_T \cdot (T_{int} - T_{ext,base})
\end{equation}

avec $H_T = \sum_{i} U_i \cdot A_i$ le coefficient de déperdition global par transmission (\si{\watt\per\kelvin}).

\subsubsection{Déperditions par ventilation et infiltrations}

\begin{equation}
  \Phi_{V} = H_V \cdot (T_{int} - T_{ext,base})
\end{equation}

avec $H_V = \rho_{air} \cdot c_{p,air} \cdot q_v$ (\si{\watt\per\kelvin}), où $q_v$ est le débit volumique d'air neuf (\si{\cubic\meter\per\second}).

\subsubsection{Puissance de chauffage de pointe}

\begin{equation}
  \Phi_{ch} = \Phi_{T} + \Phi_{V} \quad [\si{\watt}]
\end{equation}

\subsubsection{Coefficient de performance de la PAC}

\begin{equation}
  \mathrm{COP} = \frac{P_{calorifique}}{P_{electrique}}
\end{equation}

\subsubsection{Perte de charge en réseau hydraulique}

\begin{equation}
  \Delta P = \lambda \cdot \frac{L}{D} \cdot \frac{\rho v^2}{2} + \sum \xi_i \cdot \frac{\rho v^2}{2} \quad [\si{\pascal}]
\end{equation}

\subsubsection{Puissance frigorifique nécessaire}

\begin{equation}
  Q_f = Q_{apports\_solaires} + Q_{apports\_internes} + Q_{ventilation}
\end{equation}

% ============================================================
\section{Étude de cas 1 — Immeuble résidentiel RE2020}
% ============================================================

\subsection{Présentation du projet}

\begin{exemple}[title={Contexte}]
  Immeuble résidentiel collectif R+3, zone climatique H1a (Paris), altitude \SI{75}{\metre}.
  Surface habitable totale : \SI{1200}{\metre\squared} (12 appartements de \SI{100}{\metre\squared}).
  Destination : logements sociaux BBC-RE2020.
  Système CVC retenu : PAC air/eau collective + planchers chauffants + VMC double flux.
\end{exemple}

\subsection{Données architecturales}

\begin{center}
\begin{tabular}{llccc}
  \toprule
  Paroi & Description & Surface (\si{\metre\squared}) & $U$ (\si{\watt\per\metre\squared\per\kelvin}) & $U \cdot A$ (\si{\watt\per\kelvin}) \\
  \midrule
  Murs ext.    & Béton + isolation ITE 14 cm & 620 & 0,22 & 136,4 \\
  Toiture      & Terrasse accessible isolée   & 300 & 0,15 & 45,0  \\
  Plancher bas & Dalle sur terre-plein        & 300 & 0,30 & 90,0  \\
  Vitrages     & Double vitrage Uw 1,2        & 180 & 1,20 & 216,0 \\
  Portes ext.  & Porte palière isolée         & 30  & 1,40 & 42,0  \\
  \midrule
  \textbf{Total} & & & & \textbf{529,4} \\
  \bottomrule
\end{tabular}
\end{center}

\subsection{Calcul des déperditions}

\textbf{Données :} $T_{int} = \SI{19}{\celsius}$, $T_{ext,base} = \SI{-7}{\celsius}$, $\Delta T = \SI{26}{\kelvin}$.

\textbf{Renouvellement d'air :}
\begin{itemize}
  \item Débit VMC double flux : $q_v = \SI{0,5}{\per\hour} \times \SI{3600}{\metre\cubed} = \SI{1800}{\metre\cubed\per\hour}$
  \item Rendement échangeur : $\eta = 0,85$ (75\,\% récupéré)
  \item $H_V = 0,34 \times 1800 \times (1 - 0,85) = \SI{91,8}{\watt\per\kelvin}$
\end{itemize}

\textbf{Coefficient de déperdition total :}
\begin{equation}
  H_{total} = H_T + H_V = 529{,}4 + 91{,}8 = \SI{621{,}2}{\watt\per\kelvin}
\end{equation}

\textbf{Puissance de chauffage de pointe :}
\begin{equation}
  \Phi_{ch} = 621{,}2 \times 26 = \SI{16{,}15}{\kilo\watt}
\end{equation}

\begin{attention}
  La puissance de pointe calculée correspond aux déperditions en conditions hivernales extrêmes.
  Elle ne prend pas en compte les apports solaires et internes, qui sont favorables en hiver
  mais doivent être vérifiés en été pour le confort d'été (calcul des charges de refroidissement).
\end{attention}

\subsection{Sélection de la PAC}

\begin{methode}[title={Critères de sélection d'une PAC collective}]
  \begin{enumerate}
    \item Puissance calorifique $\geq \Phi_{ch}$ (avec marge de 15\,\% pour pertes réseau)
    \item COP $\geq 3{,}2$ à $-7/35$ (SCOP annuel selon EN 14825)
    \item Fluide frigorigène à faible GWP ($<$ 750 — exigences F-gas 2024)
    \item Compatibilité avec plancher chauffant (eau à \SI{35}{\celsius} max)
    \item Niveau acoustique $\leq$ 50 dB(A) à \SI{1}{\metre} (mitoyenneté)
  \end{enumerate}
\end{methode}

\textbf{Puissance installée avec marge :}
\begin{equation}
  P_{PAC} = 16{,}15 \times 1{,}15 = \SI{18{,}6}{\kilo\watt} \rightarrow \text{sélection : PAC 20 kW}
\end{equation}

\textbf{Consommation électrique annuelle estimée :}
\begin{equation}
  W_{el} = \frac{P_{calorifique} \times DJU \times 24}{SCOP \times \Delta T_{base}} = \frac{20 \times 2400 \times 24}{3{,}5 \times 26} \approx \SI{12\,615}{\kilo\watt\hour\per\year}
\end{equation}

\subsection{Dimensionnement du plancher chauffant}

Le plancher chauffant basse température (PCBT) délivre la chaleur par rayonnement depuis la dalle.

\textbf{Puissance surfacique utile :}
\begin{equation}
  q_{plancher} = \frac{\Phi_{ch}}{A_{totale}} = \frac{16\,150}{1200} = \SI{13{,}5}{\watt\per\metre\squared}
\end{equation}

\textbf{Température de départ eau :} \SI{35}{\celsius} (retour \SI{28}{\celsius}), compatible PAC.

\textbf{Pas de pose des tubes :} espacement $e = \SI{15}{\centi\metre}$ (standard pour habitat). Débit par circuit :
\begin{equation}
  \dot{m}_{circuit} = \frac{q_{plancher} \times A_{circuit}}{\rho \cdot c_p \cdot \Delta T_{eau}} = \frac{13{,}5 \times 15}{1000 \times 4186 \times 7 / 3600} \approx \SI{37{,}4}{\litre\per\hour}
\end{equation}

\subsection{Réseau hydraulique}

\begin{center}
\begin{tabular}{lcccc}
  \toprule
  Tronçon & Débit (\si{\metre\cubed\per\hour}) & $\Delta P_{lin}$ (\si{\pascal\per\metre}) & Longueur (\si{\metre}) & $\Delta P_{tot}$ (\si{\pascal}) \\
  \midrule
  Collecteur$\to$R+3 & 1,8 & 120 & 15 & 1\,800 \\
  R+3$\to$appart 1  & 0,15 & 80  & 12 & 960 \\
  Circuit PCBT       & 0,15 & 60  & 60 & 3\,600 \\
  \midrule
  \textbf{Circuit le plus défavorable} & & & & \textbf{6\,360} \\
  \bottomrule
\end{tabular}
\end{center}

\textbf{Sélection de la pompe :} $q_v = \SI{1,8}{\metre\cubed\per\hour}$, $\Delta P = \SI{6\,360}{\pascal} \approx \SI{6,4}{\metre} CE$.

\subsection{Vérification RE2020}

\begin{definition}[title={Indicateurs RE2020 à vérifier}]
  \begin{itemize}
    \item \textbf{Bbio} (Besoin bioclimatique) : efficacité de l'enveloppe — $\leq Bbio_{max}$
    \item \textbf{Cep} (Consommation d'énergie primaire) : $\leq Cep_{max}$
    \item \textbf{Cep,nr} (Énergie primaire non renouvelable) : $\leq Cep,nr_{max}$
    \item \textbf{Icénergie} (Impact carbone énergie) : $\leq$ seuil progressif
    \item \textbf{DH} (Degrés-Heures de confort d'été) : $\leq 1\,250$ DH (résidentiel)
  \end{itemize}
\end{definition}

Pour notre immeuble PAC (coefficient de conversion électricité : $f_{prim} = 2{,}3$) :
\begin{equation}
  Cep = \frac{W_{el} \times f_{prim}}{A_{ref}} = \frac{12\,615 \times 2{,}3}{1\,200} = \SI{24{,}2}{\kilo\watt\hour\per\metre\squared\per\year}
\end{equation}

Ce résultat est très inférieur au seuil RE2020 de \SI{70}{\kilo\watt\hour_{ep}\per\metre\squared\per\year} pour le chauffage en zone H1a.

% ============================================================
\section{Étude de cas 2 — Bâtiment tertiaire climatisé}
% ============================================================

\subsection{Présentation du projet}

\begin{exemple}[title={Contexte}]
  Immeuble de bureaux R+5, zone H2b (Lyon), surface utile : \SI{3\,000}{\metre\squared}.
  Occupation : 300 personnes, bureaux paysagers et salles de réunion.
  Système CVC : groupe d'eau glacée (GEG) + CTA double flux avec batterie froide
  + ventilo-convecteurs (VRV en complément), GTB centralisée.
  Objectif : confort d'été garanti ($T_{int} \leq \SI{26}{\celsius}$), labelisation HQE.
\end{exemple}

\subsection{Calcul des charges de refroidissement}

\textbf{Apports internes (été) :}
\begin{center}
\begin{tabular}{lcc}
  \toprule
  Source & Valeur unitaire & Total \\
  \midrule
  Personnes (sensible) & $70\,\si{\watt}/pers$ & $300 \times 70 = \SI{21\,000}{\watt}$ \\
  Éclairage LED & $8\,\si{\watt\per\metre\squared}$ & $3000 \times 8 = \SI{24\,000}{\watt}$ \\
  Bureautique & $150\,\si{\watt}/poste$ & $200 \times 150 = \SI{30\,000}{\watt}$ \\
  \midrule
  \textbf{Total apports internes} & & \SI{75\,000}{\watt} \\
  \bottomrule
\end{tabular}
\end{center}

\textbf{Apports solaires (façade Sud, fenêtres \SI{600}{\metre\squared}, $g = 0{,}35$, $G_{max} = \SI{600}{\watt\per\metre\squared}$) :}
\begin{equation}
  Q_{solaire} = g \times A_{vitrage} \times G_{max} = 0{,}35 \times 600 \times 600 = \SI{126\,000}{\watt}
\end{equation}

\textbf{Apports par renouvellement d'air (conditions : $T_{ext} = \SI{33}{\celsius}$, $T_{int} = \SI{26}{\celsius}$) :}
\begin{equation}
  q_v = 300\,pers \times \SI{25}{\litre\per\second\per\personne} = \SI{7{,}5}{\metre\cubed\per\second}
\end{equation}
\begin{equation}
  Q_{vent} = \rho \cdot c_p \cdot q_v \cdot (T_{ext} - T_{int}) = 1{,}2 \times 1005 \times 7{,}5 \times 7 = \SI{63\,315}{\watt}
\end{equation}

\textbf{Charge totale de refroidissement :}
\begin{equation}
  Q_f = 75\,000 + 126\,000 + 63\,315 = \SI{264\,315}{\watt} \approx \SI{265}{\kilo\watt}
\end{equation}

\subsection{Sélection du groupe d'eau glacée}

\textbf{Puissance frigorifique nécessaire (marge 10\,\%) :}
\begin{equation}
  Q_{GEG} = 265 \times 1{,}10 = \SI{292}{\kilo\watt} \rightarrow \text{GEG 300 kW}
\end{equation}

\textbf{COP du groupe froid à $5/12$\,\si{\celsius} :} $\mathrm{EER} = 3{,}0$ (classe A, norme EN 14511).

\textbf{Puissance absorbée :}
\begin{equation}
  P_{el} = \frac{Q_{GEG}}{EER} = \frac{300}{3{,}0} = \SI{100}{\kilo\watt}
\end{equation}

\textbf{Circuit eau glacée :} température $5/12\,\si{\celsius}$, débit :
\begin{equation}
  \dot{m}_{EG} = \frac{Q_{GEG}}{\rho \cdot c_p \cdot \Delta T} = \frac{300\,000}{1000 \times 4186 \times 7} = \SI{10{,}2}{\metre\cubed\per\hour}
\end{equation}

\subsection{Dimensionnement de la CTA double flux}

\begin{methode}[title={Dimensionnement d'une CTA}]
  \begin{enumerate}
    \item Déterminer le débit d'air neuf $q_{v,AN}$ (réglementation aéraulique, besoins hygiéniques).
    \item Choisir les températures de soufflage ($T_{soufflage}$) et de reprise.
    \item Dimensionner la batterie froide : $Q_{batterie} = \rho \cdot c_p \cdot q_v \cdot (T_{reprise} - T_{soufflage})$.
    \item Sélectionner le ventilateur (débit + pression disponible).
    \item Vérifier le bilan enthalpique sur le diagramme de l'air humide.
  \end{enumerate}
\end{methode}

\textbf{Débit air total CTA :} $q_{v,total} = \SI{25\,000}{\metre\cubed\per\hour}$ (8 vol/h + hygiénique).

\textbf{Température de soufflage : \SI{15}{\celsius}} (reprise \SI{26}{\celsius}).

\textbf{Puissance de la batterie froide de la CTA :}
\begin{equation}
  Q_{batterie} = \frac{1{,}2 \times 1005 \times 25\,000}{3600} \times (26 - 15) = \SI{92\,125}{\watt} \approx \SI{92}{\kilo\watt}
\end{equation}

\textbf{Perte de charge réseau aéraulique (circuit le plus défavorable) :}
\begin{center}
\begin{tabular}{lccc}
  \toprule
  Tronçon & Débit (\si{\metre\cubed\per\hour}) & $\Delta P$ (\si{\pascal}) & Vitesse (\si{\metre\per\second}) \\
  \midrule
  Gaine principale & 25\,000 & 180 & 6,0 \\
  Gaines secondaires & 5\,000 & 95 & 4,5 \\
  Diffuseurs & — & 20 & — \\
  CTA (filtres+batteries) & — & 350 & — \\
  \midrule
  \textbf{Total} & & \textbf{645} & \\
  \bottomrule
\end{tabular}
\end{center}

\textbf{Sélection du ventilateur :} $q_v = \SI{25\,000}{\metre\cubed\per\hour}$, $\Delta P_{tot} = \SI{645}{\pascal}$.
Puissance ventilateur (rendement $\eta = 0{,}75$) :
\begin{equation}
  P_{vent} = \frac{q_v \times \Delta P_{tot}}{\eta} = \frac{(25\,000/3600) \times 645}{0{,}75} = \SI{5\,972}{\watt} \approx \SI{6}{\kilo\watt}
\end{equation}

\subsection{Intégration GTB}

\begin{definition}[title={Architecture GTB pour bâtiment tertiaire}]
  La GTB centralise la supervision et la commande de toutes les installations CVC :
  \begin{itemize}
    \item \textbf{Niveau terrain} : capteurs (sondes de température, CO\textsubscript{2}, humidité, pression), actionneurs (vannes, volets, variateurs de vitesse) ;
    \item \textbf{Niveau automation} : automates DDC (Direct Digital Control) gérant les boucles de régulation locales ;
    \item \textbf{Niveau supervision} : IHM de surveillance et de pilotage, archivage des données, alarmes.
  \end{itemize}
  Protocoles courants : BACnet (IP ou MS/TP), Modbus RTU/TCP, KNX.
\end{definition}

\textbf{Fonctions GTB pour ce projet :}
\begin{itemize}
  \item Programmation horaire des CTA (démarrage anticipé, relance matinale) ;
  \item Régulation de la température d'eau glacée en fonction de la charge (reset) ;
  \item Free-cooling aéraulique ($T_{ext} < T_{int} - \SI{2}{\celsius}$) ;
  \item Comptage des consommations électriques (GEG, CTA, éclairage) ;
  \item Alarmes de dépassement de température et de pression.
\end{itemize}

% ============================================================
\section{Étude de cas 3 — Installation industrielle}
% ============================================================

\subsection{Présentation du projet}

\begin{exemple}[title={Contexte}]
  Atelier de production pharmaceutique, surface : \SI{800}{\metre\squared}.
  Exigences : température $22 \pm \SI{1}{\celsius}$, humidité relative $50 \pm 5\,\%$,
  surpression par rapport aux zones adjacentes, renouvellement d'air conforme BPF
  (Bonnes Pratiques de Fabrication — 20 volumes/heure minimum).
  Énergie : vapeur process à \SI{6}{\bar} disponible sur site, eau glacée \SI{6/12}{\celsius}.
\end{exemple}

\subsection{Bilan thermique et massique}

\textbf{Charges internes spécifiques (process) :}
\begin{center}
\begin{tabular}{lcc}
  \toprule
  Équipement & Puissance dissipée (\si{\kilo\watt}) & Fonctionnement \\
  \midrule
  Malaxeurs (4 $\times$ 15 kW) & 45 (rendement 75\,\%) & 16\,h/j \\
  Éclairage salles blanches    & 12 & 16\,h/j \\
  Personnel (50 pers $\times$ 100 W) & 5 & 8\,h/j \\
  Pertes parois (isolation $U_{moy}=0,3$) & 8 & continu \\
  \midrule
  \textbf{Total pointe} & \textbf{70} & \\
  \bottomrule
\end{tabular}
\end{center}

\textbf{Débit d'air neuf (20 vol/h) :}
\begin{equation}
  q_{v,AN} = 20 \times \SI{800}{\metre\squared} \times \SI{3}{\metre} = \SI{48\,000}{\metre\cubed\per\hour}
\end{equation}

\textbf{Charge frigorifique totale (été, $T_{ext} = \SI{33}{\celsius}$) :}
\begin{equation}
  Q_{f,total} = Q_{process} + Q_{vent} = 70 + \frac{1{,}2 \times 1005 \times (48\,000/3600) \times (33-22)}{1000} = 70 + \SI{176}{\kilo\watt} = \SI{246}{\kilo\watt}
\end{equation}

\subsection{Traitement d'air et hygrométrie}

\begin{methode}[title={Contrôle hygrométrique en salle propre}]
  \begin{enumerate}
    \item Calculer la teneur en vapeur d'eau requise : $w_{cible} = f(T_{int}, HR_{cible})$ — diagramme de l'air humide.
    \item Déshumidification si $w_{ext} > w_{cible}$ : batterie froide amenée à $T_{rosée}$, puis remontée en température.
    \item Humidification si $w_{ext} < w_{cible}$ : humidificateur vapeur (vapeur propre, sans additifs chimiques en pharmaceutique).
    \item Régulation proportionnelle-intégrale de l'humidité relative avec sonde HR (\SI{\pm 1}{\percent} de précision).
  \end{enumerate}
\end{methode}

\textbf{Calcul de la teneur en eau à $22\,\si{\celsius}$, $HR = 50\,\%$ :}
\begin{equation}
  P_{vs}(22°C) = \SI{2\,643}{\pascal}
\end{equation}
\begin{equation}
  w = 0{,}622 \times \frac{HR \times P_{vs}}{P_{atm} - HR \times P_{vs}} = 0{,}622 \times \frac{0{,}5 \times 2643}{101\,325 - 0{,}5 \times 2643} = \SI{8{,}2}{\gram\per\kilogram}
\end{equation}

\subsection{Efficacité énergétique et économies d'énergie}

\begin{definition}[title={PUE et indicateurs performance installation industrielle}]
  \begin{itemize}
    \item \textbf{SEER} (Seasonal Energy Efficiency Ratio) : EER annuel moyenné sur les conditions climatiques.
    \item \textbf{Récupération de chaleur} sur air extrait : réducteur des consommations de chauffage/humidification.
    \item \textbf{Variateurs de vitesse} sur ventilateurs et pompes : économie proportionnelle au cube de la réduction de vitesse.
    \item \textbf{Free-cooling} : utilisation de l'air extérieur frais pour le refroidissement sans compresseur.
  \end{itemize}
\end{definition}

\textbf{Économie par récupération de chaleur sur air extrait (échangeur à plaques, $\eta = 0,75$) :}
\begin{equation}
  Q_{recup} = 0{,}75 \times \frac{1{,}2 \times 1005 \times 48\,000}{3600} \times (22 - (-5)) = \SI{362}{\kilo\watt}
\end{equation}

Cette économie permet de réduire considérablement la puissance de la batterie chaude vapeur en hiver.

\textbf{Économie annuelle estimée (chauffage et humidification) :}
\begin{equation}
  W_{econo} = 362 \times 3\,000\,h/an \times 10^{-3} \approx \SI{1\,086}{\mega\watt\hour\per\year}
\end{equation}

\subsection{Sécurités spécifiques installation industrielle}

\begin{attention}
  En milieu industriel et pharmaceutique, les sécurités CVC sont critiques pour la
  production et la sécurité des personnels :
  \begin{itemize}
    \item \textbf{Surpression salle blanche} : maintenir $+15\,\si{\pascal}$ par rapport aux zones adjacentes (évite contamination).
    \item \textbf{Alarme température} : dépassement de $22 \pm \SI{2}{\celsius}$ déclenche une alarme de niveau 2 avec traçabilité.
    \item \textbf{Filtre à haute efficacité} : filtration HEPA H14 ($\geq 99{,}995\,\%$ de rétention à \SI{0,3}{\micro\metre}).
    \item \textbf{Redondance des équipements} : doublement des ventilateurs (N+1), groupe froid de secours.
    \item \textbf{Qualification IQ/OQ/PQ} : installation CVC qualifiée selon les référentiels BPF (FDA, EMA).
  \end{itemize}
\end{attention}

% ============================================================
\section{Éléments de réglementation}
% ============================================================

\subsection{Réglementation thermique RE2020}

La RE2020 (entrée en vigueur le 1\textsuperscript{er} janvier 2022) fixe des exigences de performance énergétique et environnementale pour les bâtiments neufs :

\begin{center}
\begin{tabular}{lp{7cm}l}
  \toprule
  Indicateur & Définition & Exigence type \\
  \midrule
  Bbio & Besoin bioclimatique (enveloppe seule) & $\leq Bbio_{max}$ \\
  Cep & Consommation énergie primaire totale & $\leq Cep_{max}$ \\
  Cep,nr & Part non renouvelable de Cep & $\leq Cep,nr_{max}$ \\
  Icénergie & Impact CO\textsubscript{2} des consommations & $\leq$ seuil progressif \\
  Iconstruction & Impact CO\textsubscript{2} des matériaux de construction & $\leq$ seuil progressif \\
  DH & Degrés-Heures inconfort d'été & $\leq 1\,250$ DH résidentiel \\
  \bottomrule
\end{tabular}
\end{center}

\subsection{Réglementation aéraulique et qualité de l'air intérieur (QAI)}

\begin{itemize}
  \item \textbf{Arrêté du 24 mars 1982} (modifié) : débits minimaux d'air neuf par local et par usage.
  \item \textbf{Décret n°2012-14} : obligation de vérification des systèmes de ventilation mécanique.
  \item \textbf{NF EN 16798-1} : paramètres d'entrée pour la conception des bâtiments (QAI, catégories IDA 1 à 4).
  \item \textbf{Arrêté 2011 Établissements recevant du public} : surveillance CO\textsubscript{2} et COVT obligatoire dans les ERP de catégories 1 à 4.
\end{itemize}

\subsection{Fluides frigorigènes et réglementation F-gas}

\begin{definition}[title={Règlement F-gas (UE) n°517/2014 et révision 2024}]
  Calendrier de réduction progressive des HFC à fort PRG (Potentiel de Réchauffement Global) :
  \begin{itemize}
    \item \textbf{2025} : interdiction équipements neufs charge $>$ 3 kg avec GWP $>$ 750.
    \item \textbf{2027} : interdiction GWP $>$ 150 pour équipements hermétiques.
    \item \textbf{2032} : interdiction générale HFC pour équipements neufs dans de nombreuses applications.
  \end{itemize}
  Alternatives : R-290 (propane, GWP=3), R-32 (GWP=675), R-1234yf (GWP<1), R-744 (CO\textsubscript{2}, GWP=1).
\end{definition}

\begin{attention}
  Toute intervention sur un circuit frigorifique de charge $\geq \SI{5}{\tonne}$ CO\textsubscript{2}
  équivalent nécessite une attestation de capacité délivrée par un organisme accrédité
  (Certifrigo, Qualifrigo). L'enregistrement des fluides récupérés est obligatoire.
\end{attention}

\subsection{Normes de conception}

\begin{center}
\begin{tabular}{ll}
  \toprule
  Norme & Domaine \\
  \midrule
  NF EN 12831-1 & Calcul des déperditions thermiques de base \\
  NF EN 12831-3 & Calcul des besoins de chauffage des locaux \\
  NF EN 15243 & Calcul de la charge de froid \\
  NF EN 14511 & Performances des climatiseurs et PAC \\
  NF EN 16798-3 & Conception des systèmes de ventilation \\
  DTU 65.11 & Dispositifs de sécurité installations chauffage \\
  DTU 65.12 & Réseaux de tuyauteries en acier \\
  \bottomrule
\end{tabular}
\end{center}


\end{document}
